\documentclass[11pt,a4paper,oneside]{book}

% ============================================================================
% PACKAGES
% ============================================================================
\usepackage[utf8]{inputenc}
\usepackage[T1]{fontenc}
\usepackage[french]{babel}
\usepackage{lmodern}
\usepackage{microtype}
\usepackage[top=2.5cm,bottom=2.5cm,left=2.5cm,right=2.5cm,headheight=15pt]{geometry}

\usepackage{amsmath,amsfonts,amssymb,amsthm}
\usepackage{mathtools}
\usepackage{gensymb}
\usepackage{icomma}
\usepackage{cancel}
\usepackage{multicol}
\usepackage{enumitem}
\usepackage{booktabs}
\usepackage{array}
\usepackage{graphicx}
\usepackage{xcolor}
\usepackage{tikz}
\usepackage{pgfplots}
\pgfplotsset{compat=1.18}
\usetikzlibrary{angles,quotes,calc,positioning,arrows.meta,patterns,decorations.pathreplacing,babel}

\usepackage[most]{tcolorbox}
\usepackage{fancyhdr}
\usepackage{titlesec}
\usepackage{hyperref}

% ============================================================================
% COULEURS & STYLE UCLOUVAIN
% ============================================================================
\definecolor{uclblue}{RGB}{0, 53, 107}
\definecolor{ucllightblue}{RGB}{235, 243, 250}
\definecolor{uclaccent}{RGB}{0, 114, 188}
\definecolor{defgreen}{RGB}{20, 110, 60}
\definecolor{deflightgreen}{RGB}{238, 248, 242}
\definecolor{exorange}{RGB}{180, 85, 0}
\definecolor{exlightorange}{RGB}{253, 246, 238}
\definecolor{warnred}{RGB}{170, 30, 30}
\definecolor{warnlightred}{RGB}{253, 240, 240}
\definecolor{darkgray}{RGB}{60, 60, 60}

\hypersetup{
    colorlinks=true,
    linkcolor=uclblue,
    urlcolor=uclaccent,
    citecolor=uclblue,
    pdftitle={Syllabus de Trigonométrie - Examen d'entrée Ingénieur Civil UCLouvain},
    pdfauthor={Hadrien Rattez, Jean-François Remacle et Pierre Schaus (Prompt Engineers) & IA}
}

% ============================================================================
% BOÎTES PÉDAGOGIQUES (TCOLORBOX)
% ============================================================================
\newtcolorbox{theoreme}[1][]{
    enhanced,
    breakable,
    colback=ucllightblue,
    colframe=uclblue,
    fonttitle=\bfseries,
    title={#1},
    arc=3pt,
    boxrule=0.8pt,
    left=8pt,right=8pt,top=6pt,bottom=6pt
}

\newtcolorbox{definition}[1][]{
    enhanced,
    breakable,
    colback=deflightgreen,
    colframe=defgreen,
    fonttitle=\bfseries,
    title={#1},
    arc=3pt,
    boxrule=0.8pt,
    left=8pt,right=8pt,top=6pt,bottom=6pt
}

\newtcolorbox{exemple}[1][]{
    enhanced,
    breakable,
    colback=exlightorange,
    colframe=exorange,
    fonttitle=\bfseries,
    title={#1},
    arc=3pt,
    boxrule=0.8pt,
    left=8pt,right=8pt,top=6pt,bottom=6pt
}

\newtcolorbox{attention}[1][Attention]{
    enhanced,
    breakable,
    colback=warnlightred,
    colframe=warnred,
    fonttitle=\bfseries,
    title={#1},
    arc=3pt,
    boxrule=0.8pt,
    left=8pt,right=8pt,top=6pt,bottom=6pt
}

\newtcolorbox{astuce}[1][Astuce \& Méthode]{
    enhanced,
    breakable,
    colback=gray!8,
    colframe=darkgray,
    fonttitle=\bfseries,
    title={#1},
    arc=3pt,
    boxrule=0.8pt,
    left=8pt,right=8pt,top=6pt,bottom=6pt
}

% ============================================================================
% MISE EN PAGE & EN-TÊTES
% ============================================================================
\setlength{\parindent}{0pt}
\setlength{\parskip}{0.45em}
\setlist[itemize]{topsep=3pt,itemsep=2pt}
\setlist[enumerate]{topsep=4pt,itemsep=4pt}

\pagestyle{fancy}
\fancyhf{}
\fancyhead[L]{\small\textcolor{uclblue}{\nouppercase{\leftmark}}}
\fancyhead[R]{\small\textcolor{darkgray}{Trigonométrie \& Géométrie -- UCLouvain}}
\fancyfoot[C]{\thepage}
\renewcommand{\headrulewidth}{0.4pt}

\titleformat{\chapter}[display]
  {\normalfont\huge\bfseries\color{uclblue}}
  {\filright\large\textsc{\chaptertitlename}\ \thechapter}
  {1ex}
  {\titlerule\vspace{1.2ex}\filright}
  [\vspace{1ex}\titlerule]

\titleformat{\section}
  {\normalfont\Large\bfseries\color{uclblue}}
  {\thesection}{0.8em}{}

\titleformat{\subsection}
  {\normalfont\large\bfseries\color{uclaccent}}
  {\thesubsection}{0.7em}{}

% ============================================================================
% COMMANDES MATHÉMATIQUES
% ============================================================================
\newcommand{\R}{\mathbb{R}}
\newcommand{\Z}{\mathbb{Z}}
\newcommand{\N}{\mathbb{N}}
\newcommand{\Q}{\mathbb{Q}}
\newcommand{\C}{\mathbb{C}}
\newcommand{\wideparen}[1]{\overset{\frown}{#1}}
\DeclareMathOperator{\cotg}{cotg}
\DeclareMathOperator{\tg}{tg}
\DeclareMathOperator{\cosec}{cosec}

% ============================================================================
% DOCUMENT
% ============================================================================
\begin{document}

% ----------------------------------------------------------------------------
% PAGE DE TITRE
% ----------------------------------------------------------------------------
\begin{titlepage}
    \centering
    \vspace*{1cm}
    {\Large \textsc{Université catholique de Louvain (UCLouvain)}}\\[0.3cm]
    {\large \textsc{École Polytechnique de Louvain (EPL)}}\\[1.5cm]

    \begin{tcolorbox}[colback=ucllightblue,colframe=uclblue,arc=5pt,boxrule=1.5pt,top=15pt,bottom=15pt]
        \centering
        {\huge\bfseries\color{uclblue} Syllabus de Trigonométrie\\ et de Géométrie}\\[0.4cm]
        {\Large\color{uclaccent} Préparation à l'examen d'admission\\ aux études d'ingénieur civil}
    \end{tcolorbox}

    \vspace{0.9cm}

    % Illustration centrale : cercle trigonométrique esthétique
    \begin{tikzpicture}[scale=2.1, >=Stealth]
        \draw[step=0.5, gray!20, very thin] (-1.4,-1.4) grid (1.4,1.4);
        \draw[->, thick, darkgray] (-1.4,0) -- (1.5,0) node[right] {$\cos\alpha$};
        \draw[->, thick, darkgray] (0,-1.4) -- (0,1.5) node[above] {$\sin\alpha$};
        \draw[thick, uclblue] (0,0) circle (1);
        \draw[thick, defgreen] (1,-1.3) -- (1,1.35) node[above] {$t (\tan)$};
        
        \def\ang{38}
        \coordinate (O) at (0,0);
        \coordinate (P) at ({\ang}:1);
        \coordinate (Px) at ({cos(\ang)},0);
        \coordinate (Py) at (0,{sin(\ang)});
        \coordinate (T) at (1,{tan(\ang)});

        \fill[uclaccent!15] (0,0) -- (0.35,0) arc (0:\ang:0.35) -- cycle;
        \draw[uclaccent, thick] (0.35,0) arc (0:\ang:0.35);
        \node[uclaccent] at ({\ang/2}:0.47) {$\alpha$};

        \draw[dashed, gray] (P) -- (Px);
        \draw[dashed, gray] (P) -- (Py);
        \draw[very thick, exorange] (0,0) -- (Px) node[midway, below] {\small $\cos\alpha$};
        \draw[very thick, warnred] (0,0) -- (Py) node[midway, left] {\small $\sin\alpha$};
        \draw[very thick, defgreen] (1,0) -- (T) node[midway, right] {\small $\tan\alpha$};
        \draw[thick, uclaccent] (O) -- (T);

        \fill[uclblue] (P) circle (1.2pt) node[above left] {\small $P(\cos\alpha,\sin\alpha)$};
        \fill[defgreen] (T) circle (1.2pt) node[right] {\small $T$};
        \fill[darkgray] (1,0) circle (1pt) node[below right] {\small $I$};
        \fill[darkgray] (0,1) circle (1pt) node[above left] {\small $J$};
    \end{tikzpicture}

    \vfill

    {\large \textbf{Direction artistique \& \emph{« Prompt Engineering »} :}}\\[0.2cm]
    {\Large \textbf{L'équipe trigonométrie \& géométrie de l'examen d'entrée}}\\[0.15cm]
    {\large \textbf{Hadrien Rattez}, \textbf{Jean-François Remacle} et \textbf{Pierre Schaus}}\\[0.3cm]
    {\normalsize \emph{Rédigé à la sueur des GPU par l'Intelligence Artificielle,\\ sous la supervision bienveillante (et quelques prompts bien sentis) de ses humains.}}

    \vspace{0.6cm}
    {\small \textbf{Sources fournies à l'IA :}\\
    L'excellent syllabus \emph{« Synthèse de trigonométrie »} d'Yvan Haine et Pierre Joris (ULiège, 2012)\\
    a été donné en source principale à l'IA, conjointement aux supports de cours et annales de trigonométrie et de géométrie de l'UCLouvain.}\\[0.4cm]
    {\normalsize Édition 2026--2027}
\end{titlepage}

% ----------------------------------------------------------------------------
% AVANT-PROPOS ET SOURCES
% ----------------------------------------------------------------------------
\chapter*{Avant-propos et conseils méthodologiques}
\addcontentsline{toc}{chapter}{Avant-propos et conseils méthodologiques}
\markboth{Avant-propos}{Avant-propos}

Ce syllabus a été conçu pour accompagner les étudiantes et étudiants qui préparent l'épreuve fusionnée de \textbf{trigonométrie et de géométrie} de l'examen spécial d'admission aux études d'ingénieur civil à l'\textbf{UCLouvain} (École Polytechnique de Louvain -- EPL). Il est pensé pour être agréable à lire, progressif et directement accessible à toute personne ayant suivi le cours de mathématiques de 5\ieme{} et 6\ieme{} secondaire (en particulier l'option 6h/semaine, ou 4h avec un complément d'investissement) en Communauté française de Belgique (Fédération Wallonie-Bruxelles).

\section*{Origine, auteurs et sources du document}

\begin{theoreme}[« Prompt Engineers », IA et sources du document]
\begin{itemize}
    \item \textbf{\emph{Prompt Engineers} en chef :} L'équipe trigonométrie et géométrie de l'examen d'entrée à l'UCLouvain, à savoir les professeurs \textbf{Hadrien Rattez}, \textbf{Jean-François Remacle} et \textbf{Pierre Schaus}, qui ont fourni la matière première, les idées, les annales et les instructions nécessaires pour dompter la machine.
    \item \textbf{Auteur exécutant (Intelligence Artificielle) :} Soyons honnêtes : la rédaction au kilomètre, la mise en forme \LaTeX{}, le tracé millimétré des figures TikZ et la rédaction détaillée des corrigés ont été réalisés \textbf{en très grande partie par l'intelligence artificielle} (qui, contrairement aux candidats à l'examen d'entrée, avait droit à une puissance de calcul illimitée).
    \item \textbf{Sources et remerciements :} L'excellent syllabus \textbf{\emph{« Synthèse de trigonométrie »} rédigé par Yvan Haine et Pierre Joris (Août 2012) à l'Université de Liège (ULiège)}, suite à une suggestion du Professeur E. Delhez, a été directement donné en source à l'IA pour structurer et inspirer la partie trigonométrique de cet ouvrage. Nous tenons à saluer la grande qualité pédagogique de leur travail, dont nous avons repris la progression ainsi que plusieurs exemples et exercices classiques (identifiés par la mention \emph{ULiège}).
    \item \textbf{Annales et formulaire UCLouvain :} Ce syllabus intègre également le formulaire officiel des cours préparatoires de l'UCLouvain (étendu à la géométrie synthétique et analytique plane et spatiale), les astuces présentées lors des séances de préparation (\emph{slides 2026}), ainsi que l'intégralité des questions types issues des \textbf{examens d'entrée de Trigonométrie et de Géométrie de l'UCLouvain de 2018 à 2026}. L'origine et l'année de chaque question d'examen sont systématiquement indiquées.
\end{itemize}
\end{theoreme}

\section*{Structure de l'épreuve fusionnée de Trigonométrie et Géométrie à l'UCLouvain}

Avec la fusion des épreuves de trigonométrie et de géométrie à l'examen d'admission de l'EPL (UCLouvain), l'épreuve se déroule \textbf{sans calculatrice} et couvre quatre grands piliers complémentaires :
\begin{enumerate}
    \item \textbf{Les fondements, formules et équations trigonométriques (Chapitres 1 à 3)} : cercle trigonométrique, fonctions réciproques ($\arcsin, \arccos, \arctan$), formules d'addition, de duplication, de Carnot et de Simpson, et résolution complète d'équations trigonométriques dans $\R$ (avec conditions d'existence et représentation sur le cercle).
    \item \textbf{La résolution de triangles et la géométrie synthétique plane et spatiale (Chapitres 4 et 5)} : triangles rectangles et quelconques (lois des sinus et des cosinus), angles inscrits et au centre dans le cercle, puissance d'un point par rapport à un cercle, relations métriques, triangles semblables, rapports d'aires et sections planes de solides (sphères, cônes, cylindres).
    \item \textbf{La géométrie analytique dans le plan et les coniques (Chapitre 6)} : vecteurs et produit scalaire dans $\R^2$, droites, médiatrices, distance d'un point à une droite, cercles, et étude complète des coniques (ellipse, hyperbole, parabole : réduction par complétion du carré, foyers, directrices, excentricité, asymptotes et tangentes par la condition $\Delta = 0$).
    \item \textbf{La géométrie analytique et vectorielle dans l'espace $\R^3$ (Chapitre 7)} : produit scalaire, produit vectoriel et déterminant, équations paramétriques et cartésiennes de droites et de plans, vecteur normal, point de percée, projections orthogonales, réflexion sur un miroir plan, distances et sphères.
\end{enumerate}

\section*{Notations et conventions en vigueur à l'examen}

Pour éviter toute ambiguïté le jour de l'examen, familiarisez-vous dès maintenant avec les conventions utilisées dans les énoncés officiels de l'UCLouvain :
\begin{itemize}
    \item $\sin^n x$ signifie $(\sin(x))^n$, et de même $\cos^n x = (\cos(x))^n$, $\tan^n x = (\tan(x))^n$. \textbf{Attention :} cette notation ne s'applique qu'aux exposants positifs ($n \ge 2$). La fonction réciproque se note exclusivement $\arcsin(x)$, $\arccos(x)$, $\arctan(x)$ (et non $\sin^{-1}(x)$ qui prête à confusion avec $\frac{1}{\sin x}$).
    \item La longueur d'un segment reliant les points $A$ et $B$ est notée $\overline{AB}$ ou $|AB|$.
    \item Un angle géométrique au sommet $A$ entre les segments $[AB]$ et $[AC]$ est noté $\widehat{BAC}$ ou simplement $\widehat{A}$.
    \item Dans un triangle quelconque $ABC$, les longueurs $a = \overline{BC}$, $b = \overline{AC}$ et $c = \overline{AB}$ désignent toujours les côtés opposés respectivement aux sommets $A$, $B$ et $C$.
    \item Les angles sont exprimés en \textbf{radians} par défaut, sauf présence explicite du symbole degré (par exemple $15\degree$). \textbf{Ne mélangez jamais degrés et radians dans une même formule !}
\end{itemize}

\section*{Comment utiliser ce syllabus ?}

Une étude « par cœur » des formules sans compréhension géométrique ne suffit pas pour réussir l'examen. Nous vous conseillons de :
\begin{enumerate}
    \item Comprendre comment chaque formule se visualise sur une figure (cercle trigonométrique, coupe plane, schéma vectoriel) ou se retrouve instantanément.
    \item Étudier attentivement les \textbf{exemples résolus} au fil des chapitres : ce sont d'authentiques questions d'examen de l'UCLouvain dont nous détaillons non seulement le calcul, mais surtout \emph{l'intuition qui guide le choix de la méthode}.
    \item Utiliser le \textbf{Formulaire imprimable de 4 pages} (Chapitre 8) lors de vos révisions, et chercher par vous-même les \textbf{exercices de fin de chapitre} avant de consulter les \textbf{corrigés détaillés} regroupés au Chapitre 9, en fin d'ouvrage.
\end{enumerate}

\tableofcontents

% ----------------------------------------------------------------------------
% CHAPITRES
% ----------------------------------------------------------------------------
% ============================================================================
% CHAPITRE 1 : ANGLES, ARCS ET CERCLE TRIGONOMÉTRIQUE
% ============================================================================
\chapter{Angles, Arcs et Cercle Trigonométrique}

Ce premier chapitre pose les bases indispensables de la trigonométrie : la mesure des angles orientés, la géométrie des arcs et secteurs circulaires, la définition rigoureuse des nombres trigonométriques sur le cercle de rayon unité, ainsi que les fonctions trigonométriques réciproques ($\arcsin$, $\arccos$, $\arctan$).

% ----------------------------------------------------------------------------
\section{Angles orientés et unités de mesure}
% ----------------------------------------------------------------------------

\subsection{Notion d'angle orienté}

\begin{definition}[Angle orienté et sens trigonométrique]
Un \textbf{angle orienté} de sommet $O$ est un couple ordonné de deux demi-droites de même origine, noté $([OA, [OB)$. La demi-droite $[OA$ est le \emph{premier côté} (ou côté origine) et $[OB$ est le \emph{second côté} (ou côté extrémité).

Dans le plan, on distingue deux sens de rotation :
\begin{itemize}
    \item Le \textbf{sens positif} (ou \textbf{sens trigonométrique}) est le sens \emph{anti-horlogique} (contraire à celui des aiguilles d'une montre).
    \item Le \textbf{sens négatif} est le sens \emph{horlogique} (celui des aiguilles d'une montre).
\end{itemize}
\end{definition}

\begin{center}
\begin{tikzpicture}[scale=1.1, >=Stealth]
    % Sens positif
    \begin{scope}[shift={(0,0)}]
        \coordinate (O) at (0,0);
        \coordinate (A) at (2.8,0);
        \coordinate (B) at (1.5,2.2);
        \draw[thick] (O) -- (A) node[pos=0.85, below] {$A$};
        \draw[thick] (O) -- (B) node[pos=0.85, above left] {$B$};
        \fill (O) circle (1.5pt) node[below left] {$O$};
        \fill (2.3,0) circle (1.5pt);
        \fill (1.25,1.83) circle (1.5pt);
        \draw[->, thick, uclaccent] (1.1,0) arc (0:55.7:1.1);
        \node[uclaccent, right] at (28:1.15) {$\alpha > 0$ (sens $+$)};
    \end{scope}

    % Sens négatif
    \begin{scope}[shift={(6.5,0.8)}]
        \coordinate (O) at (0,0);
        \coordinate (A) at (2.8,0);
        \coordinate (B) at (1.6,-1.8);
        \draw[thick] (O) -- (A) node[pos=0.85, above] {$A$};
        \draw[thick] (O) -- (B) node[pos=0.85, below left] {$B$};
        \fill (O) circle (1.5pt) node[above left] {$O$};
        \fill (2.3,0) circle (1.5pt);
        \fill (1.3,-1.46) circle (1.5pt);
        \draw[->, thick, warnred] (1.1,0) arc (0:-48.4:1.1);
        \node[warnred, right] at (-24:1.15) {$\beta < 0$ (sens $-$)};
    \end{scope}
\end{tikzpicture}
\end{center}

Rappelons trois angles particuliers fondamentaux :
\begin{itemize}
    \item L'\textbf{angle nul} a ses deux côtés superposés (amplitude $0$).
    \item L'\textbf{angle plat} a ses deux côtés dans le prolongement opposé l'un de l'autre.
    \item L'\textbf{angle droit} a une amplitude égale à la moitié de celle d'un angle plat.
\end{itemize}

\subsection{Le degré et le radian}

On utilise deux unités principales pour mesurer l'amplitude d'un angle :

\begin{enumerate}
    \item \textbf{Le degré ($\degree$) :} C'est l'unité telle que l'angle plat mesure $180\degree$ (et le tour complet $360\degree$, l'angle droit $90\degree$).
    \begin{itemize}
        \item En écriture \emph{décimale} (DD), la partie fractionnaire est notée après la virgule : par exemple $3,25\degree$.
        \item En écriture \emph{sexagésimale} (DMS : degrés, minutes, secondes), on subdivise le degré en $60$ minutes d'arc ($1\degree = 60'$) et la minute en $60$ secondes d'arc ($1' = 60''$). Ainsi :
        \[
            3\degree 15' = 3\degree + \frac{15}{60}\degree = 3,25\degree \qquad \text{et} \qquad 1\degree 15' = 1,25\degree = \frac{5}{4}\degree.
        \]
    \end{itemize}

    \item \textbf{Le radian ($\text{rad}$) :} C'est l'unité naturelle et incontournable en mathématiques supérieures et en physique.
\end{enumerate}

\begin{definition}[Le radian]
Le \textbf{radian} est l'amplitude d'un angle au centre d'un cercle qui intercepte sur la circonférence de ce cercle un arc de longueur égale à celle de son rayon $r$.
\end{definition}

\begin{center}
\begin{tikzpicture}[scale=1.2, >=Stealth]
    \def\R{1.6}
    \def\ang{57.2958} % 1 radian en degrés
    \draw[thick] (0,0) circle (\R);
    \draw[thick, uclblue] (0,0) -- (\R,0) node[midway, below] {$r$};
    \draw[thick, uclblue] (0,0) -- (\ang:\R) node[midway, above left] {$r$};
    \draw[very thick, warnred] (\R,0) arc (0:\ang:\R) node[midway, right=2pt] {arc $L = r$};
    \draw[->, thick, uclaccent] (0.55,0) arc (0:\ang:0.55);
    \node[uclaccent] at ({\ang/2}:0.85) {$1\text{ rad}$};
    \fill (0,0) circle (1.2pt) node[below left] {$O$};
\end{tikzpicture}
\end{center}

Puisque la circonférence complète d'un cercle de rayon $r$ vaut $2\pi r$, un tour complet ($360\degree$) correspond à $2\pi$ radians, et un angle plat ($180\degree$) correspond à $\pi$ radians.

\begin{theoreme}[Conversion degrés $\leftrightarrow$ radians]
La relation de proportionnalité fondamentale est :
\[
    \boxed{\pi \text{ radians} \longleftrightarrow 180\degree} \qquad \text{soit} \qquad 1\text{ rad} = \left(\frac{180}{\pi}\right)\degree \approx 57,3\degree \quad \text{et} \quad 1\degree = \frac{\pi}{180}\text{ rad}.
\]
\begin{center}
\renewcommand{\arraystretch}{1.4}
\begin{tabular}{@{}lccccccccc@{}}
\toprule
\textbf{Degrés} & $0\degree$ & $30\degree$ & $45\degree$ & $60\degree$ & $90\degree$ & $120\degree$ & $135\degree$ & $150\degree$ & $180\degree$ \\
\midrule
\textbf{Radians} & $0$ & $\dfrac{\pi}{6}$ & $\dfrac{\pi}{4}$ & $\dfrac{\pi}{3}$ & $\dfrac{\pi}{2}$ & $\dfrac{2\pi}{3}$ & $\dfrac{3\pi}{4}$ & $\dfrac{5\pi}{6}$ & $\pi$ \\
\bottomrule
\end{tabular}
\end{center}
\end{theoreme}

\begin{attention}[Règle d'or sur les unités]
Lorsque l'amplitude d'un angle est exprimée sous forme de fraction ou de multiple de $\pi$, on omet l'abréviation « rad ». Surtout, \textbf{on ne mélange jamais les degrés et les radians dans une même expression} ! Ainsi, lors de la résolution d'une équation :
\[
    \text{On écrit : } x = \frac{\pi}{6} + 2k\pi \quad \text{(ou } x = 30\degree + k\cdot 360\degree\text{)}, \qquad \text{mais \textbf{JAMAIS} } \cancel{\textcolor{warnred}{x = 30\degree + 2k\pi}}.
\]
\end{attention}

% ----------------------------------------------------------------------------
\section{Longueur d'un arc, aire d'un secteur et corde}
% ----------------------------------------------------------------------------

Un cercle de rayon $R$ possède une circonférence (périmètre) $P = 2\pi R$ et une aire totale $A = \pi R^2$. Par simple règle de proportionnalité par rapport à un angle au centre $\theta$ exprimé \textbf{en radians} (sur un tour complet de $2\pi$ radians), on obtient les formules suivantes, très fréquemment utilisées à l'examen d'entrée.

\begin{theoreme}[Arc, secteur, corde et segment circulaire]
Soit un cercle de centre $O$ et de rayon $R$, et soit un angle au centre $AOB$ d'amplitude $\theta$ \textbf{exprimée en radians} ($0 \le \theta \le 2\pi$) :
\begin{enumerate}
    \item \textbf{Longueur de l'arc intercepté $\overset{\frown}{AB}$ :}
    \[
        \boxed{L = R\theta}
    \]
    \item \textbf{Aire du secteur circulaire $OAB$ :}
    \[
        \boxed{A_{\text{secteur}} = \frac{1}{2} R^2 \theta}
    \]
    \item \textbf{Longueur de la corde $\overline{AB}$ (pour $0 \le \theta \le \pi$) :} En traçant la bissectrice issue de $O$ qui coupe perpendiculairement $[AB]$ en son milieu $M$, on forme deux triangles rectangles d'angle $\theta/2$, d'où :
    \[
        \boxed{c = \overline{AB} = 2R\sin\left(\frac{\theta}{2}\right)}
    \]
    \item \textbf{Aire du segment circulaire (lunule entre la corde $[AB]$ et l'arc $\overset{\frown}{AB}$) :}
    \[
        \boxed{A_{\text{segment}} = A_{\text{secteur}} - A_{\triangle OAB} = \frac{1}{2}R^2\bigl(\theta - \sin\theta\bigr)}
    \]
\end{enumerate}
\end{theoreme}

\begin{center}
\begin{tikzpicture}[scale=1.25, >=Stealth]
    \def\R{2.4}
    \def\ang{64}
    \coordinate (O) at (0,0);
    \coordinate (A) at (\R,0);
    \coordinate (B) at (\ang:\R);
    \coordinate (M) at ($(A)!0.5!(B)$);

    % Aire du segment circulaire
    \fill[uclaccent!20] (A) arc (0:\ang:\R) -- (A) -- cycle;
    % Aire du triangle OAB
    \fill[ucllightblue] (O) -- (A) -- (B) -- cycle;

    \draw[thick, uclblue] (O) -- (A) node[midway, below] {$R$};
    \draw[thick, uclblue] (O) -- (B) node[midway, above left] {$R$};
    \draw[very thick, warnred] (A) arc (0:\ang:\R) node[midway, right=4pt] {$L = R\theta$};
    \draw[thick, defgreen] (A) -- (B) node[midway, below left=-2pt] {$c = 2R\sin(\theta/2)$};
    \draw[dashed, gray] (O) -- (M);

    \draw[->, thick, uclblue] (0.65,0) arc (0:\ang:0.65);
    \node[uclblue] at ({\ang/2}:0.88) {$\theta$};

    \fill (O) circle (1.5pt) node[below left] {$O$};
    \fill (A) circle (1.5pt) node[below right] {$A$};
    \fill (B) circle (1.5pt) node[above] {$B$};
    \node[uclaccent!80!black] at ({\ang/2}:2.05) {\footnotesize Segment};
\end{tikzpicture}
\end{center}

\begin{exemple}[Application directe : UCLouvain -- Juillet 2025, Question 3c et Juillet 2026, Question 3a]
\textbf{1) Aire d'un segment circulaire (UCLouvain, Juillet 2025, Q3c) :}\\
\emph{Énoncé :} Considérons un triangle $OAB$ tel que $\overline{OA} = \overline{OB} = r$ et $\widehat{AOB} = \alpha$ (en radians). Déterminer l'aire de la région plane délimitée par le secteur circulaire de centre $O$, de rayon $r$ et d'angle $\alpha$, et le triangle $OAB$.\\
\emph{Résolution :} L'aire du secteur circulaire vaut $A_{\text{secteur}} = \frac{1}{2}r^2\alpha$. L'aire du triangle isocèle $OAB$ vaut $A_{\triangle OAB} = \frac{1}{2}\overline{OA}\cdot\overline{OB}\sin\alpha = \frac{1}{2}r^2\sin\alpha$ (voir Chapitre 4). La différence vaut donc :
\[
    \text{Aire} = \frac{1}{2}r^2(\alpha - \sin\alpha).
\]

\medskip
\textbf{2) Tangentes et arc majeur (UCLouvain, Juillet 2026, Q3a) :}\\
\emph{Énoncé :} Soit $\mathcal{C}$ un cercle de centre $O$ et de rayon $R=1$, et $P$ un point extérieur à $\mathcal{C}$. Les deux tangentes à $\mathcal{C}$ issues de $P$ touchent le cercle en $A$ et $B$. La longueur du plus grand des deux arcs délimités par $A$ et $B$ vaut $\frac{4\pi}{3}$. Déterminer la longueur commune $\overline{PA} = \overline{PB}$.\\
\emph{Résolution :} La circonférence totale du cercle de rayon $1$ vaut $2\pi$. Puisque l'arc majeur mesure $\frac{4\pi}{3}$, l'arc mineur $\overset{\frown}{AB}$ mesure $2\pi - \frac{4\pi}{3} = \frac{2\pi}{3}$. Comme $R=1$, l'angle au centre $\widehat{AOB}$ vaut $\theta = \frac{L}{R} = \frac{2\pi}{3}$.
Par symétrie, la demi-droite $[OP)$ est bissectrice de $\widehat{AOB}$, donc $\widehat{AOP} = \frac{\theta}{2} = \frac{\pi}{3}$. De plus, la tangente $(PA)$ est perpendiculaire au rayon $[OA]$, de sorte que le triangle $OAP$ est rectangle en $A$. On en déduit :
\[
    \tan\left(\widehat{AOP}\right) = \frac{\overline{PA}}{\overline{OA}} \implies \overline{PA} = 1 \cdot \tan\left(\frac{\pi}{3}\right) = \sqrt{3}.
\]
\end{exemple}

% ----------------------------------------------------------------------------
\section{Cercle trigonométrique et nombres trigonométriques}
% ----------------------------------------------------------------------------

\subsection{Le cercle trigonométrique et les quatre quadrants}

\begin{definition}[Cercle trigonométrique et point-image]
Dans le plan muni d'un repère orthonormé $(O; \vec{OI}, \vec{OJ})$, le \textbf{cercle trigonométrique} est le cercle centré à l'origine $O(0,0)$ et de rayon $R = 1$.
\begin{itemize}
    \item Un angle orienté $\alpha$ est \emph{rapporté au cercle trigonométrique} lorsque son sommet est $O$ et son premier côté est la demi-droite $[OI$ (axe horizontal positif, où $I(1,0)$ et $J(0,1)$).
    \item Le \textbf{point-image} $P$ de l'angle $\alpha$ est l'unique point d'intersection du second côté de l'angle avec le cercle trigonométrique.
\end{itemize}
\end{definition}

À chaque angle $\alpha$ correspond un unique point-image $P$, mais la réciproque est fausse : \textbf{tous les angles de la forme $\alpha + 2k\pi$ ($k \in \Z$) possèdent exactement le même point-image $P$} sur le cercle trigonométrique.

Les axes de coordonnées divisent le cercle trigonométrique en quatre régions appelées \textbf{quadrants}, numérotées de $I$ à $IV$ dans le sens trigonométrique :
\begin{itemize}
    \item \textbf{Quadrant I :} $\alpha \in \left]0, \frac{\pi}{2}\right[$ (modulo $2\pi$) ; abscisse $>0$, ordonnée $>0$.
    \item \textbf{Quadrant II :} $\alpha \in \left]\frac{\pi}{2}, \pi\right[$ (modulo $2\pi$) ; abscisse $<0$, ordonnée $>0$.
    \item \textbf{Quadrant III :} $\alpha \in \left]\pi, \frac{3\pi}{2}\right[$ (modulo $2\pi$) ; abscisse $<0$, ordonnée $<0$.
    \item \textbf{Quadrant IV :} $\alpha \in \left]\frac{3\pi}{2}, 2\pi\right[$ (ou $\left]-\frac{\pi}{2}, 0\right[$) ; abscisse $>0$, ordonnée $<0$.
\end{itemize}

\subsection{Sinus, cosinus, tangente, cotangente, sécante et cosécante}

\begin{definition}[Nombres trigonométriques d'un angle orienté $\alpha$]
Soit $P$ le point-image de l'angle orienté $\alpha$ sur le cercle trigonométrique :
\begin{enumerate}
    \item Le \textbf{cosinus} de $\alpha$, noté $\cos\alpha$, est l'\textbf{abscisse} du point $P$.
    \item Le \textbf{sinus} de $\alpha$, noté $\sin\alpha$, est l'\textbf{ordonnée} du point $P$. Ainsi, $P = (\cos\alpha, \sin\alpha)$.
    \item Soit $t$ la droite tangente au cercle au point $I(1,0)$ (d'équation $x=1$). Lorsque $\alpha \neq \frac{\pi}{2} + k\pi$ ($k\in\Z$), la droite $(OP)$ coupe $t$ en un point $T$ d'abscisse $1$ : la \textbf{tangente} de $\alpha$, notée $\tan\alpha$ (ou $\tg\alpha$), est l'\textbf{ordonnée} de $T$.
    \item Soit $t'$ la droite tangente au cercle au point $J(0,1)$ (d'équation $y=1$). Lorsque $\alpha \neq k\pi$ ($k\in\Z$), la droite $(OP)$ coupe $t'$ en un point $T'$ d'ordonnée $1$ : la \textbf{cotangente} de $\alpha$, notée $\cot\alpha$ (ou $\cotg\alpha$), est l'\textbf{abscisse} de $T'$.
\end{enumerate}
\end{definition}

\begin{center}
\begin{tikzpicture}[scale=2.3, >=Stealth]
    % Axes
    \draw[->, thick, gray] (-1.3,0) -- (1.6,0) node[right, black] {$x$};
    \draw[->, thick, gray] (0,-1.3) -- (0,1.5) node[above, black] {$y$};
    \draw[thick, uclblue] (0,0) circle (1);

    % Droites tangentes t et t'
    \draw[thick, defgreen] (1,-1.25) -- (1,1.4) node[above] {$t$};
    \draw[thick, purple] (-1.2,1) -- (1.55,1) node[right] {$t'$};

    \def\ang{38}
    \coordinate (O) at (0,0);
    \coordinate (P) at (\ang:1);
    \coordinate (Px) at ({cos(\ang)},0);
    \coordinate (Py) at (0,{sin(\ang)});
    \coordinate (T) at (1,{tan(\ang)});
    \coordinate (Tp) at ({1/tan(\ang)},1);

    % Droite OP prolongée jusqu'à T'
    \draw[thick, uclaccent] (O) -- (Tp);

    % Angle
    \fill[uclaccent!15] (0,0) -- (0.3,0) arc (0:\ang:0.3) -- cycle;
    \draw[uclaccent, thick] (0.3,0) arc (0:\ang:0.3);
    \node[uclaccent] at ({\ang/2}:0.42) {$\alpha$};

    % Projections
    \draw[dashed, gray] (P) -- (Px);
    \draw[dashed, gray] (P) -- (Py);

    % Segments colorés
    \draw[very thick, exorange] (O) -- (Px) node[midway, below] {$\cos\alpha$};
    \draw[very thick, warnred] (O) -- (Py) node[midway, left] {$\sin\alpha$};
    \draw[very thick, defgreen] (1,0) -- (T) node[midway, right] {$\tan\alpha$};
    \draw[very thick, purple] (0,1) -- (Tp) node[midway, above] {$\cot\alpha$};

    % Points
    \fill[uclblue] (P) circle (1.2pt) node[above left=-1pt] {$P$};
    \fill[defgreen] (T) circle (1.2pt) node[right] {$T(1,\tan\alpha)$};
    \fill[purple] (Tp) circle (1.2pt) node[above right] {$T'(\cot\alpha,1)$};
    \fill (1,0) circle (1pt) node[below right] {$I(1,0)$};
    \fill (0,1) circle (1pt) node[above left] {$J(0,1)$};
    \fill (0,0) circle (1pt) node[below left] {$O$};
\end{tikzpicture}
\end{center}

\begin{attention}[Abscisse et ordonnée vs longueur]
Ne dites jamais que $\cos\alpha$, $\sin\alpha$, $\tan\alpha$ ou $\cot\alpha$ sont les « longueurs » des segments correspondants ! Une longueur géométrique est toujours positive ou nulle, tandis qu'un nombre trigonométrique est une \textbf{coordonnée algébrique} (abscisse ou ordonnée) qui peut être \textbf{négative} selon le quadrant où se situe le point-image $P$.
\end{attention}

\subsection{Propriétés, signes et relations fondamentales}

Par théorème de Thalès dans les triangles semblables $OP'P$, $OIT$ et $OJT'$, et par le théorème de Pythagore appliqué au triangle rectangle $OP'P$ d'hypoténuse $\overline{OP} = 1$, on établit immédiatement toutes les relations fondamentales :

\begin{theoreme}[Relations fondamentales et conditions d'existence]
Pour tout angle orienté $\alpha \in \R$ :
\begin{enumerate}
    \item \textbf{Bornes du sinus et du cosinus :}
    \[
        \boxed{-1 \le \cos\alpha \le 1 \qquad \text{et} \qquad -1 \le \sin\alpha \le 1}
    \]
    \item \textbf{Relation fondamentale de la trigonométrie (Pythagore sur le cercle unité) :}
    \[
        \boxed{\cos^2\alpha + \sin^2\alpha = 1}
    \]
    \item \textbf{Expressions de la tangente et de la cotangente :}
    \begin{align*}
        \boxed{\tan\alpha = \frac{\sin\alpha}{\cos\alpha} \quad \left(\forall \alpha \neq \frac{\pi}{2} + k\pi,\; k\in\Z\right)} \\[4pt]
        \boxed{\cot\alpha = \frac{\cos\alpha}{\sin\alpha} = \frac{1}{\tan\alpha} \quad \left(\forall \alpha \neq k\pi,\; k\in\Z\right)}
    \end{align*}
    \item \textbf{Identités dérivées (obtenues en divisant $\cos^2\alpha + \sin^2\alpha = 1$ par $\cos^2\alpha$ ou $\sin^2\alpha$) :}
    \[
        \boxed{1 + \tan^2\alpha = \frac{1}{\cos^2\alpha} \quad \left(\alpha \neq \frac{\pi}{2}+k\pi\right)}
        \quad \text{et} \quad
        \boxed{1 + \cot^2\alpha = \frac{1}{\sin^2\alpha} \quad \left(\alpha \neq k\pi\right)}
    \]
\end{enumerate}
\end{theoreme}

À titre de complément (ces notations apparaissent parfois dans certains énoncés), on définit aussi la \textbf{sécante} et la \textbf{cosécante} comme les inverses respectifs du cosinus et du sinus :
\[
    \sec\alpha = \frac{1}{\cos\alpha} \quad \left(\alpha \neq \frac{\pi}{2}+k\pi\right)
    \qquad \text{et} \qquad
    \csc\alpha \text{ (ou } \cosec\alpha\text{)} = \frac{1}{\sin\alpha} \quad \left(\alpha \neq k\pi\right).
\]

Le tableau suivant résume le signe des nombres trigonométriques dans chacun des quatre quadrants ainsi que leurs valeurs aux frontières :

\begin{center}
\renewcommand{\arraystretch}{1.35}
\small
\begin{tabular}{@{}cccccccccc@{}}
\toprule
$\alpha$ ($\degree$) & $0\degree$ & \textit{Quad. I} & $90\degree$ & \textit{Quad. II} & $180\degree$ & \textit{Quad. III} & $270\degree$ & \textit{Quad. IV} & $360\degree$ \\
$\alpha$ (rad) & $0$ & $\left]0, \frac{\pi}{2}\right[$ & $\dfrac{\pi}{2}$ & $\left]\frac{\pi}{2}, \pi\right[$ & $\pi$ & $\left]\pi, \frac{3\pi}{2}\right[$ & $\dfrac{3\pi}{2}$ & $\left]\frac{3\pi}{2}, 2\pi\right[$ & $2\pi$ \\
\midrule
$\sin\alpha$ & $0$ & $+$ & $1$ & $+$ & $0$ & $-$ & $-1$ & $-$ & $0$ \\
$\cos\alpha$ & $1$ & $+$ & $0$ & $-$ & $-1$ & $-$ & $0$ & $+$ & $1$ \\
$\tan\alpha$ & $0$ & $+$ & $\nexists$ & $-$ & $0$ & $+$ & $\nexists$ & $-$ & $0$ \\
$\cot\alpha$ & $\nexists$ & $+$ & $0$ & $-$ & $\nexists$ & $+$ & $0$ & $-$ & $\nexists$ \\
\bottomrule
\end{tabular}
\end{center}

% ----------------------------------------------------------------------------
\section{Angles remarquables et angles associés}
% ----------------------------------------------------------------------------

\subsection{Valeurs des angles remarquables}

L'examen d'entrée se déroulant \textbf{sans calculatrice}, la connaissance instantanée des valeurs trigonométriques des angles remarquables du premier quadrant est indispensable.

\begin{theoreme}[Tableau des angles remarquables]
\begin{center}
\renewcommand{\arraystretch}{1.6}
\begin{tabular}{@{}cccccc@{}}
\toprule
\textbf{Angle $\alpha$ (degrés)} & $0\degree$ & $30\degree$ & $45\degree$ & $60\degree$ & $90\degree$ \\
\textbf{Angle $\alpha$ (radians)} & $0$ & $\dfrac{\pi}{6}$ & $\dfrac{\pi}{4}$ & $\dfrac{\pi}{3}$ & $\dfrac{\pi}{2}$ \\
\midrule
$\sin\alpha$ & $0 = \dfrac{\sqrt{0}}{2}$ & $\dfrac{1}{2} = \dfrac{\sqrt{1}}{2}$ & $\dfrac{\sqrt{2}}{2}$ & $\dfrac{\sqrt{3}}{2}$ & $1 = \dfrac{\sqrt{4}}{2}$ \\
$\cos\alpha$ & $1 = \dfrac{\sqrt{4}}{2}$ & $\dfrac{\sqrt{3}}{2}$ & $\dfrac{\sqrt{2}}{2}$ & $\dfrac{1}{2} = \dfrac{\sqrt{1}}{2}$ & $0 = \dfrac{\sqrt{0}}{2}$ \\
$\tan\alpha = \dfrac{\sin\alpha}{\cos\alpha}$ & $0$ & $\dfrac{\sqrt{3}}{3} = \dfrac{1}{\sqrt{3}}$ & $1$ & $\sqrt{3}$ & $\nexists$ \\
$\cot\alpha = \dfrac{\cos\alpha}{\sin\alpha}$ & $\nexists$ & $\sqrt{3}$ & $1$ & $\dfrac{\sqrt{3}}{3} = \dfrac{1}{\sqrt{3}}$ & $0$ \\
\bottomrule
\end{tabular}
\end{center}
\end{theoreme}

\begin{astuce}[Comment ne jamais hésiter sur le tableau des valeurs remarquables]
\begin{itemize}
    \item Pour la ligne $\sin\alpha$ de $0$ à $\frac{\pi}{2}$, il suffit de compter de $0$ à $4$ sous la racine divisée par $2$ : $\frac{\sqrt{0}}{2}, \frac{\sqrt{1}}{2}, \frac{\sqrt{2}}{2}, \frac{\sqrt{3}}{2}, \frac{\sqrt{4}}{2}$.
    \item Pour $\cos\alpha$, on lit la même suite en sens inverse (de $\frac{\sqrt{4}}{2}$ à $\frac{\sqrt{0}}{2}$).
    \item Pour $\tan\alpha$, inutile d'apprendre par cœur : on divise mentalement $\sin\alpha$ par $\cos\alpha$ (les dénominateurs $2$ se simplifient !). Par exemple, $\tan\frac{\pi}{6} = \frac{1/2}{\sqrt{3}/2} = \frac{1}{\sqrt{3}} = \frac{\sqrt{3}}{3}$.
    \item Géométriquement : les valeurs de $\frac{\pi}{4}$ ($45\degree$) proviennent d'un demi-carré de côté $1$ (diagonale $\sqrt{2}$), et celles de $\frac{\pi}{6}$ ($30\degree$) et $\frac{\pi}{3}$ ($60\degree$) proviennent d'un demi-triangle équilatéral de côté $1$ (hauteur $\frac{\sqrt{3}}{2}$).
\end{itemize}
\end{astuce}

\subsection{Angles associés sur le cercle trigonométrique}

Les \textbf{angles associés} à un angle $\alpha$ permettent de ramener le calcul des nombres trigonométriques de n'importe quel quadrant au premier quadrant. Au lieu de les mémoriser mécaniquement, visualisez toujours la symétrie du point-image sur le cercle trigonométrique (en imaginant $\alpha$ petit dans le premier quadrant, par exemple $\alpha \approx 30\degree$) !

\begin{theoreme}[Formules des angles associés]
\begin{enumerate}
    \item \textbf{Angles opposés ($\alpha$ et $-\alpha$, symétrie par rapport à l'axe $Ox$) :}
    \[
        \cos(-\alpha) = \cos\alpha, \quad \sin(-\alpha) = -\sin\alpha, \quad \tan(-\alpha) = -\tan\alpha, \quad \cot(-\alpha) = -\cot\alpha.
    \]
    \item \textbf{Angles supplémentaires ($\alpha$ et $\pi - \alpha$, symétrie par rapport à l'axe $Oy$) :}
    \[
        \cos(\pi-\alpha) = -\cos\alpha, \quad \sin(\pi-\alpha) = \sin\alpha, \quad \tan(\pi-\alpha) = -\tan\alpha, \quad \cot(\pi-\alpha) = -\cot\alpha.
    \]
    \item \textbf{Angles antisupplémentaires ($\alpha$ et $\pi + \alpha$, symétrie centrale par rapport à $O$) :}
    \[
        \cos(\pi+\alpha) = -\cos\alpha, \quad \sin(\pi+\alpha) = -\sin\alpha, \quad \tan(\pi+\alpha) = \tan\alpha, \quad \cot(\pi+\alpha) = \cot\alpha.
    \]
    \item \textbf{Angles complémentaires ($\alpha$ et $\frac{\pi}{2} - \alpha$, symétrie par rapport à la bissectrice $y=x$) :}
    \[
        \cos\left(\frac{\pi}{2}-\alpha\right) = \sin\alpha, \quad \sin\left(\frac{\pi}{2}-\alpha\right) = \cos\alpha, \quad \tan\left(\frac{\pi}{2}-\alpha\right) = \cot\alpha, \quad \cot\left(\frac{\pi}{2}-\alpha\right) = \tan\alpha.
    \]
    \item \textbf{Angles anticomplémentaires ($\alpha$ et $\frac{\pi}{2} + \alpha$, rotation d'un quart de tour $+\frac{\pi}{2}$) :}
    \[
        \cos\left(\frac{\pi}{2}+\alpha\right) = -\sin\alpha, \quad \sin\left(\frac{\pi}{2}+\alpha\right) = \cos\alpha, \quad \tan\left(\frac{\pi}{2}+\alpha\right) = -\cot\alpha, \quad \cot\left(\frac{\pi}{2}+\alpha\right) = -\tan\alpha.
    \]
\end{enumerate}
\end{theoreme}

\begin{center}
\begin{tikzpicture}[scale=1.65, >=Stealth]
    % Figure 1 : Opposés, supplémentaires, antisupplémentaires
    \begin{scope}[shift={(0,0)}]
        \draw[->, gray] (-1.3,0) -- (1.3,0) node[right, black] {\small $\cos$};
        \draw[->, gray] (0,-1.3) -- (0,1.3) node[above, black] {\small $\sin$};
        \draw[thick, uclblue] (0,0) circle (1);
        \def\a{32}
        \coordinate (P1) at (\a:1);
        \coordinate (P2) at ({180-\a}:1);
        \coordinate (P3) at ({180+\a}:1);
        \coordinate (P4) at ({-\a}:1);

        \draw[dashed, gray] (P1) -- (P2) -- (P3) -- (P4) -- cycle;
        \draw[thick, uclaccent] (0,0) -- (P1) node[above right] {\small $\alpha$};
        \draw[thick, defgreen] (0,0) -- (P2) node[above left] {\small $\pi-\alpha$};
        \draw[thick, exorange] (0,0) -- (P3) node[below left] {\small $\pi+\alpha$};
        \draw[thick, warnred] (0,0) -- (P4) node[below right] {\small $-\alpha$};

        \fill[uclaccent] (P1) circle (1.3pt);
        \fill[defgreen] (P2) circle (1.3pt);
        \fill[exorange] (P3) circle (1.3pt);
        \fill[warnred] (P4) circle (1.3pt);
        \node[below, align=center] at (0,-1.45) {\small \textbf{(a)} Symétries par rapport\\[-2pt]\small aux axes et à $O$};
    \end{scope}

    % Figure 2 : Complémentaires et anticomplémentaires
    \begin{scope}[shift={(4.8,0)}]
        \draw[->, gray] (-1.3,0) -- (1.3,0) node[right, black] {\small $\cos$};
        \draw[->, gray] (0,-1.3) -- (0,1.3) node[above, black] {\small $\sin$};
        \draw[thick, uclblue] (0,0) circle (1);
        \draw[dashed, gray!70] (-0.9,-0.9) -- (1.05,1.05) node[above right] {\tiny $y=x$};
        \def\a{25}
        \coordinate (P1) at (\a:1);
        \coordinate (P2) at ({90-\a}:1);
        \coordinate (P3) at ({90+\a}:1);

        \draw[thick, uclaccent] (0,0) -- (P1) node[above right] {\small $\alpha$};
        \draw[thick, defgreen] (0,0) -- (P2) node[above right] {\small $\frac{\pi}{2}-\alpha$};
        \draw[thick, warnred] (0,0) -- (P3) node[above left] {\small $\frac{\pi}{2}+\alpha$};

        \draw[dashed, uclaccent] (P1) -- ({cos(\a)},0);
        \draw[dashed, defgreen] (P2) -- (0,{sin(90-\a)});

        \fill[uclaccent] (P1) circle (1.3pt);
        \fill[defgreen] (P2) circle (1.3pt);
        \fill[warnred] (P3) circle (1.3pt);
        \node[below, align=center] at (0,-1.45) {\small \textbf{(b)} Angles complémentaires\\[-2pt]\small et anticomplémentaires};
    \end{scope}
\end{tikzpicture}
\end{center}

\begin{exemple}[Calcul avec un grand angle : UCLouvain -- Septembre 2023, Question 3b]
\emph{Énoncé :} Calculez la valeur exacte de $\sin\left(\frac{28\pi}{6}\right)$.\\
\emph{Résolution :} On simplifie d'abord la fraction : $\frac{28\pi}{6} = \frac{14\pi}{3}$. On extrait ensuite le plus grand multiple pair de $\pi$ (c'est-à-dire les tours complets $2k\pi$) :
\[
    \frac{14\pi}{3} = \frac{12\pi + 2\pi}{3} = 4\pi + \frac{2\pi}{3}.
\]
Par périodicité de période $2\pi$, puis par propriété des angles supplémentaires ($\frac{2\pi}{3} = \pi - \frac{\pi}{3}$) :
\[
    \sin\left(\frac{28\pi}{6}\right) = \sin\left(\frac{2\pi}{3}\right) = \sin\left(\pi - \frac{\pi}{3}\right) = \sin\left(\frac{\pi}{3}\right) = \frac{\sqrt{3}}{2}.
\]
\end{exemple}

% ----------------------------------------------------------------------------
\section{Fonctions trigonométriques et fonctions réciproques}
% ----------------------------------------------------------------------------

\subsection{Graphes et périodicité des fonctions trigonométriques}

En associant à tout réel $x \in \R$ (interprété comme une mesure d'angle en radians) son sinus, son cosinus ou sa tangente, on définit les fonctions trigonométriques :
\begin{itemize}
    \item $f(x) = \sin x$ : définie sur $\R$, à valeurs dans $[-1, 1]$, \textbf{impaire} ($\sin(-x) = -\sin x$) et \textbf{périodique de période $2\pi$} ($\sin(x+2\pi) = \sin x$).
    \item $g(x) = \cos x$ : définie sur $\R$, à valeurs dans $[-1, 1]$, \textbf{paire} ($\cos(-x) = \cos x$) et \textbf{périodique de période $2\pi$} ($\cos(x+2\pi) = \cos x$).
    \item $h(x) = \tan x$ : définie sur $\R \setminus \left\{\frac{\pi}{2} + k\pi \mid k \in \Z\right\}$, à valeurs dans $\R$, \textbf{impaire} ($\tan(-x) = -\tan x$) et \textbf{périodique de période $\pi$} ($\tan(x+\pi) = \tan x$).
\end{itemize}

\begin{center}
\begin{tikzpicture}
\begin{axis}[
    width=14cm, height=5.5cm,
    axis lines=middle,
    xmin=-6.6, xmax=6.6,
    ymin=-1.4, ymax=1.4,
    xtick={-6.283,-4.712,-3.1415,-1.5708,0,1.5708,3.1415,4.712,6.283},
    xticklabels={$-2\pi$,$-\frac{3\pi}{2}$,$-\pi$,$-\frac{\pi}{2}$,$0$,$\frac{\pi}{2}$,$\pi$,$\frac{3\pi}{2}$,$2\pi$},
    ytick={-1,1},
    grid=both,
    grid style={gray!20},
    samples=200,
    legend style={at={(0.98,0.98)},anchor=north east,font=\small}
]
    \addplot[thick, uclblue, domain=-6.4:6.4] {sin(deg(x))};
    \addlegendentry{$\sin(x)$}
    \addplot[thick, warnred, dashed, domain=-6.4:6.4] {cos(deg(x))};
    \addlegendentry{$\cos(x)$}
\end{axis}
\end{tikzpicture}
\end{center}

\subsection{Fonctions réciproques : $\arcsin$, $\arccos$, $\arctan$}

Comme les fonctions $\sin$, $\cos$ et $\tan$ sont périodiques, elles ne sont pas injectives sur leur domaine complet. Pour définir leurs \textbf{fonctions réciproques} (ou fonctions cyclométriques), on restreint chacune d'elles à un intervalle standard où elle est strictement monotone et bijective :

\begin{definition}[Définition et domaine des fonctions $\arcsin$, $\arccos$ et $\arctan$]
\begin{enumerate}
    \item \textbf{Arcsinus :} Pour tout $x \in [-1, 1]$, $\arcsin(x)$ est \textbf{l'unique angle $\theta \in \left[-\frac{\pi}{2}, \frac{\pi}{2}\right]$} (quadrants IV et I) tel que $\sin\theta = x$ :
    \[
        y = \arcsin(x) \iff \sin(y) = x \quad \text{avec } x \in [-1, 1] \text{ et } y \in \left[-\frac{\pi}{2}, \frac{\pi}{2}\right].
    \]
    \item \textbf{Arccosinus :} Pour tout $x \in [-1, 1]$, $\arccos(x)$ est \textbf{l'unique angle $\theta \in [0, \pi]$} (quadrants I et II) tel que $\cos\theta = x$ :
    \[
        y = \arccos(x) \iff \cos(y) = x \quad \text{avec } x \in [-1, 1] \text{ et } y \in [0, \pi].
    \]
    \item \textbf{Arctangente :} Pour tout $x \in \R$, $\arctan(x)$ est \textbf{l'unique angle $\theta \in \left]-\frac{\pi}{2}, \frac{\pi}{2}\right[$} (quadrants IV et I ouverts) tel que $\tan\theta = x$ :
    \[
        y = \arctan(x) \iff \tan(y) = x \quad \text{avec } x \in \R \text{ et } y \in \left]-\frac{\pi}{2}, \frac{\pi}{2}\right[.
    \]
\end{enumerate}
\end{definition}

\begin{theoreme}[Identités classiques sur les fonctions réciproques]
\begin{itemize}
    \item \textbf{Parité :} $\arcsin(-x) = -\arcsin(x)$, $\arctan(-x) = -\arctan(x)$, mais $\arccos(-x) = \pi - \arccos(x)$.
    \item \textbf{Angles complémentaires :} Pour tout $x \in [-1, 1]$, $\arcsin(x) + \arccos(x) = \dfrac{\pi}{2}$, et pour $x > 0$, $\arctan(x) + \arctan\left(\frac{1}{x}\right) = \dfrac{\pi}{2}$.
    \item \textbf{Composition d'une fonction trigonométrique et d'une fonction réciproque :} Pour tout $x \in [-1, 1]$, en posant $\theta = \arcsin(x) \in \left[-\frac{\pi}{2}, \frac{\pi}{2}\right]$, on a $\cos\theta \ge 0$, donc $\cos\theta = +\sqrt{1-\sin^2\theta} = \sqrt{1-x^2}$. De même pour $\arccos(x) \in [0, \pi]$ où $\sin\theta \ge 0$ :
    \[
        \boxed{\cos\bigl(\arcsin(x)\bigr) = \sqrt{1 - x^2}} \qquad \text{et} \qquad \boxed{\sin\bigl(\arccos(x)\bigr) = \sqrt{1 - x^2}} \qquad (\forall x \in [-1, 1]).
    \]
\end{itemize}
\end{theoreme}

\begin{astuce}[L'astuce du triangle rectangle pour simplifier $\cos(\arcsin x)$, $\sin(\arctan x)$, etc.]
Pour simplifier une expression du type $\cos(\arcsin u)$ ou $\sin(\arctan u)$ (lorsque $u > 0$), dessinez un petit \textbf{triangle rectangle} d'angle aigu $\theta$ :
\begin{itemize}
    \item Si $\theta = \arcsin(u) = \arcsin\left(\frac{u}{1}\right)$, placez le côté opposé égal à $u$ et l'hypoténuse égale à $1$. Par Pythagore, le côté adjacent vaut $\sqrt{1-u^2}$, donc $\cos\theta = \frac{\text{adjacent}}{\text{hypoténuse}} = \sqrt{1-u^2}$.
    \item Si $\theta = \arctan(u) = \arctan\left(\frac{u}{1}\right)$, placez le côté opposé égal à $u$ et l'adjacent égal à $1$. L'hypoténuse vaut $\sqrt{1+u^2}$, d'où $\cos(\arctan u) = \frac{1}{\sqrt{1+u^2}}$ et $\sin(\arctan u) = \frac{u}{\sqrt{1+u^2}}$.
\end{itemize}
\end{astuce}

\begin{exemple}[Questions d'examen UCLouvain sur les fonctions réciproques]
\textbf{1) UCLouvain -- Septembre 2025, Question 1b (et Septembre 2023, Question 3c) :}\\
\emph{Énoncé :} Simplifier $\cos(\arcsin(x))$ (Septembre 2025) et $\cos(\arcsin(x^{2023}))$ pour $|x| \le 1$ (Septembre 2023) afin qu'il n'apparaisse plus de fonction trigonométrique.\\
\emph{Résolution :} Soit $\theta = \arcsin(u)$ avec $u \in [-1, 1]$. Par définition, $\sin\theta = u$ et $\theta \in [-\frac{\pi}{2}, \frac{\pi}{2}]$. Sur cet intervalle, le cosinus est toujours positif ou nul ($\cos\theta \ge 0$). La relation fondamentale $\cos^2\theta + \sin^2\theta = 1$ donne donc :
\[
    \cos(\arcsin(x)) = \sqrt{1 - x^2} \quad \text{et, pour } u = x^{2023} : \quad \cos\bigl(\arcsin(x^{2023})\bigr) = \sqrt{1 - (x^{2023})^2} = \sqrt{1 - x^{4046}}.
\]

\medskip
\textbf{2) UCLouvain -- Mini-Blocus 2026, Question 1d :}\\
\emph{Énoncé :} Simplifier $\arcsin(\cos(2x))$ (pour $2x \in [0, \pi]$).\\
\emph{Résolution :} Par la relation des angles complémentaires, $\cos(2x) = \sin\left(\frac{\pi}{2} - 2x\right)$. Lorsque $2x \in [0, \pi]$, l'angle $\frac{\pi}{2} - 2x$ appartient bien à l'intervalle image $\left[-\frac{\pi}{2}, \frac{\pi}{2}\right]$ de la fonction $\arcsin$. Dès lors :
\[
    \arcsin\bigl(\cos(2x)\bigr) = \arcsin\left(\sin\left(\frac{\pi}{2} - 2x\right)\right) = \frac{\pi}{2} - 2x.
\]
\emph{Remarque :} Si $2x$ n'appartient pas à $[0, \pi]$, il faut d'abord ramener l'angle dans $\left[-\frac{\pi}{2}, \frac{\pi}{2}\right]$ par périodicité et supplémentarité car $\arcsin(\sin\alpha)$ vaut toujours l'unique représentant dans $\left[-\frac{\pi}{2}, \frac{\pi}{2}\right]$ ayant le même sinus que $\alpha$.
\end{exemple}

% ----------------------------------------------------------------------------
\section{Exercices du Chapitre 1}
% ----------------------------------------------------------------------------

\emph{Les corrigés complets de tous ces exercices se trouvent au Chapitre 9, en fin d'ouvrage.}

\begin{enumerate}[label=\textbf{Exercice 1.\arabic* --}]
    \item \textbf{(Inspiré d'ULiège -- La rosace hexagonale)} On considère une rosace tracée à l'intérieur d'un cercle de rayon $R = 3\text{ cm}$, dont les six pétales relient les sommets $A, B, C, D, E, F$ d'un hexagone régulier inscrit dans ce cercle (chaque arc de la rosace est centré sur un des sommets de l'hexagone et passe par le centre $O$ du cercle et les deux sommets voisins).
    \begin{enumerate}
        \item Calculer la longueur du plus petit arc $\overset{\frown}{AB}$ sur le cercle extérieur.
        \item Calculer la longueur totale des arcs formant les six pétales de la rosace.
        \item Calculer l'aire totale des six pétales de la rosace.
    \end{enumerate}

    \item \textbf{(Inspiré d'ULiège -- Mesure du rayon de la Terre par Ératosthène)} Vers 240 av. J.-C., Ératosthène savait qu'au solstice d'été à midi, le Soleil éclairait le fond des puits à Syène (à la verticale), tandis qu'à Alexandrie, située sur le même méridien à $800\text{ km}$ au nord de Syène, un obélisque vertical de $1\text{ m}$ de haut projetait une ombre de $12,6\text{ cm}$. En approximant l'angle en radians par sa tangente pour un petit angle, estimer le rayon terrestre $R$.

    \item \textbf{(UCLouvain -- Mini-Blocus 2026, Q1b)} Dans un cercle de rayon $5\text{ cm}$, une corde est située à une distance de $3\text{ cm}$ du centre. Calculez exactement la longueur $L$ de cette corde.

    \item \textbf{(Inspiré d'ULiège -- Simplifications et identités fondamentales)} Démontrer les identités suivantes en précisant leurs conditions d'existence :
    \begin{enumerate}
        \item $\tan x + \cot x = \dfrac{1}{\sin x \cos x}$
        \item $\dfrac{\tan x}{1 + \tan^2 x} = \sin x \cos x$
        \item $\sin^4 x - \cos^4 x = 1 - 2\cos^2 x = 2\sin^2 x - 1$
        \item $\sin^6 x + \cos^6 x = 1 - 3\sin^2 x \cos^2 x$
    \end{enumerate}

    \item \textbf{(UCLouvain -- Juillet 2022, Q3a)} Soit $ABC$ un triangle rectangle. On note $\alpha, \beta$ et $\gamma$ les angles aux sommets $A, B$ et $C$. Calculez la valeur exacte de :
    \[
        \sin^2(\alpha) + \sin^2(\beta) + \sin^2(\gamma).
    \]

    \item \textbf{(UCLouvain -- Juillet 2022, Q3d)} Si l'on suppose que $\sin(x) + \cos(x) = \dfrac{\sqrt{2}}{3}$ et que $\dfrac{\pi}{2} < x < \pi$, que vaut exactement $\sin(x) - \cos(x)$ ?

    \item \textbf{(Angles associés et réduction)} Simplifier au maximum l'expression suivante pour tout $x \in \R$ :
    \[
        E(x) = \cos(\pi - x) + \sin\left(\frac{\pi}{2} + x\right) + \cos(3\pi + x) - \sin\left(\frac{3\pi}{2} - x\right).
    \]

    \item \textbf{(UCLouvain -- Juillet 2026, Q3d)} Si $\sin(\theta) = x$ avec $-\dfrac{\pi}{2} \le \theta \le \dfrac{\pi}{2}$, exprimez $\sin(2\theta) = 2\sin\theta\cos\theta$ en fonction de $x$ uniquement (sans fonction trigonométrique).
\end{enumerate}

% ============================================================================
% CHAPITRE 2 : FORMULES TRIGONOMÉTRIQUES ET IDENTITÉS
% ============================================================================
\chapter{Formules Trigonométriques et Identités}

La maîtrise des formules d'addition, de duplication, de Carnot et de Simpson est le cœur technique de l'épreuve de trigonométrie. Ce chapitre présente non seulement ces formules, mais aussi les moyens mnémotechniques infaillibles pour les retrouver en quelques secondes et les appliquer aux calculs sans calculatrice et aux démonstrations d'identités.

% ----------------------------------------------------------------------------
\section{Formules d'addition}
% ----------------------------------------------------------------------------

\subsection{Sinus, cosinus et tangente d'une somme ou d'une différence}

\begin{theoreme}[Formules d'addition]
Pour tous angles réels $\alpha, \beta \in \R$ :
\begin{align}
    \cos(\alpha - \beta) &= \cos\alpha\cos\beta + \sin\alpha\sin\beta \\
    \cos(\alpha + \beta) &= \cos\alpha\cos\beta - \sin\alpha\sin\beta \\
    \sin(\alpha - \beta) &= \sin\alpha\cos\beta - \cos\alpha\sin\beta \\
    \sin(\alpha + \beta) &= \sin\alpha\cos\beta + \cos\alpha\sin\beta
\end{align}
Et pour toutes les valeurs de $\alpha$ et $\beta$ pour lesquelles les tangentes sont définies ($\alpha, \beta, \alpha\pm\beta \neq \frac{\pi}{2} + k\pi$) :
\begin{equation}
    \boxed{\tan(\alpha + \beta) = \frac{\tan\alpha + \tan\beta}{1 - \tan\alpha\tan\beta} \qquad \text{et} \qquad \tan(\alpha - \beta) = \frac{\tan\alpha - \tan\beta}{1 + \tan\alpha\tan\beta}}
\end{equation}
\end{theoreme}

\begin{astuce}[Retrouver toutes les formules d'addition grâce à la formule d'Euler]
Si vous connaissez les nombres complexes ($i^2 = -1$) et la \textbf{formule d'Euler} $e^{ix} = \cos x + i\sin x$, vous n'oublierez plus jamais une formule d'addition ou de duplication ! En effet, en utilisant la propriété des exponentielles $e^{i(\alpha+\beta)} = e^{i\alpha} \cdot e^{i\beta}$, on développe simplement le produit :
\begin{align*}
    \textcolor{uclblue}{\cos(\alpha+\beta)} + i\,\textcolor{defgreen}{\sin(\alpha+\beta)} &= e^{i(\alpha+\beta)} = e^{i\alpha} e^{i\beta} \\
    &= (\cos\alpha + i\sin\alpha)(\cos\beta + i\sin\beta) \\
    &= \textcolor{uclblue}{(\cos\alpha\cos\beta - \sin\alpha\sin\beta)} + i\,\textcolor{defgreen}{(\sin\alpha\cos\beta + \cos\alpha\sin\beta)}.
\end{align*}
En identifiant les parties réelles (en bleu) et les parties imaginaires (en vert), on retrouve instantanément $\cos(\alpha+\beta)$ et $\sin(\alpha+\beta)$ ! Les formules en $\alpha-\beta$ s'en déduisent en remplaçant $\beta$ par $-\beta$ (sachant que $\cos(-\beta)=\cos\beta$ et $\sin(-\beta)=-\sin\beta$).

Quant à $\tan(\alpha+\beta)$, elle s'obtient en divisant le numérateur et le dénominateur de $\frac{\sin(\alpha+\beta)}{\cos(\alpha+\beta)}$ par $\cos\alpha\cos\beta$.
\end{astuce}

\subsection{Application : calcul exact des nombres trigonométriques de $\frac{\pi}{12}$ ($15\degree$) et $\frac{5\pi}{12}$ ($75\degree$)}

Un grand classique des questions à boîtes fermées de l'UCLouvain consiste à calculer les valeurs exactes pour $15\degree = \frac{\pi}{12}$ ou $75\degree = \frac{5\pi}{12}$ en les décomposant à l'aide des angles remarquables $45\degree = \frac{\pi}{4}$ et $30\degree = \frac{\pi}{6}$ (ou $60\degree = \frac{\pi}{3}$) :
\[
    \frac{\pi}{12} = \frac{\pi}{4} - \frac{\pi}{6} \quad \text{(ou } \frac{\pi}{3} - \frac{\pi}{4}\text{)} \qquad \text{et} \qquad \frac{5\pi}{12} = \frac{\pi}{4} + \frac{\pi}{6}.
\]

\begin{exemple}[Calculs exacts : UCLouvain -- Juillet 2021 Q3b, Septembre 2021 Q3a, Juillet 2024 Q3d, Mini-Blocus 2026 Q1c]
\begin{enumerate}
    \item \textbf{Calcul de $\cos\left(\frac{\pi}{12}\right)$ (Septembre 2021, Q3a) :}
    \[
        \cos\left(\frac{\pi}{12}\right) = \cos\left(\frac{\pi}{4} - \frac{\pi}{6}\right) = \cos\frac{\pi}{4}\cos\frac{\pi}{6} + \sin\frac{\pi}{4}\sin\frac{\pi}{6} = \frac{\sqrt{2}}{2}\frac{\sqrt{3}}{2} + \frac{\sqrt{2}}{2}\frac{1}{2} = \boxed{\frac{\sqrt{6} + \sqrt{2}}{4}}.
    \]
    \item \textbf{Calcul de $\sin\left(\frac{\pi}{12}\right)$ (Mini-Blocus 2026, Q1c) :}
    \[
        \sin\left(\frac{\pi}{12}\right) = \sin\left(\frac{\pi}{4} - \frac{\pi}{6}\right) = \sin\frac{\pi}{4}\cos\frac{\pi}{6} - \cos\frac{\pi}{4}\sin\frac{\pi}{6} = \frac{\sqrt{2}}{2}\frac{\sqrt{3}}{2} - \frac{\sqrt{2}}{2}\frac{1}{2} = \boxed{\frac{\sqrt{6} - \sqrt{2}}{4}}.
    \]
    \item \textbf{Calcul de $\tan\left(\frac{\pi}{12}\right)$ (Juillet 2024, Q3d) et $\tan(75\degree)$ (Juillet 2021, Q3b) :}
    \begin{align*}
        \tan\left(\frac{\pi}{12}\right) &= \tan\left(\frac{\pi}{4} - \frac{\pi}{6}\right) = \frac{\tan\frac{\pi}{4} - \tan\frac{\pi}{6}}{1 + \tan\frac{\pi}{4}\tan\frac{\pi}{6}} = \frac{1 - \frac{\sqrt{3}}{3}}{1 + \frac{\sqrt{3}}{3}} \\
        &= \frac{3 - \sqrt{3}}{3 + \sqrt{3}} = \frac{(3-\sqrt{3})^2}{9 - 3} = \frac{12 - 6\sqrt{3}}{6} = \boxed{2 - \sqrt{3}}.
    \end{align*}
    De même, comme $75\degree = 90\degree - 15\degree$, ou via $\tan(45\degree+30\degree)$, on trouve immédiatement $\tan(75\degree) = \frac{3+\sqrt{3}}{3-\sqrt{3}} = \boxed{2 + \sqrt{3}}$.
\end{enumerate}
\end{exemple}

% ----------------------------------------------------------------------------
\section{Formules de duplication, de triplication et de Carnot}
% ----------------------------------------------------------------------------

\subsection{Formules de duplication (angle double $2\alpha$)}

En posant $\beta = \alpha$ dans les formules d'addition, on obtient les formules de duplication qui lient un angle $\alpha$ à son double $2\alpha$ :

\begin{theoreme}[Formules de duplication]
Pour tout $\alpha \in \R$ :
\begin{align}
    \sin(2\alpha) &= 2\sin\alpha\cos\alpha \\
    \cos(2\alpha) &= \cos^2\alpha - \sin^2\alpha = 2\cos^2\alpha - 1 = 1 - 2\sin^2\alpha \\
    \tan(2\alpha) &= \frac{2\tan\alpha}{1 - \tan^2\alpha} \qquad \left(\forall \alpha \neq \frac{\pi}{4} + k\frac{\pi}{2} \text{ et } \alpha \neq \frac{\pi}{2} + k\pi\right)
\end{align}
\end{theoreme}

\subsection{Formules de Carnot (linéarisation et angle moitié)}

En isolant $\cos^2\alpha$ et $\sin^2\alpha$ dans les deux variantes de $\cos(2\alpha)$, on obtient les \textbf{formules de Carnot}, extrêmement utiles pour abaisser le degré (linéariser un carré) ou factoriser une expression contenant $1 \pm \cos(2\alpha)$ :

\begin{theoreme}[Formules de Carnot]
Pour tout $\alpha \in \R$ :
\begin{equation}
    \boxed{1 + \cos(2\alpha) = 2\cos^2\alpha \iff \cos^2\alpha = \frac{1 + \cos(2\alpha)}{2}}
\end{equation}
\begin{equation}
    \boxed{1 - \cos(2\alpha) = 2\sin^2\alpha \iff \sin^2\alpha = \frac{1 - \cos(2\alpha)}{2}}
\end{equation}
En remplaçant $\alpha$ par $\frac{\alpha}{2}$, ces formules s'écrivent pour l'angle moitié :
\[
    1 + \cos\alpha = 2\cos^2\left(\frac{\alpha}{2}\right), \qquad 1 - \cos\alpha = 2\sin^2\left(\frac{\alpha}{2}\right), \qquad \tan\left(\frac{\alpha}{2}\right) = \frac{1 - \cos\alpha}{\sin\alpha} = \frac{\sin\alpha}{1 + \cos\alpha}.
\]
\end{theoreme}

\begin{exemple}[Applications des formules de duplication et de Carnot aux examens UCLouvain]
\textbf{1) Calcul de $\cos^4(15\degree) - \sin^4(15\degree)$ (UCLouvain, Juillet 2026, Q3c) et $\cos^4(22,5\degree) - \sin^4(22,5\degree)$ (Septembre 2026, Q3c) :}\\
\emph{Résolution :} On factorise la différence de deux carrés $a^4 - b^4 = (a^2 - b^2)(a^2 + b^2)$ :
\[
    \cos^4\alpha - \sin^4\alpha = \bigl(\cos^2\alpha - \sin^2\alpha\bigr)\underbrace{\bigl(\cos^2\alpha + \sin^2\alpha\bigr)}_{=1} = \cos(2\alpha).
\]
Dès lors :
\begin{itemize}
    \item Pour $\alpha = 15\degree$ : $\cos^4(15\degree) - \sin^4(15\degree) = \cos(30\degree) = \boxed{\dfrac{\sqrt{3}}{2}}$.
    \item Pour $\alpha = 22,5\degree$ : $\cos^4(22,5\degree) - \sin^4(22,5\degree) = \cos(45\degree) = \boxed{\dfrac{\sqrt{2}}{2}}$.
\end{itemize}

\medskip
\textbf{2) Calcul de $\tan\left(\frac{\pi}{8}\right)$ (UCLouvain, Juillet 2022 Q3e et Juillet 2023 Q3b) :}\\
\emph{Résolution :} On utilise la formule de l'angle moitié avec $\alpha = \frac{\pi}{4}$ (ou l'équation du second degré issue de $\tan\frac{\pi}{4} = \frac{2\tan(\pi/8)}{1-\tan^2(\pi/8)} = 1$) :
\[
    \tan\left(\frac{\pi}{8}\right) = \frac{1 - \cos(\pi/4)}{\sin(\pi/4)} = \frac{1 - \frac{\sqrt{2}}{2}}{\frac{\sqrt{2}}{2}} = \frac{2 - \sqrt{2}}{\sqrt{2}} = \boxed{\sqrt{2} - 1}.
\]
\end{exemple}

\subsection{Formules de triplication (angle triple $3\alpha$)}

Les expressions de $\sin(3\alpha)$ en fonction de $\sin\alpha$ et de $\cos(3\alpha)$ en fonction de $\cos\alpha$ reviennent très régulièrement dans les équations et questions courtes de l'UCLouvain (Septembre 2021 Q3c, Juillet 2023 Q1 et Q3a, Septembre 2025 Q1a et Q3, Juillet 2026 Q1, Septembre 2026 Q3d).

\begin{theoreme}[Formules de l'angle triple]
Pour tout $\alpha \in \R$ :
\begin{equation}
    \boxed{\sin(3\alpha) = 3\sin\alpha - 4\sin^3\alpha} \qquad \text{et} \qquad \boxed{\cos(3\alpha) = 4\cos^3\alpha - 3\cos\alpha}
\end{equation}
\end{theoreme}

\begin{proof}[Démonstration (deux méthodes à connaître)]
\textbf{Méthode 1 (par addition et duplication) :}
\begin{align*}
    \sin(3\alpha) &= \sin(2\alpha + \alpha) = \sin(2\alpha)\cos\alpha + \cos(2\alpha)\sin\alpha \\
    &= (2\sin\alpha\cos\alpha)\cos\alpha + (1 - 2\sin^2\alpha)\sin\alpha \\
    &= 2\sin\alpha(1 - \sin^2\alpha) + \sin\alpha - 2\sin^3\alpha = 3\sin\alpha - 4\sin^3\alpha.
\end{align*}
\textbf{Méthode 2 (par la formule d'Euler) :} En développant $(a+b)^3 = a^3 + 3a^2b + 3ab^2 + b^3$ avec $i^2 = -1$ et $i^3 = -i$ :
\begin{align*}
    \cos(3\alpha) + i\sin(3\alpha) &= (\cos\alpha + i\sin\alpha)^3 \\
    &= \bigl(\cos^3\alpha - 3\cos\alpha\sin^2\alpha\bigr) + i\bigl(3\cos^2\alpha\sin\alpha - \sin^3\alpha\bigr) \\
    &= \bigl(4\cos^3\alpha - 3\cos\alpha\bigr) + i\bigl(3\sin\alpha - 4\sin^3\alpha\bigr). \qedhere
\end{align*}
\end{proof}

\begin{exemple}[Application : UCLouvain -- Septembre 2026, Question 3d]
\emph{Énoncé :} Si $\sin(\theta) = x$ avec $-\frac{\pi}{2} \le \theta \le \frac{\pi}{2}$, exprimez $\cos(3\theta)$ en fonction de $x$ uniquement, sans fonction trigonométrique.\\
\emph{Résolution :} Comme $\theta \in [-\frac{\pi}{2}, \frac{\pi}{2}]$, on a $\cos\theta = \sqrt{1 - \sin^2\theta} = \sqrt{1-x^2}$ et $\cos^2\theta = 1-x^2$. En factorisant $\cos\theta$ dans la formule de triplication :
\[
    \cos(3\theta) = 4\cos^3\theta - 3\cos\theta = \cos\theta\bigl(4\cos^2\theta - 3\bigr) = \sqrt{1-x^2}\bigl(4(1-x^2) - 3\bigr) = \boxed{(1 - 4x^2)\sqrt{1 - x^2}}.
\]
\end{exemple}

\subsection{Expressions en fonction de $t = \tan\left(\frac{\alpha}{2}\right)$}

Il est parfois très utile de pouvoir exprimer $\cos\alpha$, $\sin\alpha$ et $\tan\alpha$ rationnellement en fonction d'une seule inconnue $t = \tan\left(\frac{\alpha}{2}\right)$ (pour $\alpha \neq (2k+1)\pi$) :

\begin{theoreme}[Formules de l'angle demi en fonction de $t = \tan(\alpha/2)$]
Pour tout $\alpha \neq (2k+1)\pi$ ($k \in \Z$), en posant $t = \tan\left(\frac{\alpha}{2}\right)$ :
\[
    \boxed{\cos\alpha = \frac{1 - t^2}{1 + t^2}, \qquad \sin\alpha = \frac{2t}{1 + t^2}, \qquad \tan\alpha = \frac{2t}{1 - t^2}}
\]
\end{theoreme}
\begin{proof}
En utilisant la duplication pour l'angle $\frac{\alpha}{2}$ et l'identité $1 + t^2 = \frac{1}{\cos^2(\alpha/2)}$ :
\[
    \sin\alpha = 2\sin\frac{\alpha}{2}\cos\frac{\alpha}{2} = 2\tan\frac{\alpha}{2}\cos^2\frac{\alpha}{2} = \frac{2t}{1+t^2}, \qquad
    \cos\alpha = 2\cos^2\frac{\alpha}{2} - 1 = \frac{2}{1+t^2} - 1 = \frac{1-t^2}{1+t^2}. \qedhere
\]
\end{proof}

% ----------------------------------------------------------------------------
\section{Formules de Simpson (Factorisation)}
% ----------------------------------------------------------------------------

Les \textbf{formules de Simpson} permettent de transformer une somme ou une différence de deux sinus ou de deux cosinus en un \textbf{produit} (ce qui est crucial pour résoudre des équations par la règle du produit nul ou pour simplifier des fractions).

\begin{theoreme}[Formules de Simpson]
Pour tous angles réels $p, q \in \R$ :
\begin{align}
    \cos p + \cos q &= 2 \cos\left(\frac{p+q}{2}\right) \cos\left(\frac{p-q}{2}\right) \\
    \cos p - \cos q &= -2 \sin\left(\frac{p+q}{2}\right) \sin\left(\frac{p-q}{2}\right) \\
    \sin p + \sin q &= 2 \sin\left(\frac{p+q}{2}\right) \cos\left(\frac{p-q}{2}\right) \\
    \sin p - \sin q &= 2 \sin\left(\frac{p-q}{2}\right) \cos\left(\frac{p+q}{2}\right)
\end{align}
Et pour la somme ou la différence de deux tangentes ($p, q \neq \frac{\pi}{2} + k\pi$) :
\begin{equation}
    \boxed{\tan p + \tan q = \frac{\sin(p+q)}{\cos p \cos q} \qquad \text{et} \qquad \tan p - \tan q = \frac{\sin(p-q)}{\cos p \cos q}}
\end{equation}
\end{theoreme}

\begin{astuce}[Astuce de mémorisation des formules de Simpson (Slides UCLouvain 2026)]
Tout repose sur deux grandeurs très simples : la \textbf{moyenne} $\text{moy} = \dfrac{p+q}{2}$ et l'\textbf{écart} (demi-différence) $\text{écart} = \dfrac{p-q}{2}$.
En effet, remarquons que $p = \text{moy} + \text{écart}$ et $q = \text{moy} - \text{écart}$. Dès lors :
\begin{center}
\renewcommand{\arraystretch}{1.35}
\small
\begin{tabular}{@{}lcc@{}}
\toprule
\textbf{Somme / Différence} & \textbf{Forme factorisée (produit)} & \textbf{Moyen mnémotechnique} \\
\midrule
$\cos p + \cos q$ & $2\cos(\text{moy})\cos(\text{écart})$ & $\cos + \cos$ est symétrique ($\cos\cdot\cos$) \\
$\cos p - \cos q$ & $-2\sin(\text{moy})\sin(\text{écart})$ & $\cos - \cos$ donne $-\sin\cdot\sin$ (signe $-$) \\
$\sin p + \sin q$ & $2\sin(\text{moy})\cos(\text{écart})$ & $\sin + \sin$ mélange $\sin(\text{moy})\cos(\text{écart})$ \\
$\sin p - \sin q$ & $2\sin(\text{écart})\cos(\text{moy})$ & remplacer $q$ par $-q$ dans $\sin p + \sin q$ \\
\bottomrule
\end{tabular}
\end{center}
Quant à $\tan p \pm \tan q$, \textbf{aucun besoin de la mémoriser} : il suffit de mettre $\frac{\sin p}{\cos p} \pm \frac{\sin q}{\cos q} = \frac{\sin p \cos q \pm \cos p \sin q}{\cos p \cos q}$ au même dénominateur et de reconnaître $\sin(p \pm q)$ au numérateur !
\end{astuce}

On peut également lire les formules de Simpson en sens inverse (formules de \textbf{Werner}, qui transforment un produit en somme) en posant $a = \frac{p+q}{2}$ et $b = \frac{p-q}{2}$ (soit $p = a+b$ et $q = a-b$) :
\begin{align*}
    2\sin a \cos b &= \sin(a+b) + \sin(a-b), \\
    2\cos a \cos b &= \cos(a+b) + \cos(a-b), \qquad 2\sin a \sin b = \cos(a-b) - \cos(a+b).
\end{align*}

\begin{exemple}[Calculs exacts par Simpson aux examens UCLouvain]
\textbf{1) UCLouvain -- Juillet 2025, Questions 3a et 3b :}\\
\begin{itemize}
    \item \emph{Calculer $2\sin(45\degree)\cos(15\degree)$ :} Par la formule produit $\to$ somme $2\sin a\cos b = \sin(a+b)+\sin(a-b)$ :
    \[
        2\sin(45\degree)\cos(15\degree) = \sin(60\degree) + \sin(30\degree) = \frac{\sqrt{3}}{2} + \frac{1}{2} = \boxed{\frac{1+\sqrt{3}}{2}}.
    \]
    \item \emph{Calculer $\dfrac{\sin(75\degree) - \sin(15\degree)}{\cos(75\degree) + \cos(15\degree)}$ :} La moyenne de $75\degree$ et $15\degree$ est $45\degree$, et leur demi-écart est $30\degree$. Par Simpson :
    \[
        \frac{\sin(75\degree) - \sin(15\degree)}{\cos(75\degree) + \cos(15\degree)} = \frac{2\sin(30\degree)\cos(45\degree)}{2\cos(45\degree)\cos(30\degree)} = \tan(30\degree) = \boxed{\frac{\sqrt{3}}{3}}.
    \]
\end{itemize}

\medskip
\textbf{2) UCLouvain -- Septembre 2025, Question 1d :}\\
\emph{Énoncé :} Calculer exactement $\cos(465\degree) + \cos(165\degree)$.\\
\emph{Résolution :} D'abord, on enlève un tour complet ($360\degree$) à $465\degree$ : $465\degree - 360\degree = 105\degree$. Ensuite, la moyenne de $165\degree$ et $105\degree$ est $\frac{165\degree+105\degree}{2} = 135\degree$ et leur demi-écart est $\frac{165\degree-105\degree}{2} = 30\degree$. Par Simpson :
\[
    \cos(165\degree) + \cos(105\degree) = 2\cos(135\degree)\cos(30\degree) = 2\left(-\frac{\sqrt{2}}{2}\right)\left(\frac{\sqrt{3}}{2}\right) = \boxed{-\frac{\sqrt{6}}{2}}.
\]

\medskip
\textbf{3) UCLouvain -- Septembre 2026, Question 3f (L'astuce du produit télescopique) :}\\
\emph{Énoncé :} Calculer exactement $P = \cos(20\degree)\cos(40\degree)\cos(80\degree)$.\\
\emph{Résolution :} On remarque que les angles doublent à chaque facteur ($20\degree \to 40\degree \to 80\degree$). L'astuce classique consiste à multiplier et diviser par $2\sin(20\degree)$ pour déclencher une réaction en chaîne de formules de duplication $\sin(2\alpha) = 2\sin\alpha\cos\alpha$ :
\begin{align*}
    P &= \frac{2\sin(20\degree)\cos(20\degree)\cos(40\degree)\cos(80\degree)}{2\sin(20\degree)} = \frac{\sin(40\degree)\cos(40\degree)\cos(80\degree)}{2\sin(20\degree)} \\
    &= \frac{\sin(80\degree)\cos(80\degree)}{4\sin(20\degree)} = \frac{\sin(160\degree)}{8\sin(20\degree)}.
\end{align*}
Or $160\degree$ et $20\degree$ sont supplémentaires ($160\degree = 180\degree - 20\degree$), donc $\sin(160\degree) = \sin(20\degree)$, ce qui donne $P = \boxed{\dfrac{1}{8}}$.
\end{exemple}

% ----------------------------------------------------------------------------
\section{Méthodologie des démonstrations d'identités}
% ----------------------------------------------------------------------------

Une \textbf{identité trigonométrique} est une égalité entre deux expressions qui est vraie pour toutes les valeurs des variables où ces expressions sont définies.

\begin{attention}[Identité vs Équation]
Ne confondez jamais \emph{« démontrer une identité $A(x) = B(x)$ »} et \emph{« résoudre une équation $A(x) = B(x)$ »}. Pour démontrer une identité, on ne part \textbf{pas} de l'égalité à prouver en faisant passer des termes d'un membre à l'autre, sauf si l'on justifie rigoureusement que chaque étape est une \textbf{équivalence} ($\iff$). La méthode la plus sûre et la plus élégante consiste à :
\begin{enumerate}
    \item Partir du membre le plus complexe (par exemple le membre de gauche, noté $\text{MG}$ ou $\text{lhs}$) et le transformer étape par étape jusqu'à obtenir le membre de droite ($\text{MD}$ ou $\text{rhs}$) ;
    \item Ou transformer indépendamment $\text{MG}$ et $\text{MD}$ jusqu'à aboutir à une même expression commune $C$.
\end{enumerate}
\end{attention}

\subsection{Identités conditionnelles (dans un triangle)}

Une \textbf{identité conditionnelle} est une égalité qui n'est vraie que sous une hypothèse liant les variables, typiquement lorsque $A, B, C$ sont les trois angles d'un triangle :
\[
    \boxed{A + B + C = \pi \quad (180\degree)}.
\]
Cette hypothèse fournit immédiatement les relations clés suivantes :
\[
    \begin{cases}
        A + B = \pi - C \implies \sin(A+B) = \sin C \quad \text{et} \quad \cos(A+B) = -\cos C, \\[4pt]
        \dfrac{A+B}{2} = \dfrac{\pi}{2} - \dfrac{C}{2} \implies \sin\left(\dfrac{A+B}{2}\right) = \cos\left(\dfrac{C}{2}\right) \quad \text{et} \quad \cos\left(\dfrac{A+B}{2}\right) = \sin\left(\dfrac{C}{2}\right).
    \end{cases}
\]

\begin{exemple}[Identités dans un triangle : UCLouvain -- Septembre 2022 Q3e et ULiège -- Septembre 2009 Q1]
\textbf{1) Calcul de $\dfrac{\sin(2A) + \sin(2B) + \sin(2C)}{\sin A \sin B \sin C}$ dans un triangle quelconque (UCLouvain, Septembre 2022, Q3e) :}\\
\emph{Résolution :} Appliquons Simpson aux deux premiers termes du numérateur et la duplication au troisième :
\[
    \sin(2A) + \sin(2B) + \sin(2C) = 2\sin(A+B)\cos(A-B) + 2\sin C \cos C.
\]
Puisque $A+B = \pi - C$, on a $\sin(A+B) = \sin C$ et $\cos C = -\cos(A+B)$. On met $2\sin C$ en évidence :
\[
    \sin(2A) + \sin(2B) + \sin(2C) = 2\sin C\bigl(\cos(A-B) - \cos(A+B)\bigr).
\]
Par Simpson (ou en développant $\cos(A\mp B)$), on sait que $\cos(A-B) - \cos(A+B) = 2\sin A \sin B$. Dès lors :
\[
    \sin(2A) + \sin(2B) + \sin(2C) = 4\sin A \sin B \sin C \implies \frac{\sin(2A) + \sin(2B) + \sin(2C)}{\sin A \sin B \sin C} = \boxed{4}.
\]

\medskip
\textbf{2) Identité des tangentes dans un triangle non rectangle :}\\
Dans tout triangle $ABC$ non rectangle, on a $A+B = \pi - C$, d'où $\tan(A+B) = \tan(\pi-C) = -\tan C$. En utilisant la formule d'addition de la tangente :
\[
    \frac{\tan A + \tan B}{1 - \tan A \tan B} = -\tan C \iff \tan A + \tan B = -\tan C + \tan A \tan B \tan C,
\]
ce qui donne la célèbre identité conditionnelle (souvent utilisée à l'examen, par ex. ULiège Juillet 2006 Q2 et Septembre 2009 Q1) :
\begin{align*}
    \boxed{\tan A + \tan B + \tan C = \tan A \tan B \tan C} \\[4pt]
    \text{et, en divisant par } \tan A \tan B \tan C : \qquad \boxed{\cot A \cot B + \cot B \cot C + \cot C \cot A = 1}.
\end{align*}
\end{exemple}

\subsection{Identités comportant des fonctions réciproques ($\arcsin, \arccos, \arctan$)}

Lorsqu'on doit calculer ou démontrer une relation du type $\alpha + \beta = \gamma$ où $\alpha, \beta$ sont des $\arcsin$, $\arccos$ ou $\arctan$, la méthode comporte \textbf{deux étapes obligatoires} :
\begin{enumerate}
    \item Calculer un nombre trigonométrique de $\alpha + \beta$ (par exemple $\tan(\alpha+\beta)$ ou $\sin(\alpha+\beta)$) à l'aide des formules d'addition.
    \item \textbf{Encadrer soigneusement l'angle $\alpha + \beta$} à partir des intervalles de définition de $\alpha$ et $\beta$ afin de déterminer à quel multiple de $\pi$ ou $2\pi$ près se situe la valeur exacte !
\end{enumerate}

\begin{exemple}[Calculs avec $\arctan$ et $\arcsin$ : UCLouvain -- Juillet 2024 Q3e et Septembre 2023 Q3e]
\textbf{1) Calculer $\arctan(2) - \arctan(-3)$ (UCLouvain, Juillet 2024, Q3e) :}\\
\emph{Résolution :} Posons $\alpha = \arctan(2)$ et $\beta = \arctan(-3) = -\arctan(3)$.
\begin{itemize}
    \item \emph{Étape 1 (Formule d'addition) :} Comme $\tan\alpha = 2$ et $\tan\beta = -3$, on a :
    \[
        \tan(\alpha - \beta) = \frac{\tan\alpha - \tan\beta}{1 + \tan\alpha\tan\beta} = \frac{2 - (-3)}{1 + 2(-3)} = \frac{5}{-5} = -1.
    \]
    Donc $\alpha - \beta = -\frac{\pi}{4} + k\pi$ pour un certain entier $k \in \Z$.
    \item \emph{Étape 2 (Encadrement de l'angle) :} Comme $2 > 1$, on a $\alpha = \arctan(2) \in \left]\frac{\pi}{4}, \frac{\pi}{2}\right[$. Comme $-3 < -1$, on a $\beta = \arctan(-3) \in \left]-\frac{\pi}{2}, -\frac{\pi}{4}\right[$, soit $-\beta \in \left]\frac{\pi}{4}, \frac{\pi}{2}\right[$. En sommant ces deux encadrements :
    \[
        \frac{\pi}{2} < \alpha - \beta < \pi.
    \]
    L'unique angle de la famille $-\frac{\pi}{4} + k\pi$ appartenant à $\left]\frac{\pi}{2}, \pi\right[$ correspond à $k=1$, soit :
    \[
        \arctan(2) - \arctan(-3) = -\frac{\pi}{4} + \pi = \boxed{\frac{3\pi}{4}}.
    \]
\end{itemize}

\medskip
\textbf{2) Calculer $12\sin\left(\arcsin\left(\frac{1}{4}\right) + \arcsin\left(\frac{1}{3}\right)\right)$ (UCLouvain, Septembre 2023, Q3e) :}\\
\emph{Résolution :} Posons $\alpha = \arcsin\left(\frac{1}{4}\right)$ et $\beta = \arcsin\left(\frac{1}{3}\right)$. On a $\sin\alpha = \frac{1}{4}$, $\cos\alpha = \sqrt{1 - \frac{1}{16}} = \frac{\sqrt{15}}{4}$, $\sin\beta = \frac{1}{3}$ et $\cos\beta = \sqrt{1 - \frac{1}{9}} = \frac{\sqrt{8}}{3} = \frac{2\sqrt{2}}{3}$. Par la formule d'addition du sinus :
\[
    12\sin(\alpha + \beta) = 12\bigl(\sin\alpha\cos\beta + \cos\alpha\sin\beta\bigr) = 12\left(\frac{1}{4}\cdot\frac{2\sqrt{2}}{3} + \frac{\sqrt{15}}{4}\cdot\frac{1}{3}\right) = \boxed{2\sqrt{2} + \sqrt{15}}.
\]
\end{exemple}

% ----------------------------------------------------------------------------
\section{Exercices du Chapitre 2}
% ----------------------------------------------------------------------------

\emph{Les corrigés complets de tous ces exercices se trouvent au Chapitre 9, en fin d'ouvrage.}

\begin{enumerate}[label=\textbf{Exercice 2.\arabic* --}]
    \item \textbf{(UCLouvain -- Juillet 2022, Q3f et Juillet 2023, Q3d)} En utilisant les formules de Simpson (ou de Werner), calculez la valeur exacte de :
    \[
        E = \sin(70\degree)\cos(50\degree) + \sin(260\degree)\cos(280\degree).
    \]

    \item \textbf{(UCLouvain -- Septembre 2022, Q3a et Q3b)} Calculez la valeur exacte des deux sommes suivantes :
    \[
        S_1 = \sin(50\degree) + \sin(10\degree) - \sin(70\degree) \qquad \text{et} \qquad S_2 = \cos\left(\frac{\pi}{9}\right) + \cos\left(\frac{5\pi}{9}\right) + \cos\left(\frac{7\pi}{9}\right).
    \]

    \item \textbf{(UCLouvain -- Septembre 2022, Q3d et Juillet 2024, Q3b)}
    \begin{enumerate}
        \item Si $\cos x = \frac{1}{3}$ et que $0 < x < \pi$, que vaut exactement $7 \times \tan(2x)$ ?
        \item Sachant que $\tan(10\degree) \simeq \frac{9}{50}$, donnez la valeur correspondante de $\sin(20\degree)$ sous forme d'une fraction irréductible.
    \end{enumerate}

    \item \textbf{(UCLouvain -- Septembre 2024, Q3a, Q3b et Q3d)}
    \begin{enumerate}
        \item Sachant que $\sin(x) = \dfrac{\sqrt{5}-1}{4}$ (qui est la valeur exacte de $\sin(18\degree) = \sin\frac{\pi}{10}$), calculez la valeur exacte de $\cos(4x) + \dfrac{1}{4}$.
        \item Sachant que $\cos(36\degree) = \dfrac{1+\sqrt{5}}{4}$, simplifiez au maximum l'expression :
        \[
            E(x) = \sin(36\degree + x) - \sin(36\degree - x) + \sin(72\degree - x) - \sin(72\degree + x).
        \]
        \item Calculez la valeur exacte de $4 - \sqrt{2}(\sqrt{3}-1)\tan\left(\dfrac{\pi}{24}\right)$.
    \end{enumerate}

    \item \textbf{(UCLouvain -- Juillet 2021, Q3d et Septembre 2023, Q3d)} Calculez la valeur exacte des produits suivants :
    \begin{align*}
        P_1 &= \sin\left(\frac{\pi}{24}\right)\sin\left(\frac{5\pi}{24}\right)\sin\left(\frac{7\pi}{24}\right)\sin\left(\frac{11\pi}{24}\right), \\
        P_2 &= \sin\left(\frac{\pi}{9}\right)\sin\left(\frac{2\pi}{9}\right)\sin\left(\frac{3\pi}{9}\right)\sin\left(\frac{4\pi}{9}\right).
    \end{align*}

    \item \textbf{(UCLouvain -- Juillet 2021, Q3e)} Sachant que $\tan(x) + \tan(y) + \tan(x)\tan(y) = 1$, que vaut l'angle $(x+y) \in [0, \pi]$ exprimé en radians ?

    \item \textbf{(ULiège -- Juillet 2002 Q1b, Juillet 2003 Q1 et Juillet 2004 Q1)} Démontrer les identités suivantes (en précisant les conditions d'existence pour les deux premières) :
    \begin{enumerate}
        \item $\dfrac{\sin(2x) + \sin(2y)}{\cos(2x) + \cos(2y)} = \tan(x+y)$
        \item $\dfrac{1 - \cos(2a) + \sin(2a)}{1 + \cos(2a) + \sin(2a)} = \tan a$
        \item $\sin^2 x + \sin^2\left(x + \dfrac{2\pi}{3}\right) + \sin^2\left(x - \dfrac{2\pi}{3}\right) = \dfrac{3}{2}$
    \end{enumerate}

    \item \textbf{(ULiège -- Septembre 2002 Q1 et Juillet 2010 Q4)}
    \begin{enumerate}
        \item Démontrer l'identité $\csc a = \cot\left(\dfrac{a}{2}\right) - \cot a$ (pour $a \neq k\pi$), et en déduire une expression simplifiée de la somme télescopique $\csc a + \csc(2a) + \csc(4a) + \csc(8a)$.
        \item Calculer sans calculatrice la valeur exacte de $E = \tan 9\degree - \tan 27\degree - \tan 63\degree + \tan 81\degree$.
    \end{enumerate}

    \item \textbf{(Identités conditionnelles dans un triangle : ULiège -- Septembre 1999 Q2, Septembre 2004 Q1 et Juillet 2008 Q1)} Soient $A, B, C$ les trois angles d'un triangle ($A+B+C = \pi$).
    \begin{enumerate}
        \item Démontrer que $\sin A + \sin B + \sin C = 4\cos\left(\dfrac{A}{2}\right)\cos\left(\dfrac{B}{2}\right)\cos\left(\dfrac{C}{2}\right)$.
        \item Démontrer que $\cos(2A) + \cos(2B) + \cos(2C) = -1 - 4\cos A \cos B \cos C$.
        \item Démontrer que si les angles du triangle vérifient $\sin A = \dfrac{\sin B + \sin C}{\cos B + \cos C}$, alors le triangle $ABC$ est rectangle en $A$.
    \end{enumerate}

    \item \textbf{(Identité et réciproque : ULiège -- Septembre 1983 Q2 / Septembre 2010 Q3)}
    \begin{enumerate}
        \item Démontrer que si $a, b, c$ sont trois réels en progression arithmétique (c'est-à-dire $a = b-r$ et $c = b+r$), alors :
        \[
            \frac{\sin a + \sin b + \sin c}{\cos a + \cos b + \cos c} = \tan b.
        \]
        Pour quelles valeurs de la raison $r$ et de $b$ cette propriété est-elle en défaut ?
        \item La réciproque est-elle toujours vraie ? Déterminer la condition exacte liant $a, b, c$.
    \end{enumerate}
\end{enumerate}

% ============================================================================
% CHAPITRE 3 : ÉQUATIONS ET INÉQUATIONS TRIGONOMÉTRIQUES
% ============================================================================
\chapter{Équations et Inéquations Trigonométriques}

La première grande question de l'examen d'entrée en trigonométrie (commune aux quatre facultés polytechniques de la Fédération Wallonie-Bruxelles) porte systématiquement sur la \textbf{résolution complète d'une équation trigonométrique dans $\R$}, souvent accompagnée de la représentation de ses solutions sur le cercle trigonométrique. Ce chapitre expose de manière structurée toutes les techniques permettant de ramener n'importe quelle équation d'examen aux trois équations élémentaires de base.

% ----------------------------------------------------------------------------
\section{Principes d'équivalence fondamentaux et équations élémentaires}
% ----------------------------------------------------------------------------

Une équation trigonométrique admet généralement une \textbf{infinité de solutions} réparties en une ou plusieurs familles périodiques (indexées par un entier $k \in \Z$). Les solutions appartenant à l'intervalle $[0, 2\pi[$ (ou $[-\pi, \pi[$ selon l'énoncé) sont appelées \textbf{solutions fondamentales}.

À partir des propriétés des angles associés vues au Chapitre 1 (deux angles ont le même cosinus ssi ils sont égaux ou opposés à $2k\pi$ près ; le même sinus ssi ils sont égaux ou supplémentaires à $2k\pi$ près ; la même tangente ssi ils sont égaux à $k\pi$ près), on établit les trois principes d'équivalence fondamentaux :

\begin{theoreme}[Les trois équations trigonométriques élémentaires ($k \in \Z$)]
\begin{enumerate}
    \item \textbf{Égalité de deux cosinus :}
    \[
        \boxed{\cos u = \cos a \iff u = a + 2k\pi \quad \text{ou} \quad u = -a + 2k\pi \qquad (k \in \Z)}
    \]
    \emph{Écriture condensée :} $u = \pm a + 2k\pi$ ($k \in \Z$).
    \begin{itemize}
        \item \emph{Cas particuliers (une seule famille) :} $\cos u = 1 \iff u = 2k\pi$ ; \quad $\cos u = -1 \iff u = \pi + 2k\pi$ ; \quad $\cos u = 0 \iff u = \dfrac{\pi}{2} + k\pi$.
    \end{itemize}

    \item \textbf{Égalité de deux sinus :}
    \[
        \boxed{\sin u = \sin a \iff u = a + 2k\pi \quad \text{ou} \quad u = \pi - a + 2k\pi \qquad (k \in \Z)}
    \]
    \begin{itemize}
        \item \emph{Cas particuliers (une seule famille) :} $\sin u = 1 \iff u = \dfrac{\pi}{2} + 2k\pi$ ; \quad $\sin u = -1 \iff u = -\dfrac{\pi}{2} + 2k\pi$ ; \quad $\sin u = 0 \iff u = k\pi$.
    \end{itemize}

    \item \textbf{Égalité de deux tangentes (ou cotangentes) :}
    \[
        \boxed{\tan u = \tan a \iff \begin{cases} u = a + k\pi \quad (k \in \Z) \\ \text{CE : } u \neq \frac{\pi}{2} + m\pi,\; a \neq \frac{\pi}{2} + m\pi \end{cases}}
    \]
    \begin{itemize}
        \item \emph{De même pour la cotangente :} $\cot u = \cot a \iff u = a + k\pi$ ($k \in \Z$), sous les CE $u, a \neq m\pi$.
    \end{itemize}
\end{enumerate}
\end{theoreme}

\begin{center}
\begin{tikzpicture}[scale=1.55, >=Stealth]
    % Cosinus
    \begin{scope}[shift={(0,0)}]
        \draw[->, gray] (-1.25,0) -- (1.3,0) node[right, black] {\small $\cos$};
        \draw[->, gray] (0,-1.25) -- (0,1.3) node[above, black] {\small $\sin$};
        \draw[thick, uclblue] (0,0) circle (1);
        \def\a{50}
        \draw[dashed, exorange, thick] ({cos(\a)},-1.05) -- ({cos(\a)},1.05);
        \draw[thick, uclaccent] (0,0) -- (\a:1) node[above right] {\small $a$};
        \draw[thick, uclaccent] (0,0) -- (-\a:1) node[below right] {\small $-a$};
        \fill[warnred] (\a:1) circle (1.5pt);
        \fill[warnred] (-\a:1) circle (1.5pt);
        \node[below] at (0,-1.35) {\small \textbf{(a)} $\cos u = \cos a$};
    \end{scope}

    % Sinus
    \begin{scope}[shift={(3.4,0)}]
        \draw[->, gray] (-1.25,0) -- (1.3,0) node[right, black] {\small $\cos$};
        \draw[->, gray] (0,-1.25) -- (0,1.3) node[above, black] {\small $\sin$};
        \draw[thick, uclblue] (0,0) circle (1);
        \def\a{35}
        \draw[dashed, warnred, thick] (-1.05,{sin(\a)}) -- (1.05,{sin(\a)});
        \draw[thick, uclaccent] (0,0) -- (\a:1) node[above right] {\small $a$};
        \draw[thick, uclaccent] (0,0) -- ({180-\a}:1) node[above left] {\small $\pi-a$};
        \fill[warnred] (\a:1) circle (1.5pt);
        \fill[warnred] ({180-\a}:1) circle (1.5pt);
        \node[below] at (0,-1.35) {\small \textbf{(b)} $\sin u = \sin a$};
    \end{scope}

    % Tangente
    \begin{scope}[shift={(6.8,0)}]
        \draw[->, gray] (-1.25,0) -- (1.3,0) node[right, black] {\small $\cos$};
        \draw[->, gray] (0,-1.25) -- (0,1.3) node[above, black] {\small $\sin$};
        \draw[thick, uclblue] (0,0) circle (1);
        \draw[thick, defgreen] (1,-1.15) -- (1,1.2) node[above] {\small $t$};
        \def\a{40}
        \draw[thick, uclaccent] ({180+\a}:1) -- (1,{tan(\a)});
        \fill[defgreen] (1,{tan(\a)}) circle (1.4pt) node[right] {\small $\tan a$};
        \fill[warnred] (\a:1) circle (1.5pt) node[above left] {\small $a$};
        \fill[warnred] ({180+\a}:1) circle (1.5pt) node[below left] {\small $a+\pi$};
        \node[below] at (0,-1.35) {\small \textbf{(c)} $\tan u = \tan a$};
    \end{scope}
\end{tikzpicture}
\end{center}

\begin{attention}[Deux réflexes vitaux à l'examen : Conditions d'Existence (CE) et non-simplification]
\begin{enumerate}
    \item \textbf{Conditions d'existence (CE) :} Dès qu'une équation contient une fraction, une tangente ($\tan u \implies \cos u \neq 0 \iff u \neq \frac{\pi}{2}+k\pi$), une cotangente ($\cot u \implies \sin u \neq 0 \iff u \neq k\pi$), une sécante ou une cosécante, \textbf{écrivez les CE en tout premier lieu} (cela vaut souvent 1 à 2 points sur le barème officiel !) et vérifiez à la fin que vos solutions respectent ces CE.
    \item \textbf{Ne divisez jamais par un facteur dépendant de $x$ sans vérifier s'il peut s'annuler :} Simplifier $2\sin x = \sin x$ par $\sin x$ fait perdre toute la famille de solutions $\sin x = 0 \iff x = k\pi$ ! \textbf{Mettez toujours le facteur commun en évidence} : $\sin x(2-1) = 0 \iff \sin x = 0$.
\end{enumerate}
\end{attention}

% ----------------------------------------------------------------------------
\section{Équations réductibles et comptage des solutions sur le cercle}
% ----------------------------------------------------------------------------

\subsection{Équations du type $\sin(nx) = c$, $\sin(\alpha x) = \sin(\beta x)$ et $\sin(\alpha x) = \cos(\beta x)$}

\begin{itemize}
    \item \textbf{Pour $\sin(nx) = \sin a$ (ou $\cos(nx) = \cos a$, $\tan(nx) = \tan a$) :} On applique d'abord le principe d'équivalence à l'argument $nx$, \textbf{puis} on divise toute l'égalité (y compris la période $2k\pi$ ou $k\pi$ !) par $n$ :
    \begin{align*}
        \sin(3x) = \frac{\sqrt{3}}{2} = \sin\frac{\pi}{3} &\iff 3x = \frac{\pi}{3} + 2k\pi \quad \text{ou} \quad 3x = \frac{2\pi}{3} + 2k\pi \\
        &\iff x = \frac{\pi}{9} + \frac{2k\pi}{3} \quad \text{ou} \quad x = \frac{2\pi}{9} + \frac{2k\pi}{3} \qquad (k\in\Z).
    \end{align*}
    Pour obtenir toutes les \textbf{solutions fondamentales dans $[0, 2\pi[$}, il suffit de donner à l'entier $k$ les $n$ valeurs consécutives $k = 0, 1, 2, \dots, n-1$ dans chaque famille !
    \item \textbf{Pour $\sin(u) = \cos(v)$ :} On transforme le cosinus en sinus (ou l'inverse) grâce aux angles complémentaires :
    \[
        \boxed{\sin(u) = \cos(v) \iff \cos\left(\frac{\pi}{2} - u\right) = \cos(v) \quad \text{ou} \quad \sin(u) = \sin\left(\frac{\pi}{2} - v\right)}.
    \]
    De même, pour $-\sin(u) = \cos(v)$, on utilise $\cos\left(\frac{\pi}{2} + u\right) = -\sin(u)$ : l'équation devient $\cos\left(\frac{\pi}{2} + u\right) = \cos(v)$.
\end{itemize}

\begin{exemple}[Équations réductibles et dénombrement : Questions d'examen UCLouvain]
\textbf{1) UCLouvain -- Septembre 2022, Question 3c :}\\
\emph{Énoncé :} Combien y a-t-il de solutions sur le cercle trigonométrique à l'équation $\tan(3x) = -\sqrt{3}$ ?\\
\emph{Résolution :} $\tan(3x) = -\sqrt{3} = \tan\left(-\frac{\pi}{3}\right) \iff 3x = -\frac{\pi}{3} + k\pi \iff x = -\frac{\pi}{9} + k\frac{\pi}{3}$ ($k \in \Z$).
Puisque le pas entre deux solutions consécutives est $\frac{\pi}{3}$, il y a exactement $\frac{2\pi}{\pi/3} = \boxed{6 \text{ solutions}}$ distinctes sur le cercle trigonométrique (pour $k = 0, 1, 2, 3, 4, 5$).

\medskip
\textbf{2) UCLouvain -- Juillet 2023, Question 3c :}\\
\emph{Énoncé :} Combien y a-t-il de solutions pour $x \in [0, 2\pi[$ à l'équation $\cos(x/3) = \sin(2x)$ ?\\
\emph{Résolution :} En écrivant $\sin(2x) = \cos\left(\frac{\pi}{2} - 2x\right)$, l'équation devient $\cos\left(\frac{x}{3}\right) = \cos\left(\frac{\pi}{2} - 2x\right)$, ce qui équivaut à :
\begin{itemize}
    \item \emph{Famille 1 :} $\dfrac{x}{3} = \dfrac{\pi}{2} - 2x + 2k\pi \iff \dfrac{7x}{3} = \dfrac{\pi}{2} + 2k\pi \iff x = \dfrac{(3 + 12k)\pi}{14}$,
    \item \emph{Famille 2 :} $\dfrac{x}{3} = -\dfrac{\pi}{2} + 2x + 2k\pi \iff -\dfrac{5x}{3} = -\dfrac{\pi}{2} + 2k\pi \iff x = \dfrac{(3 + 12k)\pi}{10}$.
\end{itemize}
Cherchons les entiers $k \in \Z$ tels que $x \in [0, 2\pi[$ :
\begin{itemize}
    \item Dans la Famille 1 ($0 \le 3 + 12k < 28$) : $k = 0, 1, 2$ donnent $x \in \left\{\frac{3\pi}{14}, \frac{15\pi}{14}, \frac{27\pi}{14}\right\}$ (3 solutions).
    \item Dans la Famille 2 ($0 \le 3 + 12k < 20$) : $k = 0, 1$ donnent $x \in \left\{\frac{3\pi}{10}, \frac{15\pi}{10} = \frac{3\pi}{2}\right\}$ (2 solutions).
\end{itemize}
Aucune solution n'est commune aux deux familles (dénominateurs irréductibles $14$ vs $10$ et $2$), soit un total de $3 + 2 = \boxed{5 \text{ solutions}}$.

\medskip
\textbf{3) UCLouvain -- Juillet 2025 Q3d, Juillet 2026 Q3b et Septembre 2026 Q3b :}\\
\begin{itemize}
    \item \emph{Solutions de $\sin(3x) = -\frac{\sqrt{2}}{2}$ dans $[0, \pi]$ (Juillet 2025, Q3d) :} $3x = -\frac{\pi}{4} + 2k\pi$ ou $3x = \frac{5\pi}{4} + 2k\pi$, soit $x = -\frac{\pi}{12} + \frac{2k\pi}{3}$ ou $x = \frac{5\pi}{12} + \frac{2k\pi}{3}$. Dans $[0, \pi] = [0, \frac{12\pi}{12}]$, la première famille ($k=1$) donne $x = \frac{7\pi}{12}$ et la seconde ($k=0$) donne $x = \frac{5\pi}{12}$. D'où $S = \boxed{\left\{\frac{5\pi}{12}, \frac{7\pi}{12}\right\}}$.
    \item \emph{Solutions de $\sin(2x) = -\frac{1}{2}$ dans $[0, \pi]$ (Juillet 2026, Q3b) :} $2x = \frac{7\pi}{6} + 2k\pi$ ou $2x = \frac{11\pi}{6} + 2k\pi \iff x = \frac{7\pi}{12} + k\pi$ ou $x = \frac{11\pi}{12} + k\pi$. Dans $[0, \pi]$, on obtient $S = \boxed{\left\{\frac{7\pi}{12}, \frac{11\pi}{12}\right\}}$.
    \item \emph{Solutions de $\sin(2x) = \cos x$ dans $[0, 2\pi[$ (Septembre 2026, Q3b) :} On développe $\sin(2x) = 2\sin x\cos x$ et on factorise : $\cos x(2\sin x - 1) = 0 \iff \cos x = 0$ ou $\sin x = \frac{1}{2}$. Dans $[0, 2\pi[$, cela donne immédiatement $S = \boxed{\left\{\frac{\pi}{6}, \frac{\pi}{2}, \frac{5\pi}{6}, \frac{3\pi}{2}\right\}}$.
\end{itemize}
\end{exemple}

% ----------------------------------------------------------------------------
\section{Les 6 grandes familles d'équations de l'examen d'entrée}
% ----------------------------------------------------------------------------

À l'examen d'admission, l'équation trigonométrique principale se ramène toujours, après quelques transformations, à l'une des six grandes structures suivantes.

\subsection{Famille 1 : Équations polynomiales en $\sin x$, $\cos x$ ou $\tan x$ (Substitution)}

Lorsqu'une équation peut s'exprimer à l'aide d'un seul nombre trigonométrique (par exemple $\cos x$, $\sin x$, $\cos^2 x$, $\cos(2x)$ ou $\tan x$), on pose une variable auxiliaire $y = \cos x$ (ou $y = \sin x$, $y = \sin^2 x$, etc.) pour obtenir une \textbf{équation algébrique polynomiale $P(y) = 0$}.
\begin{itemize}
    \item Si l'équation mélange $\sin^2 x$ et $\cos x$, on remplace $\sin^2 x$ par $1 - \cos^2 x$ pour n'avoir que du $\cos x$.
    \item Si $P(y)$ est de degré 3, cherchez toujours une \textbf{racine évidente} parmi $y \in \{1, -1, 2, -2, \frac{1}{2}\}$ pour factoriser $P(y)$ par $(y - y_0)$ (division polynomiale ou schéma de Horner).
    \item \textbf{Filtre de validité :} Si $y = \cos x$ ou $y = \sin x$, toute racine $y$ telle que $|y| > 1$ est impossible dans $\R$ et doit être rejetée ! Si $y = \sin^2 x$ ou $\cos^2 x$, on doit avoir $0 \le y \le 1$.
\end{itemize}

\begin{exemple}[Polynômes trigonométriques : UCLouvain -- Juillet 2023 Q1, Juillet 2025 Q1, Juillet 2026 Q1, Septembre 2026 Q1]
\textbf{1) UCLouvain -- Juillet 2023, Question 1 :}\\
\emph{Énoncé :} Résoudre dans $\R$ l'équation $2\cos(3x) - 14\cos(2x) + 34\cos(x) - 22 = 0$.\\
\emph{Résolution :} Exprimons tout en fonction de $\cos x$ grâce aux formules de duplication et de triplication :
\[
    \cos(2x) = 2\cos^2 x - 1 \qquad \text{et} \qquad \cos(3x) = 4\cos^3 x - 3\cos x.
\]
En substituant dans l'équation :
\begin{align*}
    2(4\cos^3 x - 3\cos x) - 14(2\cos^2 x - 1) + 34\cos x - 22 = 0 \\
    \iff 8\cos^3 x - 28\cos^2 x + 28\cos x - 8 = 0.
\end{align*}
Divisons par $4$ et posons $y = \cos x$ :
\[
    P(y) = 2y^3 - 7y^2 + 7y - 2 = 0.
\]
On remarque que la somme des coefficients est nulle ($2 - 7 + 7 - 2 = 0$), donc $y = 1$ est racine évidente ! En factorisant par $(y-1)$ :
\[
    P(y) = (y - 1)(2y^2 - 5y + 2) = (y - 1)(y - 2)(2y - 1) = 0.
\]
Les racines sont $y_1 = 1$, $y_2 = \frac{1}{2}$ et $y_3 = 2$. Comme $y = \cos x \in [-1, 1]$, la racine $y_3 = 2 > 1$ est rejetée. Il reste :
\[
    \cos x = 1 \iff x = 2k\pi \qquad \text{ou} \qquad \cos x = \frac{1}{2} \iff x = \pm\frac{\pi}{3} + 2k\pi \qquad (k \in \Z).
\]
D'où l'ensemble des solutions : $S = \boxed{\{2k\pi \mid k \in \Z\} \cup \left\{\pm\frac{\pi}{3} + 2k\pi \;\middle|\; k \in \Z\right\}}$.

\medskip
\textbf{2) UCLouvain -- Juillet 2026, Question 1 :}\\
\emph{Énoncé :} Résoudre dans $\R$ l'équation $\sin(3x) - \cos(2x) + 3\sin(x) + 1 = 0$ et représenter les solutions appartenant à $[-\pi, \pi[$ sur le cercle trigonométrique.\\
\emph{Résolution :} Exprimons tout en fonction de $\sin x$ à l'aide de $\sin(3x) = 3\sin x - 4\sin^3 x$ et de la formule de Carnot $1 - \cos(2x) = 2\sin^2 x$ :
\[
    (3\sin x - 4\sin^3 x) + 2\sin^2 x + 3\sin x = 0 \iff -4\sin^3 x + 2\sin^2 x + 6\sin x = 0.
\]
En posant $y = \sin x$ et en divisant par $-2$, on factorise immédiatement :
\[
    y(2y^2 - y - 3) = 0 \iff y(y + 1)\left(2y - 3\right) = 0.
\]
La racine $y = \frac{3}{2} > 1$ est impossible pour un sinus. Il reste :
\begin{itemize}
    \item $\sin x = 0 \iff x = k\pi$ ($k \in \Z$),
    \item $\sin x = -1 \iff x = -\dfrac{\pi}{2} + 2k\pi$ ($k \in \Z$).
\end{itemize}
Dans l'intervalle $[-\pi, \pi[$, il y a exactement \textbf{3 solutions} : $\boxed{x \in \left\{-\pi, -\frac{\pi}{2}, 0\right\}}$.

\begin{center}
\begin{tikzpicture}[scale=1.25, >=Stealth]
    \draw[->, thick, gray] (-1.3,0) -- (1.35,0) node[right, black] {\small $\cos x$};
    \draw[->, thick, gray] (0,-1.35) -- (0,1.3) node[above, black] {\small $\sin x$};
    \draw[thick, uclblue] (0,0) circle (1);
    \fill[warnred] (1,0) circle (2.2pt) node[above right] {\small $0$};
    \fill[warnred] (0,-1) circle (2.2pt) node[below right] {\small $-\frac{\pi}{2}$};
    \fill[warnred] (-1,0) circle (2.2pt) node[above left] {\small $-\pi$};
    \node[below left] at (0,0) {\small $O$};
\end{tikzpicture}
\end{center}

\medskip
\textbf{3) UCLouvain -- Juillet 2025, Question 1 (Une équation, quatre méthodes !) :}\\
\emph{Énoncé :} Résoudre dans $\R$ l'équation $\sin^6 x - \cos^6 x = 3\cos^2 x \sin^2 x$.\\
\emph{Résolution (Méthode la plus directe : substitution $y = \sin^2 x$) :} Comme seuls des exposants pairs apparaissent, posons $y = \sin^2 x$, de sorte que $\cos^2 x = 1 - y$. L'équation devient :
\begin{align*}
    y^3 - (1 - y)^3 = 3y(1 - y) &\iff y^3 - (1 - 3y + 3y^2 - y^3) = 3y - 3y^2 \\
    &\iff 2y^3 - 3y^2 + 3y - 1 = 3y - 3y^2.
\end{align*}
Les termes en $y^2$ et en $y$ s'annulent miraculeusement ! Il reste :
\[
    2y^3 - 1 = 0 \iff y = \sin^2 x = \sqrt[3]{\frac{1}{2}} \in [0, 1] \iff \sin x = \pm \sqrt{\sqrt[3]{\frac{1}{2}}} = \pm 2^{-1/6}.
\]
Puisque $\sin x$ peut prendre la valeur positive ou négative, les quatre points sur le cercle sont deux à deux opposés et supplémentaires, ce qui se résume par :
\[
    S = \boxed{\left\{ \pm\arcsin\left(\sqrt{\sqrt[3]{\frac{1}{2}}}\right) + k\pi \;\middle|\; k \in \Z \right\}}.
\]
\emph{Note :} En posant $u = \cos(2x)$ via Carnot ($\sin^2 x = \frac{1-u}{2}$, $\cos^2 x = \frac{1+u}{2}$), on trouve de façon équivalente $(u+1)^3 = 4 \iff \cos(2x) = \sqrt[3]{4} - 1$, soit $x = \pm\frac{1}{2}\arccos(\sqrt[3]{4}-1) + k\pi$.

\medskip
\textbf{4) UCLouvain -- Septembre 2026, Question 1 :}\\
\emph{Énoncé :} Résoudre dans $\R$ l'équation $\sin^2(4x) + \cos^4(2x) = 1$.\\
\emph{Résolution :} Par duplication, $\sin(4x) = 2\sin(2x)\cos(2x)$, donc $\sin^2(4x) = 4\sin^2(2x)\cos^2(2x)$. Posons $u = \cos^2(2x)$ (d'où $\sin^2(2x) = 1 - u$) :
\[
    4u(1-u) + u^2 = 1 \iff 3u^2 - 4u + 1 = 0 \iff (3u - 1)(u - 1) = 0 \iff u = 1 \text{ ou } u = \frac{1}{3}.
\]
\begin{itemize}
    \item $u = 1 \iff \cos^2(2x) = 1 \iff \sin(2x) = 0 \iff 2x = k\pi \iff x = \dfrac{k\pi}{2}$ ($k \in \Z$).
    \item $u = \frac{1}{3} \iff \cos(2x) = \pm\frac{1}{\sqrt{3}} \iff 2x = \pm\arccos\left(\frac{1}{\sqrt{3}}\right) + k\pi \iff x = \pm\frac{1}{2}\arccos\left(\frac{1}{\sqrt{3}}\right) + \frac{k\pi}{2}$ ($k \in \Z$).
\end{itemize}
\end{exemple}

\subsection{Famille 2 : Mise en évidence et identités algébriques remarquables}

Avant de développer aveuglément une équation, observez toujours si un facteur commun peut être \textbf{mis en évidence}, éventuellement après avoir utilisé une identité algébrique comme :
\[
    a^2 - b^2 = (a-b)(a+b), \qquad a^3 - b^3 = (a-b)(a^2 + ab + b^2), \qquad a^3 + b^3 = (a+b)(a^2 - ab + b^2).
\]

\begin{exemple}[Factorisation : UCLouvain -- Septembre 2021 Q1 et Septembre 2025 Q3]
\textbf{1) UCLouvain -- Septembre 2021, Question 1 :}\\
\emph{Énoncé :} Résoudre dans $\R$ l'équation $\cos^3(x) - \sin^3(x) = \dfrac{3}{5}\bigl(\cos(x) - \sin(x)\bigr)$.\\
\emph{Résolution :} Utilisons $a^3 - b^3 = (a-b)(a^2 + ab + b^2)$ avec $a = \cos x$ et $b = \sin x$, et regroupons tout dans le membre de gauche pour mettre $(\cos x - \sin x)$ en évidence :
\[
    \bigl(\cos x - \sin x\bigr)\Bigl(\underbrace{\cos^2 x + \sin^2 x}_{=1} + \cos x \sin x - \frac{3}{5}\Bigr) = 0 \iff \bigl(\cos x - \sin x\bigr)\Bigl(\cos x \sin x + \frac{2}{5}\Bigr) = 0.
\]
Par la règle du produit nul, l'équation se scinde en deux équations simples :
\begin{itemize}
    \item \emph{1\up{re} famille :} $\cos x - \sin x = 0 \iff \tan x = 1 \iff \boxed{x = \frac{\pi}{4} + k\pi}$ ($k \in \Z$).
    \item \emph{2\up{e} famille :} $\cos x \sin x = -\frac{2}{5} \iff \sin(2x) = -\frac{4}{5}$. Comme $\left|-\frac{4}{5}\right| \le 1$, on obtient :
    \[
        2x = -\arcsin\left(\frac{4}{5}\right) + 2k\pi \quad \text{ou} \quad 2x = \pi + \arcsin\left(\frac{4}{5}\right) + 2k\pi,
    \]
    soit $\boxed{x = -\frac{1}{2}\arcsin\left(\frac{4}{5}\right) + k\pi \quad \text{ou} \quad x = \frac{\pi}{2} + \frac{1}{2}\arcsin\left(\frac{4}{5}\right) + k\pi}$ ($k \in \Z$).
\end{itemize}

\medskip
\textbf{2) UCLouvain -- Septembre 2025, Question 3 :}\\
\emph{Énoncé :} Résoudre dans $\R$ l'équation $\sin(3x) - 4\sin(x) + \tan(x) = 0$.\\
\emph{Résolution :}
\begin{itemize}
    \item \textbf{Conditions d'existence (CE) :} La présence de $\tan x$ impose $\cos x \neq 0 \iff x \neq \frac{\pi}{2} + k\pi$ ($k \in \Z$).
    \item \textbf{Transformation et factorisation :} Remplaçons $\sin(3x) = 3\sin x - 4\sin^3 x$ et $\tan x = \frac{\sin x}{\cos x}$ :
    \[
        (3\sin x - 4\sin^3 x) - 4\sin x + \frac{\sin x}{\cos x} = 0 \iff \sin x\left(-1 - 4\sin^2 x + \frac{1}{\cos x}\right) = 0.
    \]
    \item \emph{1\up{er} facteur :} $\sin x = 0 \iff \boxed{x = k\pi}$ ($k \in \Z$), compatible avec les CE.
    \item \emph{2\up{e} facteur :} Sous la CE $\cos x \neq 0$, multiplions la parenthèse par $\cos x$ et remplaçons $\sin^2 x = 1 - \cos^2 x$ :
    \[
        -\cos x - 4(1 - \cos^2 x)\cos x + 1 = 0 \iff 4\cos^3 x - 5\cos x + 1 = 0.
    \]
    En posant $y = \cos x$, le polynôme $4y^3 - 5y + 1 = 0$ admet la racine évidente $y = 1$ (qui donne $\cos x = 1 \iff x = 2k\pi$, déjà incluse dans $x = k\pi$). En factorisant par $(y-1)$ :
    \[
        (y - 1)(4y^2 + 4y - 1) = 0 \implies y = \frac{-4 \pm \sqrt{16 + 16}}{8} = \frac{-1 \pm \sqrt{2}}{2}.
    \]
    La racine $\frac{-1-\sqrt{2}}{2} \approx -1,207 < -1$ est rejetée. La racine $\frac{-1+\sqrt{2}}{2} \in [-1, 1]$ donne la seconde famille :
    \[
        \boxed{x = \pm\arccos\left(\frac{\sqrt{2}-1}{2}\right) + 2k\pi \qquad (k \in \Z)}.
    \]
\end{itemize}
\end{exemple}

\subsection{Famille 3 : Équations linéaires en $\sin x$ et $\cos x$ ($a\sin x + b\cos x = c$)}

Il s'agit des équations de la forme :
\begin{equation}
    a\sin x + b\cos x = c \qquad (a, b \in \R^*)
\end{equation}
Deux méthodes équivalentes permettent de la ramener en deux lignes à une équation élémentaire :

\begin{theoreme}[Résolution de l'équation linéaire $a\sin x + b\cos x = c$]
\textbf{Condition d'existence de solutions réelles :} L'équation admet des solutions dans $\R$ si et seulement si :
\[
    \boxed{c^2 \le a^2 + b^2 \iff -\sqrt{a^2+b^2} \le c \le \sqrt{a^2+b^2}}.
\]
\begin{enumerate}
    \item \textbf{Méthode 1 (Méthode de l'angle auxiliaire $\tan\phi = \frac{b}{a}$) :}
    On divise par $a$ et on pose $\tan\phi = \frac{b}{a}$ (c'est-à-dire $\phi = \arctan(b/a) \in \left]-\frac{\pi}{2}, \frac{\pi}{2}\right[$) :
    \[
        \sin x + \frac{b}{a}\cos x = \frac{c}{a} \iff \sin x + \frac{\sin\phi}{\cos\phi}\cos x = \frac{c}{a} \iff \sin x\cos\phi + \cos x\sin\phi = \frac{c}{a}\cos\phi
    \]
    ce qui se réduit par la formule d'addition à :
    \[
        \boxed{\sin(x + \phi) = \frac{c}{a}\cos\phi}.
    \]
    \item \textbf{Méthode 2 (Méthode de l'amplitude $r = \sqrt{a^2+b^2}$) :}
    On divise toute l'équation par $r = \sqrt{a^2 + b^2}$. Comme $\left(\frac{a}{r}\right)^2 + \left(\frac{b}{r}\right)^2 = 1$, il existe un unique angle $\theta \in [0, 2\pi[$ tel que $\cos\theta = \frac{a}{r}$ et $\sin\theta = \frac{b}{r}$ (ou $\sin\theta = \frac{a}{r}$ et $\cos\theta = \frac{b}{r}$). L'équation devient :
    \[
        \frac{a}{r}\sin x + \frac{b}{r}\cos x = \frac{c}{r} \iff \boxed{\sin(x + \theta) = \frac{c}{\sqrt{a^2+b^2}} \quad \text{ou} \quad \cos(x - \theta') = \frac{c}{\sqrt{a^2+b^2}}}.
    \]
\end{enumerate}
\end{theoreme}

\begin{exemple}[Équations linéaires : UCLouvain -- Juillet 2024 Q1 et Septembre 2024 Q1]
\textbf{1) UCLouvain -- Septembre 2024, Question 1 :}\\
\emph{Énoncé :} (a) Pour quelles valeurs du réel $m$ l'équation $\sqrt{3}\cos(2x) - \sin(2x) = m$ admet-elle des solutions dans $\R$ ? (b) Résoudre cette équation dans $\R$ lorsque $m = \sqrt{3}$.\\
\emph{Résolution :}
\begin{enumerate}[label=(\alph*)]
    \item Calculons l'amplitude $r = \sqrt{(\sqrt{3})^2 + (-1)^2} = \sqrt{3 + 1} = 2$. En divisant les deux membres par $2$, on fait apparaître les valeurs remarquables de $\frac{\pi}{6}$ :
    \[
        \frac{\sqrt{3}}{2}\cos(2x) - \frac{1}{2}\sin(2x) = \frac{m}{2} \iff \cos\left(\frac{\pi}{6}\right)\cos(2x) - \sin\left(\frac{\pi}{6}\right)\sin(2x) = \frac{m}{2},
    \]
    soit $\cos\left(2x + \frac{\pi}{6}\right) = \frac{m}{2}$. Cette équation admet des solutions réelles ssi $\left|\frac{m}{2}\right| \le 1 \iff \boxed{-2 \le m \le 2}$.
    \item Pour $m = \sqrt{3}$, on obtient $\cos\left(2x + \frac{\pi}{6}\right) = \frac{\sqrt{3}}{2} = \cos\left(\frac{\pi}{6}\right)$, soit :
    \[
        2x + \frac{\pi}{6} = \pm\frac{\pi}{6} + 2k\pi \iff \boxed{x = k\pi \quad \text{ou} \quad x = -\frac{\pi}{6} + k\pi} \quad (k\in\Z).
    \]
\end{enumerate}

\medskip
\textbf{2) UCLouvain -- Juillet 2024, Question 1 :}\\
\emph{Énoncé :} Résoudre dans $\R$ l'équation $2\sin(2x) + \cos(2x) = -2$.\\
\emph{Résolution (Méthode de l'angle auxiliaire) :} C'est une équation linéaire en $\sin(2x)$ et $\cos(2x)$. Posons $\tan y = 2$ (c'est-à-dire $y = \arctan(2)$). L'équation s'écrit :
\[
    \tan(y)\sin(2x) + \cos(2x) = -\tan(y) \iff \sin(y)\sin(2x) + \cos(y)\cos(2x) = -\sin(y).
\]
Le membre de gauche est $\cos(2x - y)$ et le membre de droite est $\sin(-y) = \cos\left(\frac{\pi}{2} + y\right)$. Dès lors :
\[
    \cos(2x - y) = \cos\left(\frac{\pi}{2} + y\right) \iff 2x - y = \frac{\pi}{2} + y + 2k\pi \quad \text{ou} \quad 2x - y = -\frac{\pi}{2} - y + 2k\pi,
\]
ce qui donne :
\[
    \boxed{x = \frac{\pi}{4} + \arctan(2) + k\pi \quad \text{ou} \quad x = -\frac{\pi}{4} + k\pi \qquad (k \in \Z)}.
\]
\emph{Méthode alternative (équation homogène) :} En développant $\sin(2x) = 2\sin x\cos x$, $\cos(2x) = \cos^2 x - \sin^2 x$ et $-2 = -2(\cos^2 x + \sin^2 x)$, on obtient $\tan^2 x + 4\tan x + 3 = 0 \iff (\tan x + 1)(\tan x + 3) = 0$, soit $x = -\frac{\pi}{4} + k\pi$ ou $x = \arctan(-3) + k\pi$ (qui est identique à $\frac{\pi}{4} + \arctan(2) + k\pi$ comme démontré à l'Exemple du Chapitre 2 !).
\end{exemple}

\subsection{Famille 4 : Équations homogènes en $\sin x$ et $\cos x$}

\begin{definition}[Équation homogène et méthode de résolution]
Une équation est \textbf{homogène de degré $d$} en $\sin x$ et $\cos x$ si elle peut s'écrire sous la forme $P(\cos x, \sin x) = 0$ où chaque terme non nul est du type $a_{m,n} \sin^m x \cos^n x$ avec une \textbf{somme des exposants $m + n = d$ constante}.

\textbf{Méthode de résolution :}
\begin{enumerate}
    \item Mettre d'abord en évidence tout facteur éventuel $\sin x$ ou $\cos x$ commun à tous les termes.
    \item Vérifier que $\cos x = 0$ n'est pas solution de l'équation restante, puis \textbf{diviser tous les termes par $\cos^d x$} afin d'obtenir une équation polynomiale de degré $d$ en l'inconnue $t = \tan x$.
\end{enumerate}
\emph{Astuce d'homogénéisation :} Si une équation contient des termes de degré $2$ (ou $4, 6$) et une constante $c$ au second membre, on peut la rendre homogène en multipliant la constante $c$ par $(\cos^2 x + \sin^2 x)^k = 1$ !
\end{definition}

\begin{exemple}[Équations homogènes : UCLouvain -- Juillet 2018 et Mini-Blocus 2026 Q3]
\textbf{1) UCLouvain -- Mini-Blocus 2026, Question 3 :}\\
\emph{Énoncé :} Résoudre dans $\R$ l'équation $3\sin(2x)\cos(2x) - \cos^2(2x) = 2\sin^2(2x)$.\\
\emph{Résolution :} Tous les termes sont de degré $2$ en $\sin(2x)$ et $\cos(2x)$. Comme $\cos(2x) = 0$ impliquerait $\sin(2x) = 0$ (impossible), on divise par $\cos^2(2x)$ :
\[
    2\tan^2(2x) - 3\tan(2x) + 1 = 0 \iff (2\tan(2x) - 1)(\tan(2x) - 1) = 0.
\]
On a donc $\tan(2x) = \frac{1}{2}$ ou $\tan(2x) = 1$, d'où :
\[
    \boxed{x = \frac{1}{2}\arctan\left(\frac{1}{2}\right) + \frac{k\pi}{2} \quad \text{ou} \quad x = \frac{\pi}{8} + \frac{k\pi}{2} \qquad (k \in \Z)}.
\]

\medskip
\textbf{2) UCLouvain -- Juillet 2018, Question 1 :}\\
\emph{Énoncé :} Résoudre dans $\R$ l'équation homogène $\sin^3 x + 2\cos^3 x = 3\sin^2 x \cos x$.\\
\emph{Résolution :} L'équation est homogène de degré $3$. Comme $\cos x = 0$ donnerait $\sin^3 x = 0$ (impossible), divisons les deux membres par $\cos^3 x$ et posons $t = \tan x$ :
\[
    \tan^3 x + 2 = 3\tan^2 x \iff t^3 - 3t^2 + 2 = 0.
\]
On remarque que $t = 1$ est racine évidente ($1 - 3 + 2 = 0$). En factorisant par $(t-1)$ :
\[
    (t - 1)(t^2 - 2t - 2) = 0 \iff t = 1 \quad \text{ou} \quad t = 1 \pm \sqrt{3}.
\]
D'où les trois familles de solutions : $x = \frac{\pi}{4} + k\pi$, $x = \arctan(1+\sqrt{3}) + k\pi = \frac{5\pi}{12} + k\pi$, et $x = \arctan(1-\sqrt{3}) + k\pi = -\frac{\pi}{12} + k\pi$ ($k \in \Z$).
\end{exemple}

\subsection{Famille 5 : Équations symétriques en $\sin x$ et $\cos x$}

\begin{definition}[Équation symétrique et changement de variable $x = y + \frac{\pi}{4}$]
Une équation est \textbf{symétrique en $\sin x$ et $\cos x$} si elle reste inchangée lorsqu'on permute $\sin x$ et $\cos x$ (elle ne dépend que de la somme $S = \sin x + \cos x$ et du produit $P = \sin x \cos x$).

\textbf{Méthode 1 (Changement de variable universel $x = y + \frac{\pi}{4}$) :}
En posant $\boxed{x = y + \dfrac{\pi}{4}}$ (soit $y = x - \frac{\pi}{4}$), les formules d'addition donnent :
\[
    \boxed{\sin x = \frac{\sqrt{2}}{2}(\cos y + \sin y) \qquad \text{et} \qquad \cos x = \frac{\sqrt{2}}{2}(\cos y - \sin y)}.
\]
En particulier, la somme et le produit deviennent des expressions en $\cos y$ uniquement :
\[
    \sin x + \cos x = \sqrt{2}\cos y \qquad \text{et} \qquad \sin x \cos x = \frac{1}{2}(\cos^2 y - \sin^2 y) = \cos^2 y - \frac{1}{2}.
\]

\textbf{Méthode 2 (Variable somme $u = \sin x + \cos x$) :}
En posant $u = \sin x + \cos x = \sqrt{2}\sin\left(x + \frac{\pi}{4}\right)$ (avec $|u| \le \sqrt{2}$), on a $u^2 = \sin^2 x + \cos^2 x + 2\sin x\cos x = 1 + 2\sin x\cos x$, d'où $\sin x \cos x = \dfrac{u^2 - 1}{2}$.
\end{definition}

\begin{exemple}[Équation symétrique : UCLouvain -- Juillet 2022, Question 1]
\emph{Énoncé :} Résoudre dans $\R$ l'équation $\cos(x) + \sin(x) = \sqrt{2}\cos(x)\sin(x)$.\\
\emph{Résolution :} L'équation est symétrique en $\sin x$ et $\cos x$. Posons $x = y + \frac{\pi}{4}$. On a :
\[
    \cos x + \sin x = \sqrt{2}\cos y \qquad \text{et} \qquad \cos x \sin x = \frac{1}{2}(2\cos^2 y - 1).
\]
En substituant dans l'équation :
\[
    \sqrt{2}\cos y = \frac{\sqrt{2}}{2}(2\cos^2 y - 1) \iff 2\cos^2 y - 2\cos y - 1 = 0.
\]
Les racines de ce trinôme du second degré en $\cos y$ sont :
\[
    \cos y = \frac{2 \pm \sqrt{4 + 8}}{4} = \frac{1 \pm \sqrt{3}}{2}.
\]
Comme $\frac{1+\sqrt{3}}{2} \approx 1,366 > 1$, cette racine est rejetée. En revanche, $\frac{1-\sqrt{3}}{2} \approx -0,366 \in [-1, 1]$ est valide et donne :
\[
    y = \pm\arccos\left(\frac{1-\sqrt{3}}{2}\right) + 2k\pi \implies \boxed{x = \frac{\pi}{4} \pm \arccos\left(\frac{1-\sqrt{3}}{2}\right) + 2k\pi \qquad (k \in \Z)}.
\]
\end{exemple}

\subsection{Famille 6 : Équations résolubles par Simpson ou par mise au même dénominateur}

Lorsqu'une équation contient une somme ou une différence de sinus/cosinus d'arguments multiples ($x, 2x, 3x, 4x, 5x\dots$), il faut regrouper astucieusement deux termes dont la demi-somme (moyenne) ou la demi-différence fait apparaître l'argument d'un troisième terme, afin de pouvoir \textbf{factoriser grâce aux formules de Simpson}.

\begin{exemple}[Simpson et fractions en tangentes : UCLouvain -- Juillet 2021 Q1, Septembre 2022 Q1, Septembre 2023 Q1]
\textbf{1) UCLouvain -- Juillet 2021, Question 1 :}\\
\emph{Énoncé :} Résoudre dans $\R$ l'équation $\cos(x) - \cos(2x) - \sin(3x) = 0$.\\
\emph{Résolution :} Appliquons Simpson à la différence $\cos(x) - \cos(2x)$ (moyenne $\frac{3x}{2}$, demi-écart $-\frac{x}{2}$) et la formule de duplication à $\sin(3x) = \sin\left(2\cdot\frac{3x}{2}\right) = 2\sin\left(\frac{3x}{2}\right)\cos\left(\frac{3x}{2}\right)$ :
\begin{align*}
    -2\sin\left(\frac{3x}{2}\right)\sin\left(-\frac{x}{2}\right) - 2\sin\left(\frac{3x}{2}\right)\cos\left(\frac{3x}{2}\right) = 0 \\
    \iff 2\sin\left(\frac{3x}{2}\right)\left[\sin\left(\frac{x}{2}\right) - \cos\left(\frac{3x}{2}\right)\right] = 0.
\end{align*}
\begin{itemize}
    \item \emph{1\up{er} facteur :} $\sin\left(\frac{3x}{2}\right) = 0 \iff \frac{3x}{2} = k\pi \iff \boxed{x = \frac{2k\pi}{3}}$ ($k \in \Z$).
    \item \emph{2\up{e} facteur :} $\cos\left(\frac{3x}{2}\right) = \sin\left(\frac{x}{2}\right) = \cos\left(\frac{\pi}{2} - \frac{x}{2}\right)$, ce qui équivaut à :
    \[
        \frac{3x}{2} = \pm\left(\frac{\pi}{2} - \frac{x}{2}\right) + 2k\pi \iff \boxed{x = \frac{\pi}{4} + k\pi \quad \text{ou} \quad x = -\frac{\pi}{2} + 2k\pi} \quad (k\in\Z).
    \]
\end{itemize}

\medskip
\textbf{2) UCLouvain -- Septembre 2022, Question 1 :}\\
\emph{Énoncé :} Résoudre dans $\R$ l'équation $\dfrac{\tan(x)}{\tan(2x)} = \dfrac{\tan(2x)}{\tan(x)}$ et donner les solutions dans $[0, 2\pi[$.\\
\emph{Résolution :}
\begin{itemize}
    \item \textbf{Conditions d'existence (CE) :} Il faut que $\tan x$ et $\tan(2x)$ existent et soient non nulles (car toutes deux figurent au dénominateur) :
    \[
        \begin{cases}
            \cos x \neq 0 \implies x \neq \frac{\pi}{2} + k\pi \\
            \cos(2x) \neq 0 \implies x \neq \frac{\pi}{4} + k\frac{\pi}{2} \\
            \sin x \neq 0 \implies x \neq k\pi \\
            \sin(2x) \neq 0 \implies x \neq k\frac{\pi}{2}
        \end{cases}
        \iff \boxed{\text{CE : } x \neq k\frac{\pi}{4} \quad (k \in \Z)}.
    \]
    \item \textbf{Résolution :} Sous ces CE, l'équation équivaut à $\tan^2(2x) = \tan^2(x) \iff \tan(2x) = \pm\tan(x) = \tan(\pm x)$.
    \begin{itemize}
        \item $\tan(2x) = \tan(x) \iff 2x = x + k\pi \iff x = k\pi$, \textbf{à rejeter intégralement} car exclu par les CE !
        \item $\tan(2x) = \tan(-x) \iff 2x = -x + k\pi \iff x = \dfrac{k\pi}{3}$ ($k \in \Z$). En retirant les multiples de $\pi$ ($k$ multiple de $3$) qui violent la CE $x \neq m\frac{\pi}{4}$, il reste dans $[0, 2\pi[$ les \textbf{4 solutions} :
        \[
            \boxed{S_{[0, 2\pi[} = \left\{\frac{\pi}{3}, \frac{2\pi}{3}, \frac{4\pi}{3}, \frac{5\pi}{3}\right\}}.
        \]
    \end{itemize}
\end{itemize}

\medskip
\textbf{3) UCLouvain -- Septembre 2023, Question 1 :}\\
\emph{Énoncé :} Résoudre dans $\R$ l'équation $\tan(x) - \tan(3x) = 2$.\\
\emph{Résolution :}
\begin{itemize}
    \item \textbf{CE :} $\cos x \neq 0$ et $\cos(3x) \neq 0 \iff x \neq \frac{\pi}{2} + k\pi$ et $x \neq \frac{\pi}{6} + k\frac{\pi}{3}$ ($k \in \Z$).
    \item \textbf{Transformation :} Utilisons la formule de la différence de deux tangentes $\tan p - \tan q = \frac{\sin(p-q)}{\cos p \cos q}$ :
    \[
        \tan x - \tan(3x) = \frac{\sin(x - 3x)}{\cos x \cos(3x)} = \frac{\sin(-2x)}{\cos x \cos(3x)} = \frac{-2\sin x \cos x}{\cos x \cos(3x)}.
    \]
    Puisque $\cos x \neq 0$ par les CE, on simplifie par $\cos x$ :
    \[
        \frac{-2\sin x}{\cos(3x)} = 2 \iff \cos(3x) = -\sin(x) = \cos\left(\frac{\pi}{2} + x\right).
    \]
    On résout enfin cette équation élémentaire en cosinus :
    \[
        3x = \pm\left(\frac{\pi}{2} + x\right) + 2k\pi \iff \boxed{x = \frac{\pi}{4} + k\pi \quad \text{ou} \quad x = -\frac{\pi}{8} + \frac{k\pi}{2}} \quad (k\in\Z),
    \]
    toutes compatibles avec les CE.
\end{itemize}
\end{exemple}

% ----------------------------------------------------------------------------
\section{Inéquations trigonométriques}
% ----------------------------------------------------------------------------

La résolution d'une \textbf{inéquation trigonométrique} repose toujours sur l'étude du signe sur le \textbf{cercle trigonométrique} :
\begin{enumerate}
    \item Préciser les éventuelles conditions d'existence (CE).
    \item Ramener l'inéquation à l'étude du signe d'un facteur simple (du type $\sin(u) \ge a$, $\cos(u) \le a$, $\tan(u) > a$) ou d'un produit/quotient de facteurs simples.
    \item Résoudre l'équation associée pour trouver les bornes sur le cercle trigonométrique, puis lire l'arc (ou les arcs) solution dans le \textbf{sens trigonométrique positif} en veillant à ce que la borne inférieure de chaque intervalle soit toujours strictement inférieure à la borne supérieure !
\end{enumerate}

\begin{center}
\begin{tikzpicture}[scale=1.35, >=Stealth]
    % sin u >= 1/2
    \begin{scope}[shift={(0,0)}]
        \draw[->, gray] (-1.25,0) -- (1.3,0) node[right, black] {\small $\cos$};
        \draw[->, gray] (0,-1.25) -- (0,1.3) node[above, black] {\small $\sin$};
        \draw[thick, gray!60] (0,0) circle (1);
        \draw[dashed, warnred] (-1.1,0.5) -- (1.1,0.5);
        \draw[ultra thick, defgreen] (30:1) arc (30:150:1);
        \fill[defgreen] (30:1) circle (1.5pt) node[right] {\small $\frac{\pi}{6}$};
        \fill[defgreen] (150:1) circle (1.5pt) node[left] {\small $\frac{5\pi}{6}$};
        \node[below, align=center] at (0,-1.35) {\small $\sin x \ge \frac{1}{2}$\\[-2pt]\footnotesize $\frac{\pi}{6}+2k\pi \le x \le \frac{5\pi}{6}+2k\pi$};
    \end{scope}

    % cos u >= 1/2
    \begin{scope}[shift={(3.8,0)}]
        \draw[->, gray] (-1.25,0) -- (1.3,0) node[right, black] {\small $\cos$};
        \draw[->, gray] (0,-1.25) -- (0,1.3) node[above, black] {\small $\sin$};
        \draw[thick, gray!60] (0,0) circle (1);
        \draw[dashed, exorange] (0.5,-1.1) -- (0.5,1.1);
        \draw[ultra thick, defgreen] (-60:1) arc (-60:60:1);
        \fill[defgreen] (60:1) circle (1.5pt) node[above right] {\small $\frac{\pi}{3}$};
        \fill[defgreen] (-60:1) circle (1.5pt) node[below right] {\small $-\frac{\pi}{3}$};
        \node[below, align=center] at (0,-1.35) {\small $\cos x \ge \frac{1}{2}$\\[-2pt]\footnotesize $-\frac{\pi}{3}+2k\pi \le x \le \frac{\pi}{3}+2k\pi$};
    \end{scope}

    % sin u >= -1/2
    \begin{scope}[shift={(7.6,0)}]
        \draw[->, gray] (-1.25,0) -- (1.3,0) node[right, black] {\small $\cos$};
        \draw[->, gray] (0,-1.25) -- (0,1.3) node[above, black] {\small $\sin$};
        \draw[thick, gray!60] (0,0) circle (1);
        \draw[dashed, warnred] (-1.1,-0.5) -- (1.1,-0.5);
        \draw[ultra thick, defgreen] (-30:1) arc (-30:210:1);
        \fill[defgreen] (-30:1) circle (1.5pt) node[below right] {\small $-\frac{\pi}{6}$};
        \fill[defgreen] (210:1) circle (1.5pt) node[below left] {\small $\frac{7\pi}{6}$};
        \node[below, align=center] at (0,-1.35) {\small $\sin x \ge -\frac{1}{2}$\\[-2pt]\footnotesize $-\frac{\pi}{6}+2k\pi \le x \le \frac{7\pi}{6}+2k\pi$};
    \end{scope}
\end{tikzpicture}
\end{center}

\begin{exemple}[Inéquation trigonométrique : ULiège -- Juillet 2004, Question 2]
\emph{Énoncé :} Résoudre dans $\R$ l'inéquation $\cot x - \tan x - 2\tan(2x) - 4\tan(4x) > 8\sqrt{3}$.\\
\emph{Résolution :}
\begin{itemize}
    \item \textbf{CE :} $\sin x \neq 0$, $\cos x \neq 0$, $\cos(2x) \neq 0$, $\cos(4x) \neq 0$, ce qui se résume par $\sin(8x) \neq 0 \iff x \neq \frac{k\pi}{8}$ ($k\in\Z$).
    \item \textbf{Simplification télescopique :} Remarquons l'identité classique :
    \[
        \cot x - \tan x = \frac{1}{\tan x} - \tan x = \frac{1 - \tan^2 x}{\tan x} = 2\cdot\frac{1 - \tan^2 x}{2\tan x} = \frac{2}{\tan(2x)} = 2\cot(2x).
    \]
    En appliquant cette identité trois fois en cascade :
    \begin{align*}
        \underbrace{\cot x - \tan x}_{=2\cot(2x)} - 2\tan(2x) - 4\tan(4x) &= 2\bigl(\cot(2x) - \tan(2x)\bigr) - 4\tan(4x) \\
        &= 4\bigl(\cot(4x) - \tan(4x)\bigr) = 8\cot(8x).
    \end{align*}
    \item \textbf{Résolution de l'inéquation de base :} L'inéquation se réduit donc à $8\cot(8x) > 8\sqrt{3} \iff \cot(8x) > \sqrt{3}$. Sur le cercle trigonométrique (où la cotangente est de période $\pi$), $\cot u > \sqrt{3}$ ssi $u$ est compris strictement entre $k\pi$ et $\frac{\pi}{6} + k\pi$ ($k \in \Z$) :
    \[
        k\pi < 8x < \frac{\pi}{6} + k\pi \iff \boxed{\frac{k\pi}{8} < x < \frac{\pi}{48} + \frac{k\pi}{8} \qquad (k \in \Z)},
    \]
    qui respecte automatiquement la CE $x \neq \frac{k\pi}{8}$.
\end{itemize}
\end{exemple}

% ----------------------------------------------------------------------------
\section{Exercices du Chapitre 3}
% ----------------------------------------------------------------------------

\emph{Les corrigés complets de tous ces exercices se trouvent au Chapitre 9, en fin d'ouvrage.}

\begin{enumerate}[label=\textbf{Exercice 3.\arabic* --}, align=left, leftmargin=*]
    \item \textbf{(UCLouvain -- Mini-Blocus 2026, Q1a et Juillet 2024, Q3a)}
    \begin{enumerate}
        \item Trouvez toutes les solutions dans $[0, 2\pi[$ de l'équation $\sin x + \cos x = \sqrt{2}$.
        \item Trouvez toutes les solutions dans l'intervalle $]0, \pi[$ de l'équation $\sin(x) + \sin(2x) + \sin(3x) = 0$.
    \end{enumerate}

    \item \textbf{(UCLouvain -- Septembre 2019, Q1)} Résoudre dans $\R$ les deux équations trigonométriques suivantes (issues des deux variantes de la session de Septembre 2019) :
    \begin{enumerate}
        \item $\sin(x) - 2\sin(3x) + \sin(5x) = -\cos(2x) + \cos(4x)$
        \item $\dfrac{\cos(x) - \cos(3x)}{\sin(x)} = \cos(x)$
    \end{enumerate}

    \item \textbf{(UCLouvain -- Juillet 2019, Q1 et Juillet 2020, Q1)} Résoudre dans $\R$ les équations symétriques suivantes :
    \begin{enumerate}
        \item $4\bigl(\sin(x) + \cos(x)\bigr) - 8\sin(x)\cos(x) = 5$
        \item $3\bigl(\sin(x) + \cos(x)\bigr) - 4\bigl(\sin^3(x) + \cos^3(x)\bigr) = 0$
    \end{enumerate}

    \item \textbf{(UCLouvain -- Septembre 2020, Q1)} Résoudre dans $\R$ l'équation suivante en précisant ses conditions d'existence :
    \[
        \tan^2(x) - 3\frac{\tan(x)}{\cos(x)} - \frac{1}{\cos^2(x)} = 1.
    \]

    \item \textbf{(ULiège -- Septembre 2002 Q2, Juillet 2003 Q3 et Juillet 2008 Q2)} Résoudre dans $\R$ les équations suivantes :
    \begin{enumerate}
        \item $\dfrac{\cot x - \cos x}{\cot x + \cos x} = 2(1 - \sin x)$
        \item $\cos(2x) + \cos(6x) = 1 + \cos(8x)$
        \item $\sin x + \sin(2x) + \sin(3x) = \sqrt{2}\bigl(1 + \cos x + \cos(2x)\bigr)$
    \end{enumerate}

    \item \textbf{(ULiège -- Septembre 2005 Q2, Juillet 2006 Q4 et Septembre 2009 Q2)} Résoudre dans $\R$ les équations suivantes en précisant les conditions d'existence :
    \begin{enumerate}
        \item $\tan^3 x - 3\tan^2 x - 3\tan x + 1 = 0$
        \item $\left(\dfrac{1 + \tan(x/2)}{1 - \tan(x/2)}\right)^2 + \sin x = 0$
        \item $\tan x + \tan(2x) + \tan(3x) = 0$
    \end{enumerate}

    \item \textbf{(ULiège -- Juillet 2005 Q2, Juillet 2007 Q2 et Juillet 2010 Q2)} Résoudre dans $\R$ les équations suivantes :
    \begin{enumerate}
        \item $\sin^6 x + \cos^6 x = \dfrac{5}{8}$
        \item $\cos^6 x + \sin^6 x = \dfrac{7}{16}$
        \item $\sin^4 x + \cos^4 x = \sin x \cos x$
    \end{enumerate}

    \item \textbf{(Inéquations trigonométriques -- Inspiré d'ULiège)} Résoudre dans $\R$ les inéquations suivantes et représenter les solutions sur le cercle trigonométrique :
    \begin{enumerate}
        \item $\sin(2x) \le -\dfrac{\sqrt{3}}{2}$
        \item $\sin\left(x - \dfrac{\pi}{3}\right) > \sin x$
        \item $2\cos^2 x - \cos x - 1 \ge 0$
    \end{enumerate}
\end{enumerate}

% ============================================================================
% CHAPITRE 4 : RÉSOLUTION DE TRIANGLES ET PROBLÈMES GÉOMÉTRIQUES
% ============================================================================
\chapter[Résolution de Triangles et Problèmes]{Résolution de Triangles et Problèmes Géométriques}

La deuxième grande question ouverte de l'examen d'entrée à l'UCLouvain (ainsi que plusieurs questions courtes) consiste en un \textbf{problème géométrique et trigonométrique}. Qu'il s'agisse d'un système mécanique, d'un satellite en orbite, d'un pliage de papier, de cercles tangents ou d'un quadrilatère, la clé consiste toujours à identifier les bons triangles (rectangles ou quelconques) et à leur appliquer méthodiquement les théorèmes de ce chapitre.

Dans tout ce chapitre, pour un triangle $ABC$, nous notons $a = \overline{BC}$, $b = \overline{AC}$ et $c = \overline{AB}$ les longueurs des côtés opposés aux sommets $A, B, C$, et $\alpha = \widehat{A}$, $\beta = \widehat{B}$, $\gamma = \widehat{C}$ les amplitudes des angles correspondants. N'oubliez jamais la relation fondamentale de la géométrie euclidienne plane :
\[
    \boxed{\alpha + \beta + \gamma = \pi \text{ radians} = 180\degree}.
\]

% ----------------------------------------------------------------------------
\section{Triangles rectangles, triangles semblables et théorème de Thalès}
% ----------------------------------------------------------------------------

\subsection{Relations dans le triangle rectangle (Pythagore et SOH-CAH-TOA)}

Considérons un triangle $ABC$ rectangle en $B$ ($\beta = 90\degree = \frac{\pi}{2}$), de sorte que le côté opposé $b = \overline{AC}$ est l'\textbf{hypoténuse}, et les angles aigus $\alpha$ et $\gamma$ sont complémentaires ($\alpha + \gamma = 90\degree$).

\begin{theoreme}[Théorème de Pythagore et trigonométrie du triangle rectangle]
Dans un triangle rectangle d'hypoténuse $H$, pour tout angle aigu $\theta$ :
\begin{enumerate}
    \item \textbf{Théorème de Pythagore :} Le carré de la longueur de l'hypoténuse est égal à la somme des carrés des longueurs des deux côtés de l'angle droit (cathètes).
    \item \textbf{Nombres trigonométriques (moyen mnémotechnique \textcolor{warnred}{SOH -- CAH -- TOA}) :}
    \[
        \boxed{\textbf{\textcolor{warnred}{S}OH : } \sin\theta = \frac{\text{\textcolor{warnred}{O}pp.}}{\text{\textcolor{warnred}{H}yp.}}, \qquad \textbf{\textcolor{warnred}{C}AH : } \cos\theta = \frac{\text{\textcolor{warnred}{A}dj.}}{\text{\textcolor{warnred}{H}yp.}}, \qquad \textbf{\textcolor{warnred}{T}OA : } \tan\theta = \frac{\text{\textcolor{warnred}{O}pp.}}{\text{\textcolor{warnred}{A}dj.}}}.
    \]
\end{enumerate}
\end{theoreme}

\begin{center}
\begin{tikzpicture}[scale=1.05, >=Stealth]
    \coordinate (A) at (0,0);
    \coordinate (B) at (4.8,0);
    \coordinate (C) at (4.8,2.4);

    \draw[very thick, defgreen] (A) -- (B) node[midway, below=3pt] {Côté \textbf{A}djacent ($c = b\cos\theta$)};
    \draw[very thick, uclaccent] (B) -- (C) node[midway, right=4pt] {Côté \textbf{O}pposé ($a = b\sin\theta$)};
    \draw[very thick, warnred] (A) -- (C) node[midway, sloped, above=3pt] {\textbf{H}ypoténuse ($b$)};

    \draw[fill=defgreen!20, draw=defgreen] (4.8,0) -- (4.45,0) -- (4.45,0.35) -- (4.8,0.35) -- cycle;
    \pic [draw, thick, uclblue, "$\theta$", angle radius=0.9cm, angle eccentricity=1.25] {angle = B--A--C};

    \fill (A) circle (1.5pt) node[below left] {$A$};
    \fill (B) circle (1.5pt) node[below right] {$B$};
    \fill (C) circle (1.5pt) node[above right] {$C$};
\end{tikzpicture}
\end{center}

\subsection{Théorème de Thalès et triangles semblables}

Deux triangles sont \textbf{semblables} s'ils ont les mêmes angles deux à deux (il suffit que deux angles soient égaux, par exemple un angle droit et un angle aigu commun ou alterne-interne). Leurs côtés homologues sont alors proportionnels.

\begin{theoreme}[Théorème de Thalès]
Soit un triangle $ABC$, et soient $D \in (AB)$ et $E \in (AC)$ deux points tels que la droite $(DE)$ soit parallèle à la droite $(BC)$. Alors les triangles $ADE$ et $ABC$ sont semblables et :
\[
    \boxed{\frac{\overline{AD}}{\overline{AB}} = \frac{\overline{AE}}{\overline{AC}} = \frac{\overline{DE}}{\overline{BC}}}.
\]
\end{theoreme}

\begin{exemple}[Problèmes d'examen UCLouvain sur les triangles rectangles]
\textbf{1) Topographie et double visée (UCLouvain -- Juillet 2021, Question 2 et Septembre 2024, Question 3e) :}\\
\emph{Énoncé (Juillet 2021, Q2) :} On désire mesurer la hauteur $H$ d'un grand arbre dont le pied $P$ n'est pas accessible. Deux graphomètres $A$ et $B$, placés au niveau du sol horizontal, alignés avec le pied de l'arbre et séparés d'une distance connue $\overline{AB} = d$, visent le sommet $S$ de l'arbre sous des angles d'élévation respectifs $\widehat{A}$ et $\widehat{B}$ (avec $A$ plus proche de l'arbre que $B$, donc $\widehat{A} > \widehat{B}$). Exprimer $H$ en fonction de $d = \overline{AB}$, $\widehat{A}$ et $\widehat{B}$.\\
\emph{Résolution :} Notons $x = \overline{PA}$ la distance inconnue entre le pied de l'arbre et le point $A$. Dans les triangles rectangles $SPA$ et $SPB$ (rectangles en $P$), on a :
\[
    \tan\widehat{A} = \frac{H}{x} \implies x = \frac{H}{\tan\widehat{A}} = H\cot\widehat{A}, \quad \tan\widehat{B} = \frac{H}{x + d} \implies x + d = \frac{H}{\tan\widehat{B}} = H\cot\widehat{B}.
\]
En soustrayant la première équation de la seconde pour éliminer l'inconnue intermédiaire $x$ :
\[
    d = H\left(\frac{1}{\tan\widehat{B}} - \frac{1}{\tan\widehat{A}}\right) = H\,\frac{\tan\widehat{A} - \tan\widehat{B}}{\tan\widehat{A}\tan\widehat{B}} \implies \boxed{H = \frac{d\,\tan\widehat{A}\tan\widehat{B}}{\tan\widehat{A} - \tan\widehat{B}} = \frac{d}{\cot\widehat{B} - \cot\widehat{A}}}.
\]
\emph{Application numérique (Septembre 2024, Q3e) :} Pour un château entouré de douves avec $d = 30\text{ m}$, $\widehat{A} = 60\degree$, $\widehat{B} = 30\degree$ et $\sqrt{3} \simeq 1,7$ : $H = \frac{30}{\cot 30\degree - \cot 60\degree} = \frac{30}{\sqrt{3} - \frac{1}{\sqrt{3}}} = 15\sqrt{3} \simeq \boxed{25,5\text{ m}}$.

\medskip
\textbf{2) Carré inscrit dans un triangle rectangle et optimisation (UCLouvain -- Septembre 2021, Question 2) :}\\
\emph{Énoncé :} Soit $ABC$ un triangle rectangle en $A$, d'hypoténuse $a = \overline{BC}$ et d'angle aigu $\gamma = \widehat{ACB}$. On inscrit un carré $XYPQ$ de côté $x = \overline{XY}$ avec $X, Y \in [BC]$, $P \in [AC]$ et $Q \in [AB]$. (a) Identifier les triangles semblables à $ABC$ et démontrer que $\overline{BX}\cdot\overline{CY} = x^2$. (b) Exprimer $x$ en fonction de $a$ et $\gamma$. (c) Pour quelle valeur de $\gamma$ l'aire du carré est-elle maximale ?

\begin{center}
\begin{tikzpicture}[scale=1.1, >=Stealth]
    \def\angC{32}
    \def\hyp{6}
    \pgfmathsetmacro{\b}{\hyp*cos(\angC)}
    \pgfmathsetmacro{\c}{\hyp*sin(\angC)}
    \coordinate (C) at (0,0);
    \coordinate (B) at (-\hyp,0);
    \coordinate (A) at ({-\b*cos(\angC)},{\b*sin(\angC)});

    % Carré de côté x = a / (1 + tan + cot)
    \pgfmathsetmacro{\x}{\hyp / (1 + tan(\angC) + 1/tan(\angC))}
    \pgfmathsetmacro{\CY}{\x / tan(\angC)}
    \pgfmathsetmacro{\BX}{\x * tan(\angC)}
    \coordinate (Y) at (-\CY, 0);
    \coordinate (X) at ({-\CY-\x}, 0);
    \coordinate (P) at (-\CY, \x);
    \coordinate (Q) at ({-\CY-\x}, \x);

    \draw[thick, uclblue] (B) -- (C) -- (A) -- cycle;
    \draw[thick, defgreen, fill=deflightgreen] (X) -- (Y) -- (P) -- (Q) -- cycle;

    % Angle droit en A
    \draw ($(A)!0.25cm!(B)$) -- ($(A)!0.25cm!(B)!0.25cm!90:(B)$) -- ($(A)!0.25cm!(C)$);
    \pic [draw, uclaccent, "$\gamma$", angle radius=0.8cm, angle eccentricity=1.3] {angle = A--C--B};

    \fill (A) circle (1.5pt) node[above] {$A$};
    \fill (B) circle (1.5pt) node[left] {$B$};
    \fill (C) circle (1.5pt) node[right] {$C$};
    \fill (X) circle (1.3pt) node[below] {$X$};
    \fill (Y) circle (1.3pt) node[below] {$Y$};
    \fill (P) circle (1.3pt) node[above right] {$P$};
    \fill (Q) circle (1.3pt) node[above left] {$Q$};
    \node[defgreen!80!black] at ($(X)!0.5!(P)$) {$x$};
\end{tikzpicture}
\end{center}

\emph{Résolution :}
\begin{enumerate}[label=(\alph*)]
    \item Les triangles rectangles $AQP$ (rectangle en $A$), $XQB$ (rectangle en $X$) et $YPC$ (rectangle en $Y$) ont les mêmes angles aigus ($\gamma$ et $\frac{\pi}{2}-\gamma$) que $ABC$ : ils lui sont tous trois semblables ! La similitude entre $XQB$ et $YPC$ donne $\frac{\overline{BX}}{\overline{XQ}} = \frac{\overline{YP}}{\overline{CY}}$, et comme $\overline{XQ} = \overline{YP} = x$, il vient $\overline{BX}\cdot\overline{CY} = x^2$.
    \item Dans $YPC$ ($\widehat{YCP}=\gamma$), on a $\tan\gamma = \frac{x}{\overline{CY}} \implies \overline{CY} = \frac{x}{\tan\gamma}$. Dans $XQB$ ($\widehat{XQB}=\gamma$), on a $\tan\gamma = \frac{\overline{BX}}{x} \implies \overline{BX} = x\tan\gamma$. Or $a = \overline{BX} + \overline{XY} + \overline{CY} = x\left(\tan\gamma + 1 + \frac{1}{\tan\gamma}\right)$, d'où :
    \[
        \boxed{x = \frac{a}{1 + \tan\gamma + \dfrac{1}{\tan\gamma}} = \frac{a}{1 + \dfrac{2}{\sin(2\gamma)}}}.
    \]
    \item Pour maximiser l'aire $x^2$, il faut maximiser $x$, donc minimiser le dénominateur $1 + \tan\gamma + \frac{1}{\tan\gamma} = 1 + \frac{2}{\sin(2\gamma)}$. Celui-ci est minimal lorsque $\sin(2\gamma)$ est maximal, c'est-à-dire $\sin(2\gamma) = 1 \iff \boxed{\gamma = \frac{\pi}{4} = 45\degree}$ (triangle rectangle isocèle, donnant $x_{\max} = \frac{a}{3}$).
\end{enumerate}

\medskip
\textbf{3) Pliage d'une bande de papier (UCLouvain -- Juillet 2026, Question 2) :}\\
\emph{Énoncé :} Une bande de papier rectangulaire de largeur $w$ est pliée de telle sorte que le coin supérieur droit $B$ vienne toucher le bord inférieur opposé au point $E$. La ligne de pli $[AC]$ fait un angle $\theta \in \left]0\degree, 45\degree\right[$ avec le bord supérieur $[AB]$. Exprimer la longueur du pli $L = \overline{AC}$ en fonction de $w$ et de $\theta$.

\begin{center}
\begin{tikzpicture}[scale=1.05, >=Stealth]
    \def\myW{2.4}
    \pgfmathsetmacro{\myd}{\myW/sin(60)}
    \def\myL{6.2}
    \coordinate (B) at (\myL, \myW);
    \coordinate (A) at ({\myL-\myd}, \myW);
    \coordinate (C) at (\myL, {\myW - \myd*tan(30)});
    \pgfmathsetmacro{\Ex}{\myL-\myd + \myd*cos(60)}
    \pgfmathsetmacro{\Ey}{\myW - \myd*sin(60)}
    \coordinate (E) at (\Ex, \Ey);

    \draw[thick, gray!50] (0, \myW) -- (\myL, \myW) -- (\myL, 0) -- (0, 0);
    \draw[thick] (0, \myW) -- (A);
    \draw[thick] (0, 0) -- (\myL, 0);
    \draw[thick, uclblue, fill=ucllightblue] (A) -- (C) -- (E) -- cycle;
    \draw[very thick, warnred, dashed] (A) -- (C) node[midway, above right=-2pt] {$L$};
    \draw[thick, dashed, gray] (A) -- (B) -- (C);

    \fill (A) circle (1.5pt) node[above left] {$A$};
    \fill[gray] (B) circle (1.5pt) node[above right] {$B$};
    \fill (C) circle (1.5pt) node[right] {$C$};
    \fill (E) circle (1.5pt) node[below] {$E$};

    \draw[uclaccent, thick] (A) ++(0.85,0) arc (0:-30:0.85);
    \path (A) ++(-15:1.1) node[uclaccent] {$\theta$};
    \draw[uclaccent, thick] (A) ++(-30:0.7) arc (-30:-60:0.7);
    \path (A) ++(-45:0.95) node[uclaccent] {$\theta$};
    \draw[<->, thick, defgreen] (0.6, 0) -- (0.6, \myW) node[midway, left] {$w$};
\end{tikzpicture}
\end{center}

\emph{Résolution :}
\begin{enumerate}
    \item Par symétrie axiale du pliage autour de la droite $(AC)$, le triangle plié $AEC$ est isométrique au triangle initial $ABC$ (rectangle en $B$). On a donc $\overline{AE} = \overline{AB}$ et $\widehat{EAC} = \widehat{BAC} = \theta$.
    \item Dans le triangle rectangle $ABC$, le côté adjacent à $\theta$ vaut $\overline{AB} = L\cos\theta$, d'où $\overline{AE} = L\cos\theta$.
    \item Le segment $[AE]$ forme avec le bord supérieur horizontal $[AB]$ un angle total $\widehat{BAE} = \theta + \theta = 2\theta$. Or la projection verticale de $[AE]$ entre le bord supérieur ($A$) et le bord inférieur ($E$) est exactement la largeur $w$ de la bande ! Dans le triangle rectangle projetant $E$ verticalement sur $[AB]$, on a donc :
    \[
        w = \overline{AE}\sin(2\theta) = L\cos\theta\sin(2\theta) \implies \boxed{L = \frac{w}{\cos\theta\sin(2\theta)} = \frac{w}{2\sin\theta\cos^2\theta}}.
    \]
\end{enumerate}

\medskip
\textbf{4) Trois cercles mutuellement tangents posés sur le sol (UCLouvain -- Septembre 2025, Question 1c) :}\\
\emph{Énoncé :} Trois cercles de centres $A, B, C$ et de rayons respectifs $x, y, z$ sont chacun tangents à une même droite horizontale ainsi que deux à deux entre eux ($C$ étant coincé entre $A$ et $B$). Exprimer $z$ en fonction de $x$ et $y$.

\begin{center}
\begin{tikzpicture}[scale=0.045, >=Stealth]
    \pgfmathsetmacro{\x}{85}
    \pgfmathsetmacro{\y}{42}
    \pgfmathsetmacro{\z}{\x*\y/(sqrt(\x)+sqrt(\y))^2}
    \pgfmathsetmacro{\dxAC}{2*sqrt(\x*\z)}
    \pgfmathsetmacro{\dxCB}{2*sqrt(\y*\z)}
    \pgfmathsetmacro{\dxAB}{\dxAC+\dxCB}

    \coordinate (Ap) at (0,0);
    \coordinate (A) at (0,\x);
    \coordinate (Cp) at (\dxAC,0);
    \coordinate (C) at (\dxAC,\z);
    \coordinate (Bp) at (\dxAB,0);
    \coordinate (B) at (\dxAB,\y);
    \coordinate (H) at (0,\y);

    \draw[thick, darkgray] (-60,0) -- ({\dxAB+55},0);
    \draw[thick, uclblue] (A) circle (\x);
    \draw[thick, uclblue] (B) circle (\y);
    \draw[thick, warnred] (C) circle (\z);

    \draw[dashed, gray] (A)--(Ap) node[midway, left] {\small $x$};
    \draw[dashed, gray] (B)--(Bp) node[midway, right] {\small $y$};
    \draw[dashed, gray] (C)--(Cp);
    \draw[thick, defgreen] (A)--(B) node[midway, above right] {\small $x+y$};
    \draw[dashed, defgreen] (A)--(H)--(B);

    \draw[<->, thick, uclaccent] (0,-18) -- (\dxAC,-18) node[midway, below] {\small $2\sqrt{xz}$};
    \draw[<->, thick, uclaccent] (\dxAC,-18) -- (\dxAB,-18) node[midway, below] {\small $2\sqrt{yz}$};

    \fill (A) circle (2) node[above left] {\small $A$};
    \fill (B) circle (2) node[above right] {\small $B$};
    \fill (C) circle (2) node[above] {\small $C$};
\end{tikzpicture}
\end{center}

\emph{Résolution :} Calculons d'abord la distance horizontale $d_{AB}$ entre les points de contact au sol de deux cercles tangents de rayons $x$ et $y$. Le segment $[AB]$ reliant leurs centres a pour longueur $x + y$ (somme des rayons au point de tangence), et la différence de hauteur entre $A$ et $B$ vaut $|x - y|$. Par le théorème de Pythagore dans le trapèze rectangle formé par les rayons verticaux :
\[
    d_{AB}^2 + (x - y)^2 = (x + y)^2 \implies d_{AB}^2 = (x+y)^2 - (x-y)^2 = 4xy \implies \boxed{d_{AB} = 2\sqrt{xy}}.
\]
En appliquant exactement la même formule aux paires de cercles tangents $(A,C)$ et $(C,B)$, leurs distances horizontales au sol valent $d_{AC} = 2\sqrt{xz}$ et $d_{CB} = 2\sqrt{yz}$. Comme le point de contact de $C$ est situé entre ceux de $A$ et de $B$ :
\[
    d_{AB} = d_{AC} + d_{CB} \iff \sqrt{xy} = \sqrt{z}(\sqrt{x} + \sqrt{y}) \iff \boxed{z = \frac{xy}{(\sqrt{x} + \sqrt{y})^2}}.
\]
\end{exemple}

% ----------------------------------------------------------------------------
\section{Triangles quelconques : lois des sinus, des cosinus et formules d'aire}
% ----------------------------------------------------------------------------

Dans un triangle quelconque $ABC$ (qui n'est pas nécessairement rectangle), trois relations majeures permettent de déterminer tous les côtés et angles à partir de trois données indépendantes (dont au moins une longueur).

\begin{theoreme}[Loi des cosinus, loi des sinus et formules de l'aire d'un triangle quelconque]
Soit un triangle quelconque $ABC$ de côtés $a, b, c$, d'angles $\alpha = \widehat{A}$, $\beta = \widehat{B}$, $\gamma = \widehat{C}$, de demi-périmètre $p = \dfrac{a+b+c}{2}$, de rayon du cercle inscrit $r$ et de rayon du cercle circonscrit $R$ :
\begin{enumerate}
    \item \textbf{Loi des cosinus (Théorème de Pythagore généralisé ou d'Al-Kashi) :}
    \begin{align}
        \boxed{a^2 = b^2 + c^2 - 2bc\cos\alpha} \\
        b^2 = a^2 + c^2 - 2ac\cos\beta \\
        c^2 = a^2 + b^2 - 2ab\cos\gamma
    \end{align}
    \item \textbf{Loi des sinus (et lien avec le rayon du cercle circonscrit $R$) :}
    \begin{equation}
        \boxed{\frac{a}{\sin\alpha} = \frac{b}{\sin\beta} = \frac{c}{\sin\gamma} = 2R}
    \end{equation}
    \item \textbf{Formules de l'aire $S$ du triangle $ABC$ :}
    \begin{itemize}
        \item \emph{En fonction de deux côtés et de l'angle compris entre eux :}
        \[
            \boxed{S = \frac{1}{2}bc\sin\alpha = \frac{1}{2}ac\sin\beta = \frac{1}{2}ab\sin\gamma}
        \]
        \item \emph{En fonction du demi-périmètre $p = \frac{P}{2} = \frac{a+b+c}{2}$ et du rayon du cercle inscrit $r$ :} En décomposant $ABC$ en trois triangles $IAB, IBC, ICA$ de hauteur $r$ autour du centre $I$ du cercle inscrit :
        \[
            \boxed{S = \frac{1}{2}(a+b+c)r = p \cdot r \iff r = \frac{S}{p} = \frac{2S}{a+b+c}}
        \]
        \item \emph{Formule de Héron (en fonction des trois côtés uniquement) :}
        \[
            \boxed{S = \sqrt{p(p-a)(p-b)(p-c)}}
        \]
    \end{itemize}
\end{enumerate}
\end{theoreme}

\begin{center}
\begin{tikzpicture}[scale=1.15, >=Stealth]
    % Figure 1 : Triangle quelconque et hauteur h = b sin alpha
    \begin{scope}[shift={(0,0)}]
        \coordinate (A) at (0,0);
        \coordinate (B) at (4.5,0);
        \coordinate (C) at (1.6,2.5);
        \coordinate (H) at (1.6,0);

        \draw[thick, uclblue, fill=ucllightblue!50] (A) -- (B) -- (C) -- cycle;
        \draw[dashed, warnred, thick] (C) -- (H) node[midway, right] {\small $h = b\sin\alpha$};
        \draw[warnred] (1.6,0.22) -- (1.82,0.22) -- (1.82,0);

        \pic [draw, uclaccent, "$\alpha$", angle radius=0.6cm, angle eccentricity=1.3] {angle = B--A--C};
        \pic [draw, uclaccent, "$\beta$", angle radius=0.6cm, angle eccentricity=1.3] {angle = C--B--A};
        \pic [draw, uclaccent, "$\gamma$", angle radius=0.5cm, angle eccentricity=1.35] {angle = A--C--B};

        \node[below] at ($(A)!0.5!(B)$) {$c$};
        \node[above left] at ($(A)!0.5!(C)$) {$b$};
        \node[above right] at ($(B)!0.5!(C)$) {$a$};
        \fill (A) circle (1.3pt) node[below left] {$A$};
        \fill (B) circle (1.3pt) node[below right] {$B$};
        \fill (C) circle (1.3pt) node[above] {$C$};
        \node[below] at (2.25,-0.55) {\small \textbf{(a)} $S = \frac{1}{2}c h = \frac{1}{2}bc\sin\alpha$};
    \end{scope}

    % Figure 2 : Cercle inscrit et S = p r
    \begin{scope}[shift={(6.8,0)}]
        \coordinate (A) at (0,0);
        \coordinate (B) at (4.5,0);
        \coordinate (C) at (1.6,2.5);
        % Incenter computation
        \pgfmathsetmacro{\c}{4.5}
        \pgfmathsetmacro{\b}{sqrt(1.6^2+2.5^2)}
        \pgfmathsetmacro{\a}{sqrt((4.5-1.6)^2+2.5^2)}
        \pgfmathsetmacro{\p}{(\a+\b+\c)/2}
        \pgfmathsetmacro{\Area}{0.5*4.5*2.5}
        \pgfmathsetmacro{\rin}{\Area/\p}
        \pgfmathsetmacro{\Ix}{(\a*0 + \b*4.5 + \c*1.6)/(\a+\b+\c)}
        \pgfmathsetmacro{\Iy}{(\a*0 + \b*0 + \c*2.5)/(\a+\b+\c)}
        \coordinate (I) at (\Ix,\Iy);

        \draw[thick, uclblue] (A) -- (B) -- (C) -- cycle;
        \draw[dashed, gray] (I) -- (A) (I) -- (B) (I) -- (C);
        \draw[thick, defgreen, fill=deflightgreen!60] (I) circle (\rin);
        \draw[thick, warnred] (I) -- (\Ix,0) node[midway, right] {\small $r$};
        \fill (I) circle (1.3pt) node[above] {\small $I$};
        \fill (A) circle (1.3pt) node[below left] {$A$};
        \fill (B) circle (1.3pt) node[below right] {$B$};
        \fill (C) circle (1.3pt) node[above] {$C$};
        \node[below] at (2.25,-0.55) {\small \textbf{(b)} Cercle inscrit : $S = \frac{a+b+c}{2}\cdot r = p\cdot r$};
    \end{scope}
\end{tikzpicture}
\end{center}

\subsection{Aire d'un quadrilatère quelconque}

Deux formules sur les quadrilatères ont fait l'objet de questions complètes aux examens de Septembre 2023 (Q2), Juillet 2025 (Q2) et Septembre 2026 (Q3a) :

\begin{theoreme}[Aire d'un quadrilatère en fonction de ses diagonales]
Soit $ABCD$ un quadrilatère quelconque (convexe ou concave) dont les diagonales mesurent $x = \overline{AC}$ et $y = \overline{BD}$, et dont les droites $(AC)$ et $(BD)$ forment un angle $\theta$. Alors son aire vaut :
\[
    \boxed{S_{ABCD} = \frac{1}{2} x y \sin\theta}.
\]
\end{theoreme}
\begin{proof}[Démonstration (UCLouvain -- Septembre 2023 Q2a et Juillet 2025 Q2b-2c)]
\mbox{}\\
Dans le cas convexe, soit $O$ le point d'intersection des diagonales $[AC]$ et $[BD]$, et posons $a = \overline{OA}$, $c = \overline{OC}$ (avec $a+c = x$), $b = \overline{OB}$, $d = \overline{OD}$ (avec $b+d = y$). Les quatre angles en $O$ sont $\theta$, $\pi-\theta$, $\theta$, $\pi-\theta$, qui ont tous le même sinus $\sin(\pi-\theta) = \sin\theta$. En additionnant les aires des quatre triangles $OAB, OBC, OCD, ODA$ :
\begin{align*}
    S &= \frac{1}{2}ab\sin\theta + \frac{1}{2}bc\sin\theta + \frac{1}{2}cd\sin\theta + \frac{1}{2}da\sin\theta \\
      &= \frac{1}{2}\sin\theta\bigl(a(b+d) + c(b+d)\bigr) = \frac{1}{2}(a+c)(b+d)\sin\theta = \frac{1}{2}xy\sin\theta.
\end{align*}
Dans le cas concave, la même factorisation s'obtient par différence des aires des triangles $ABD$ et $CBD$.
\end{proof}

% ----------------------------------------------------------------------------
\section{Problèmes complets résolus des examens UCLouvain}
% ----------------------------------------------------------------------------

Étudions maintenant six problèmes complets tirés des annales récentes de l'UCLouvain pour illustrer la méthode de résolution attendue le jour de l'examen.

\begin{exemple}[Problème 1 : Résolution d'un triangle avec $\widehat{A} = 2\widehat{B}$ (UCLouvain -- Septembre 2022, Question 2)]
\emph{Énoncé :} On considère un triangle $ABC$ de côtés $a = \overline{BC}$, $b = \overline{AC}$ et $c = \overline{AB}$. Les longueurs $a$ et $b$ sont connues et on suppose que $\widehat{A} = 2\widehat{B}$.
\begin{enumerate}
    \item Calculez l'angle $\widehat{B}$ en fonction de $a$ et $b$ uniquement.
    \item Trouvez la longueur $c = \overline{AB}$ en fonction de $a$ et $b$ uniquement.
    \item Si $\widehat{B} \to 0$, calculez la limite du rapport $\dfrac{b}{a}$.
\end{enumerate}
\emph{Résolution :}
\begin{enumerate}
    \item Appliquons la \textbf{loi des sinus} aux côtés $a$ et $b$ en utilisant l'hypothèse $\widehat{A} = 2\widehat{B}$ et la formule de duplication $\sin(2\widehat{B}) = 2\sin\widehat{B}\cos\widehat{B}$ :
    \[
        \frac{b}{\sin\widehat{B}} = \frac{a}{\sin\widehat{A}} = \frac{a}{\sin(2\widehat{B})} = \frac{a}{2\sin\widehat{B}\cos\widehat{B}}.
    \]
    Puisque dans un triangle non dégénéré $\sin\widehat{B} \neq 0$, on simplifie par $\sin\widehat{B}$ :
    \[
        \cos\widehat{B} = \frac{a}{2b} \implies \boxed{\widehat{B} = \arccos\left(\frac{a}{2b}\right)}.
    \]
    \item Appliquons la \textbf{loi des cosinus} à l'angle $\widehat{B}$ et remplaçons $\cos\widehat{B}$ par $\frac{a}{2b}$ :
    \[
        b^2 = a^2 + c^2 - 2ac\cos\widehat{B} = a^2 + c^2 - 2ac\left(\frac{a}{2b}\right) = a^2 + c^2 - \frac{a^2 c}{b}.
    \]
    Regroupons les termes pour factoriser par $(b - c)$ :
    \[
        b^2 - c^2 = a^2\left(1 - \frac{c}{b}\right) \iff (b - c)(b + c) = \frac{a^2}{b}(b - c).
    \]
    Lorsque $\widehat{B}$ est quelconque ($\widehat{A} = 2\widehat{B} \neq \widehat{B} \implies a \neq b \implies b \neq c$), on divise par $(b-c)$ pour obtenir :
    \[
        b + c = \frac{a^2}{b} \implies \boxed{c = \frac{a^2}{b} - b = \frac{a^2 - b^2}{b}}.
    \]
    \item Lorsque $\widehat{B} \to 0$, on a $\cos\widehat{B} = \frac{a}{2b} \to \cos(0) = 1$, d'où immédiatement $\boxed{\lim_{\widehat{B}\to 0} \frac{b}{a} = \frac{1}{2}}$.
\end{enumerate}
\end{exemple}

\begin{exemple}[Problème 2 : Partage d'un côté d'un triangle (UCLouvain -- Juillet 2023, Question 2)]
\emph{Énoncé :} On considère un triangle $ABC$ et un point $M$ sur le segment $[BC]$ tel que $\dfrac{\overline{BC}}{\overline{BM}} = 3$.
\begin{enumerate}
    \item \textbf{(Sous-question 2a)} On suppose que $ABC$ est un triangle équilatéral. Calculez la valeur exacte de l'angle $\beta = \widehat{BAM}$.
    \item \textbf{(Sous-question 2b)} Le triangle $ABC$ n'est plus équilatéral mais $\frac{\overline{BC}}{\overline{BM}} = 3$ reste vrai. On suppose que $\overline{AM} = \overline{BM} = y$ et $\overline{AB} = \sqrt{2}y$. Calculez la longueur $\overline{AC}$ en fonction de $y$.
\end{enumerate}

\begin{center}
\begin{tikzpicture}[scale=1.1, >=Stealth]
    % 2a : Equilatéral
    \begin{scope}[shift={(0,0)}]
        \coordinate (C) at (0,0);
        \coordinate (B) at (3.6,0);
        \coordinate (A) at (1.8,{3.6*sin(60)});
        \coordinate (M) at (2.4,0);

        \draw[thick, uclblue] (C) -- (B) -- (A) -- cycle;
        \draw[thick, warnred] (A) -- (M) node[midway, left] {$h$};
        \pic [draw, uclaccent, "$\beta$", angle radius=0.85cm, angle eccentricity=1.25] {angle = M--A--B};

        \node[below] at (1.2,0) {$2x$};
        \node[below] at (3.0,0) {$x$};
        \node[above right] at ($(A)!0.5!(B)$) {$3x$};
        \fill (A) circle (1.4pt) node[above] {$A$};
        \fill (B) circle (1.4pt) node[right] {$B$};
        \fill (C) circle (1.4pt) node[left] {$C$};
        \fill (M) circle (1.4pt) node[below] {$M$};
        \node[below] at (1.8,-0.55) {\small \textbf{(a)} Triangle équilatéral};
    \end{scope}

    % 2b : Triangle ABM isocèle rectangle
    \begin{scope}[shift={(6.2,0)}]
        \coordinate (C) at (0,0);
        \coordinate (M) at (3.2,0);
        \coordinate (B) at (4.8,0);
        \coordinate (A) at (3.2,1.6);

        \draw[thick, uclblue] (C) -- (B) -- (A) -- cycle;
        \draw[thick, warnred] (A) -- (M) node[midway, left] {$y$};
        \draw[warnred] (3.2,0.25) -- (2.95,0.25) -- (2.95,0);

        \node[below] at (1.6,0) {$2y$};
        \node[below] at (4.0,0) {$y$};
        \node[above right] at ($(A)!0.5!(B)$) {$\sqrt{2}y$};
        \node[above left] at ($(A)!0.5!(C)$) {$y\sqrt{5}$};
        \fill (A) circle (1.4pt) node[above] {$A$};
        \fill (B) circle (1.4pt) node[right] {$B$};
        \fill (C) circle (1.4pt) node[left] {$C$};
        \fill (M) circle (1.4pt) node[below] {$M$};
        \node[below] at (2.4,-0.55) {\small \textbf{(b)} Cas $\overline{AM}=\overline{BM}=y$, $\overline{AB}=\sqrt{2}y$};
    \end{scope}
\end{tikzpicture}
\end{center}

\emph{Résolution :}
\begin{enumerate}
    \item \textbf{Attention au piège classique !} Beaucoup d'étudiants écrivent $\beta = \frac{1}{3}\cdot 60\degree = 20\degree$ sous prétexte que $\overline{BM} = \frac{1}{3}\overline{BC}$. \textbf{C'est totalement faux !} Une cévienne qui coupe un côté au tiers ne coupe pas l'angle au tiers.
    Posons $x = \overline{BM}$, de sorte que $\overline{MC} = 2x$, $\overline{AB} = \overline{AC} = 3x$, $\widehat{B} = \widehat{C} = \frac{\pi}{3}$, $\widehat{BAM} = \beta$ et $\widehat{CAM} = \frac{\pi}{3} - \beta$. Notons $h = \overline{AM}$. En appliquant la \textbf{loi des sinus} dans les triangles $ABM$ et $ACM$ :
    \[
        \text{Dans } ABM : \frac{x}{\sin\beta} = \frac{h}{\sin(\pi/3)}, \qquad \text{et dans } ACM : \frac{2x}{\sin(\pi/3 - \beta)} = \frac{h}{\sin(\pi/3)}.
    \]
    En égalant les deux expressions de $\frac{h}{\sin(\pi/3)}$ et en développant $\sin(\frac{\pi}{3}-\beta)$ :
    \[
        \sin\left(\frac{\pi}{3} - \beta\right) = 2\sin\beta \iff \frac{\sqrt{3}}{2}\cos\beta - \frac{1}{2}\sin\beta = 2\sin\beta \iff \frac{5}{2}\sin\beta = \frac{\sqrt{3}}{2}\cos\beta,
    \]
    soit $\tan\beta = \frac{\sqrt{3}}{5}$. Comme $\beta \in \left]0, \frac{\pi}{3}\right[$, on conclut que $\boxed{\beta = \arctan\left(\frac{\sqrt{3}}{5}\right)}$ (ou $\arccos\left(\frac{5}{2\sqrt{7}}\right)$ par la loi des cosinus).
    \item Dans le triangle $ABM$, on constate que $\overline{AM}^2 + \overline{BM}^2 = y^2 + y^2 = 2y^2 = (\sqrt{2}y)^2 = \overline{AB}^2$ ! Par la réciproque du théorème de Pythagore (ou la loi des cosinus), le triangle $ABM$ est \textbf{rectangle isocèle en $M$} ($\widehat{AMB} = \frac{\pi}{2}$). Dès lors, son angle supplémentaire $\widehat{AMC} = \pi - \frac{\pi}{2} = \frac{\pi}{2}$ est aussi droit ! Dans le triangle $AMC$ rectangle en $M$, avec $\overline{AM} = y$ et $\overline{MC} = 2\overline{BM} = 2y$, Pythagore donne immédiatement :
    \[
        \overline{AC} = \sqrt{\overline{AM}^2 + \overline{MC}^2} = \sqrt{y^2 + (2y)^2} = \boxed{y\sqrt{5}}.
    \]
\end{enumerate}
\end{exemple}

\begin{exemple}[Problème 3 : Cercle inscrit dans un triangle isocèle (UCLouvain -- Juillet 2024, Question 2)]
\emph{Énoncé :} Soit un triangle isocèle $ABC$ avec $\overline{AB} = \overline{AC}$. Le cercle inscrit $\mathcal{C}$ à ce triangle est de rayon $r$ et l'angle à la base vaut $\widehat{ABC} = 2\theta$. Calculez l'aire $S$ du triangle $ABC$ en fonction de $r$ et de $\theta$.

\begin{center}
\begin{tikzpicture}[scale=1.15, >=Stealth]
    \coordinate (P) at (0,0);
    \coordinate (B) at (1.8,0);
    \coordinate (C) at (-1.8,0);
    \coordinate (A) at (0,3.2);
    \pgfmathsetmacro{\angB}{atan(3.2/1.8)}
    \pgfmathsetmacro{\th}{\angB/2}
    \pgfmathsetmacro{\rin}{1.8*tan(\th)}
    \coordinate (O) at (0,\rin);
    \coordinate (Q) at ({1.8 - 1.8*cos(\angB)}, {1.8*sin(\angB)});

    \draw[thick, uclblue] (C) -- (B) -- (A) -- cycle;
    \draw[thick, defgreen, fill=deflightgreen!50] (O) circle (\rin);
    \draw[dashed, gray] (A) -- (P);
    \draw[thick, uclaccent] (O) -- (B);
    \draw[thick, warnred] (O) -- (P) node[midway, left] {$r$};
    \draw[thick, warnred] (O) -- (Q) node[midway, above left=-2pt] {$r$};

    \node[below] at (-0.9,0) {$y$};
    \node[below] at (0.9,0) {$y$};
    \node[above left] at ($(A)!0.5!(C)$) {$x$};
    \fill (A) circle (1.4pt) node[above] {$A$};
    \fill (B) circle (1.4pt) node[right] {$B$};
    \fill (C) circle (1.4pt) node[left] {$C$};
    \fill (P) circle (1.4pt) node[below] {$P$};
    \fill (O) circle (1.4pt) node[left] {$O$};
    \fill (Q) circle (1.4pt) node[above right] {$Q$};
    \node[uclaccent] at (1.15,0.22) {$\theta$};
    \node[uclaccent] at (1.28,0.52) {$\theta$};
\end{tikzpicture}
\end{center}

\emph{Résolution :} Soit $P$ le milieu de la base $[BC]$ et $O$ le centre du cercle inscrit $\mathcal{C}$. Par symétrie du triangle isocèle, la hauteur $[AP]$ est perpendiculaire à $[BC]$ en $P$ et passe par $O$ ($[OP]$ est un rayon de longueur $r$).
Posons $y = \overline{BP} = \overline{PC}$ (demi-base) et $h = \overline{AP}$ (hauteur).
\begin{enumerate}
    \item Le centre $O$ du cercle inscrit est à l'intersection des bissectrices des angles du triangle, donc la demi-droite $[BO)$ divise l'angle $\widehat{ABC} = 2\theta$ en deux angles égaux à $\theta$ : $\widehat{OBP} = \theta$.
    \item Dans le triangle $OPB$ rectangle en $P$, on a $\tan\theta = \dfrac{\overline{OP}}{\overline{BP}} = \dfrac{r}{y} \implies y = r\cot\theta = \dfrac{r}{\tan\theta}$.
    \item Dans le triangle $APB$ rectangle en $P$, avec $\widehat{ABP} = 2\theta$, on a $\tan(2\theta) = \dfrac{\overline{AP}}{\overline{BP}} = \dfrac{h}{y} \implies h = y\tan(2\theta)$.
    \item L'aire $S$ du triangle $ABC$ vaut $\frac{1}{2} \cdot \text{base} \cdot \text{hauteur} = \frac{1}{2}(2y)h = y h = y^2\tan(2\theta)$, d'où :
    \[
        \boxed{S = r^2\cot^2(\theta)\tan(2\theta) = \frac{2r^2}{\tan\theta(1 - \tan^2\theta)}}.
    \]
\end{enumerate}
\end{exemple}

\begin{exemple}[Problème 4 : Satellite et station spatiale en orbite (UCLouvain -- Septembre 2025, Question 2)]
\emph{Énoncé :} Un satellite $S$ et une station spatiale $T$ sont en orbite circulaire autour de la Terre (de centre $O$), dans un même plan, sur des cercles centrés en $O$ de rayons respectifs $R$ et $r$ ($r < R$). Un vaisseau spatial $V$ est situé en un point tel que : (i) les segments $[VS]$ et $[VT]$ sont tangents respectivement aux deux orbites en $S$ et $T$, et (ii) $\overline{OS} = \overline{SV}$.
\begin{enumerate}
    \item Déterminez une formule exprimant l'angle $\alpha = \widehat{TVO}$ en fonction de $r$ et $R$ uniquement.
    \item En supposant $\alpha$ connu, exprimez la distance $\overline{ST}$ en fonction de $R$, $r$ et $\alpha$.
\end{enumerate}

\begin{center}
\begin{tikzpicture}[scale=1.6, >=Stealth]
    \def\R{1.6}
    \def\r{0.9}
    \coordinate (O) at (0,0);
    \coordinate (S) at (0,\R);
    \coordinate (T) at ({\r*cos(65)},{\r*sin(65)});
    \coordinate (V) at (-\R,\R);

    \draw[thick, uclblue] (O) circle (\R);
    \draw[thick, uclaccent] (O) circle (\r);

    \draw[thick, uclblue] (O) -- (S) node[midway, left] {$R$};
    \draw[thick, uclaccent] (O) -- (T) node[midway, below right=-2pt] {$r$};
    \draw[thick, defgreen] (V) -- (S) node[midway, above] {$R$};
    \draw[thick, defgreen] (V) -- (T);
    \draw[thick, darkgray] (V) -- (O) node[midway, below left] {$R\sqrt{2}$};
    \draw[very thick, warnred] (S) -- (T);

    % Angle droit en S
    \draw (-0.15,\R) -- (-0.15,{\R-0.15}) -- (0,{\R-0.15});
    \pic [draw, uclaccent, "$\alpha$", angle radius=0.9cm, angle eccentricity=1.25] {angle = O--V--T};

    \fill (O) circle (1.2pt) node[below] {$O$};
    \fill (S) circle (1.2pt) node[above] {$S$};
    \fill (T) circle (1.2pt) node[right] {$T$};
    \fill (V) circle (1.2pt) node[left] {$V$};
\end{tikzpicture}
\end{center}

\emph{Résolution :}
\begin{enumerate}
    \item Comme $[VS]$ est tangent au cercle de rayon $R$ en $S$, le rayon $[OS]$ est perpendiculaire à $[VS]$ ($\widehat{OSV} = \frac{\pi}{2}$). De plus, $\overline{OS} = \overline{SV} = R$, donc le triangle $OSV$ est \textbf{rectangle isocèle en $S$}. Par Pythagore, son hypoténuse vaut $\overline{OV} = \sqrt{R^2 + R^2} = R\sqrt{2}$, et ses angles aigus valent $\widehat{SOV} = \widehat{SVO} = \frac{\pi}{4}$.
    De même, comme $[VT]$ est tangent au cercle de rayon $r$ en $T$, le triangle $OTV$ est \textbf{rectangle en $T$} ($\widehat{OTV} = \frac{\pi}{2}$), d'hypoténuse $\overline{OV} = R\sqrt{2}$ et de côté opposé à $\alpha = \widehat{TVO}$ égal à $\overline{OT} = r$. On en tire :
    \[
        \sin\alpha = \frac{\overline{OT}}{\overline{OV}} = \frac{r}{R\sqrt{2}} \implies \boxed{\alpha = \arcsin\left(\frac{r}{R\sqrt{2}}\right)}.
    \]
    \item Dans le triangle rectangle $OTV$, l'angle en $O$ vaut $\widehat{TOV} = \frac{\pi}{2} - \alpha$. Or $\widehat{SOV} = \frac{\pi}{4}$, donc l'angle $\beta = \widehat{SOT}$ du triangle $SOT$ vaut $\beta = \widehat{TOV} - \widehat{SOV} = \left(\frac{\pi}{2} - \alpha\right) - \frac{\pi}{4} = \frac{\pi}{4} - \alpha$. En appliquant la \textbf{loi des cosinus} dans le triangle $SOT$ (où $\overline{OS} = R$, $\overline{OT} = r$ et $\widehat{SOT} = \frac{\pi}{4} - \alpha$) :
    \[
        \boxed{\overline{ST} = \sqrt{R^2 + r^2 - 2Rr\cos\left(\frac{\pi}{4} - \alpha\right)}}.
    \]
\end{enumerate}
\end{exemple}

\begin{exemple}[Problème 5 : Cercle inscrit dans un secteur circulaire (UCLouvain -- Septembre 2026, Question 2)]
\emph{Énoncé :} Un secteur circulaire de centre $O$, de rayon $R$ et d'angle au centre $\theta \in \left]0, \pi\right[$ est donné. On inscrit un cercle de centre $C$ et de rayon $r$ à l'intérieur de ce secteur, de sorte qu'il soit tangent à l'arc du secteur ainsi qu'aux deux rayons qui le délimitent.
\begin{enumerate}
    \item Exprimez le rayon $r$ du cercle inscrit en fonction de $R$ et de $\theta$.
    \item Déterminez la valeur de $\theta$ (en radians et en degrés) pour laquelle $\dfrac{r}{R} = \dfrac{1}{3}$.
\end{enumerate}

\begin{center}
\begin{tikzpicture}[scale=1.15, >=Stealth]
    \def\R{3.4}
    \def\ang{30}
    \pgfmathsetmacro{\r}{\R*sin(\ang)/(1+sin(\ang))}
    \pgfmathsetmacro{\d}{\R - \r}

    \coordinate (O) at (0,0);
    \coordinate (C) at (\d,0);
    \coordinate (A) at ({\R*cos(\ang)}, {\R*sin(\ang)});
    \coordinate (B) at ({\R*cos(-\ang)}, {\R*sin(-\ang)});
    \coordinate (T) at ({\d*cos(\ang)*cos(\ang)}, {\d*cos(\ang)*sin(\ang)});

    \draw[thick, uclblue] (O) -- (A) arc (\ang:-\ang:\R) -- cycle;
    \draw[thick, defgreen, fill=deflightgreen!60] (C) circle (\r);
    \draw[dashed, uclaccent, thick] (O) -- (C) node[midway, below] {$R-r$};
    \draw[thick, warnred] (C) -- (T) node[midway, above left=-2pt] {$r$};

    \begin{scope}[shift={(T)}, rotate=\ang]
        \draw[warnred] (-0.18,0) -- (-0.18,-0.18) -- (0,-0.18);
    \end{scope}

    \draw[uclaccent] (0.7,0) arc (0:\ang:0.7);
    \node[uclaccent] at (0.95, 0.2) {$\theta/2$};
    \draw[<->, gray] (0, -1.85) -- (\R, -1.85) node[midway, below] {$R$};

    \fill (O) circle (1.3pt) node[left] {$O$};
    \fill (C) circle (1.3pt) node[below right] {$C$};
    \fill (T) circle (1.3pt) node[above left] {$T$};
\end{tikzpicture}
\end{center}

\emph{Résolution :}
\begin{enumerate}
    \item Par symétrie, le centre $C$ du cercle inscrit se situe sur la bissectrice de l'angle au centre $\theta$, qui partage l'angle en deux angles égaux à $\frac{\theta}{2}$ et rencontre l'arc du secteur au point de tangence entre le cercle et l'arc. La distance entre $O$ et $C$ vaut donc $\overline{OC} = R - r$.
    Soit $T$ le point de tangence du cercle avec l'un des deux rayons du secteur. Le rayon $[CT]$ de longueur $r$ est perpendiculaire à ce rayon en $T$, de sorte que le triangle $OTC$ est \textbf{rectangle en $T$}. Dans ce triangle rectangle :
    \[
        \sin\left(\frac{\theta}{2}\right) = \frac{\overline{CT}}{\overline{OC}} = \frac{r}{R - r} \iff r\bigl(1 + \sin(\theta/2)\bigr) = R\sin(\theta/2) \iff \boxed{r = \frac{R\sin(\theta/2)}{1 + \sin(\theta/2)}}.
    \]
    \item Pour que $\dfrac{r}{R} = \dfrac{1}{3}$, on résout :
    \[
        \frac{\sin(\theta/2)}{1 + \sin(\theta/2)} = \frac{1}{3} \iff 3\sin\left(\frac{\theta}{2}\right) = 1 + \sin\left(\frac{\theta}{2}\right) \iff \sin\left(\frac{\theta}{2}\right) = \frac{1}{2}.
    \]
    Comme $0 < \theta < \pi$, on a $0 < \frac{\theta}{2} < \frac{\pi}{2}$, d'où $\frac{\theta}{2} = \frac{\pi}{6}$, soit $\boxed{\theta = \frac{\pi}{3} \text{ rad} = 60\degree}$.
\end{enumerate}
\end{exemple}

\begin{exemple}[Problème 6 : Point de Brocard d'un triangle (UCLouvain -- Mini-Blocus 2026, Question 2a)]
\emph{Énoncé :} Soit $ABC$ un triangle non aplati de côtés $a, b, c$ et d'angles $\widehat{A}, \widehat{B}, \widehat{C}$. On considère un point $M$ intérieur au triangle tel que $\widehat{MAB} = \widehat{MBC} = \widehat{MCA} = \theta \in \left]0, \frac{\pi}{3}\right[$. On pose $x = \overline{MA}$, $y = \overline{MB}$ et $z = \overline{MC}$. Dans les triangles $MAB$, $MBC$ et $MCA$, exprimer $x, y, z$ en fonction de $a, b, c$, des angles $\widehat{A}, \widehat{B}, \widehat{C}$ et de $\theta$.\\
\emph{Résolution :} Dans le triangle $MAB$, l'angle en $A$ vaut $\widehat{MAB} = \theta$ et l'angle en $B$ vaut $\widehat{ABM} = \widehat{B} - \widehat{MBC} = \widehat{B} - \theta$. Par conséquent, le troisième angle du triangle $MAB$ (en $M$) vaut :
\[
    \widehat{AMB} = \pi - \bigl(\theta + (\widehat{B} - \theta)\bigr) = \pi - \widehat{B}.
\]
En appliquant la \textbf{loi des sinus} dans le triangle $MAB$ (sachant que $\sin(\pi - \widehat{B}) = \sin\widehat{B}$) :
\[
    \frac{c}{\sin(\pi - \widehat{B})} = \frac{x}{\sin(\widehat{B} - \theta)} = \frac{y}{\sin\theta} \implies \boxed{x = \frac{c\sin(\widehat{B} - \theta)}{\sin\widehat{B}}}.
\]
Par permutation circulaire sur les triangles $MBC$ (où $\widehat{BMC} = \pi - \widehat{C}$) et $MCA$ (où $\widehat{CMA} = \pi - \widehat{A}$), on obtient immédiatement :
\[
    \boxed{y = \frac{a\sin(\widehat{C} - \theta)}{\sin\widehat{C}} \quad \text{(ou } y = \frac{c\sin\theta}{\sin\widehat{B}}\text{)} \qquad \text{et} \qquad z = \frac{b\sin(\widehat{A} - \theta)}{\sin\widehat{A}} \quad \text{(ou } z = \frac{a\sin\theta}{\sin\widehat{C}}\text{)}}.
\]
\end{exemple}

% ----------------------------------------------------------------------------
\section{Exercices du Chapitre 4}
% ----------------------------------------------------------------------------

\emph{Les corrigés complets de tous ces exercices se trouvent au Chapitre 9, en fin d'ouvrage.}

\begin{enumerate}[label=\textbf{Exercice 4.\arabic* --}, align=left, leftmargin=*]
    \item \textbf{(UCLouvain -- Septembre 2024, Q2 : Théorème de la bissectrice)} Soit un triangle quelconque $ABC$ de côtés $a = \overline{BC}$, $b = \overline{AC}$ et $c = \overline{AB}$. Soit $X$ le point du segment $[BC]$ tel que la demi-droite $[AX)$ soit la bissectrice de l'angle $\widehat{BAC}$ (c'est-à-dire $\widehat{BAX} = \widehat{XAC} = \alpha$). En appliquant la loi des sinus dans les triangles $ABX$ et $ACX$, exprimez la longueur $x = \overline{BX}$ en fonction de $a, b$ et $c$ uniquement.

    \item \textbf{(UCLouvain -- Mini-Blocus 2026 Q1e, Septembre 2026 Q3e et Septembre 2023 Q3f : Cercles inscrit et circonscrit d'un triangle rectangle isocèle)} Soit un triangle rectangle isocèle dont l'hypoténuse mesure $c$ (et les deux côtés de l'angle droit mesurent $\frac{c}{\sqrt{2}}$). On note $r$ le rayon de son cercle inscrit et $R$ le rayon de son cercle circonscrit.
    \begin{enumerate}
        \item Calculez exactement le rapport $\dfrac{r}{R}$.
        \item Si l'hypoténuse mesure $c = 1\text{ m}$, calculez la distance $d$ entre le centre du cercle inscrit et le centre du cercle circonscrit à ce triangle.
    \end{enumerate}

    \item \textbf{(UCLouvain -- Juillet 2026, Q3e et Septembre 2026, Q3a : Trapèzes à diagonales orthogonales)}
    \begin{enumerate}
        \item Soit $ABCD$ un trapèze \textbf{rectangle} en $A$ et $D$ dont les bases parallèles $[AB]$ et $[CD]$ mesurent respectivement $\overline{AB} = x$ et $\overline{CD} = y$. Les diagonales $[AC]$ et $[BD]$ se coupent perpendiculairement ($90\degree$). Déterminez la hauteur $h = \overline{AD}$ de ce trapèze en fonction de $x$ et $y$.
        \item Soit $ABCD$ un trapèze \textbf{isocèle} dont les bases parallèles $[AB]$ et $[CD]$ mesurent $\overline{AB} = x$ et $\overline{CD} = y$, et dont les diagonales $[AC]$ et $[BD]$ se coupent à angle droit ($90\degree$). Déterminez l'aire de ce trapèze en fonction de $x$ et $y$.
    \end{enumerate}

    \item \textbf{(UCLouvain -- Juillet 2022, Q3b : Trois cordes dans un cercle)} Dans un cercle, trois cordes de longueurs $2$, $3$ et $4$ déterminent respectivement des angles au centre $\alpha$, $\beta$ et $\alpha+\beta$ (avec $\alpha+\beta < \pi$). Calculez la valeur exacte de $\cos(\alpha)$.

    \item \textbf{(UCLouvain -- Juillet 2021, Q3a, Q3c et Q3f)}
    \begin{enumerate}
        \item Soient deux polygones réguliers à $n$ côtés, l'un ($I$) inscrit dans un cercle de rayon $R$ et l'autre ($C$) circonscrit à ce même cercle. Calculez le rapport de leurs périmètres $\dfrac{P_C}{P_I}$ en fonction de $n$ uniquement.
        \item Dans un triangle $ABC$, si $\cos^2(\widehat{A} - \widehat{B}) - \cos(2\widehat{A})\cos(2\widehat{B}) = \dfrac{1}{2}$ et que $\widehat{C}$ est un angle obtus, calculez l'angle $\widehat{C}$ en degrés.
        \item Soit un triangle $ABC$ rectangle en $A$ dont le cercle circonscrit a pour rayon $R = 2$ et tel que $\overline{AB} = \sqrt{7}$. Calculez la valeur exacte de $\cos(\widehat{B} - \widehat{C})$.
    \end{enumerate}

    \item \textbf{(UCLouvain -- Juillet 2024, Q3c et Septembre 2024, Q3c)}
    \begin{enumerate}
        \item Dans un triangle quelconque $ABC$, l'angle $\widehat{BAC} = \frac{\pi}{6}$, le côté $\overline{AC} = 2\text{ cm}$ et l'aire du triangle vaut $3\text{ cm}^2$. Calculez la valeur exacte de la longueur du côté $\overline{BC}$.
        \item Dans un triangle $ABC$, sachant que $\widehat{A} = 60\degree$ et $\widehat{B} = 45\degree$, calculez la valeur exacte de $\sqrt{2}(\sqrt{3}+1)\dfrac{a+b}{c}$.
    \end{enumerate}

    \item \textbf{(UCLouvain -- Juillet 2025 / Recueil 2025, Q3.5)} Considérons le cercle $\mathcal{C}$ de centre $O$ et de rayon $1$, passant par deux points $A$ et $B$ tels que $\widehat{AOB} = \dfrac{\pi}{3}$. Sur la droite $d$ passant par $A$ et $B$, on place le point $C$ (du côté de $B$) tel que $\overline{AB} = \overline{BC}$. Déterminez la valeur exacte de l'angle $\widehat{BOC}$.

    \item \textbf{(UCLouvain -- Juillet 2022, Q2 : Carrés et rectangle inscrits dans un cercle)} On considère un cercle $\mathcal{C}$ de centre $O$ et de rayon $1$, et un carré $ABCD$ inscrit dans ce cercle.
    \begin{enumerate}
        \item Calculez la longueur $\overline{AB}$ du côté du carré $ABCD$.
        \item On construit un petit carré $EFGH$ à l'extérieur du côté $[AB]$ mais à l'intérieur du cercle $\mathcal{C}$, avec $E, F \in [AB]$ et $G, H \in \mathcal{C}$. Calculez la valeur exacte du rapport $\dfrac{\overline{AB}}{\overline{HE}}$.
        \item On construit de même un rectangle $IJKL$ de l'autre côté, avec $I, J \in [CD]$ et $K, L \in \mathcal{C}$. Exprimez l'aire $A$ du rectangle $IJKL$ en fonction de $y = \overline{JK}$, et déterminez la valeur du rapport $\dfrac{\overline{DC}}{\overline{JK}}$ qui maximise cette aire.
    \end{enumerate}

    \item \textbf{(ULiège -- Juillet 2008, Q3 : Trois cordes équidistantes dans un demi-cercle)} Dans un demi-cercle de rayon $R$, on trace trois cordes parallèles au diamètre de base, de longueurs respectives $C_1 = 8\text{ m}$, $C_2 = 16\text{ m}$ et $C_3 = 20\text{ m}$, telles que la distance $h$ entre $C_1$ et $C_2$ soit égale à la distance $h$ entre $C_2$ et $C_3$. Calculez la valeur exacte de $h$ ainsi que le rayon $R$ du demi-cercle.

    \item \textbf{(ULiège -- Septembre 2006, Q2 : Théorème de Ménélaüs par la loi des sinus)} Soit un triangle $ABC$ et soit une droite $\mathcal{D}$ qui coupe le côté $[AB]$ en $F$, le côté $[AC]$ en $E$ et le prolongement du côté $[BC]$ (au-delà de $C$) en un point $D$. Sachant que $\dfrac{\overline{FA}}{\overline{FB}} = \dfrac{1}{3}$ et $\dfrac{\overline{EA}}{\overline{EC}} = \dfrac{2}{3}$, calculez la valeur exacte du rapport $\dfrac{\overline{DC}}{\overline{DB}}$.
\end{enumerate}

% ============================================================================
% CHAPITRE 5 : GÉOMÉTRIE SYNTHÉTIQUE PLANE ET DANS L'ESPACE
% ============================================================================
\chapter[Géométrie Synthétique Plane et Spatiale]{Géométrie Synthétique Plane et dans l'Espace}

Avec la fusion des épreuves de trigonométrie et de géométrie à l'examen d'admission de l'École Polytechnique de Louvain (UCLouvain), la maîtrise de la \textbf{géométrie synthétique} (c'est-à-dire la géométrie pure, sans recours obligatoire à un système d'axes cartésiens) constitue un atout décisif. Les questions de géométrie synthétique de l'examen évaluent votre capacité à observer une figure, à tracer une droite auxiliaire judicieuse (hauteur, médiatrice, rayon au point de tangence, parallèle) et à enchaîner des déductions rigoureuses appuyées sur les théorèmes classiques du secondaire (angles dans le cercle, puissance d'un point, triangles semblables, découpages d'aires et sections planes de solides).

% ----------------------------------------------------------------------------
\section{Géométrie du cercle : angles, cordes, tangentes et quadrilatères inscriptibles}
% ----------------------------------------------------------------------------

\subsection{Cordes, médiatrices et tangentes à un cercle}

Commençons par trois propriétés élémentaires mais omniprésentes dans les questions d'examen (notamment en Juillet 2024, Juillet 2026 et Septembre 2026) :

\begin{theoreme}[Propriétés fondamentales des cordes et des tangentes]
Soit $\mathcal{C}$ un cercle de centre $O$ et de rayon $r$ :
\begin{enumerate}
    \item \textbf{Médiatrice d'une corde :} Toute corde $[AB]$ forme avec le centre $O$ un triangle isocèle $OAB$ ($|OA| = |OB| = r$). Par conséquent, la droite reliant le centre $O$ au milieu $M$ de la corde $[AB]$ est la \textbf{médiatrice} de $[AB]$ (elle est perpendiculaire à $[AB]$ en $M$) et la \textbf{bissectrice} de l'angle au centre $\widehat{AOB}$ :
    \[
        M \text{ milieu de } [AB] \implies (OM) \perp (AB) \quad \text{et} \quad |OM|^2 + |AM|^2 = r^2.
    \]
    \item \textbf{Égalité des cordes et des arcs :} Dans un même cercle, deux arcs ont la même longueur si et seulement si les cordes qui les sous-tendent ont la même longueur :
    \[
        \text{longueur}(\wideparen{AB}) = \text{longueur}(\wideparen{CD}) \iff |AB| = |CD| \iff |\widehat{AOB}| = |\widehat{COD}|.
    \]
    \item \textbf{Tangente à un cercle et cercles tangents :}
    \begin{itemize}
        \item La tangente à $\mathcal{C}$ en un point $T \in \mathcal{C}$ est \textbf{perpendiculaire} au rayon $[OT]$.
        \item Si d'un point extérieur $P$ on mène les deux tangentes au cercle $\mathcal{C}$ aux points $T_1$ et $T_2$, alors $|PT_1| = |PT_2|$ et la demi-droite $[PO)$ est bissectrice de $\widehat{T_1 P T_2}$.
        \item Si deux cercles de centres $O_1, O_2$ et de rayons $r_1, r_2$ sont tangents en un point $T$, alors les trois points $O_1, O_2, T$ sont \textbf{alignés}, avec $|O_1 O_2| = r_1 + r_2$ (tangence extérieure) ou $|O_1 O_2| = |r_1 - r_2|$ (tangence intérieure).
    \end{itemize}
\end{enumerate}
\end{theoreme}

\subsection{Théorème de l'angle inscrit, de l'angle au centre et quadrilatère inscriptible}

\begin{theoreme}[Théorème de l'angle inscrit et de l'angle au centre]
Soient $A, B, M$ trois points distincts sur un cercle $\mathcal{C}$ de centre $O$ :
\begin{enumerate}
    \item \textbf{Relation angle inscrit -- angle au centre :} L'amplitude d'un angle inscrit $\widehat{AMB}$ est égale à la \textbf{moitié} de l'amplitude de l'angle au centre $\widehat{AOB}$ qui intercepte le même arc $\wideparen{AB}$ :
    \[
        \boxed{|\widehat{AMB}| = \frac{1}{2}|\widehat{AOB}|}.
    \]
    \item \textbf{Invariance de l'angle inscrit :} Deux angles inscrits $\widehat{AMB}$ et $\widehat{ANB}$ qui interceptent le même arc $\wideparen{AB}$ (situés du même côté de la corde $[AB]$) ont la \textbf{même amplitude} : $|\widehat{AMB}| = |\widehat{ANB}|$.
    \item \textbf{Angle tangentiel :} L'angle formé par une corde $[AB]$ et la tangente $(AT)$ au cercle en $A$ a la même amplitude que tout angle inscrit $\widehat{AMB}$ interceptant l'arc $\wideparen{AB}$.
    \item \textbf{Triangle inscrit dans un demi-cercle (Théorème de Thalès du cercle) :} Tout angle inscrit qui intercepte un diamètre $[AB]$ est un \textbf{angle droit} ($90\degree = \frac{\pi}{2}$) :
    \[
        [AB] \text{ est un diamètre de } \mathcal{C} \text{ et } M \in \mathcal{C} \setminus \{A,B\} \iff \widehat{AMB} = 90\degree.
    \]
    \item \textbf{Quadrilatère inscriptible (cocyclique) :} Un quadrilatère convexe $ABCD$ est inscriptible dans un cercle si et seulement si ses \textbf{angles opposés sont supplémentaires} :
    \[
        \boxed{|\widehat{ABC}| + |\widehat{ADC}| = 180\degree \iff |\widehat{DAB}| + |\widehat{DCB}| = 180\degree}.
    \]
\end{enumerate}
\end{theoreme}

\begin{center}
\begin{tikzpicture}[scale=1.15]
    % Figure 1 : Angle inscrit et angle au centre
    \begin{scope}[shift={(0,0)}]
        \draw[thick, uclblue] (0,0) circle (1.6);
        \coordinate (O) at (0,0);
        \coordinate (A) at (-130:1.6);
        \coordinate (B) at (-30:1.6);
        \coordinate (M) at (105:1.6);
        \coordinate (N) at (50:1.6);

        \draw[thick, warnred] (A) -- (O) -- (B);
        \draw[thick, uclaccent] (A) -- (M) -- (B);
        \draw[thick, defgreen, dashed] (A) -- (N) -- (B);

        \pic[draw=warnred, fill=warnred!15, angle radius=0.35cm, "$2\alpha$", angle eccentricity=1.5] {angle = A--O--B};
        \pic[draw=uclaccent, fill=uclaccent!15, angle radius=0.45cm, "$\alpha$", angle eccentricity=1.4] {angle = A--M--B};
        \pic[draw=defgreen, fill=defgreen!15, angle radius=0.45cm, "$\alpha$", angle eccentricity=1.4] {angle = A--N--B};

        \fill (O) circle (1.3pt) node[above] {\small $O$};
        \fill (A) circle (1.3pt) node[below left] {\small $A$};
        \fill (B) circle (1.3pt) node[below right] {\small $B$};
        \fill (M) circle (1.3pt) node[above left] {\small $M$};
        \fill (N) circle (1.3pt) node[above right] {\small $N$};
        \node[below, align=center] at (0,-1.85) {\small \textbf{(a)} Angle inscrit $|\widehat{AMB}| = |\widehat{ANB}| = \alpha$\\[-2pt]\small et angle au centre $|\widehat{AOB}| = 2\alpha$};
    \end{scope}

    % Figure 2 : Quadrilatère inscriptible
    \begin{scope}[shift={(5.2,0)}]
        \draw[thick, uclblue] (0,0) circle (1.6);
        \coordinate (A) at (145:1.6);
        \coordinate (B) at (-135:1.6);
        \coordinate (C) at (-25:1.6);
        \coordinate (D) at (55:1.6);

        \draw[thick, uclaccent] (A) -- (B) -- (C) -- (D) -- cycle;
        \pic[draw=uclaccent, fill=uclaccent!15, angle radius=0.4cm, "$\beta$", angle eccentricity=1.45] {angle = C--B--A};
        \pic[draw=warnred, fill=warnred!15, angle radius=0.35cm, "$\pi-\beta$", angle eccentricity=1.65] {angle = A--D--C};

        \fill (A) circle (1.3pt) node[above left] {\small $A$};
        \fill (B) circle (1.3pt) node[below left] {\small $B$};
        \fill (C) circle (1.3pt) node[below right] {\small $C$};
        \fill (D) circle (1.3pt) node[above right] {\small $D$};
        \node[below, align=center] at (0,-1.85) {\small \textbf{(b)} Quadrilatère inscriptible :\\[-2pt]\small $|\widehat{B}| + |\widehat{D}| = 180\degree = \pi$};
    \end{scope}
\end{tikzpicture}
\end{center}

\begin{exemple}[Problèmes d'examen UCLouvain sur les cercles et angles inscrits]
\textbf{1) Rayon d'un cercle à partir d'une corde partagée (UCLouvain -- Septembre 2026, Question 3) :}\\
\emph{Énoncé :} Dans un cercle de centre $O$ et de rayon $r$, on trace une corde $[AB]$ sur laquelle se situe un point $P$ tel que $|AP| = 3\text{ cm}$, $|PB| = 7\text{ cm}$ et $|OP| = 6\text{ cm}$. Déterminer le rayon $r$ du cercle.\\
\emph{Résolution :} Soit $M$ le milieu de la corde $[AB]$.
\begin{enumerate}
    \item La longueur totale de la corde vaut $|AB| = |AP| + |PB| = 3 + 7 = 10\text{ cm}$, donc le demi-segment mesure $|AM| = |MB| = 5\text{ cm}$, et la distance entre $P$ et le milieu $M$ vaut $|PM| = |AM| - |AP| = 5 - 3 = 2\text{ cm}$.
    \item Puisque $M$ est le milieu de la corde $[AB]$ et que $|OA| = |OB| = r$, la droite $(OM)$ est la médiatrice de $[AB]$, donc $(OM) \perp (AB)$.
    \item Dans le triangle $OMP$ rectangle en $M$, le théorème de Pythagore donne $|OM|^2 = |OP|^2 - |PM|^2 = 6^2 - 2^2 = 36 - 4 = 32\text{ cm}^2$.
    \item Dans le triangle $OMA$ rectangle en $M$, Pythagore donne enfin $r^2 = |OA|^2 = |AM|^2 + |OM|^2 = 5^2 + 32 = 25 + 32 = 57$, d'où $r = \boxed{\sqrt{57}\text{ cm}}$.
\end{enumerate}

\medskip
\textbf{2) Le rayon du triple arc (UCLouvain -- Juillet 2026, Question 4) :}\\
\emph{Énoncé :} Sur un demi-cercle de centre $O$ et de diamètre $[PQ]$, les trois arcs $\wideparen{PA}$, $\wideparen{AB}$ et $\wideparen{BC}$ ont la même longueur. Le point $T$ est le pied de la perpendiculaire à la corde $[PC]$ abaissée depuis $A$. Sachant que $|PT| = 7\text{ cm}$ et $|TC| = 15\text{ cm}$, déterminer : (a) la longueur de la corde $|PB|$, (b) l'égalité des angles $|\widehat{BPC}| = |\widehat{BOM}|$ où $M$ est le milieu de $[BC]$, et (c) le rayon $r$ du demi-cercle.
\begin{center}
\begin{tikzpicture}[scale=1.1]
    \draw[thick, uclblue] (-2.5,0) -- (2.5,0);
    \draw[thick, uclblue] (2.5,0) arc (0:180:2.5);
    \coordinate (O) at (0,0);
    \coordinate (P) at (-2.5,0);
    \coordinate (Q) at (2.5,0);
    \coordinate (A) at (130:2.5);
    \coordinate (B) at (80:2.5);
    \coordinate (C) at (30:2.5);
    \coordinate (M) at ($0.5*(B)+0.5*(C)$);

    \draw[thick] (P) -- (A) -- (B) -- (C);
    \draw[thick, uclaccent] (P) -- (C);
    \draw[thick, warnred] (P) -- (B);
    \draw[thick, defgreen] (O) -- (B) (O) -- (M);

    \coordinate (T) at ($(P)!(A)!(C)$);
    \coordinate (N) at ($(P)!(B)!(C)$);
    \draw[dashed] (A) -- (T) node[below] {\small $T$};
    \draw[dashed, uclaccent] (B) -- (N) node[below right] {\small $N$};
    \pic[draw, angle radius=0.18cm] {right angle = A--T--C};
    \pic[draw, angle radius=0.18cm] {right angle = B--N--P};

    \fill (O) circle (1.3pt) node[below] {\small $O$};
    \fill (P) circle (1.3pt) node[below left] {\small $P$};
    \fill (Q) circle (1.3pt) node[below right] {\small $Q$};
    \fill (A) circle (1.3pt) node[above left] {\small $A$};
    \fill (B) circle (1.3pt) node[above] {\small $B$};
    \fill (C) circle (1.3pt) node[right] {\small $C$};
    \fill (M) circle (1.3pt) node[above right] {\small $M$};
\end{tikzpicture}
\end{center}
\emph{Résolution :}
\begin{enumerate}[label=(\alph*)]
    \item Puisque les trois arcs $\wideparen{PA}$, $\wideparen{AB}$ et $\wideparen{BC}$ ont la même longueur, les trois cordes ont la même longueur $c = |PA| = |AB| = |BC|$. Soit $N$ le pied de la perpendiculaire à $[PC]$ passant par $B$. Par symétrie orthogonale de la trapèze isocèle $PABC$ par rapport à la médiatrice de $[PC]$, on a $|NC| = |PT| = 7\text{ cm}$ et $ABNT$ est un rectangle. Dès lors, $c = |AB| = |TN| = |TC| - |NC| = 15 - 7 = 8\text{ cm}$, et $|BC| = 8\text{ cm}$.
    \begin{itemize}
        \item Dans le triangle $BNC$ rectangle en $N$ : $|BN|^2 = |BC|^2 - |NC|^2 = 8^2 - 7^2 = 64 - 49 = 15 \implies |BN| = \sqrt{15}\text{ cm}$.
        \item Dans le triangle $PBN$ rectangle en $N$ (où $|PN| = |PT| + |TN| = 7 + 8 = 15\text{ cm}$) : $|PB|^2 = |PN|^2 + |BN|^2 = 15^2 + 15 = 240 \implies |PB| = \boxed{4\sqrt{15}\text{ cm}}$.
    \end{itemize}
    \item L'angle inscrit $\widehat{BPC}$ et l'angle au centre $\widehat{BOC}$ interceptent le même arc $\wideparen{BC}$, donc $|\widehat{BOC}| = 2|\widehat{BPC}|$. Or $M$ est le milieu de la corde $[BC]$ dans le triangle isocèle $OBC$ ($|OB| = |OC| = r$), donc $(OM)$ est la médiatrice de $[BC]$ et la bissectrice de $\widehat{BOC}$, c'est-à-dire $|\widehat{BOC}| = 2|\widehat{BOM}|$. On conclut que $|\widehat{BPC}| = |\widehat{BOM}|$.
    \item Dans le triangle $PBN$ rectangle en $N$, $\sin\widehat{BPC} = \frac{|BN|}{|PB|} = \frac{\sqrt{15}}{4\sqrt{15}} = \frac{1}{4}$. Dans le triangle $BOM$ rectangle en $M$ (où $|BM| = \frac{1}{2}|BC| = 4\text{ cm}$ et $|OB| = r$), $\sin\widehat{BOM} = \frac{|BM|}{|OB|} = \frac{4}{r}$. Puisque $\widehat{BPC} = \widehat{BOM}$, on a $\frac{4}{r} = \frac{1}{4}$, soit $r = \boxed{16\text{ cm}}$.
\end{enumerate}

\medskip
\textbf{3) Un théorème (modifié) d'Archimède (UCLouvain -- Septembre 2022, Question 3) :}\\
\emph{Énoncé :} Soit $A$ le milieu d'un demi-cercle de diamètre $[BC]$. Une demi-droite d'origine $B$ coupe le segment $[AC]$ en $D$ et le demi-cercle en $Q$ ($Q \neq B$). Soit $P$ le pied de la perpendiculaire à $(BQ)$ issue de $A$, et soit $R$ le point extérieur au demi-cercle sur la demi-droite $[BQ)$ tel que $|QR| = |QC|$. (a) Calculer l'amplitude de $\widehat{AQC}$. (b) Montrer que $|AC| = |AR|$. (c) En déduire que $|BP| = |PQ| + |QC|$ (indépendamment de la position de $D$ !).\\
\emph{Résolution :}
\begin{enumerate}[label=(\alph*)]
    \item Le triangle $ABC$ est inscrit dans le demi-cercle de diamètre $[BC]$ et $A$ est le milieu de l'arc $\wideparen{BC}$, donc $ABC$ est rectangle isocèle en $A$, ce qui donne $|\widehat{ABC}| = 45\degree$. Or le quadrilatère $ABCQ$ est inscrit dans le cercle, donc ses angles opposés $\widehat{ABC}$ et $\widehat{AQC}$ sont supplémentaires : $|\widehat{AQC}| = 180\degree - 45\degree = \boxed{135\degree}$.
    \item L'angle inscrit $\widehat{BQC}$ intercepte le diamètre $[BC]$, donc $|\widehat{BQC}| = 90\degree$, et son supplément sur la droite $(BR)$ vaut $|\widehat{CQR}| = 90\degree$. Comme le tour complet autour de $Q$ vaut $360\degree$, on a $|\widehat{AQR}| = 360\degree - |\widehat{AQC}| - |\widehat{CQR}| = 360\degree - 135\degree - 90\degree = 135\degree = |\widehat{AQC}|$. Les triangles $AQC$ et $AQR$ ont le côté $[AQ]$ commun, $|QC| = |QR|$ par construction, et l'angle compris égal ($135\degree$) : ils sont donc \textbf{isométriques} (cas CAC), d'où $|AC| = |AR|$.
    \item Comme $|AB| = |AC|$ (triangle rectangle isocèle $ABC$) et $|AC| = |AR|$, on a $|AB| = |AR|$ : le triangle $ABR$ est donc \textbf{isocèle en $A$}. Sa hauteur $(AP) \perp (BR)$ est aussi la médiatrice de la base $[BR]$! Par conséquent, $P$ est le milieu de $[BR]$, soit $|BP| = |PR| = |PQ| + |QR| = |PQ| + |QC|$.
\end{enumerate}
\end{exemple}

% ----------------------------------------------------------------------------
\section{Puissance d'un point par rapport à un cercle et relations métriques}
% ----------------------------------------------------------------------------

Les deux sessions de l'examen d'entrée 2025 (Juillet 2025 Questions 3 et 4, Septembre 2025 Question 3) ont porté directement sur un théorème classique de la géométrie du cercle : la \textbf{puissance d'un point par rapport à un cercle}.

\begin{theoreme}[Puissance d'un point par rapport à un cercle]
Soit $\mathcal{C}$ un cercle de centre $O$ et de rayon $r$, et soit $P$ un point quelconque du plan. Si l'on trace par $P$ une droite sécante qui coupe $\mathcal{C}$ en deux points $A$ et $B$, alors le produit des distances $|PA| \cdot |PB|$ est \textbf{constant} (indépendant de l'orientation de la sécante issue de $P$) :
\begin{enumerate}
    \item \textbf{Si $P$ est extérieur au cercle ($|PO| > r$) :} Pour deux sécantes $PAB$ et $PA'B'$, et pour une tangente $(PT)$ issue de $P$ (où $T$ est le point de tangence) :
    \[
        \boxed{|PA| \cdot |PB| = |PA'| \cdot |PB'| = |PT|^2 = |PO|^2 - r^2}.
    \]
    \item \textbf{Si $P$ est intérieur au cercle ($|PO| < r$) :} Pour deux cordes $[AB]$ et $[A'B']$ qui se croisent en $P$ :
    \[
        \boxed{|PA| \cdot |PB| = |PA'| \cdot |PB'| = r^2 - |PO|^2}.
    \]
\end{enumerate}
\end{theoreme}

\begin{center}
\begin{tikzpicture}[scale=1.1]
    % Figure 1 : Point extérieur
    \begin{scope}[shift={(0,0)}]
        \draw[thick, uclblue] (0,0) circle (1.4);
        \coordinate (O) at (0,0);
        \coordinate (P) at (-3.0,0);
        \coordinate (A) at (152:1.4);
        \coordinate (B) at (28:1.4);
        % recalc B on line PA:
        % Let's place A and B cleanly on a secant from P(-3,0):
        % Line y = 0.38*(x+3) -> at circle x^2 + y^2 = 1.96
        \coordinate (A) at (-1.28, 0.56);
        \coordinate (B) at (0.86, 1.10);
        \coordinate (Ap) at (-1.15, -0.80);
        \coordinate (Bp) at (1.15, -0.80);
        \coordinate (T) at (-0.653, 1.238);

        \draw[thick, uclaccent] (P) -- (B);
        \draw[thick, uclaccent] (P) -- (Bp);
        \draw[thick, warnred] (P) -- (T);
        \draw[dashed, gray] (A) -- (Bp) (Ap) -- (B);
        \draw[dashed, warnred] (O) -- (T);
        \draw[dotted] (P) -- (O);

        \fill (O) circle (1.3pt) node[above right] {\small $O$};
        \fill (P) circle (1.3pt) node[left] {\small $P$};
        \fill (A) circle (1.3pt) node[above left] {\small $A$};
        \fill (B) circle (1.3pt) node[above right] {\small $B$};
        \fill (Ap) circle (1.3pt) node[below left] {\small $A'$};
        \fill (Bp) circle (1.3pt) node[below right] {\small $B'$};
        \fill (T) circle (1.3pt) node[above] {\small $T$};
        \node[below, align=center] at (-0.8,-1.65) {\small \textbf{(a)} Point extérieur :\\[-2pt]\small $|PA|\cdot|PB| = |PT|^2 = |PO|^2 - r^2$};
    \end{scope}

    % Figure 2 : Point intérieur
    \begin{scope}[shift={(5.0,0)}]
        \draw[thick, uclblue] (0,0) circle (1.4);
        \coordinate (O) at (0,0);
        \coordinate (A) at (150:1.4);
        \coordinate (B) at (-20:1.4);
        \coordinate (Ap) at (-125:1.4);
        \coordinate (Bp) at (55:1.4);
        \coordinate (P) at (intersection of A--B and Ap--Bp);

        \draw[thick, uclaccent] (A) -- (B);
        \draw[thick, warnred] (Ap) -- (Bp);
        \draw[dashed, gray] (A) -- (Ap) (B) -- (Bp);
        \draw[dotted] (O) -- (P);

        \fill (O) circle (1.3pt) node[below right] {\small $O$};
        \fill (P) circle (1.3pt) node[above left] {\small $P$};
        \fill (A) circle (1.3pt) node[above left] {\small $A$};
        \fill (B) circle (1.3pt) node[below right] {\small $B$};
        \fill (Ap) circle (1.3pt) node[below left] {\small $A'$};
        \fill (Bp) circle (1.3pt) node[above right] {\small $B'$};
        \node[below, align=center] at (0,-1.65) {\small \textbf{(b)} Point intérieur :\\[-2pt]\small $|PA|\cdot|PB| = |PA'|\cdot|PB'| = r^2 - |PO|^2$};
    \end{scope}
\end{tikzpicture}
\end{center}

\begin{proof}[Démonstration (UCLouvain -- Juillet 2025 Q3 et Septembre 2025 Q3)]
\mbox{}
\begin{enumerate}
    \item \textbf{Cas extérieur ($|PO| > r$) :} Considérons deux sécantes $(PAB)$ et $(PA'B')$. Les triangles $PAB'$ et $PA'B$ possèdent l'angle $\widehat{P}$ en commun, et leurs angles inscrits $\widehat{AB'P}$ et $\widehat{A'BP}$ interceptent le même arc $\wideparen{AA'}$, donc $|\widehat{AB'P}| = |\widehat{A'BP}|$. Les triangles $PAB'$ et $PA'B$ sont donc \textbf{semblables}, ce qui entraîne :
    \[
        \frac{|PA|}{|PA'|} = \frac{|PB'|}{|PB|} \iff |PA| \cdot |PB| = |PA'| \cdot |PB'|.
    \]
    En choisissant comme sécante particulière la droite $(PO)$ qui coupe le cercle en $A''$ (du côté de $P$) et $B''$ (de l'autre côté), on a $|PA''| = |PO| - r$ et $|PB''| = |PO| + r$, d'où :
    \[
        |PA| \cdot |PB| = |PA''| \cdot |PB''| = (|PO| - r)(|PO| + r) = |PO|^2 - r^2.
    \]
    Enfin, pour la tangente $(PT)$ en $T$, le triangle $POT$ est rectangle en $T$, donc par Pythagore $|PT|^2 = |PO|^2 - |OT|^2 = |PO|^2 - r^2$.
    \item \textbf{Cas intérieur ($|PO| < r$) :} Pour deux cordes $[AB]$ et $[A'B']$ sécantes en $P$, les triangles $PAA'$ et $PB'B$ ont leurs angles en $P$ opposés par le sommet ($|\widehat{APA'}| = |\widehat{B'PB}|$) et leurs angles inscrits $\widehat{PAA'}$ et $\widehat{PB'B}$ interceptent le même arc $\wideparen{A'B}$, donc $|\widehat{PAA'}| = |\widehat{PB'B}|$. Par similitude, $\frac{|PA|}{|PB'|} = \frac{|PA'|}{|PB|} \iff |PA|\cdot|PB| = |PA'|\cdot|PB'|$. Sur le diamètre passant par $P$ et $O$, les distances aux deux extrémités valent $r - |PO|$ et $r + |PO|$, dont le produit vaut $(r - |PO|)(r + |PO|) = r^2 - |PO|^2$.
\end{enumerate}
\end{proof}

\begin{exemple}[Application de la puissance d'un point : limite géométrique (UCLouvain -- Juillet 2025, Question 4)]
\emph{Énoncé :} L'hypoténuse $[AB]$ d'un triangle $ABO$ rectangle isocèle en $O$ (avec $|OA| = |OB| = 1$) est divisée en $2n+1$ petits segments de même longueur $[P_i P_{i+1}]$ pour $i = 0, \dots, 2n$ (où $A = P_0$ et $B = P_{2n+1}$). On considère le cercle de rayon $r_n$ tangent aux deux côtés de l'angle droit $[OA]$ et $[OB]$ qui passe par les deux extrémités $P_n$ et $P_{n+1}$ du segment central de l'hypoténuse. (a) Calculer $r_n$ en fonction de $n$. (b) En déduire $r = \lim_{n\to\infty} r_n$.\\
\emph{Résolution :}
\begin{enumerate}[label=(\alph*)]
    \item Soit $Q$ le point de tangence du cercle avec le côté $[OA]$. Comme le cercle de rayon $r_n$ est tangent aux deux côtés perpendiculaires $[OA]$ et $[OB]$, son centre est à coordonnées $(r_n, r_n)$ et le point de tangence $Q$ est à distance $|OQ| = r_n$ de l'origine $O$. La distance entre le sommet $A$ et le point de tangence $Q$ vaut donc $|AQ| = |OA| - |OQ| = 1 - r_n$.\\
    Or la droite $(AB)$ est une sécante issue de $A$ qui coupe le cercle aux points $P_n$ et $P_{n+1}$, tandis que $(AQ)$ est tangente au cercle en $Q$ ! Par le théorème de la \textbf{puissance du point $A$ par rapport au cercle} :
    \[
        |AQ|^2 = |AP_n| \cdot |AP_{n+1}|.
    \]
    L'hypoténuse $[AB]$ a pour longueur $\sqrt{1^2+1^2} = \sqrt{2}$, donc chaque petit segment mesure $\frac{\sqrt{2}}{2n+1}$. Ainsi $|AP_n| = \frac{n\sqrt{2}}{2n+1}$ et $|AP_{n+1}| = \frac{(n+1)\sqrt{2}}{2n+1}$, ce qui donne :
    \[
        (1 - r_n)^2 = \frac{n\sqrt{2}}{2n+1} \cdot \frac{(n+1)\sqrt{2}}{2n+1} = \frac{2n(n+1)}{(2n+1)^2}.
    \]
    Comme $0 \le r_n \le 1$, on prend la racine carrée positive : $r_n = \boxed{1 - \frac{\sqrt{2n(n+1)}}{2n+1}}$.
    \item Lorsque $n \to \infty$, le segment central $[P_n P_{n+1}]$ se réduit au point milieu $P$ de l'hypoténuse $[AB]$ (à distance $\frac{\sqrt{2}}{2}$ de $A$) et le cercle devient tangent à l'hypoténuse en $P$ (cercle inscrit!). Par égalité des deux tangentes issues de $A$, $|AQ| = |AP| \iff 1 - r = \frac{\sqrt{2}}{2}$, soit $r = \boxed{1 - \frac{\sqrt{2}}{2}}$.
\end{enumerate}
\end{exemple}

\subsection{Relations métriques dans le triangle rectangle}

Lorsque l'on abaisse la hauteur $[CH]$ relative à l'hypoténuse $[AB]$ d'un triangle $ABC$ rectangle en $C$, on découpe le triangle rectangle en deux petits triangles rectangles $ACH$ et $CBH$ qui sont tous deux \textbf{semblables} au grand triangle $ABC$.

\begin{theoreme}[Relations métriques dans le triangle rectangle]
Soit $ABC$ un triangle rectangle en $C$, et soit $H$ le pied de la hauteur issue de $C$ sur l'hypoténuse $[AB]$. Notons $c = |AB|$, $a = |BC|$, $b = |AC|$, $h = |CH|$, et $p = |AH|$, $q = |BH|$ les projections orthogonales des deux cathètes sur l'hypoténuse ($p + q = c$). Alors :
\begin{enumerate}
    \item \textbf{Théorème de la hauteur :} Le carré de la hauteur relative à l'hypoténuse est égal au produit des deux segments qu'elle détermine sur l'hypoténuse :
    \[
        \boxed{h^2 = p \cdot q \qquad \left(\text{c'est-à-dire } |CH|^2 = |AH| \cdot |HB|\right)}.
    \]
    \item \textbf{Théorème des cathètes (d'Euclide) :} Le carré de chaque côté de l'angle droit est égal au produit de l'hypoténuse entière par sa projection orthogonale sur l'hypoténuse :
    \[
        \boxed{b^2 = c \cdot p \quad \text{et} \quad a^2 = c \cdot q}.
    \]
    \item \textbf{Produit des cathètes :} En exprimant le double de l'aire du triangle de deux façons, on a $a \cdot b = c \cdot h$.
\end{enumerate}
\end{theoreme}

\begin{center}
\begin{tikzpicture}[scale=1.1]
    \coordinate (A) at (0,0);
    \coordinate (B) at (5,0);
    \coordinate (C) at (1.8, 2.4);
    \coordinate (H) at (1.8, 0);

    \draw[thick, uclblue] (A) -- (B) -- (C) -- cycle;
    \draw[thick, warnred] (C) -- (H) node[midway, right] {\small $h$};
    \draw[dashed, gray] (2.5,0) arc (0:180:2.5);
    \pic[draw, angle radius=0.22cm] {right angle = A--C--B};
    \pic[draw=warnred, angle radius=0.22cm] {right angle = C--H--B};

    \node[below] at (0.9, 0) {\small $p = |AH|$};
    \node[below] at (3.4, 0) {\small $q = |HB|$};
    \node[above left] at (0.9, 1.2) {\small $b$};
    \node[above right] at (3.4, 1.2) {\small $a$};

    \fill (A) circle (1.3pt) node[below left] {\small $A$};
    \fill (B) circle (1.3pt) node[below right] {\small $B$};
    \fill (C) circle (1.3pt) node[above] {\small $C$};
    \fill (H) circle (1.3pt) node[below] {\small $H$};
\end{tikzpicture}
\end{center}

\begin{exemple}[Un problème géométrique de Fermat (UCLouvain -- Juillet 2022, Question 3)]
\emph{Énoncé :} Soit $ABCD$ un rectangle de côtés $L = |AB| = |DC|$ et $\ell = |BC| = |AD|$ tels que $L^2 = 2\ell^2$. À l'extérieur du rectangle, on trace un demi-cercle de diamètre $[AB]$ sur lequel on choisit un point quelconque $P$. Les segments $[PC]$ et $[PD]$ coupent $[AB]$ respectivement en $X$ et $Y$. Enfin, la perpendiculaire à $[DC]$ issue de $P$ coupe $[AB]$ en $Q$ et $[DC]$ en $R$. On pose $h = |PQ|$.
\begin{enumerate}[label=(\alph*)]
    \item Exprimer $|XY|^2$ en fonction de $\ell$ et $h$.
    \item Exprimer le produit $|DR| \cdot |RC|$ en fonction de $h$.
    \item Démontrer que $|XY|^2 = 2\,|AY|\cdot|BX|$ (ce qui équivaut à $|AB|^2 = |AX|^2 + |BY|^2$, indépendamment de la position de $P$ sur le demi-cercle!).
\end{enumerate}
\emph{Résolution :}
\begin{enumerate}[label=(\alph*)]
    \item Comme $(AB) \parallel (DC)$, le triangle $PXY$ (de hauteur $h = |PQ|$) est semblable au triangle $PDC$ (de base $|DC| = L$ et de hauteur $|PR| = |PQ| + |QR| = h + \ell$). Par Thalès :
    \[
        \frac{|XY|}{|DC|} = \frac{h}{h+\ell} \implies |XY|^2 = \frac{h^2 L^2}{(h+\ell)^2} = \boxed{\frac{2h^2\ell^2}{(h+\ell)^2}} \quad (\text{car } L^2 = 2\ell^2).
    \]
    \item Le point $P$ appartient au demi-cercle de diamètre $[AB]$, donc le triangle $APB$ est \textbf{rectangle en $P$}, et $[PQ]$ est la hauteur relative à l'hypoténuse $[AB]$. Par le théorème de la hauteur :
    \[
        h^2 = |AQ| \cdot |QB|.
    \]
    Or $AQRD$ et $QRCB$ sont des rectangles, donc $|DR| = |AQ|$ et $|RC| = |QB|$. Il vient immédiatement $|DR| \cdot |RC| = \boxed{h^2}$.
    \item Par Thalès dans le triangle $PRD$ (où $(AY) \parallel (DR)$) et dans le triangle $PRC$ (où $(BX) \parallel (RC)$) :
    \[
        \frac{|AY|}{|DR|} = \frac{\ell}{h+\ell} \implies |AY| = \frac{\ell\,|DR|}{h+\ell}, \qquad \text{et} \qquad \frac{|BX|}{|RC|} = \frac{\ell}{h+\ell} \implies |BX| = \frac{\ell\,|RC|}{h+\ell}.
    \]
    En multipliant ces deux égalités et en utilisant $|DR|\cdot|RC| = h^2$ ainsi que (a) :
    \[
        2\,|AY|\cdot|BX| = \frac{2\ell^2\,|DR|\cdot|RC|}{(h+\ell)^2} = \frac{2\ell^2 h^2}{(h+\ell)^2} = |XY|^2.
    \]
\end{enumerate}
\end{exemple}

% ----------------------------------------------------------------------------
\section{Similitudes, rapports d'aires, polygones et découpages géométriques}
% ----------------------------------------------------------------------------

\subsection{Rapports de similitude, rapports d'aires et propriété du parallélogramme}

\begin{theoreme}[Principes fondamentaux pour le calcul de longueurs et d'aires]
\begin{enumerate}
    \item \textbf{Rapport des aires de deux figures semblables :} Si deux triangles (ou polygones) sont semblables avec un rapport de similitude $\rho = \frac{L_2}{L_1}$ sur les longueurs homologues, alors le rapport de leurs \textbf{aires} est égal au \textbf{carré} du rapport de similitude :
    \[
        \boxed{\frac{\text{Aire}(\triangle_2)}{\text{Aire}(\triangle_1)} = \rho^2 = \left(\frac{L_2}{L_1}\right)^2}.
    \]
    \item \textbf{Triangles partageant la même hauteur (ou la même base) :} Si deux triangles ont un sommet commun et leurs bases situées sur une même droite, ils possèdent la \textbf{même hauteur} $h$. Le rapport de leurs aires est alors simplement égal au rapport de leurs bases :
    \[
        \boxed{\text{Même hauteur } h \implies \frac{\text{Aire}(\triangle_1)}{\text{Aire}(\triangle_2)} = \frac{\text{base}_1}{\text{base}_2}}.
    \]
    En particulier, une \textbf{médiane} d'un triangle partage toujours celui-ci en deux triangles d'aires égales !
    \item \textbf{Identité du parallélogramme (UCLouvain -- Juillet 2024, Question 4a) :} Dans tout parallélogramme $ABCD$ de côtés $a = |AB| = |CD|$ et $b = |BC| = |DA|$, la somme des carrés des longueurs des deux diagonales est égale à la somme des carrés des longueurs des quatre côtés :
    \[
        \boxed{|AC|^2 + |BD|^2 = 2\bigl(|AB|^2 + |BC|^2\bigr) = 2(a^2 + b^2)}.
    \]
    \emph{Démonstration :} Par la loi des cosinus dans $ABC$ (d'angle $\beta$) et $BAD$ (d'angle $\alpha = \pi - \beta$, d'où $\cos\alpha = -\cos\beta$), on a $|BD|^2 = a^2 + b^2 - 2ab\cos\alpha$ et $|AC|^2 = a^2 + b^2 - 2ab\cos\beta$. En additionnant, les termes en $\cos\alpha + \cos\beta = 0$ s'annulent.
\end{enumerate}
\end{theoreme}

\begin{exemple}[Cinq problèmes d'examen UCLouvain sur les similitudes et les aires]
\textbf{1) Trois carrés inscrits dans un triangle (UCLouvain -- Septembre 2026, Question 4) :}\\
\emph{Énoncé :} Dans un triangle $ABC$, trois carrés $M_1 R_2 P_2 P_1$ (d'aire $4\text{ cm}^2$, donc de côté $2\text{ cm}$), $M_2 N_2 Q_2 P_2$ (d'aire $36\text{ cm}^2$, donc de côté $6\text{ cm}$) et $S_2 N_1 Q_1 Q_2$ (d'aire $9\text{ cm}^2$, donc de côté $3\text{ cm}$) sont posés côte à côte sur la base $[BC]$, avec $M_1, M_2 \in [AC]$ et $N_1, N_2 \in [AB]$. Calculer : (a) la longueur de la base $|BC|$, (b) la hauteur $|AO|$ du petit triangle supérieur $AM_2 N_2$, et (c) l'aire totale du triangle $ABC$.
\begin{center}
\begin{tikzpicture}[scale=0.45]
    \coordinate (C) at (0,0);
    \coordinate (P1) at (1,0);
    \coordinate (P2) at (3,0);
    \coordinate (Q2) at (9,0);
    \coordinate (Q1) at (12,0);
    \coordinate (B) at (15,0);
    \coordinate (A) at (5,10);

    \draw[fill=ucllightblue, thick, uclblue] (P1) rectangle (3,2);
    \draw[fill=deflightgreen, thick, defgreen] (P2) rectangle (9,6);
    \draw[fill=exlightorange, thick, exorange] (Q2) rectangle (12,3);
    \draw[thick, uclblue] (C) -- (B) -- (A) -- cycle;
    \draw[dashed, warnred] (A) -- (5,6);
    \fill[warnred] (5,6) circle (3pt) node[below right] {\small $O$};

    \node at (2,1) {\small $2\times 2$};
    \node at (6,3) {\small $6\times 6$};
    \node at (10.5,1.5) {\small $3\times 3$};
    \node[below] at (C) {\small $C$};
    \node[below] at (P1) {\small $P_1$};
    \node[below] at (P2) {\small $P_2$};
    \node[below] at (Q2) {\small $Q_2$};
    \node[below] at (Q1) {\small $Q_1$};
    \node[below] at (B) {\small $B$};
    \node[above] at (A) {\small $A$};
    \node[above left] at (1,2) {\small $M_1$};
    \node[above left] at (3,6) {\small $M_2$};
    \node[above right] at (9,6) {\small $N_2$};
    \node[above right] at (12,3) {\small $N_1$};
\end{tikzpicture}
\end{center}
\emph{Résolution :}
\begin{enumerate}[label=(\alph*)]
    \item Les côtés des carrés étant parallèles et perpendiculaires à $(BC)$, les triangles rectangles $C P_1 M_1$ et $M_1 R_2 M_2$ sont \textbf{semblables} (angles correspondants égaux). Comme $|P_1 M_1| = |M_1 R_2| = 2\text{ cm}$ et $|R_2 M_2| = 6 - 2 = 4\text{ cm}$, on a :
    \[
        \frac{|CP_1|}{|P_1 M_1|} = \frac{|M_1 R_2|}{|R_2 M_2|} \implies |CP_1| = \frac{2 \times 2}{4} = 1\text{ cm}.
    \]
    De même, à droite, les triangles rectangles $Q_1 B N_1$ et $S_2 N_1 N_2$ sont semblables, avec $|Q_1 N_1| = |S_2 N_1| = 3\text{ cm}$ et $|S_2 N_2| = 6 - 3 = 3\text{ cm}$, d'où :
    \[
        \frac{|Q_1 B|}{|Q_1 N_1|} = \frac{|S_2 N_1|}{|S_2 N_2|} \implies |Q_1 B| = \frac{3 \times 3}{3} = 3\text{ cm}.
    \]
    La base totale vaut donc $|BC| = |CP_1| + |P_1 P_2| + |P_2 Q_2| + |Q_2 Q_1| + |Q_1 B| = 1 + 2 + 6 + 3 + 3 = \boxed{15\text{ cm}}$.
    \item Soit $O$ le pied de la hauteur issue de $A$ sur le côté supérieur $[M_2 N_2]$ du grand carré (de longueur $6\text{ cm}$), et posons $x = |M_2 O|$ (de sorte que $|O N_2| = 6 - x$).
    \begin{itemize}
        \item Par similitude entre $M_2 O A$ et $M_1 R_2 M_2$ : $\frac{|AO|}{|M_2 O|} = \frac{|R_2 M_2|}{|M_1 R_2|} = \frac{4}{2} = 2 \implies |AO| = 2x$.
        \item Par similitude entre $A O N_2$ et $N_2 S_2 N_1$ : $\frac{|AO|}{|O N_2|} = \frac{|N_2 S_2|}{|S_2 N_1|} = \frac{3}{3} = 1 \implies |AO| = 6 - x$.
    \end{itemize}
    En égalant ces deux expressions : $2x = 6 - x \implies x = 2\text{ cm}$, d'où $|AO| = 4\text{ cm}$.
    \item La hauteur totale du triangle $ABC$ vaut $H = |AO| + 6 = 4 + 6 = 10\text{ cm}$, et son aire vaut :
    \[
        \text{Aire}(\triangle ABC) = \frac{1}{2} \times |BC| \times H = \frac{1}{2} \times 15 \times 10 = \boxed{75\text{ cm}^2}.
    \]
\end{enumerate}

\medskip
\textbf{2) Le « triangle arrondi » et égalité d'aires (UCLouvain -- Septembre 2025, Question 4) :}\\
\emph{Énoncé :} Dans un demi-cercle de centre $O$, de rayon $r = 4\text{ cm}$ et de diamètre $[AB]$, on trace une corde $[AQ]$ telle que $|\widehat{BAQ}| = \frac{\pi}{6}\text{ rad}$. Par le milieu $P$ de $[AQ]$, on trace le segment $[PB]$. Calculer l'aire $S$ de la région délimitée par les segments $[PQ]$, $[PB]$ et l'arc $\wideparen{QB}$.\\
\emph{Résolution :}
\begin{enumerate}
    \item Comparons l'aire des triangles $PQB$ et $OQB$ à celle du grand triangle $AQB$ :
    \begin{itemize}
        \item Les triangles $PQB$ et $AQB$ ont le sommet $B$ en commun et leurs bases $[PQ]$ et $[AQ]$ alignées avec $|PQ| = \frac{1}{2}|AQ|$ ($P$ milieu de $[AQ]$), donc $\text{Aire}(\triangle PQB) = \frac{1}{2}\text{Aire}(\triangle AQB)$.
        \item Les triangles $OQB$ et $AQB$ ont le sommet $Q$ en commun et leurs bases $[OB]$ et $[AB]$ alignées avec $|OB| = \frac{1}{2}|AB|$ ($O$ milieu du diamètre $[AB]$), donc $\text{Aire}(\triangle OQB) = \frac{1}{2}\text{Aire}(\triangle AQB)$.
    \end{itemize}
    Par conséquent, $\text{Aire}(\triangle PQB) = \text{Aire}(\triangle OQB)$!
    \item Comme la région cherchée d'aire $S$ est la réunion du triangle $PQB$ et du segment circulaire sous l'arc $\wideparen{QB}$, tandis que le secteur circulaire $OQB$ est la réunion du triangle $OQB$ et de ce même segment circulaire, on a exactement :
    \[
        S = \text{Aire}(\text{secteur } OQB) = \frac{1}{2}r^2 \cdot |\widehat{BOQ}|.
    \]
    Or l'angle au centre $\widehat{BOQ}$ intercepte le même arc $\wideparen{BQ}$ que l'angle inscrit $\widehat{BAQ} = \frac{\pi}{6}$, donc $|\widehat{BOQ}| = 2 \cdot \frac{\pi}{6} = \frac{\pi}{3}\text{ rad}$. D'où $S = \frac{1}{2} \cdot 4^2 \cdot \frac{\pi}{3} = \boxed{\frac{8\pi}{3}\text{ cm}^2}$.
\end{enumerate}

\medskip
\textbf{3) La lunule d'Hippocrate (UCLouvain -- Septembre 2024, Question 4) :}\\
\emph{Énoncé :} Soit $ABO$ un triangle rectangle isocèle en $O$ d'aire $T = 8\text{ cm}^2$. On construit la lunule délimitée par l'arc $\wideparen{AB}$ du cercle de centre $O$ passant par $A$ et $B$, et le demi-cercle de diamètre $[AB]$ extérieur à $O$. Calculer l'aire $L$ de cette lunule.\\
\emph{Résolution :} Notons $R = |OA| = |OB|$ le rayon du grand cercle de centre $O$, et $r = \frac{|AB|}{2}$ le rayon du demi-cercle de diamètre $[AB]$. Par Pythagore dans $ABO$, $|AB|^2 = 2R^2 \implies (2r)^2 = 2R^2 \implies r^2 = \frac{R^2}{2}$.\\
L'aire du demi-disque de diamètre $[AB]$ vaut $D = \frac{1}{2}\pi r^2 = \frac{1}{4}\pi R^2$, ce qui est \textbf{exactement égal} à l'aire $Q = \frac{1}{4}\pi R^2$ du quart de disque de centre $O$ et de rayon $R$ ! Dès lors :
\[
    L = (T + D) - Q = T + (D - Q) = T = \boxed{8\text{ cm}^2}.
\]

\medskip
\textbf{4) Ligne brisée orthogonale dans un rectangle (UCLouvain -- Septembre 2023, Question 4 et Juillet 2023, Question 4a) :}\\
\emph{Énoncé (Septembre 2023, Q4) :} Dans un rectangle $ABCD$ de largeur $\ell = |AD| = 15$, une ligne brisée $A M N O P C$ reliant les sommets opposés $A$ et $C$ possède uniquement des angles droits (alternés), avec $|AM| = 10$, $|MN| = 3$, $|NO| = 8$, $|OP| = 4$ et $|PC| = 6$. Calculer la longueur $L = |DC|$ du rectangle.\\
\emph{Résolution :} Les segments $[AM]$, $[NO]$ et $[PC]$ sont tous trois parallèles entre eux (car perpendiculaires à $(MN)$ et $(OP)$), et les segments $[MN]$ et $[OP]$ leur sont perpendiculaires. En projetant toute la ligne brisée sur ces deux directions orthogonales, la diagonale $[AC]$ du rectangle $ABCD$ est aussi l'hypoténuse d'un grand triangle rectangle dont les deux cathètes mesurent :
\[
    X = |AM| + |NO| + |PC| = 10 + 8 + 6 = 24 \qquad \text{et} \qquad Y = |MN| + |OP| = 3 + 4 = 7.
\]
Par Pythagore, le carré de la diagonale vaut $|AC|^2 = 24^2 + 7^2 = 576 + 49 = 625 = 25^2$. Or, dans le rectangle $ABCD$, $|AC|^2 = \ell^2 + L^2$, d'où $L = \sqrt{25^2 - 15^2} = \sqrt{625 - 225} = \sqrt{400} = \boxed{20}$.
\end{exemple}

% ----------------------------------------------------------------------------
\section{Géométrie synthétique dans l'espace : solides et coupes planes}
% ----------------------------------------------------------------------------

Les problèmes de géométrie synthétique dans l'espace portent principalement sur des assemblages de \textbf{sphères, cônes, cylindres ou pyramides}. La règle méthodologique essentielle consiste toujours à \textbf{se ramener à une figure plane} en effectuant une \textbf{coupe bien choisie} (section par un plan de symétrie ou un plan contenant l'axe de révolution, appelée \emph{coupe méridienne}).

\begin{theoreme}[Formulaire des volumes et aires des solides classiques]
\begin{enumerate}
    \item \textbf{Cylindre droit} de rayon de base $R$ et de hauteur $h$ :
    \[
        V_{\text{cylindre}} = \pi R^2 h, \qquad A_{\text{latérale}} = 2\pi R h, \qquad A_{\text{totale}} = 2\pi R h + 2\pi R^2.
    \]
    \item \textbf{Cône droit de révolution} de rayon de base $R$, de hauteur $h$ et d'apothème (génératrice) $a = \sqrt{R^2 + h^2}$ :
    \[
        \boxed{V_{\text{cône}} = \frac{1}{3}\pi R^2 h}, \qquad A_{\text{latérale}} = \pi R a = \pi R\sqrt{R^2+h^2}.
    \]
    \item \textbf{Pyramide} quelconque d'aire de base $\mathcal{B}$ et de hauteur $h$ :
    \[
        V_{\text{pyramide}} = \frac{1}{3}\,\mathcal{B}\cdot h.
    \]
    \item \textbf{Sphère et boule} de rayon $r$ :
    \[
        \boxed{V_{\text{sphère}} = \frac{4}{3}\pi r^3}, \qquad \boxed{A_{\text{sphère}} = 4\pi r^2}.
    \]
\end{enumerate}
\end{theoreme}

\begin{exemple}[Sphère inscrite dans un cône droit (UCLouvain -- Septembre 2023, Question 3)]
\emph{Énoncé :} Dans un cône droit de sommet $A$, de hauteur $h$ et de base de centre $M$ et de rayon $R$, est inscrite une sphère de centre $O$ et de rayon $r$, tangente à la base et à la surface latérale du cône.
\begin{enumerate}[label=(\alph*)]
    \item Représenter une coupe passant par l'axe de rotation $(AM)$ et exprimer le rayon $r$ de la sphère en fonction de $h$ et $R$.
    \item Calculer le rapport $\dfrac{V_{\text{cône}}}{V_{\text{sphère}}}$ lorsque $3h = 4R$.
\end{enumerate}
\begin{center}
\begin{tikzpicture}[scale=0.9]
    \coordinate (M) at (0,0);
    \coordinate (B) at (-2.25,0);
    \coordinate (C) at (2.25,0);
    \coordinate (A) at (0,3);
    % For R=2.25, h=3, a = 3.75, r = R*(a-R)/h = 2.25*1.5/3 = 1.125
    \coordinate (O) at (0,1.125);
    \coordinate (P) at ($(A)!(O)!(B)$);
    \coordinate (N) at ($(A)!(O)!(C)$);

    \draw[thick, uclblue] (B) -- (C) -- (A) -- cycle;
    \draw[thick, uclaccent, fill=ucllightblue!50] (O) circle (1.125);
    \draw[dashed, uclblue] (A) -- (M);
    \draw[thick, warnred] (O) -- (P) node[midway, above right] {\small $r$};
    \draw[thick, warnred] (O) -- (M) node[midway, right] {\small $r$};
    \pic[draw=warnred, angle radius=0.18cm] {right angle = A--P--O};
    \pic[draw, angle radius=0.18cm] {right angle = A--M--C};

    \fill (A) circle (1.3pt) node[above] {\small $A$};
    \fill (B) circle (1.3pt) node[below left] {\small $B$};
    \fill (C) circle (1.3pt) node[below right] {\small $C$};
    \fill (M) circle (1.3pt) node[below] {\small $M$};
    \fill (O) circle (1.3pt) node[left] {\small $O$};
    \fill (P) circle (1.3pt) node[above left] {\small $P$};
\end{tikzpicture}
\end{center}
\emph{Résolution :}
\begin{enumerate}[label=(\alph*)]
    \item La coupe méridienne du cône par un plan contenant l'axe $(AM)$ est un triangle isocèle $ABC$ de base $|BC| = 2R$ (milieu $M$) et de hauteur $|AM| = h$, dans lequel la sphère apparaît comme le \textbf{cercle inscrit} de centre $O$ et de rayon $r$. Soit $P$ le point de tangence du cercle avec la génératrice $[AB]$.
    \begin{itemize}
        \item Dans le triangle $AMB$ rectangle en $M$, l'apothème vaut $|AB| = \sqrt{h^2 + R^2}$.
        \item Les segments $[BM]$ et $[BP]$ sont tous deux tangents au cercle issus du point $B$, donc $|BP| = |BM| = R$, et par suite $|AP| = |AB| - |BP| = \sqrt{h^2+R^2} - R$ (et $|AO| = h - r$).
        \item Les triangles rectangles $APO$ (rectangle en $P$) et $AMB$ (rectangle en $M$) partagent l'angle aigu en $A$ : ils sont donc \textbf{semblables} ! Dès lors :
        \[
            \frac{|OP|}{|AP|} = \frac{|BM|}{|AM|} \iff \frac{r}{\sqrt{h^2+R^2}-R} = \frac{R}{h} \implies r = \boxed{\frac{R\bigl(\sqrt{h^2+R^2}-R\bigr)}{h}}.
        \]
    \end{itemize}
    \item Lorsque $3h = 4R$, posons $h = 4a$ et $R = 3a$ (triangle pythagoricien $3-4-5$ : $\sqrt{h^2+R^2} = 5a$). La formule (a) donne $r = \frac{3a(5a-3a)}{4a} = \frac{3}{2}a$. Le rapport des volumes vaut alors :
    \[
        \frac{V_{\text{cône}}}{V_{\text{sphère}}} = \frac{\frac{1}{3}\pi R^2 h}{\frac{4}{3}\pi r^3} = \frac{R^2 h}{4r^3} = \frac{(3a)^2(4a)}{4\left(\frac{3}{2}a\right)^3} = \frac{36a^3}{\frac{27}{2}a^3} = \boxed{\frac{8}{3}}.
    \]
\end{enumerate}
\end{exemple}

% ----------------------------------------------------------------------------
\section{Exercices du Chapitre 5}
% ----------------------------------------------------------------------------

\emph{Les corrigés complets de tous ces exercices se trouvent au Chapitre 9, en fin d'ouvrage.}

\begin{enumerate}[label=\textbf{Exercice 5.\arabic* --}, align=left, leftmargin=*]

    \item \textbf{(UCLouvain -- Juillet 2026, Q3 : Carré, arc de cercle et milieu d'un segment)}
    Dans un carré $ABCD$ de côté $c$, on trace le quart de cercle de centre $A$ passant par $B$ et $D$. On trace ensuite un segment $[DY]$ où $Y \in [BC]$, qui coupe l'arc de cercle en un point $X$. Sachant que $X$ est exactement le milieu du segment $[DY]$, calculez l'amplitude (en degrés) de l'angle $\widehat{ADY}$. \emph{(Indice : tracez la parallèle à $(AD)$ passant par $X$.)}

    \item \textbf{(UCLouvain -- Juillet 2024, Q3 : Tangence inscrite dans un demi-cercle)}
    Un cercle $\mathcal{C}$ de centre $C$ et de rayon $r$ est inscrit dans un demi-cercle $\mathcal{D}$ de centre $O$, de diamètre $[AB]$ et de rayon $R$ (de sorte que $\mathcal{C}$ est tangent au diamètre $[AB]$ en un point $S$ et à l'arc du demi-cercle en un point $T$). On construit les deux segments tangents à $\mathcal{C}$ perpendiculaires à $[AB]$, qui coupent $[AB]$ en $P$ (du côté de $A$) et $Q$ (du côté de $B$). Sachant que $|AP| = 2$ et $|QB| = 8$ :
    \begin{enumerate}[label=(\alph*)]
        \item Exprimez le rayon $R$ du demi-cercle en fonction du rayon $r$ du cercle $\mathcal{C}$.
        \item En utilisant le triangle rectangle $OSC$, déterminez la valeur exacte de $r$ puis celle de $R$.
    \end{enumerate}

    \item \textbf{(UCLouvain -- Juillet 2023, Q3 : La surface cachée de l'hexagone)}
    Un hexagone régulier $ABCDEF$ de côté $c$ est inscrit dans un rectangle $MNPQ$ de telle sorte que $A \in [MQ]$, $B \in [MN]$ et $D \in [NP]$.
    \begin{enumerate}[label=(\alph*)]
        \item Démontrez que l'angle $\widehat{ABD}$ est droit ($90\degree$) et que les triangles rectangles $MAB$ et $NBD$ sont semblables.
        \item Calculez le rapport $r = \dfrac{\text{Aire}(\triangle NBD)}{\text{Aire}(\triangle MAB)}$.
        \item Sachant que $\text{Aire}(\triangle MAB) = 20$ et que l'aire du quadrilatère concave $NBCD$ vaut $24$, déduisez-en l'aire totale de l'hexagone régulier $ABCDEF$.
    \end{enumerate}

    \item \textbf{(UCLouvain -- Septembre 2024, Q3 : Rapport d'aires paramétré)}
    Dans un carré $ABCD$ de côté $c$, le point $N$ du côté $[AD]$ est situé à distance $|AN| = \lambda c$ du sommet $A$ (avec $0 < \lambda < 1$). On construit un rectangle $BPDM$ (avec $M$ sur la droite $(BN)$ au-delà de $N$, et $P$ de l'autre côté) de telle sorte que le segment $[BN]$ soit inclus dans le côté $[BM]$ du rectangle et que le côté opposé $[DP]$ lui soit parallèle et coupe $[BC]$ en un point $O$.
    \begin{enumerate}[label=(\alph*)]
        \item Calculez l'aire $\mathcal{Q}$ du parallélogramme $NBOD$ en fonction de $c$ et $\lambda$.
        \item Montrez que les triangles rectangles $ABN$ et $MDN$ sont semblables et calculez leur rapport de similitude $\rho(\lambda) = \dfrac{|MD|}{|AB|}$.
        \item Déduisez-en le rapport $R(\lambda) = \dfrac{\text{Aire}(BPDM)}{\text{Aire}(ABCD)}$ en fonction de $\lambda$ uniquement.
    \end{enumerate}

    \item \textbf{(UCLouvain -- Juillet 2024, Q4b : Carré, quarts de cercle et optimisation)}
    Dans un carré $ABCD$ de côté $c$, on trace les deux quarts de cercle de centres respectifs $B$ et $D$ reliant les sommets $A$ et $C$. Soit $P$ un point du quart de cercle centré en $B$, situé à une distance $d = |DP|$ du sommet $D$. Symétriquement, on place le point $Q$ sur le quart de cercle centré en $D$ tel que $(DP) \parallel (QB)$ et $|QB| = d$.
    \begin{enumerate}[label=(\alph*)]
        \item En appliquant l'identité du parallélogramme au quadrilatère $DQBP$, exprimez la longueur $\ell = |PQ|$ en fonction de $d$ uniquement.
        \item Déterminez la valeur minimale $m(c)$ de la longueur $\ell$ en fonction de $c$.
    \end{enumerate}

    \item \textbf{(UCLouvain -- Juillet 2022, Q4 : Sélection « Top Cheffe » et cylindres tangents)}
    On entoure quatre cylindres identiques de rayon $r$ et de longueur $\ell$ (quatre asperges tangentes deux à deux disposées en carré) à l'aide d'une pâte rectangulaire de longueur $\ell$ et de largeur $L$ qui fait exactement \emph{deux fois} le tour du grand cylindre extérieur de rayon $R$ circonscrit aux quatre petits cylindres.
    \begin{enumerate}[label=(\alph*)]
        \item Calculez le rapport $\dfrac{r}{R}$.
        \item Exprimez le volume $V$ de l'espace vide à l'intérieur du grand cylindre (entre le grand cylindre et les quatre petits cylindres) en fonction de $\ell$ et $L$.
    \end{enumerate}

    \item \textbf{(UCLouvain -- Septembre 2022, Q4 : Tangence équilatérale)}
    À l'extérieur d'un petit triangle équilatéral de côté $2a$, on trace trois demi-cercles ayant chacun un des côtés du triangle pour diamètre. Un plus grand triangle équilatéral, homothétique au premier, a ses trois côtés tangents à chaque paire de demi-cercles. Calculez en fonction de $a$ :
    \begin{enumerate}[label=(\alph*)]
        \item La longueur $L$ d'un côté du grand triangle équilatéral.
        \item L'aire totale $S$ des trois régions comprises entre chaque sommet du grand triangle et le demi-cercle le plus proche.
    \end{enumerate}

    \item \textbf{(UCLouvain -- Juillet 2021, Q1 et Q2 : Hexagone et petits calculs d'aires)}
    \begin{enumerate}[label=(\alph*)]
        \item Dans un hexagone régulier $ABCDEF$, on construit un rectangle $AGDH$ ayant la diagonale $[AD]$ de l'hexagone pour diagonale (avec $|AG| \le |AB|$). Que vaut l'angle $\alpha = \widehat{DGC}$ ?
        \item Soit un anneau compris entre deux cercles concentriques de rayons $r < R$. Une corde $[AB]$ du grand cercle, de longueur $2L$, est tangente au petit cercle. Exprimez l'aire $\mathcal{A}$ de l'anneau en fonction de $L$.
        \item Dans un carré $ABCD$ de côté $4$ et de centre $M$, on note $N$ le milieu de $[AD]$, $E$ le milieu de $[NM]$ et $F$ le milieu de $[EB]$. Calculez l'aire $\mathcal{T}$ du triangle $AEF$.
    \end{enumerate}

    \item \textbf{(UCLouvain -- Septembre 2021, Q3 et Q4 : Carrés adjacents et partage d'un carré)}
    \begin{enumerate}[label=(\alph*)]
        \item Dans un triangle équilatéral $ABC$ de côté $|AB| = 2$, on trace la hauteur $[AM]$ (où $M \in [BC]$), que l'on partage en son milieu $F$ pour construire deux petits carrés adjacents $ADEF$ et $FEGM$ (avec $G \in [MB]$). La diagonale $[DH]$ du rectangle $ADGM$ coupe le côté $[AB]$ en $H$. Calculez $|AM|$, $|AH|$ et l'angle $\alpha = \widehat{MHA}$.
        \item Dans un carré $ABCD$ de centre $M$, on place les milieux $E$ de $[AM]$ et $F$ de $[MC]$. La droite $(BE)$ coupe $[AD]$ en $G$ et la droite $(DF)$ coupe $[BC]$ en $H$. Calculez le rapport $r = \frac{\text{Aire}(\triangle AEG)}{\text{Aire}(\triangle GEM)}$, le rapport $\frac{|BH|}{|HC|}$ et la fraction $f$ de l'aire du carré $ABCD$ occupée par le parallélogramme $GBHD$.
    \end{enumerate}

    \item \textbf{(UCLouvain -- Juillet 2023, Q4a et Q4b : Carré et égalité d'aires dans un rectangle)}
    \begin{enumerate}[label=(\alph*)]
        \item Dans un carré $ABCD$, deux segments $[AX]$ et $[CY]$ sont tous deux perpendiculaires à un même segment $[XY]$ situé à l'intérieur du carré (avec $A$ et $C$ de part et d'autre de $(XY)$). Sachant que $|AX| = 13$, $|XY| = 3$ et $|YC| = 8$, calculez le côté $c$ du carré $ABCD$.
        \item Dans un rectangle de largeur $\ell$ et de longueur $L$, on trace les deux quarts de cercle de rayon $\ell$ centrés sur les deux sommets supérieurs du rectangle. Soit $A$ la région commune aux deux quarts de disque et $B$ la région du rectangle extérieure aux deux quarts de disque. Sachant que $\text{Aire}(A) = \text{Aire}(B)$, exprimez $L$ en fonction de $\ell$.
    \end{enumerate}

\end{enumerate}

% ============================================================================
% CHAPITRE 6 : GÉOMÉTRIE ANALYTIQUE PLANE : DROITES, CERCLES ET CONIQUES
% ============================================================================
\chapter[Géométrie Analytique Plane et Coniques]{Géométrie Analytique Plane : Droites, Cercles et Coniques}

La première question de la partie géométrie de l'examen d'admission à l'UCLouvain porte systématiquement sur la \textbf{géométrie analytique dans le plan} muni d'un repère orthonormé $(O;\vec{i},\vec{j})$. Ce chapitre couvre l'ensemble du programme de 5\ieme{} et 6\ieme{} secondaire : le calcul vectoriel et le produit scalaire dans le plan, les droites, les médiatrices, la distance d'un point à une droite, les cercles, ainsi que l'étude complète des \textbf{coniques} (ellipse, hyperbole, parabole) et de leurs tangentes.

% ----------------------------------------------------------------------------
\section{Vecteurs, produit scalaire, droites et distances dans le plan}
% ----------------------------------------------------------------------------

\subsection{Vecteurs, milieu, distance et produit scalaire dans $\R^2$}

Dans le plan muni d'un repère orthonormé $(O;\vec{i},\vec{j})$, soient deux points $A(x_A, y_A)$ et $B(x_B, y_B)$, et deux vecteurs $\vec{u} = (u_x, u_y)$ et $\vec{v} = (v_x, v_y)$ :

\begin{theoreme}[Formulaire vectoriel et métrique dans le plan]
\begin{enumerate}
    \item \textbf{Vecteur $\overrightarrow{AB}$, distance et point milieu $M$ de $[AB]$ :}
    \[
        \overrightarrow{AB} = (x_B - x_A,\; y_B - y_A), \qquad |AB| = \|\overrightarrow{AB}\| = \sqrt{(x_B-x_A)^2 + (y_B-y_A)^2},
    \]
    \[
        M = \left(\frac{x_A+x_B}{2},\; \frac{y_A+y_B}{2}\right).
    \]
    \item \textbf{Colinéarité (alignement ou parallélisme) :} Deux vecteurs non nuls $\vec{u} = (u_x, u_y)$ et $\vec{v} = (v_x, v_y)$ sont colinéaires ($\vec{u} \parallel \vec{v}$) ssi leur déterminant est nul :
    \[
        \vec{u} \parallel \vec{v} \iff \exists k \in \R : \vec{v} = k\vec{u} \iff \det(\vec{u},\vec{v}) = \begin{vmatrix} u_x & v_x \\ u_y & v_y \end{vmatrix} = u_x v_y - u_y v_x = 0.
    \]
    \item \textbf{Produit scalaire et angle $\theta$ entre deux vecteurs :}
    \[
        \boxed{\vec{u} \cdot \vec{v} = u_x v_x + u_y v_y = \|\vec{u}\|\,\|\vec{v}\|\cos\theta} \implies \cos\theta = \frac{u_x v_x + u_y v_y}{\sqrt{u_x^2+u_y^2}\sqrt{v_x^2+v_y^2}}.
    \]
    \item \textbf{Condition d'orthogonalité ($\vec{u} \perp \vec{v}$) :}
    \[
        \boxed{\vec{u} \perp \vec{v} \iff \vec{u} \cdot \vec{v} = u_x v_x + u_y v_y = 0}.
    \]
    \emph{Astuce immédiate dans le plan :} Pour obtenir un vecteur orthogonal à $\vec{u} = (a, b)$, il suffit d'échanger ses deux composantes et de changer un signe : $\vec{u}_\perp = (-b, a)$ (ou $(b, -a)$).
\end{enumerate}
\end{theoreme}

\subsection{Équations d'une droite, médiatrice et distance d'un point à une droite}

\begin{theoreme}[Équations d'une droite et distance point--droite dans le plan]
\begin{enumerate}
    \item \textbf{Équations paramétriques} de la droite $d$ passant par $A(x_A, y_A)$ et de vecteur directeur $\vec{d} = (u, v)$ :
    \[
        d \equiv \begin{cases} x = x_A + t\,u \\ y = y_A + t\,v \end{cases} \quad (t \in \R).
    \]
    \item \textbf{Équation cartésienne générale :}
    \[
        \boxed{d \equiv ax + by + c = 0} \qquad \text{avec} \quad \begin{cases} \vec{n} = (a, b) & \text{\textbf{vecteur normal} (perpendiculaire à $d$)} \\ \vec{d} = (-b, a) & \text{\textbf{vecteur directeur} (parallèle à $d$)} \end{cases}.
    \]
    \item \textbf{Forme réduite (pente -- ordonnée à l'origine)} lorsque $b \neq 0$ (droite non verticale) :
    \[
        y = mx + p \qquad \text{où } m = -\frac{a}{b} = \frac{y_B - y_A}{x_B - x_A} = \tan\alpha \text{ est le \textbf{coefficient directeur (pente)}}.
    \]
    \begin{itemize}
        \item Deux droites $y = m_1 x + p_1$ et $y = m_2 x + p_2$ sont \textbf{parallèles} ssi $m_1 = m_2$.
        \item Elles sont \textbf{perpendiculaires} ssi $\boxed{m_1 \cdot m_2 = -1}$ (soit $m_2 = -\frac{1}{m_1}$).
    \end{itemize}
    \item \textbf{Médiatrice $m$ d'un segment $[AB]$ :} C'est la droite passant par le milieu $M = \frac{A+B}{2}$ et ayant le vecteur $\overrightarrow{AB}$ comme \textbf{vecteur normal} (ou, de façon équivalente, le lieu des points $P(x,y)$ tels que $|PA|^2 = |PB|^2$).
    \item \textbf{Distance d'un point $P(x_0, y_0)$ à une droite $d \equiv ax + by + c = 0$ :}
    \[
        \boxed{d(P, d) = \frac{|ax_0 + by_0 + c|}{\sqrt{a^2 + b^2}}}.
    \]
\end{enumerate}
\end{theoreme}

\begin{exemple}[Calcul d'aire par distance point--droite (UCLouvain -- Septembre 2024, Question 1b)]
\emph{Énoncé :} Dans un repère orthonormé, on donne les points $A(4, 5)$ et $B(7, 2)$. Sachant que l'aire du triangle $ABC$ vaut $6$ et que l'ordonnée du point $C$ vaut $4$, déterminer toutes les valeurs possibles de l'abscisse $x$ du point $C(x, 4)$.\\
\emph{Résolution :}
\begin{enumerate}
    \item \textbf{Longueur de la base $[AB]$ et équation de la droite $(AB)$ :} On a $\overrightarrow{AB} = (7-4, 2-5) = (3, -3)$, donc la base mesure $|AB| = \sqrt{3^2 + (-3)^2} = 3\sqrt{2}$. Un vecteur normal à $(AB)$ est $\vec{n} = (1, 1)$, et comme $A(4,5) \in (AB)$ ($4+5=9$), la droite $(AB)$ a pour équation cartésienne :
    \[
        (AB) \equiv x + y - 9 = 0.
    \]
    \item \textbf{Hauteur $h_C = d(C, (AB))$ :} La distance du point $C(x, 4)$ à la droite $(AB)$ vaut :
    \[
        h_C = d\bigl(C, (AB)\bigr) = \frac{|x + 4 - 9|}{\sqrt{1^2 + 1^2}} = \frac{|x - 5|}{\sqrt{2}}.
    \]
    \item \textbf{Aire du triangle :} $\text{Aire}(\triangle ABC) = \frac{1}{2}|AB|\cdot h_C = \frac{1}{2}\cdot 3\sqrt{2} \cdot \frac{|x-5|}{\sqrt{2}} = \frac{3}{2}|x-5|$. En imposant que cette aire vaut $6$ :
    \[
        \frac{3}{2}|x - 5| = 6 \iff |x - 5| = 4 \iff x - 5 = \pm 4 \iff x = 9 \quad\text{ou}\quad x = 1.
    \]
    Il y a donc deux points possibles : $\boxed{C(1, 4) \text{ ou } C(9, 4)}$.
\end{enumerate}
\end{exemple}

% ----------------------------------------------------------------------------
\section{Le cercle en géométrie analytique}
% ----------------------------------------------------------------------------

\begin{theoreme}[Équation cartésienne d'un cercle]
Le cercle $\mathcal{C}$ de centre $\Omega(x_0, y_0)$ et de rayon $R > 0$ est le lieu des points $M(x,y)$ situés à distance $R$ de $\Omega$ :
\[
    \boxed{\mathcal{C} \equiv (x - x_0)^2 + (y - y_0)^2 = R^2}.
\]
Sous forme développée, une équation $x^2 + y^2 + \alpha x + \beta y + \gamma = 0$ se ramène à la forme canonique en \textbf{complétant les carrés} : $\left(x + \frac{\alpha}{2}\right)^2 + \left(y + \frac{\beta}{2}\right)^2 = \frac{\alpha^2+\beta^2}{4} - \gamma$ (c'est un cercle réel ssi le membre de droite est strictement positif).
\end{theoreme}

\begin{astuce}[Trois problèmes classiques sur les cercles aux examens UCLouvain]
\begin{enumerate}
    \item \textbf{Cercle circonscrit à trois points $A, B, C$ (Juillet 2024 Q1b et Septembre 2022 Q1) :} Le centre $\Omega$ du cercle passant par $A, B, C$ est équidistant des trois sommets : c'est donc le \textbf{point d'intersection des médiatrices} de deux des segments $[AB]$ et $[AC]$ ! Une fois $\Omega$ trouvé, le rayon vaut $R = |\Omega A|$.
    \item \textbf{Cercle de centre connu $\Omega$ tangent à une droite $d$ (Juillet 2021 Q3b) :} Le rayon $R$ d'un cercle centré en $\Omega(x_0, y_0)$ et tangent à une droite $d \equiv ax + by + c = 0$ est égal à la distance du centre à la droite : $R = d(\Omega, d) = \frac{|ax_0+by_0+c|}{\sqrt{a^2+b^2}}$.
    \item \textbf{Cercle de rayon connu $R$ tangent à deux droites $d_1$ et $d_2$ (Septembre 2021 Q1c) :} Le centre $\Omega(x_0, y_0)$ vérifie le système $d(\Omega, d_1) = R$ et $d(\Omega, d_2) = R$ (il se situe sur les bissectrices des deux droites).
\end{enumerate}
\end{astuce}

\begin{exemple}[Cercle passant par trois points (UCLouvain -- Juillet 2024, Question 1b)]
\emph{Énoncé :} Déterminer les coordonnées du centre $\Omega$ ainsi que l'aire du cercle $\mathcal{C}$ contenant les trois points $A(6, -4)$, $B(9, -5)$ et $C(6, -8)$.\\
\emph{Résolution :}
\begin{itemize}
    \item \textbf{Médiatrice $m_1$ de $[AC]$ :} Les points $A(6, -4)$ et $C(6, -8)$ ont la même abscisse $x = 6$ (segment vertical). Leur milieu est $M_1(6, -6)$, donc la médiatrice $m_1$ est la droite horizontale d'équation $\mathbf{y = -6}$.
    \item \textbf{Médiatrice $m_2$ de $[AB]$ :} Le milieu de $A(6, -4)$ et $B(9, -5)$ est $M_2\left(\frac{15}{2}, -\frac{9}{2}\right)$, et le vecteur normal est $\overrightarrow{AB} = (3, -1)$. L'équation cartésienne de $m_2$ est donc $3\left(x - \frac{15}{2}\right) - 1\left(y + \frac{9}{2}\right) = 0 \iff \mathbf{3x - y - 18 = 0}$.
    \item \textbf{Intersection $\Omega = m_1 \cap m_2$ :} En remplaçant $y = -6$ dans $3x - y - 18 = 0$, on obtient $3x + 6 - 18 = 0 \implies x = 4$... Attention! Vérifions le calcul : $3(x - \frac{15}{2}) - (y + \frac{9}{2}) = 3x - y - \frac{45}{2} - \frac{9}{2} = 3x - y - 27 = 0$ ! En remplaçant $y = -6$ dans $3x - y - 27 = 0$, on trouve $3x + 6 - 27 = 0 \implies 3x = 21 \implies x = 7$.\\
    Le centre du cercle est donc $\Omega = \boxed{(7, -6)}$.
    \item \textbf{Rayon et aire :} Le carré du rayon vaut $R^2 = |\Omega A|^2 = (7-6)^2 + (-6 - (-4))^2 = 1^2 + (-2)^2 = 5$. L'aire du disque vaut donc $\pi R^2 = \boxed{5\pi}$.
\end{itemize}
\end{exemple}

% ----------------------------------------------------------------------------
\section{Les coniques : l'ellipse, l'hyperbole et la parabole}
% ----------------------------------------------------------------------------

En 6\ieme{} secondaire, les \textbf{coniques} (courbes obtenues en coupant un cône de révolution par un plan, ou lieux géométriques du second degré en $x$ et $y$) sont étudiées sous deux angles complémentaires : leur \textbf{définition géométrique par foyer et directrice} (ou bifocale) et leur \textbf{équation cartésienne réduite}.

\subsection{Définition monofocale par foyer, directrice et excentricité}

\begin{definition}[Définition par foyer, directrice et excentricité]
Soit $F$ un point fixe appelé \textbf{foyer}, $d$ une droite fixe (ne contenant pas $F$) appelée \textbf{directrice}, et $e > 0$ un réel strictement positif appelé \textbf{excentricité}. Le lieu géométrique des points $M(x,y)$ du plan dont le rapport de la distance au foyer $F$ à la distance à la directrice $d$ est constant et égal à $e$ :
\[
    \boxed{\frac{d(M, F)}{d(M, d)} = e \iff d(M, F) = e \cdot d(M, d)}
\]
est une \textbf{conique} :
\begin{itemize}
    \item Si $\mathbf{0 < e < 1}$ : la conique est une \textbf{ellipse} (et si $e = 0$, c'est un cercle).
    \item Si $\mathbf{e = 1}$ : la conique est une \textbf{parabole} ($d(M, F) = d(M, d)$ : points équidistants du foyer et de la directrice).
    \item Si $\mathbf{e > 1}$ : la conique est une \textbf{hyperbole}.
\end{itemize}
\end{definition}

\begin{exemple}[Conique définie par foyer et directrice (UCLouvain -- Septembre 2024, Question 1a)]
\emph{Énoncé :} Donner l'équation cartésienne et le nom de la conique rassemblant tous les points $M(x,y)$ dont la distance à la droite $d \equiv x = 1$ est la moitié de la distance au point $F(4, 0)$.\\
\emph{Résolution :} La distance à la droite verticale $x = 1$ est $d_1 = d(M, d) = |x - 1|$, et la distance au foyer $F(4,0)$ est $d_2 = d(M, F) = \sqrt{(x-4)^2 + y^2}$. La condition $d_1 = \frac{1}{2}d_2$ s'écrit $d_2 = 2\,d_1$ (excentricité $e = 2 > 1$, c'est donc une \textbf{hyperbole} !). En élevant au carré :
\[
    (x - 4)^2 + y^2 = 4(x - 1)^2 \iff x^2 - 8x + 16 + y^2 = 4x^2 - 8x + 4 \iff 3x^2 - y^2 = 12 \iff \boxed{\frac{x^2}{4} - \frac{y^2}{12} = 1}.
\]
\end{exemple}

\subsection{Étude complète de l'ellipse, de l'hyperbole et de la parabole}

Lorsqu'une équation du second degré sans terme rectangle en $xy$ ($Ax^2 + By^2 + Cx + Dy + E = 0$) vous est donnée à l'examen, la première étape consiste toujours à \textbf{compléter les carrés} en $x$ et en $y$ pour identifier le centre (ou sommet) $(x_0, y_0)$ et ramener l'équation à l'une des trois formes réduites suivantes.

\begin{theoreme}[1. L'Ellipse d'axes parallèles aux axes du repère]
Soit l'ellipse $\mathcal{E}$ de \textbf{centre $C(x_0, y_0)$} et de demi-axes $a > 0$ (horizontal) et $b > 0$ (vertical) :
\[
    \boxed{\mathcal{E} \equiv \frac{(x - x_0)^2}{a^2} + \frac{(y - y_0)^2}{b^2} = 1}.
\]
Elle possède 4 sommets : $S_{1,2}(x_0 \pm a, y_0)$ et $S_{3,4}(x_0, y_0 \pm b)$. La position des deux foyers $F_1, F_2$ dépend de celui des deux demi-axes qui est le plus grand (\textbf{grand axe = axe focal}) :
\begin{enumerate}
    \item \textbf{Cas 1 : Grand axe horizontal ($a > b$)} (axe focal $y = y_0$, longueur du grand axe $2a$, petit axe $2b$) :
    \begin{itemize}
        \item \textbf{Distance focale $c$ :} $\boxed{c = \sqrt{a^2 - b^2}}$ \quad (relation pythagoricienne $a^2 = b^2 + c^2$).
        \item \textbf{Foyers :} $F_1(x_0 - c,\; y_0)$ et $F_2(x_0 + c,\; y_0)$. \quad (\emph{Propriété bifocale :} $|MF_1| + |MF_2| = 2a$.)
        \item \textbf{Excentricité et directrices :} $e = \dfrac{c}{a} < 1$ \quad et \quad $d_{1,2} \equiv x = x_0 \pm \dfrac{a^2}{c} = x_0 \pm \dfrac{a}{e}$.
    \end{itemize}
    \item \textbf{Cas 2 : Grand axe vertical ($b > a$)} (axe focal $x = x_0$, longueur du grand axe $2b$, petit axe $2a$) :
    \begin{itemize}
        \item \textbf{Distance focale $c$ :} $\boxed{c = \sqrt{b^2 - a^2}}$ \quad (ici $b^2 = a^2 + c^2$).
        \item \textbf{Foyers :} $F_1(x_0,\; y_0 - c)$ et $F_2(x_0,\; y_0 + c)$.
        \item \textbf{Excentricité et directrices :} $e = \dfrac{c}{b} < 1$ \quad et \quad $d_{1,2} \equiv y = y_0 \pm \dfrac{b^2}{c} = y_0 \pm \dfrac{b}{e}$.
    \end{itemize}
\end{enumerate}
\end{theoreme}

\begin{theoreme}[2. L'Hyperbole d'axes parallèles aux axes du repère]
Soit l'hyperbole $\mathcal{H}$ de \textbf{centre $C(x_0, y_0)$} :
\begin{enumerate}
    \item \textbf{Axe focal horizontal (le signe $+$ est devant $(x-x_0)^2$) :}
    \[
        \boxed{\mathcal{H} \equiv \frac{(x - x_0)^2}{a^2} - \frac{(y - y_0)^2}{b^2} = 1}.
    \]
    \begin{itemize}
        \item \textbf{Sommets :} $S_1(x_0 - a, y_0)$ et $S_2(x_0 + a, y_0)$ (sur l'axe focal $y = y_0$).
        \item \textbf{Distance focale $c$ :} $\boxed{c = \sqrt{a^2 + b^2}}$ \quad (\emph{Attention au signe $+$ : pour l'hyperbole, $c > a$!}).
        \item \textbf{Foyers :} $F_1(x_0 - c,\; y_0)$ et $F_2(x_0 + c,\; y_0)$. \quad (\emph{Propriété bifocale :} $\bigl||MF_1| - |MF_2|\bigr| = 2a$.)
        \item \textbf{Excentricité et directrices :} $e = \dfrac{c}{a} > 1$ \quad et \quad $\boxed{d_{1,2} \equiv x = x_0 \pm \dfrac{a^2}{c}}$.
        \item \textbf{Asymptotes :} Les deux droites passant par le centre $C(x_0,y_0)$ d'équations :
        \[
            \boxed{y - y_0 = \pm \frac{b}{a}(x - x_0)} \qquad \left(\text{obtenues en remplaçant $1$ par $0$ dans l'équation réduite!}\right).
        \]
    \end{itemize}
    \item \textbf{Axe focal vertical (le signe $+$ est devant $(y-y_0)^2$) :} $-\dfrac{(x-x_0)^2}{a^2} + \dfrac{(y-y_0)^2}{b^2} = 1$.\\
    Les sommets sont $(x_0, y_0 \pm b)$, la distance focale reste $c = \sqrt{a^2+b^2}$, les foyers sont $(x_0, y_0 \pm c)$, les directrices sont $y = y_0 \pm \frac{b^2}{c}$ et les asymptotes restent $y - y_0 = \pm \frac{b}{a}(x - x_0)$.
\end{enumerate}
\end{theoreme}

\begin{theoreme}[3. La Parabole d'axe parallèle à l'un des axes du repère]
Soit une parabole $\mathcal{P}$ de \textbf{sommet $S(x_0, y_0)$} et de paramètre focal $p > 0$ (distance foyer--directrice) :
\begin{enumerate}
    \item \textbf{Axe de symétrie vertical $x = x_0$ ($y$ en fonction de $x^2$) :}
    \[
        \boxed{(x - x_0)^2 = \pm 2p(y - y_0) \quad \Longleftrightarrow \quad y - y_0 = a(x - x_0)^2} \quad \left(\text{avec } |a| = \frac{1}{2p}\right).
    \]
    \begin{itemize}
        \item Si $a > 0$ (concave vers les $y$ croissants, signe $+$) : \textbf{Foyer} $F\left(x_0,\; y_0 + \frac{p}{2}\right) = \left(x_0,\; y_0 + \frac{1}{4a}\right)$ et \textbf{Directrice} $d \equiv y = y_0 - \frac{p}{2} = y_0 - \frac{1}{4a}$.
        \item Si $a < 0$ (concave vers les $y$ décroissants, signe $-$) : Foyer $F\left(x_0,\; y_0 - \frac{p}{2}\right)$ et Directrice $d \equiv y = y_0 + \frac{p}{2}$.
    \end{itemize}
    \item \textbf{Axe de symétrie horizontal $y = y_0$ ($x$ en fonction de $y^2$) :}
    \[
        \boxed{(y - y_0)^2 = \pm 2p(x - x_0) \quad \Longleftrightarrow \quad x - x_0 = a(y - y_0)^2}.
    \]
    Foyer $F\left(x_0 \pm \frac{p}{2},\; y_0\right)$ et Directrice $d \equiv x = x_0 \mp \frac{p}{2}$.
\end{enumerate}
\end{theoreme}

\begin{center}
\begin{tikzpicture}[scale=0.82]
    % Figure 1 : Ellipse
    \begin{scope}[shift={(0,0)}]
        \draw[->, gray] (-2.6,0) -- (2.6,0) node[right, black] {\small $x$};
        \draw[->, gray] (0,-1.8) -- (0,1.8) node[above, black] {\small $y$};
        \draw[thick, uclblue] (0,0) ellipse (2.0 and 1.2);
        \draw[dashed, defgreen] (2.5,-1.7) -- (2.5,1.7) node[above] {\tiny $x=\frac{a^2}{c}$};
        \draw[dashed, defgreen] (-2.5,-1.7) -- (-2.5,1.7);
        \fill[warnred] (-1.6,0) circle (1.8pt) node[below] {\small $F_1$};
        \fill[warnred] (1.6,0) circle (1.8pt) node[below] {\small $F_2$};
        \fill[uclblue] (2.0,0) circle (1.5pt) node[below right] {\tiny $a$};
        \fill[uclblue] (0,1.2) circle (1.5pt) node[above left] {\tiny $b$};
        \node[below, align=center] at (0,-2.0) {\small \textbf{(a)} Ellipse ($e = \frac{c}{a} < 1$)\\[-2pt]\small $c = \sqrt{a^2-b^2}$};
    \end{scope}

    % Figure 2 : Hyperbole
    \begin{scope}[shift={(6.2,0)}]
        \draw[->, gray] (-2.6,0) -- (2.6,0) node[right, black] {\small $x$};
        \draw[->, gray] (0,-1.8) -- (0,1.8) node[above, black] {\small $y$};
        \draw[dashed, gray] (-2.3,-1.61) -- (2.3,1.61) node[above right, black] {\tiny $y=\frac{b}{a}x$};
        \draw[dashed, gray] (-2.3,1.61) -- (2.3,-1.61);
        \draw[thick, uclblue, domain=-1.15:1.15, samples=40] plot ({1.1*cosh(\x)}, {0.77*sinh(\x)});
        \draw[thick, uclblue, domain=-1.15:1.15, samples=40] plot ({-1.1*cosh(\x)}, {0.77*sinh(\x)});
        \fill[warnred] (-1.55,0) circle (1.8pt) node[below] {\small $F_1$};
        \fill[warnred] (1.55,0) circle (1.8pt) node[below] {\small $F_2$};
        \fill[uclblue] (1.1,0) circle (1.5pt) node[below right] {\tiny $a$};
        \node[below, align=center] at (0,-2.0) {\small \textbf{(b)} Hyperbole ($e = \frac{c}{a} > 1$)\\[-2pt]\small $c = \sqrt{a^2+b^2}$};
    \end{scope}

    % Figure 3 : Parabole
    \begin{scope}[shift={(12.2,0)}]
        \draw[->, gray] (-2.2,0) -- (2.2,0) node[right, black] {\small $x$};
        \draw[->, gray] (0,-1.8) -- (0,1.8) node[above, black] {\small $y$};
        \draw[thick, uclblue, domain=-1.8:1.8, samples=40] plot (\x, {0.5*\x*\x - 0.4});
        \draw[dashed, defgreen] (-2.0,-0.9) -- (2.0,-0.9) node[right] {\tiny $d : y = y_0 - \frac{p}{2}$};
        \fill[warnred] (0,0.1) circle (1.8pt) node[above right] {\small $F$};
        \fill[uclblue] (0,-0.4) circle (1.5pt) node[below right] {\small $S$};
        \node[below, align=center] at (0,-2.0) {\small \textbf{(c)} Parabole ($e = 1$)\\[-2pt]\small $(x-x_0)^2 = 2p(y-y_0)$};
    \end{scope}
\end{tikzpicture}
\end{center}

\begin{exemple}[Quatre problèmes d'examen UCLouvain sur les coniques]
\textbf{1) Réduction d'une hyperbole, directrices et aire (UCLouvain -- Juillet 2026, Question 1) :}\\
\emph{Énoncé :} On considère la conique $\mathcal{C} \equiv x^2 - y^2 + 2y = 4$. (a) Mettre cette équation sous forme réduite, déterminer sa nature et ses directrices. (b) Calculer l'aire du triangle formé par les deux foyers $F_1, F_2$ et l'origine $O(0,0)$.\\
\emph{Résolution :}
\begin{enumerate}[label=(\alph*)]
    \item Complétons le carré en $y$ : $x^2 - (y^2 - 2y) = 4 \iff x^2 - \bigl((y-1)^2 - 1\bigr) = 4 \iff x^2 - (y-1)^2 = 3$. En divisant par $3$, on obtient l'équation réduite :
    \[
        \boxed{\frac{x^2}{3} - \frac{(y-1)^2}{3} = 1}.
    \]
    C'est une \textbf{hyperbole} (équilatère) de centre $C(0, 1)$, d'axe focal horizontal $y = 1$, avec $a^2 = b^2 = 3$. La distance focale vaut $c = \sqrt{a^2 + b^2} = \sqrt{3 + 3} = \sqrt{6}$. Les deux directrices (verticales) ont pour équations :
    \[
        x = x_0 \pm \frac{a^2}{c} = 0 \pm \frac{3}{\sqrt{6}} = \boxed{\pm \frac{\sqrt{6}}{2}}.
    \]
    \item Les deux foyers sont $F_1(-\sqrt{6}, 1)$ et $F_2(\sqrt{6}, 1)$. Le segment $[F_1 F_2]$ est horizontal (sur la droite $y = 1$), de longueur $|F_1 F_2| = 2c = 2\sqrt{6}$, et sa hauteur par rapport à l'origine $O(0,0)$ vaut $h = 1$. L'aire du triangle $O F_1 F_2$ vaut donc $\frac{1}{2} \times 2\sqrt{6} \times 1 = \boxed{\sqrt{6}}$.
\end{enumerate}

\medskip
\textbf{2) Réduction d'une ellipse et intersection avec une droite (UCLouvain -- Septembre 2026, Question 1) :}\\
\emph{Énoncé :} On considère la conique $\mathcal{C} \equiv 9x^2 + 4y^2 - 36x + 8y + 4 = 0$. (a) Déterminer sa forme réduite, sa nature et ses foyers $F_1, F_2$. (b) Déterminer ses points d'intersection $A$ et $B$ avec la droite $y = x - 1$.\\
\emph{Résolution :}
\begin{enumerate}[label=(\alph*)]
    \item Factorisons par $9$ les termes en $x$ et par $4$ les termes en $y$ pour compléter les carrés :
    \begin{align*}
        9(x^2 - 4x) + 4(y^2 + 2y) + 4 = 0 &\iff 9\bigl((x-2)^2 - 4\bigr) + 4\bigl((y+1)^2 - 1\bigr) + 4 = 0 \\
        &\iff 9(x-2)^2 + 4(y+1)^2 = 36.
    \end{align*}
    En divisant par $36$, on obtient :
    \[
        \boxed{\frac{(x-2)^2}{4} + \frac{(y+1)^2}{9} = 1}.
    \]
    C'est une \textbf{ellipse} de centre $C(2, -1)$, avec $a^2 = 4$ ($a=2$) et $b^2 = 9$ ($b=3$). Comme $b > a$, le \textbf{grand axe est vertical} ($x = 2$), et $c = \sqrt{b^2 - a^2} = \sqrt{9 - 4} = \sqrt{5}$. Les foyers sont donc :
    \[
        \boxed{F_1(2,\; -1 - \sqrt{5}) \quad \text{et} \quad F_2(2,\; -1 + \sqrt{5})}.
    \]
    \item Remplaçons $y = x - 1$ (c'est-à-dire $y + 1 = x$) dans $9(x-2)^2 + 4(y+1)^2 = 36$ :
    \[
        9(x^2 - 4x + 4) + 4x^2 = 36 \iff 13x^2 - 36x = 0 \implies x = 0 \;\text{ ou }\; x = \frac{36}{13}.
    \]
    Avec $y = x - 1$, on obtient les deux points d'intersection : $\boxed{A(0, -1) \text{ et } B\left(\frac{36}{13}, \frac{23}{13}\right)}$.
\end{enumerate}

\medskip
\textbf{3) Hyperbole à partir de ses asymptotes (UCLouvain -- Septembre 2025, Question 1a) :}\\
\emph{Énoncé :} Une hyperbole $\mathcal{C}$ centrée à l'origine $O(0,0)$ a pour asymptotes les droites $y = \pm \frac{x}{2}$ et passe par le point $A(2, 0)$. Donner son équation réduite et ses foyers.\\
\emph{Résolution :} Puisque le sommet $A(2,0)$ est sur l'axe $Ox$, l'équation réduite est $\frac{x^2}{a^2} - \frac{y^2}{b^2} = 1$ avec $a = 2$ ($a^2 = 4$). La pente des asymptotes vaut $\frac{b}{a} = \frac{1}{2}$, d'où $b = \frac{a}{2} = 1$ ($b^2 = 1$). L'équation est donc $\boxed{\frac{x^2}{4} - y^2 = 1}$, et avec $c = \sqrt{a^2+b^2} = \sqrt{4+1} = \sqrt{5}$, les foyers sont $\boxed{F_1(-\sqrt{5}, 0) \text{ et } F_2(\sqrt{5}, 0)}$.

\medskip
\textbf{4) Trajectoire parabolique d'un projectile (UCLouvain -- Septembre 2023, Question 1a) :}\\
\emph{Énoncé :} Un projectile lancé depuis $(0, 11)$ décrit une parabole $\mathcal{P}$ dont le point culminant est à $36\text{ m}$ au-dessus du sol ($y_{\max} = 36$) et qui retombe au sol en $(11, 0)$. Déterminer l'équation cartésienne de $\mathcal{P}$.\\
\emph{Résolution :} Le sommet de la parabole a pour ordonnée $y_0 = 36$, donc son équation canonique s'écrit $y = 36 - a(x - x_0)^2$ avec $a > 0$. En exprimant qu'elle passe par $(0, 11)$ et $(11, 0)$ :
\[
    \begin{cases} 11 = 36 - a x_0^2 \iff a x_0^2 = 25 \\ 0 = 36 - a(11 - x_0)^2 \iff a(11 - x_0)^2 = 36 \end{cases} \implies \frac{11 - x_0}{x_0} = \frac{6}{5} \implies x_0 = 5 \text{ et } a = 1.
\]
D'où l'équation de la parabole : $\boxed{y = 36 - (x-5)^2 = -x^2 + 10x + 11}$.
\end{exemple}

% ----------------------------------------------------------------------------
\section{Tangentes à une conique et angles d'intersection}
% ----------------------------------------------------------------------------

\begin{theoreme}[Condition de tangence droite--conique ($\Delta = 0$) et vecteur tangent]
\begin{enumerate}
    \item \textbf{Méthode du discriminant nul ($\Delta = 0$) pour la tangence :} Pour déterminer si une droite $d$ (d'équation $y = mx + p$) coupe ou est tangente à une conique $\mathcal{C}$, on substitue $y = mx + p$ dans l'équation de $\mathcal{C}$. On obtient une équation du second degré $Ax^2 + Bx + C = 0$ de discriminant $\Delta = B^2 - 4AC$ :
    \begin{itemize}
        \item $\Delta > 0 \iff$ la droite est \textbf{sécante} (2 points d'intersection distincts) ;
        \item $\boxed{\Delta = 0 \iff \text{la droite est \textbf{tangente} à la conique (1 point de contact double)}}$ ;
        \item $\Delta < 0 \iff$ la droite est \textbf{extérieure} (aucun point d'intersection).
    \end{itemize}
    \item \textbf{Équation directe de la tangente en un point $M_1(x_1, y_1) \in \mathcal{C}$ (Règle de dédoublement) :}
    \begin{itemize}
        \item Pour l'ellipse $\frac{(x-x_0)^2}{a^2} + \frac{(y-y_0)^2}{b^2} = 1$ : la tangente en $(x_1, y_1)$ est $\dfrac{(x_1-x_0)(x-x_0)}{a^2} + \dfrac{(y_1-y_0)(y-y_0)}{b^2} = 1$.
        \item Pour l'hyperbole $\frac{(x-x_0)^2}{a^2} - \frac{(y-y_0)^2}{b^2} = 1$ : la tangente en $(x_1, y_1)$ est $\dfrac{(x_1-x_0)(x-x_0)}{a^2} - \dfrac{(y_1-y_0)(y-y_0)}{b^2} = 1$.
        \item Pour une courbe $y = f(x)$ (par exemple une parabole $y = ax^2 + bx + c$) : la pente de la tangente au point d'abscisse $x_1$ est la dérivée $f'(x_1)$, et un \textbf{vecteur directeur de la tangente} est $\boxed{\vec{t} = \bigl(1,\; f'(x_1)\bigr)}$.
    \end{itemize}
\end{enumerate}
\end{theoreme}

\begin{exemple}[Problèmes de tangence aux examens UCLouvain (Juillet 2025 Q1b et Juillet 2023 Q1c)]
\textbf{1) Pentes limites d'une droite coupant une ellipse (UCLouvain -- Juillet 2025, Question 1) :}\\
\emph{Énoncé :} Soit l'ellipse $4x^2 - 32x + 64 + y^2 = 4 \iff \frac{(x-4)^2}{1} + \frac{y^2}{4} = 1$. Une droite passant par l'origine, d'équation $y = mx$, coupe cette ellipse. Déterminer les coefficients directeurs maximum $m_{\max}$ et minimum $m_{\min}$ permettant que la droite coupe l'ellipse.\\
\emph{Résolution :} Substituons $y = mx$ dans l'équation de l'ellipse $4(x-4)^2 + y^2 = 4$ :
\[
    4(x^2 - 8x + 16) + m^2 x^2 = 4 \iff (4 + m^2)x^2 - 32x + 60 = 0.
\]
Cette équation du second degré en $x$ admet au moins une solution réelle (les cas limites correspondant aux deux tangentes issues de l'origine) si et seulement si son discriminant $\Delta$ est positif ou nul :
\[
    \Delta = (-32)^2 - 4(4 + m^2)(60) \ge 0 \iff 64 - 240m^2 \ge 0 \iff m^2 \le \frac{64}{240} = \frac{4}{15}.
\]
D'où les deux pentes limites : $\boxed{m_{\min} = -\frac{2}{\sqrt{15}} \quad \text{et} \quad m_{\max} = \frac{2}{\sqrt{15}}}$.

\medskip
\textbf{2) Angles d'intersection entre une droite et une parabole (UCLouvain -- Juillet 2023, Question 1) :}\\
\emph{Énoncé :} Soit la parabole $\mathcal{L} \equiv (x+2)^2 - 2y = 8 \iff y = \frac{1}{2}(x+2)^2 - 4$ et la droite $d \equiv x - 2y = 0$. (a) Trouver les points d'intersection de $\mathcal{L}$ et $d$. (b) Calculer les angles $\varphi_1$ et $\varphi_2$ entre la droite $d$ et les tangentes à $\mathcal{L}$ en ces deux points.\\
\emph{Résolution :}
\begin{enumerate}[label=(\alph*)]
    \item En remplaçant $2y = x$ dans $(x+2)^2 - 2y = 8$, on obtient $x^2 + 4x + 4 - x - 8 = 0 \iff x^2 + 3x - 4 = 0 \iff (x+4)(x-1) = 0$. Les deux points d'intersection sont donc $\boxed{P_1(-4, -2) \text{ et } P_2\left(1, \frac{1}{2}\right)}$.
    \item La droite $d \equiv y = \frac{1}{2}x$ a pour vecteur directeur $\vec{d} = (2, 1)$ (de norme $\|\vec{d}\| = \sqrt{5}$). La dérivée de $y(x) = \frac{1}{2}(x+2)^2 - 4$ est $y'(x) = x + 2$, donc le vecteur directeur de la tangente au point d'abscisse $x$ est $\vec{t} = (1, x+2)$ :
    \begin{itemize}
        \item En $P_1$ ($x = -4$) : $\vec{t}_1 = (1, -2)$. On constate que $\vec{d} \cdot \vec{t}_1 = 2(1) + 1(-2) = 0$, donc $\vec{d} \perp \vec{t}_1$ et $\boxed{\varphi_1 = 90\degree}$.
        \item En $P_2$ ($x = 1$) : $\vec{t}_2 = (1, 3)$ (de norme $\|\vec{t}_2\| = \sqrt{10}$). Par le produit scalaire :
        \[
            \cos\varphi_2 = \frac{\vec{d} \cdot \vec{t}_2}{\|\vec{d}\|\,\|\vec{t}_2\|} = \frac{2(1) + 1(3)}{\sqrt{5}\sqrt{10}} = \frac{5}{5\sqrt{2}} = \frac{\sqrt{2}}{2} \implies \boxed{\varphi_2 = 45\degree}.
        \]
    \end{itemize}
\end{enumerate}
\end{exemple}

% ----------------------------------------------------------------------------
\section{Exercices du Chapitre 6}
% ----------------------------------------------------------------------------

\emph{Les corrigés complets de tous ces exercices se trouvent au Chapitre 9, en fin d'ouvrage.}

\begin{enumerate}[label=\textbf{Exercice 6.\arabic* --}, leftmargin=*]

    \item \textbf{(UCLouvain -- Septembre 2023, Q1b : Alignement et orthogonalité dans le plan)}
    On donne les points $A(4, -3)$ et $B(1, -2)$, ainsi que la droite $d \equiv y = 1 + x$. Déterminez les coordonnées des deux points $C \in d$ et $D \in d$ tels que :
    \begin{enumerate}[label=(\alph*)]
        \item les trois points $A, B, C$ soient alignés ;
        \item les segments $[AB]$ et $[BD]$ soient perpendiculaires.
    \end{enumerate}

    \item \textbf{(UCLouvain -- Septembre 2021, Q1 : Droites perpendiculaires et cercle tangent)}
    Soit la droite $\mathcal{A} \equiv 3x + y + 12 = 0$ et le point $P(2, -2)$.
    \begin{enumerate}[label=(\alph*)]
        \item Déterminez l'équation cartésienne de la droite $\mathcal{B}$ perpendiculaire à $\mathcal{A}$ et passant par $P$.
        \item Un cercle $\mathcal{C}$ de rayon $R = \sqrt{10}$ est tangent simultanément aux deux droites $\mathcal{A}$ et $\mathcal{B}$, et son centre possède une abscisse négative et une ordonnée positive (situé dans le quadrant supérieur/droit par rapport au point d'intersection de $\mathcal{A}$ et $\mathcal{B}$, du côté de l'origine). Calculez les coordonnées exactes du centre de $\mathcal{C}$.
    \end{enumerate}

    \item \textbf{(UCLouvain -- Septembre 2022, Q1 : Médiatrice, cercle circonscrit et cercle tangent)}
    Dans un repère orthonormé, on considère le triangle de sommets $E(-2, -4)$, $F(3, 1)$ et $G(-5, 5)$.
    \begin{enumerate}[label=(\alph*)]
        \item Déterminez l'équation cartésienne de la médiatrice $m$ du segment $[EG]$.
        \item Déterminez l'équation cartésienne du cercle $\mathcal{C}$ circonscrit au triangle $EFG$.
        \item Calculez le rayon $r$ du plus grand cercle $\mathcal{B}$ que l'on puisse inscrire dans la région comprise à l'extérieur du triangle $EFG$ (du côté du segment $[EG]$) et à l'intérieur du cercle circonscrit $\mathcal{C}$.
    \end{enumerate}

    \item \textbf{(UCLouvain -- Juillet 2021, Q3 : Cercle tangent à une droite et tangente à une parabole)}
    On considère la droite $d \equiv \dfrac{x}{3} + \dfrac{y}{4} = 1$, le cercle $\mathcal{C} \equiv \dfrac{x^2}{a} + \dfrac{y^2}{a} = 1$ (avec $a > 0$) et la parabole $\mathcal{P} \equiv \dfrac{x}{6} + \dfrac{y^2}{27} = -\dfrac{1}{2}$.
    \begin{enumerate}[label=(\alph*)]
        \item Calculez la valeur exacte du paramètre $a$ qui rend le cercle $\mathcal{C}$ tangent à la droite $d$.
        \item Déterminez l'équation cartésienne de la droite $m$ qui est perpendiculaire à la droite $d$ et tangente à la parabole $\mathcal{P}$, ainsi que les coordonnées du point de tangence.
    \end{enumerate}

    \item \textbf{(UCLouvain -- Juillet 2022, Q2 : Médiatrice et ellipse tangente à une droite)}
    On donne les points $C\left(-\frac{12}{5}, \frac{6}{5}\right)$ et $D(5, -3)$.
    \begin{enumerate}[label=(\alph*)]
        \item Donnez une équation paramétrique (et cartésienne) de la médiatrice $m$ du segment $[CD]$.
        \item Déterminez l'équation cartésienne de l'ellipse $\mathcal{E}$ centrée à l'origine, d'axes parallèles aux axes du repère, sachant qu'elle passe par le point $C$ et qu'elle est tangente à la droite $d \equiv 8x + 9y = 30$.
    \end{enumerate}

    \item \textbf{(UCLouvain -- Juillet 2024, Q1a : Réduction et foyers d'une ellipse)}
    Mettez l'équation cartésienne $4x^2 - 8x + y^2 = 0$ sous sa forme réduite, précisez la nature de la conique, son centre, les longueurs de son grand axe et de son petit axe, ses foyers, son excentricité et ses directrices.

    \item \textbf{(UCLouvain -- Septembre 2025, Q1b : Médiatrice et intersection avec une hyperbole)}
    On considère l'hyperbole $\mathcal{C} \equiv \dfrac{x^2}{4} - y^2 = 1$ et les deux points $A(2, 0)$ et $B(2, 2)$. Déterminez l'équation cartésienne de la médiatrice du segment $[AB]$, ainsi que les coordonnées de ses deux points d'intersection $C$ et $D$ avec l'hyperbole $\mathcal{C}$.

    \item \textbf{(UCLouvain -- Juillet 2023, Q1a : Éléments caractéristiques d'une parabole)}
    Pour la parabole $\mathcal{L} \equiv (x+2)^2 - 2y = 8$, déterminez les coordonnées de son sommet $S$, de son foyer $F$, l'équation de sa directrice $d$, ainsi que ses points d'intersection avec l'axe des abscisses $Ox$.

    \item \textbf{(Classification et réduction de coniques)}
    Pour chacune des deux équations suivantes, complétez les carrés pour obtenir l'équation réduite, identifiez la nature de la conique, son centre (ou sommet), ses foyers, son excentricité et, le cas échéant, ses asymptotes :
    \begin{enumerate}[label=(\alph*)]
        \item $4x^2 - 9y^2 - 16x - 18y - 29 = 0$
        \item $y^2 - 6y - 8x + 17 = 0$
    \end{enumerate}

    \item \textbf{(Propriété optique et lieu géométrique bifocal)}
    Déterminez l'équation cartésienne réduite de l'ellipse dont les foyers sont $F_1(-3, 0)$ et $F_2(3, 0)$ et dont la somme des distances de tout point $M$ aux deux foyers vaut $|MF_1| + |MF_2| = 10$. Déterminez ensuite les équations des deux tangentes à cette ellipse qui sont parallèles à la droite $4x + 3y = 0$.

\end{enumerate}

% ============================================================================
% CHAPITRE 7 : GÉOMÉTRIE ANALYTIQUE ET VECTORIELLE DANS L'ESPACE
% ============================================================================
\chapter[Géométrie Analytique dans l'Espace]{Géométrie Analytique et Vectorielle dans l'Espace}

La deuxième question de la partie géométrie de l'examen d'admission à l'UCLouvain est consacrée à la \textbf{géométrie analytique dans l'espace à trois dimensions} muni d'un repère orthonormé $(O;\vec{i},\vec{j},\vec{k})$. Qu'il s'agisse de programmer le bras d'un robot charpentier (Juillet 2026), de creuser des galeries minières (Septembre 2026), de suivre un éboulement rocheux sur un versant (Juillet 2025), de calculer la réflexion d'un rayon lumineux sur une verrière (Septembre 2025) ou la descente en rappel d'une alpiniste (Juillet 2021), tous ces problèmes se résolvent à l'aide d'un petit nombre d'outils vectoriels fondamentaux.

% ----------------------------------------------------------------------------
\section{Calcul vectoriel dans l'espace : produit scalaire, produit vectoriel et déterminant}
% ----------------------------------------------------------------------------

Dans l'espace muni d'un repère orthonormé $(O;\vec{i},\vec{j},\vec{k})$, soient deux points $A(x_A, y_A, z_A)$ et $B(x_B, y_B, z_B)$, et trois vecteurs $\vec{u} = (u_x, u_y, u_z)$, $\vec{v} = (v_x, v_y, v_z)$ et $\vec{w} = (w_x, w_y, w_z)$.

\begin{theoreme}[Outils vectoriels fondamentaux dans $\R^3$]
\begin{enumerate}
    \item \textbf{Vecteur $\overrightarrow{AB}$, distance $|AB|$ et point milieu $M$ de $[AB]$ :}
    \[
        \overrightarrow{AB} = (x_B - x_A,\; y_B - y_A,\; z_B - z_A), \qquad |AB| = \|\overrightarrow{AB}\| = \sqrt{(x_B-x_A)^2 + (y_B-y_A)^2 + (z_B-z_A)^2},
    \]
    \[
        M = \left(\frac{x_A+x_B}{2},\; \frac{y_A+y_B}{2},\; \frac{z_A+z_B}{2}\right).
    \]
    \item \textbf{Produit scalaire $\vec{u} \cdot \vec{v}$ et angle $\theta \in [0,\pi]$ entre deux vecteurs :}
    \[
        \boxed{\vec{u} \cdot \vec{v} = u_x v_x + u_y v_y + u_z v_z = \|\vec{u}\|\,\|\vec{v}\|\cos\theta} \implies \cos\theta = \frac{u_x v_x + u_y v_y + u_z v_z}{\sqrt{u_x^2+u_y^2+u_z^2}\sqrt{v_x^2+v_y^2+v_z^2}}.
    \]
    En particulier, deux vecteurs non nuls sont \textbf{orthogonaux} ssi leur produit scalaire est nul :
    \[
        \boxed{\vec{u} \perp \vec{v} \iff \vec{u} \cdot \vec{v} = u_x v_x + u_y v_y + u_z v_z = 0}.
    \]
    \item \textbf{Produit vectoriel $\vec{u} \wedge \vec{v}$ (pour construire un vecteur perpendiculaire à $\vec{u}$ et à $\vec{v}$) :}
    Le produit vectoriel de $\vec{u}$ et $\vec{v}$ est le vecteur $\vec{n} = \vec{u} \wedge \vec{v}$ simultanément orthogonal à $\vec{u}$ et à $\vec{v}$ (de norme $\|\vec{u}\wedge\vec{v}\| = \|\vec{u}\|\,\|\vec{v}\|\sin\theta$, égale à l'aire du parallélogramme construit sur $\vec{u}$ et $\vec{v}$) :
    \[
        \boxed{\vec{u} \wedge \vec{v} = \begin{vmatrix} \vec{i} & u_x & v_x \\ \vec{j} & u_y & v_y \\ \vec{k} & u_z & v_z \end{vmatrix} = \bigl(u_y v_z - u_z v_y,\;\; u_z v_x - u_x v_z,\;\; u_x v_y - u_y v_x\bigr)}.
    \]
    \item \textbf{Produit mixte (déterminant $3 \times 3$) et condition de coplanarité :}
    Trois vecteurs $\vec{u}, \vec{v}, \vec{w}$ sont \textbf{coplanaires} (appartiennent à un même plan vectoriel) ssi leur déterminant est nul :
    \[
        \boxed{\vec{u}, \vec{v}, \vec{w} \text{ coplanaires} \iff \det(\vec{u}, \vec{v}, \vec{w}) = \begin{vmatrix} u_x & v_x & w_x \\ u_y & v_y & w_y \\ u_z & v_z & w_z \end{vmatrix} = (\vec{u} \wedge \vec{v}) \cdot \vec{w} = 0}.
    \]
\end{enumerate}
\end{theoreme}

% ----------------------------------------------------------------------------
\section{Droites et plans dans l'espace : équations et représentations}
% ----------------------------------------------------------------------------

\subsection{Équations paramétriques et cartésiennes d'une droite dans $\R^3$}

\begin{theoreme}[Représentations d'une droite $d$ dans l'espace]
Soit $d$ la droite passant par le point $A(x_A, y_A, z_A)$ et de \textbf{vecteur directeur} $\vec{u} = (a, b, c) \neq \vec{0}$ :
\begin{enumerate}
    \item \textbf{Équations paramétriques (la forme la plus pratique pour tous les calculs !) :}
    \[
        \boxed{d \equiv \begin{cases} x = x_A + t\,a \\ y = y_A + t\,b \\ z = z_A + t\,c \end{cases} \quad (t \in \R) \qquad \Longleftrightarrow \qquad \overrightarrow{OM}(t) = \overrightarrow{OA} + t\,\vec{u}}.
    \]
    \emph{Remarque importante :} Le paramètre $t$ mesure le déplacement le long de la droite en unités de $\|\vec{u}\|$ ! Ainsi, la distance entre le point initial $A$ ($t=0$) et le point $M(t)$ vaut directement :
    \[
        \boxed{|AM(t)| = |t| \cdot \|\vec{u}\| = |t|\sqrt{a^2 + b^2 + c^2}}.
    \]
    \item \textbf{Équations cartésiennes (intersection de deux plans sécants) :} Si $a,b,c \neq 0$, en isolant $t$ :
    \[
        \frac{x - x_A}{a} = \frac{y - y_A}{b} = \frac{z - z_A}{c}, \qquad \text{ou sous forme d'un système : } d \equiv \begin{cases} a_1 x + b_1 y + c_1 z + d_1 = 0 \\ a_2 x + b_2 y + c_2 z + d_2 = 0 \end{cases}.
    \]
    \emph{Astuce :} Lorsqu'une droite $d$ est donnée comme l'intersection de deux plans de vecteurs normaux $\vec{n}_1$ et $\vec{n}_2$, son vecteur directeur est $\vec{u} = \vec{n}_1 \wedge \vec{n}_2$ (ou on pose simplement $z = t$ et on exprime $x$ et $y$ en fonction de $t$, comme en Septembre 2022 Q2a!).
\end{enumerate}
\end{theoreme}

\subsection{Équation cartésienne d'un plan et plan médiateur}

\begin{theoreme}[Équation cartésienne et paramétrique d'un plan $\Pi$]
\begin{enumerate}
    \item \textbf{Par un point $A(x_A, y_A, z_A)$ et un vecteur normal $\vec{n} = (a, b, c) \neq \vec{0}$ :}
    Un point $M(x,y,z)$ appartient au plan $\Pi$ ssi $\overrightarrow{AM} \perp \vec{n} \iff \vec{n} \cdot \overrightarrow{AM} = 0$, ce qui donne l'\textbf{équation cartésienne du plan} :
    \[
        \boxed{\Pi \equiv a(x - x_A) + b(y - y_A) + c(z - z_A) = 0 \quad \Longleftrightarrow \quad ax + by + cz + d = 0}
    \]
    où les coefficients $\mathbf{(a, b, c)}$ de $x, y, z$ forment toujours un \textbf{vecteur normal $\vec{n}$} (perpendiculaire) au plan $\Pi$ !
    \item \textbf{Par un point $A$ et deux vecteurs directeurs non colinéaires $\vec{u}$ et $\vec{v}$ (ou par 3 points $A, B, C$) :}
    \begin{itemize}
        \item \emph{Équation paramétrique :} $\overrightarrow{OM} = \overrightarrow{OA} + r\,\vec{u} + s\,\vec{v}$ ($r, s \in \R$).
        \item \emph{Passage à l'équation cartésienne :} Soit on calcule le vecteur normal $\vec{n} = \vec{u} \wedge \vec{v}$ (ou on résout $\vec{n}\cdot\vec{u} = 0$ et $\vec{n}\cdot\vec{v} = 0$), soit on annule le déterminant de coplanarité :
        \[
            M(x,y,z) \in \Pi \iff \det(\overrightarrow{AM}, \vec{u}, \vec{v}) = \begin{vmatrix} x - x_A & u_x & v_x \\ y - y_A & u_y & v_y \\ z - z_A & u_z & v_z \end{vmatrix} = 0.
        \]
    \end{itemize}
    \item \textbf{Plan médiateur d'un segment $[AB]$ (Septembre 2021 Q2a et Septembre 2024 Q2b) :}
    Le lieu géométrique des points de l'espace équidistants de deux points $A$ et $B$ est le \textbf{plan médiateur} de $[AB]$ : c'est le plan qui passe par le \textbf{milieu $M = \frac{A+B}{2}$} et qui a pour \textbf{vecteur normal $\vec{n} = \overrightarrow{AB}$}.
\end{enumerate}
\end{theoreme}

\begin{exemple}[Équations de plans et droites d'appui (UCLouvain -- Juillet 2024 Q2 et Septembre 2022 Q2)]
\textbf{1) Plan contenant une droite et droite s'appuyant sur deux droites (UCLouvain -- Juillet 2024, Question 2) :}\\
\emph{Énoncé :} Dans l'espace muni d'un repère orthonormé, on considère les deux droites :
\[
    d_1 \equiv \begin{cases} x = 5r + 1 \\ y = -r - 1 \\ z = 4 \end{cases} \quad (r \in \R) \qquad \text{et} \qquad d_2 \equiv \begin{cases} x = s \\ y = 3s + 2 \\ z = -2s - 2 \end{cases} \quad (s \in \R).
\]
(a) Déterminer une équation cartésienne du plan $\Pi$ qui contient l'origine $O(0,0,0)$ ainsi que la droite $d_1$.\\
(b) Déterminer un vecteur directeur de la droite $d$ qui passe par l'origine $O$ et qui s'appuie à la fois sur $d_1$ et sur $d_2$ (c'est-à-dire qui coupe $d_1$ et $d_2$).\\
\emph{Résolution :}
\begin{enumerate}[label=(\alph*)]
    \item La droite $d_1$ passe par le point $A(1, -1, 4)$ (pour $r=0$) et a pour vecteur directeur $\vec{u}_1 = (5, -1, 0)$. Comme le plan $\Pi$ contient $O(0,0,0)$ et toute la droite $d_1$, il passe par $O$ et est dirigé par les deux vecteurs $\vec{u}_1 = (5, -1, 0)$ et $\overrightarrow{OA} = (1, -1, 4)$. Leur produit vectoriel donne un vecteur normal $\vec{n}$ à $\Pi$ :
    \[
        \vec{n} = \vec{u}_1 \wedge \overrightarrow{OA} = \begin{vmatrix} \vec{i} & 5 & 1 \\ \vec{j} & -1 & -1 \\ \vec{k} & 0 & 4 \end{vmatrix} = (-4,\; -20,\; -4) = -4(1,\; 5,\; 1).
    \]
    Comme $\Pi$ passe par l'origine $O(0,0,0)$, son équation cartésienne est immédiatement $\boxed{\Pi \equiv x + 5y + z = 0}$.
    \item \emph{Astuce géométrique lumineuse :} Toute droite $d$ passant par l'origine $O$ et coupant la droite $d_1$ est nécessairement \textbf{contenue dans le plan $\Pi$} (qui est justement le plan formé par $O$ et $d_1$)! Pour que cette droite $d \subset \Pi$ coupe également la droite $d_2$, elle doit passer par le \textbf{point de percée $P$ de $d_2$ dans le plan $\Pi$}.\\
    Substituons les équations paramétriques de $d_2$ dans $\Pi \equiv x + 5y + z = 0$ :
    \[
        s + 5(3s + 2) + (-2s - 2) = 0 \iff 14s + 8 = 0 \implies s = -\frac{4}{7}.
    \]
    Le point d'intersection $P \in d_2 \cap \Pi$ a donc pour coordonnées $P\left(-\frac{4}{7},\; \frac{2}{7},\; -\frac{6}{7}\right) = -\frac{2}{7}(2,\; -1,\; 3)$. La droite $d = (OP)$ a donc pour vecteur directeur $\boxed{\vec{v} = (2,\; -1,\; 3)}$.
\end{enumerate}

\medskip
\textbf{2) Plan parallèle à une droite et angle droit sur une droite (UCLouvain -- Septembre 2022, Question 2) :}\\
\emph{Énoncé :} On donne les points $A(-4, 2, 1)$ et $B(-1, 0, 3)$, ainsi que la droite $d \equiv \begin{cases} x - y + z = 0 \\ x + 2z - 5 = 0 \end{cases}$.\\
(a) Donner l'équation cartésienne du plan $\Pi$ contenant $A$ et $B$ et parallèle à $d$.\\
(b) Déterminer les coordonnées du point $C \in d$ tel que l'angle $\widehat{ACB}$ soit droit ($90\degree$).\\
\emph{Résolution :}
\begin{enumerate}[label=(\alph*)]
    \item Trouvons d'abord une représentation paramétrique de $d$ en posant $z = k$ : la seconde équation donne $x = 5 - 2k$, et la première donne $y = x + z = 5 - k$. Ainsi :
    \[
        d \equiv \begin{cases} x = 5 - 2k \\ y = 5 - k \\ z = k \end{cases} \quad (k \in \R), \qquad \text{de vecteur directeur } \vec{u} = (-2, -1, 1).
    \]
    Le plan $\Pi$ contient $A$ et $B$, donc le vecteur $\overrightarrow{AB} = (3, -2, 2)$ est parallèle à $\Pi$, tout comme $\vec{u} = (-2, -1, 1)$. Leur produit vectoriel fournit la normale à $\Pi$ :
    \[
        \vec{n} = \overrightarrow{AB} \wedge \vec{u} = \begin{vmatrix} \vec{i} & 3 & -2 \\ \vec{j} & -2 & -1 \\ \vec{k} & 2 & 1 \end{vmatrix} = (0,\; -7,\; -7) = -7(0, 1, 1).
    \]
    L'équation de $\Pi$ est donc de la forme $y + z + d = 0$. Comme $A(-4, 2, 1) \in \Pi$, on a $2 + 1 + d = 0 \implies d = -3$, soit $\boxed{\Pi \equiv y + z - 3 = 0}$.
    \item Un point quelconque $C \in d$ a pour coordonnées $C(5-2k,\; 5-k,\; k)$. Formons les vecteurs $\overrightarrow{AC}$ et $\overrightarrow{BC}$ :
    \[
        \overrightarrow{AC} = (9-2k,\; 3-k,\; k-1) \qquad \text{et} \qquad \overrightarrow{BC} = (6-2k,\; 5-k,\; k-3).
    \]
    L'angle $\widehat{ACB}$ est droit ssi $\overrightarrow{AC} \cdot \overrightarrow{BC} = 0$ :
    \[
        (9-2k)(6-2k) + (3-k)(5-k) + (k-1)(k-3) = 0 \iff 6k^2 - 42k + 72 = 0 \iff k^2 - 7k + 12 = 0,
    \]
    dont les racines sont $k = 3$ et $k = 4$, donnant les deux points $\boxed{C_1(-1, 2, 3) \text{ et } C_2(-3, 1, 4)}$.
\end{enumerate}
\end{exemple}

% ----------------------------------------------------------------------------
\section{Positions relatives, projections, distances, sphère et réflexion}
% ----------------------------------------------------------------------------

\subsection{Parallélisme, orthogonalité, intersections et distances dans l'espace}

\begin{theoreme}[Positions relatives et distances dans l'espace]
Soient une droite $d$ (de vecteur directeur $\vec{u}$) et deux plans $\Pi_1$ (de normale $\vec{n}_1 \equiv (a,b,c)$) et $\Pi_2$ (de normale $\vec{n}_2$) :
\begin{enumerate}
    \item \textbf{Attention à l'inversion entre droite et plan !} Une droite est définie par son \emph{vecteur directeur} $\vec{u}$, tandis qu'un plan est défini par son \emph{vecteur normal} $\vec{n}$. Par conséquent :
    \begin{itemize}
        \item \textbf{Droite parallèle à un plan ($d \parallel \Pi_1$) :} $\boxed{\vec{u} \perp \vec{n}_1 \iff \vec{u} \cdot \vec{n}_1 = 0}$ \quad (Juillet 2025 Q2b).
        \item \textbf{Droite perpendiculaire à un plan ($d \perp \Pi_1$) :} $\boxed{\vec{u} \parallel \vec{n}_1 \iff \exists k \in \R^* : \vec{u} = k\,\vec{n}_1}$ \quad (Septembre 2023 Q2a).
        \item \textbf{Deux plans parallèles ($\Pi_1 \parallel \Pi_2$) :} $\vec{n}_1 \parallel \vec{n}_2$.
        \item \textbf{Deux plans perpendiculaires ($\Pi_1 \perp \Pi_2$) :} $\boxed{\vec{n}_1 \perp \vec{n}_2 \iff \vec{n}_1 \cdot \vec{n}_2 = 0}$ \quad (Juillet 2022 Q1c).
    \end{itemize}
    \item \textbf{Angle entre une droite $d$ et la normale $\vec{n}$ au plan (angle d'incidence $\theta$) :}
    \[
        \cos\theta = \frac{|\vec{u} \cdot \vec{n}|}{\|\vec{u}\|\,\|\vec{n}\|} \qquad \left(\text{et l'angle $\alpha = \frac{\pi}{2}-\theta$ entre la droite et le plan vérifie } \sin\alpha = \frac{|\vec{u}\cdot\vec{n}|}{\|\vec{u}\|\,\|\vec{n}\|}\right).
    \]
    \item \textbf{Distance d'un point $A(x_A, y_A, z_A)$ à un plan $\Pi \equiv ax + by + cz + d = 0$ :}
    \[
        \boxed{d(A, \Pi) = \frac{|ax_A + by_A + cz_A + d|}{\sqrt{a^2 + b^2 + c^2}}}.
    \]
    \item \textbf{Sphère $\mathcal{S}$ de centre $A(x_A, y_A, z_A)$ et de rayon $R$ :}
    \[
        \boxed{\mathcal{S} \equiv (x - x_A)^2 + (y - y_A)^2 + (z - z_A)^2 = R^2} \quad \text{(et boule fermée : $\le R^2$, ouverte : $< R^2$)}.
    \]
    \textbf{Intersection d'une sphère $\mathcal{S}(A, R)$ et d'un plan $\Pi$ (Juillet 2022 Q1b) :} Si la distance $d = d(A, \Pi)$ du centre de la sphère au plan est inférieure à $R$ ($d < R$), alors l'intersection $\mathcal{S} \cap \Pi$ est un \textbf{cercle} (ou un disque pour la boule) dont le rayon $r$ est donné par le théorème de Pythagore :
    \[
        \boxed{r = \sqrt{R^2 - d(A, \Pi)^2} \implies \text{Aire du disque de section} = \pi r^2 = \pi\bigl(R^2 - d(A,\Pi)^2\bigr)}.
    \]
\end{enumerate}
\end{theoreme}

\subsection{Projection orthogonale, translation de droite et réflexion sur un miroir plan}

\begin{theoreme}[Trois constructions géométriques récurrentes à l'examen]
\begin{enumerate}
    \item \textbf{Pied $P$ de la perpendiculaire abaissée d'un point $C$ sur une droite $d \equiv A + t\vec{u}$ (Projeté orthogonal, Juillet 2023 Q2a et Septembre 2024 Q2a) :}
    On écrit les coordonnées génériques d'un point de la droite $P(t) = A + t\vec{u}$ et on détermine l'unique valeur de $t \in \R$ telle que le vecteur $\overrightarrow{CP(t)}$ soit orthogonal au vecteur directeur $\vec{u}$ :
    \[
        \boxed{P(t) = A + t\vec{u} \quad \text{avec} \quad \overrightarrow{CP(t)} \cdot \vec{u} = 0}.
    \]
    La distance du point $C$ à la droite $d$ est alors simplement la norme $|CP|$.
    \item \textbf{Droite $d_2$ parallèle à $d_1 \equiv A + t\vec{u}$, située à une distance $D$ dans une direction orthogonale $\vec{v} \perp \vec{u}$ (Juillet 2026 Q2a, Septembre 2026 Q2a, Juillet 2023 Q2c) :}
    On norme le vecteur de direction transversale $\vec{v}$ et on le multiplie par la distance $D$ souhaitée pour obtenir un point $B$ de la seconde droite :
    \[
        \boxed{B = A + D\,\frac{\vec{v}}{\|\vec{v}\|} \qquad \text{et} \qquad d_2 \equiv B + t\,\vec{u} \quad (t \in \R)}.
    \]
    \item \textbf{Réflexion d'un rayon lumineux sur un miroir plan $\Pi$ de vecteur normal $\vec{n}$ (Septembre 2025 Q2b) :}
    Lorsqu'un rayon incident de vecteur directeur $\vec{v}_{\text{inc}}$ frappe un miroir plan en un point $B$, sa composante tangentielle (parallèle au miroir) est conservée tandis que sa composante normale $\vec{v}_{\perp} = \frac{\vec{v}_{\text{inc}}\cdot\vec{n}}{\|\vec{n}\|^2}\vec{n}$ est inversée ($-\vec{v}_\perp$). Le vecteur directeur du \textbf{rayon réfléchi} est donc donné par la formule de Householder :
    \[
        \boxed{\vec{v}_{\text{réf}} = \vec{v}_{\text{inc}} - 2\,\frac{\vec{v}_{\text{inc}} \cdot \vec{n}}{\|\vec{n}\|^2}\,\vec{n}}.
    \]
\end{enumerate}
\end{theoreme}

\begin{center}
\begin{tikzpicture}[scale=1.05]
    % Figure 1 : Section sphère - plan
    \begin{scope}[shift={(0,0)}]
        \draw[thick, uclblue] (0,0) circle (1.5);
        \draw[thick, uclaccent, fill=ucllightblue!70] (0,0.75) ellipse (1.3 and 0.35);
        \coordinate (A) at (0,0);
        \coordinate (H) at (0,0.75);
        \coordinate (M) at (1.3,0.75);
        \draw[thick, warnred] (A) -- (H) node[midway, left] {\small $d$};
        \draw[thick, defgreen] (H) -- (M) node[midway, above] {\small $r$};
        \draw[thick, uclblue] (A) -- (M) node[midway, below right] {\small $R$};
        \pic[draw=warnred, angle radius=0.18cm] {right angle = A--H--M};
        \fill (A) circle (1.3pt) node[below left] {\small $A$};
        \fill (H) circle (1.3pt) node[above left] {\small $H$};
        \fill (M) circle (1.3pt) node[right] {\small $M$};
        \node[below, align=center] at (0,-1.75) {\small \textbf{(a)} Section d'une sphère par un plan :\\[-2pt]\small $r = \sqrt{R^2 - d^2}$};
    \end{scope}

    % Figure 2 : Réflexion d'un rayon sur un plan
    \begin{scope}[shift={(6.2,-0.5)}]
        \draw[thick, uclblue, fill=ucllightblue!40] (-2.4,0) -- (1.8,0) -- (2.6,0.9) -- (-1.6,0.9) -- cycle;
        \node[below right, uclblue] at (1.3,0.35) {\small $\Pi$};
        \coordinate (B) at (0,0.45);
        \coordinate (A) at (-1.5, 2.1);
        \coordinate (C) at (1.5, 2.1);
        \coordinate (N) at (0, 2.2);
        \draw[->, very thick, exorange] (A) -- (B) node[midway, above left] {\small $\vec{v}_{\text{inc}}$};
        \draw[->, very thick, defgreen] (B) -- (C) node[midway, above right] {\small $\vec{v}_{\text{réf}}$};
        \draw[->, thick, warnred] (B) -- (N) node[above] {\small $\vec{n}$};
        \draw[dashed, gray] (0,0.45) -- (0,-0.6);
        \pic[draw=exorange, angle radius=0.55cm, "$\theta$", angle eccentricity=1.35] {angle = N--B--A};
        \fill (B) circle (1.5pt) node[below] {\small $B$};
        \node[below, align=center] at (0,-1.25) {\small \textbf{(b)} Réflexion sur un miroir plan :\\[-2pt]\small $\vec{v}_{\text{réf}} = \vec{v}_{\text{inc}} - 2\frac{\vec{v}_{\text{inc}}\cdot\vec{n}}{\|\vec{n}\|^2}\vec{n}$};
    \end{scope}
\end{tikzpicture}
\end{center}

\begin{exemple}[Trois problèmes complets d'examen UCLouvain dans l'espace]
\textbf{1) Robot charpentier et toiture en 3D (UCLouvain -- Juillet 2026, Question 2) :}\\
\emph{Énoncé :} Dans un repère orthonormé $(O;\vec{i},\vec{j},\vec{k})$ (unité : le mètre, sol en $z=0$), l'arête supérieure du premier mur d'une maison est portée par la droite $d_1$ passant par $A(2, -1, 3)$ et de vecteur directeur $\vec{u} = (4, 2, 0)$. Le second mur est parallèle au premier, à la même altitude ($z=3$), situé du côté des $y$ positifs à une distance de $2\sqrt{5}\text{ m}$ du premier mur. Le bâtiment forme un rectangle $ABCD$ de longueur $|AB| = 4\sqrt{5}\text{ m}$ (avec $B \in d_1$ dans le sens de $\vec{u}$, et $C, D \in d_2$ en face de $B, A$).
\begin{enumerate}[label=(\alph*)]
    \item Déterminer une équation paramétrique de la droite $d_2$ portant l'arête supérieure du second mur.
    \item La poutre de faîtage $[EF]$ a ses extrémités $E$ et $F$ situées à la verticale des milieux des segments $[AD]$ et $[BC]$. Sachant que les chevrons $[AE]$ et $[DE]$ ont une longueur maximale de $5\text{ m}$, déterminer la hauteur maximale $H$ du faîtage par rapport au sol ($z=0$) et les coordonnées du point $E$.
    \item Le bras du robot est fixé au point $G(4, -2, 3)$. Déterminer la longueur minimale $L = |GC|$ du bras pour atteindre le point le plus éloigné $C$.
\end{enumerate}
\emph{Résolution :}
\begin{enumerate}[label=(\alph*)]
    \item Cherchons un vecteur horizontal ($z=0$) perpendiculaire à $\vec{u} = (4, 2, 0)$ et orienté vers les $y > 0$ : le vecteur $\vec{v} = (-2, 4, 0)$ vérifie $\vec{u} \cdot \vec{v} = -8 + 8 + 0 = 0$ et sa norme vaut justement $\|\vec{v}\| = \sqrt{(-2)^2 + 4^2} = \sqrt{20} = 2\sqrt{5}\text{ m}$ (qui est exactement la distance entre les deux murs!). Le point $D$ en face de $A$ sur le second mur est donc :
    \[
        D = A + \vec{v} = (2, -1, 3) + (-2, 4, 0) = (0, 3, 3).
    \]
    Comme $d_2$ passe par $D(0,3,3)$ et a le même vecteur directeur $\vec{u} = (4,2,0)$, son équation paramétrique est :
    \[
        \boxed{d_2 \equiv \begin{cases} x = 4t \\ y = 3 + 2t \\ z = 3 \end{cases} \quad (t \in \R)}.
    \]
    \item Le projeté horizontal $M$ de $E$ à l'altitude des murs ($z=3$) est le milieu de $A(2,-1,3)$ et $D(0,3,3)$, soit $M = \left(\frac{2+0}{2}, \frac{-1+3}{2}, 3\right) = (1, 1, 3)$. La demi-distance entre les deux murs vaut $|AM| = \frac{1}{2}|AD| = \sqrt{5}\text{ m}$.\\
    Dans le triangle rectangle formé par $A, M$ et $E$ (de hauteur $h = |ME|$ au-dessus des murs et d'hypoténuse $|AE| = 5\text{ m}$), Pythagore donne :
    \[
        h = \sqrt{5^2 - (\sqrt{5})^2} = \sqrt{20} = 2\sqrt{5}\text{ m}.
    \]
    Comme les murs sont à l'altitude $z=3$, la hauteur maximale par rapport au sol vaut $H = \boxed{3 + 2\sqrt{5}\text{ m}}$ et $\boxed{E(1,\; 1,\; 3 + 2\sqrt{5})}$.
    \item Le vecteur directeur $\vec{u} = (4, 2, 0)$ a pour norme $\|\vec{u}\| = 2\sqrt{5}\text{ m}$. Comme la longueur du mur est $|AB| = 4\sqrt{5} = 2\|\vec{u}\|$, on a $\overrightarrow{AB} = 2\vec{u} = (8, 4, 0)$, d'où $B = A + 2\vec{u} = (10, 3, 3)$ et $C = B + \overrightarrow{AD} = (10, 3, 3) + (-2, 4, 0) = (8, 7, 3)$. La distance entre $G(4, -2, 3)$ et $C(8, 7, 3)$ vaut donc :
    \[
        L = |GC| = \sqrt{(8-4)^2 + (7 - (-2))^2 + (3-3)^2} = \sqrt{16 + 81 + 0} = \boxed{\sqrt{97}\text{ m}}.
    \]
\end{enumerate}

\medskip
\textbf{2) Galeries minières et filon de minerai (UCLouvain -- Septembre 2026, Question 2) :}\\
\emph{Énoncé :} Une galerie principale est modélisée par la droite $d_1$ passant par son point d'entrée $A(0, 0, -5)$ et de vecteur directeur $\vec{u} = (2, 1, 2)$. Une galerie secondaire $d_2$ est creusée parallèlement à $d_1$, à une distance de $2\sqrt{5}\text{ m}$ dans la direction du vecteur $\vec{v} = (1, -2, 0)$. Le filon de minerai est modélisé par le plan $\pi \equiv x + y + z = 4$. (a) Donner une équation paramétrique de la galerie secondaire $d_2$. (b) Déterminer le point de rencontre $C$ de la galerie principale $d_1$ avec le filon $\pi$, et la longueur $L = |AC|$ à creuser.\\
\emph{Résolution :}
\begin{enumerate}[label=(\alph*)]
    \item Vérifions que $\vec{v} = (1, -2, 0)$ est bien orthogonal à $\vec{u} = (2, 1, 2)$ : $\vec{u}\cdot\vec{v} = 2 - 2 + 0 = 0$. Comme $\|\vec{v}\| = \sqrt{1^2+(-2)^2+0^2} = \sqrt{5}$, le déplacement de distance $D = 2\sqrt{5}$ dans la direction de $\vec{v}$ est $2\vec{v} = (2, -4, 0)$. La galerie secondaire passe donc par $B = A + 2\vec{v} = (2, -4, -5)$ avec le vecteur directeur $\vec{u} = (2, 1, 2)$ :
    \[
        \boxed{d_2 \equiv \begin{cases} x = 2 + 2t \\ y = -4 + t \\ z = -5 + 2t \end{cases} \quad (t \in \R)}.
    \]
    \item La galerie principale a pour équations $x = 2t, y = t, z = -5 + 2t$. En substituant dans $\pi \equiv x + y + z = 4$ :
    \[
        2t + t + (-5 + 2t) = 4 \iff 5t = 9 \implies t = \frac{9}{5}.
    \]
    D'où le point d'intersection $\boxed{C\left(\frac{18}{5},\; \frac{9}{5},\; -\frac{7}{5}\right)}$. Comme $\|\vec{u}\| = \sqrt{2^2+1^2+2^2} = \sqrt{9} = 3$, la longueur à creuser depuis $A$ ($t=0$) vaut directement $L = |t|\cdot\|\vec{u}\| = \frac{9}{5} \times 3 = \boxed{\frac{27}{5}\text{ m}}$.
\end{enumerate}

\medskip
\textbf{3) Réflexion d'un rayon lumineux dans un bâtiment passif (UCLouvain -- Septembre 2025, Question 2) :}\\
\emph{Énoncé :} Un rayon lumineux part d'un puits de lumière en $A(2, 0, 4)$ et frappe une paroi vitrée plane $\Pi \equiv x + 2y + 2z = 14$ au point $B(6, 2, 2)$. (a) Calculer le cosinus de l'angle $\theta$ entre le rayon incident et la normale au plan. (b) Déterminer l'équation paramétrique du rayon réfléchi issu de $B$.\\
\emph{Résolution :}
\begin{enumerate}[label=(\alph*)]
    \item Le vecteur incident est $\vec{v}_{\text{inc}} = \overrightarrow{AB} = (6-2,\; 2-0,\; 2-4) = (4, 2, -2)$, de norme $\|\vec{v}_{\text{inc}}\| = \sqrt{16+4+4} = \sqrt{24} = 2\sqrt{6}$. Le vecteur normal au plan $\Pi$ est $\vec{n} = (1, 2, 2)$, de norme $\|\vec{n}\| = \sqrt{1+4+4} = 3$. Leur produit scalaire vaut $\vec{v}_{\text{inc}} \cdot \vec{n} = 4(1) + 2(2) + (-2)(2) = 4$. Dès lors :
    \[
        \cos\theta = \frac{\vec{v}_{\text{inc}} \cdot \vec{n}}{\|\vec{v}_{\text{inc}}\|\,\|\vec{n}\|} = \frac{4}{3\sqrt{24}} = \boxed{\frac{2}{3\sqrt{6}}}.
    \]
    \item Par la formule de réflexion :
    \begin{align*}
        \vec{v}_{\text{réf}} &= \vec{v}_{\text{inc}} - 2\,\frac{\vec{v}_{\text{inc}}\cdot\vec{n}}{\|\vec{n}\|^2}\,\vec{n} = (4, 2, -2) - 2\cdot\frac{4}{9}\cdot(1, 2, 2) \\
        &= \left(4 - \frac{8}{9},\; 2 - \frac{16}{9},\; -2 - \frac{16}{9}\right) = \frac{1}{9}(28,\; 2,\; -34).
    \end{align*}
    Le rayon réfléchi partant de $B(6, 2, 2)$ a donc pour équation paramétrique (avec $s = \frac{t}{9}$) :
    \[
        \boxed{\begin{cases} x = 6 + \frac{28}{9}t \\ y = 2 + \frac{2}{9}t \\ z = 2 - \frac{34}{9}t \end{cases} \quad \text{ou} \quad \begin{cases} x = 6 + 28s \\ y = 2 + 2s \\ z = 2 - 34s \end{cases}}.
    \]
\end{enumerate}
\end{exemple}

% ----------------------------------------------------------------------------
\section{Exercices du Chapitre 7}
% ----------------------------------------------------------------------------

\emph{Les corrigés complets de tous ces exercices se trouvent au Chapitre 9, en fin d'ouvrage.}

\begin{enumerate}[label=\textbf{Exercice 7.\arabic* --}, leftmargin=*]

    \item \textbf{(UCLouvain -- Juillet 2025, Q2 : Éboulement rocheux et condition de non-décrochage)}
    La pente d'un versant montagneux est modélisée par le plan $\Pi \equiv z = -0,6x + 0,8y + 85$. Un rocher se détache du point $A(10, 20, 100)$ et suit une trajectoire rectiligne de vecteur directeur $\vec{v} = (4, -2, -4)$.
    \begin{enumerate}[label=(\alph*)]
        \item Un village se situe à l'altitude $z = 80$. Déterminez les coordonnées du point $M$ où la trajectoire du rocher atteint cette altitude, ainsi que la distance $|AM|$ parcourue par le rocher.
        \item Pour une trajectoire quelconque de vecteur directeur $\vec{w} = (a, b, c)$, quelle relation algébrique doivent vérifier $a, b, c$ pour que la trajectoire suive le plan $\Pi$ sans décrochage ?
        \item Observe-t-on un décrochage (décollage par rapport au versant $\Pi$) pour la trajectoire observée $\vec{v} = (4, -2, -4)$ ?
    \end{enumerate}

    \item \textbf{(UCLouvain -- Septembre 2024, Q2 : Projeté orthogonal et point équidistant sur une droite)}
    On considère dans l'espace les trois points $A(-7, 10, 5)$, $B(2, 4, 2)$ et $C(3, 10, 7)$.
    \begin{enumerate}[label=(\alph*)]
        \item Déterminez les coordonnées du point $P$ appartenant à la droite $(AB)$ tel que le vecteur $\overrightarrow{CP}$ soit perpendiculaire au vecteur $\overrightarrow{AB}$.
        \item Déterminez les coordonnées du point $Q$ appartenant à la droite $(BC)$ qui est équidistant des points $A$ et $C$.
    \end{enumerate}

    \item \textbf{(UCLouvain -- Septembre 2023, Q2 : Point de percée et distance entre deux droites parallèles)}
    Dans l'espace orthonormé $Oxyz$, on considère le plan $\Pi \equiv -2x + y - 3z + 4 = 0$ et le point $A(5, -1, 7)$.
    \begin{enumerate}[label=(\alph*)]
        \item Donnez des équations paramétriques de la droite $d_1$ passant par $A$ et perpendiculaire au plan $\Pi$, ainsi que les coordonnées de son point de percée $B$ dans le plan $\Pi$.
        \item Calculez la distance qui sépare la droite $d_1$ de la droite $d_2$ parallèle à $d_1$ et passant par le point $C(-6, 5, -7)$.
    \end{enumerate}

    \item \textbf{(UCLouvain -- Juillet 2023, Q2 : Droites perpendiculaires, plan et distance dans un plan)}
    Soit la droite $d_1 \equiv \begin{cases} x = t - 3 \\ y = 3t - 4 \\ z = -t + 5 \end{cases}$ ($t \in \R$) et le point $A(1, 1, 2)$.
    \begin{enumerate}[label=(\alph*)]
        \item Déterminez une équation paramétrique de la droite $d_2$ contenant $A$ et perpendiculaire à $d_1$ en leur point d'intersection $B$ (dont on déterminera les coordonnées).
        \item Donnez une équation cartésienne du plan $\mathcal{P}$ contenant les droites $d_1$ et $d_2$.
        \item Déterminez les coordonnées des deux points $C$ possibles situés à l'intersection du plan $\mathcal{P}$ et du plan $\mathcal{S} \equiv x = 2z$, tels que la distance de $C$ à la droite $d_1$ vaille $\sqrt{6}$.
    \end{enumerate}

    \item \textbf{(UCLouvain -- Juillet 2022, Q1 : Sphère, section par un plan et plans inclinés)}
    Soit le point $A(7, -2, 8)$ (coordonnées en mètres) et le plan $p \equiv 3x - 6y + 2z = -7$.
    \begin{enumerate}[label=(\alph*)]
        \item Donnez l'inéquation cartésienne de la boule ouverte $\mathcal{L}$ rassemblant tous les points situés à moins de $10\text{ m}$ du point $A$.
        \item Calculez la distance du centre $A$ au plan $p$, et déduisez-en l'aire exacte du disque obtenu par l'intersection de $\mathcal{L}$ avec le plan $p$.
        \item Déterminez l'équation cartésienne $x + ay + bz = c$ du plan $q$ contenant $A$, perpendiculaire au plan $p$, formant un angle de $45\degree$ avec tout plan horizontal ($z = \text{cte}$), et tel que le point $B(5,5,5)$ se situe en dessous de $q$.
    \end{enumerate}

    \item \textbf{(UCLouvain -- Juillet 2021, Q4 : Descente en rappel d'une alpiniste)}
    Une alpiniste descend en rappel le long d'une trajectoire rectiligne ($t$ en minutes, coordonnées en mètres) définie par :
    \[
        x(t) = -t - 50, \qquad y(t) = 6t - 100, \qquad z(t) = -18t + 144.
    \]
    \begin{enumerate}[label=(\alph*)]
        \item Quelle est la distance totale parcourue par l'alpiniste entre $t = 0$ et l'instant où elle touche le sol ($z = 0$) ?
        \item Déterminez un vecteur directeur d'une droite $h$ qui est à la fois horizontale et perpendiculaire à la trajectoire de l'alpiniste.
        \item La paroi de la montagne est le plan qui contient la trajectoire de l'alpiniste et qui est parallèle à la droite $h$. Donnez une équation cartésienne de ce plan.
    \end{enumerate}

    \item \textbf{(UCLouvain -- Septembre 2021, Q2 : Les sommets d'une pyramide)}
    Les points $A, B, C, D$ sont les quatre sommets d'une pyramide tels que :
    \begin{itemize}
        \item la base $ABC$ est contenue dans le plan horizontal $z = 0$, avec $B(-3, -1, 0)$ et $C(4, 0, 0)$ ;
        \item le sommet $D$ appartient au plan horizontal $z = 5$ et est équidistant des trois sommets $A, B$ et $C$ ($|DA| = |DB| = |DC|$) ;
        \item l'une des faces latérales est contenue dans le plan $\Pi \equiv -3x + y + z - 8 = 0$.
    \end{itemize}
    \begin{enumerate}[label=(\alph*)]
        \item Donnez l'équation cartésienne du plan médiateur $\mathcal{L}$ rassemblant tous les points équidistants de $B$ et $C$.
        \item En remarquant que $C \notin \Pi$ (donc $D \in \Pi$ et $A \in \Pi$), calculez les coordonnées du sommet $D$.
        \item Déduisez-en les coordonnées du sommet $A$.
    \end{enumerate}

    \item \textbf{(Symétrique d'un point par rapport à un plan et distance point--plan)}
    Soit le plan $\Pi \equiv 2x - y + 2z - 6 = 0$ et le point $A(5, -3, 7)$.
    \begin{enumerate}[label=(\alph*)]
        \item Calculez la distance $d(A, \Pi)$ à l'aide de la formule de distance point--plan.
        \item Déterminez les coordonnées du projeté orthogonal $H$ de $A$ sur $\Pi$, ainsi que celles du point $A'$ symétrique de $A$ par rapport au plan $\Pi$.
    \end{enumerate}

    \item \textbf{(Intersection de deux plans et distance d'un point à une droite)}
    On considère les deux plans $\Pi_1 \equiv x + y - z = 1$ et $\Pi_2 \equiv 2x - y + z = 2$, ainsi que le point $P(3, 1, 0)$.
    \begin{enumerate}[label=(\alph*)]
        \item Déterminez une représentation paramétrique de la droite $d = \Pi_1 \cap \Pi_2$.
        \item Déterminez le projeté orthogonal $H$ de $P$ sur $d$ et la distance $d(P, d)$.
    \end{enumerate}

    \item \textbf{(Sphère tangente à un plan et plan tangent à une sphère)}
    \begin{enumerate}[label=(\alph*)]
        \item Déterminez l'équation cartésienne de la sphère $\mathcal{S}$ de centre $\Omega(1, -2, 3)$ qui est tangente au plan $\Pi \equiv 2x + 3y - 6z + 5 = 0$.
        \item Soit la sphère $\mathcal{S}' \equiv (x-2)^2 + (y-1)^2 + (z+1)^2 = 14$. Vérifiez que le point $T(3, 3, 2)$ appartient à $\mathcal{S}'$ et déterminez l'équation cartésienne du plan tangent à $\mathcal{S}'$ en $T$.
    \end{enumerate}

\end{enumerate}

% ============================================================================
% CHAPITRE 8 : FORMULAIRE SYNTHÉTIQUE ET AIDE-MÉMOIRE (VERSION IMPRIMABLE 4 PAGES)
% ============================================================================
\cleardoublepage
\newgeometry{top=1.3cm,bottom=1.4cm,left=1.4cm,right=1.4cm}
\refstepcounter{chapter}
\addcontentsline{toc}{chapter}{\protect\numberline{\thechapter}Formulaire Synthétique et Aide-Mémoire (Imprimable -- 4 pages)}
\chaptermark{Formulaire Synthétique et Aide-Mémoire}

% ============================================================================
% PAGE 1/4 : FORMULAIRE DE TRIGONOMÉTRIE (1 PAGE A4 IMPRIMABLE)
% ============================================================================
\thispagestyle{plain}
\refstepcounter{section}
\addcontentsline{toc}{section}{\protect\numberline{\thesection}Formulaire de Trigonométrie}

\begin{center}
    {\LARGE\bfseries\color{uclblue} Formulaire de Trigonométrie}\\[2pt]
    {\small\color{darkgray} Examen d'entrée Ingénieur Civil -- UCLouvain (Fiche récapitulative imprimable 1/4)}
    \vspace{2pt}
    \hrule height 1pt
\end{center}
\vspace{-6pt}

{\fontsize{9.5pt}{11.5pt}\selectfont
\setlength{\columnsep}{18pt}
\begin{multicols}{2}

\noindent{\large\bfseries\color{uclblue} 1. Identités Fondamentales}
\begin{itemize}[leftmargin=1.2em, itemsep=1.5pt, topsep=2pt]
    \item $\cos^2\alpha + \sin^2\alpha = 1$
    \item $\tan\alpha = \dfrac{\sin\alpha}{\cos\alpha} \;\left(\alpha \neq \frac{\pi}{2}+k\pi\right)$, \quad $\cot\alpha = \dfrac{\cos\alpha}{\sin\alpha} \;(\alpha \neq k\pi)$
    \item $1 + \tan^2\alpha = \dfrac{1}{\cos^2\alpha}$ \quad et \quad $1 + \cot^2\alpha = \dfrac{1}{\sin^2\alpha}$
\end{itemize}

\vspace{3pt}
\noindent{\large\bfseries\color{uclblue} 2. Angles Associés}
\begin{itemize}[leftmargin=1.2em, itemsep=1.5pt, topsep=2pt]
    \item \textbf{Opposés ($-\alpha$) :} $\cos(-\alpha) = \cos\alpha$, \; $\sin(-\alpha) = -\sin\alpha$, \; $\tan(-\alpha) = -\tan\alpha$
    \item \textbf{Supplémentaires ($\pi-\alpha$) :} $\cos(\pi-\alpha) = -\cos\alpha$, \; $\sin(\pi-\alpha) = \sin\alpha$, \; $\tan(\pi-\alpha) = -\tan\alpha$
    \item \textbf{Antisupplémentaires ($\pi+\alpha$) :} $\cos(\pi+\alpha) = -\cos\alpha$, \; $\sin(\pi+\alpha) = -\sin\alpha$, \; $\tan(\pi+\alpha) = \tan\alpha$
    \item \textbf{Complémentaires ($\frac{\pi}{2}-\alpha$) :} $\cos\left(\frac{\pi}{2}-\alpha\right) = \sin\alpha$, \; $\sin\left(\frac{\pi}{2}-\alpha\right) = \cos\alpha$, \; $\tan\left(\frac{\pi}{2}-\alpha\right) = \cot\alpha$
    \item \textbf{Anticomplémentaires ($\frac{\pi}{2}+\alpha$) :} $\cos\left(\frac{\pi}{2}+\alpha\right) = -\sin\alpha$, \; $\sin\left(\frac{\pi}{2}+\alpha\right) = \cos\alpha$, \; $\tan\left(\frac{\pi}{2}+\alpha\right) = -\cot\alpha$
\end{itemize}

\vspace{3pt}
\noindent{\large\bfseries\color{uclblue} 3. Formules d'Addition}
\begin{itemize}[leftmargin=1.2em, itemsep=1.5pt, topsep=2pt]
    \item $\cos(\alpha \pm \beta) = \cos\alpha\cos\beta \mp \sin\alpha\sin\beta$
    \item $\sin(\alpha \pm \beta) = \sin\alpha\cos\beta \pm \cos\alpha\sin\beta$
    \item $\tan(\alpha \pm \beta) = \dfrac{\tan\alpha \pm \tan\beta}{1 \mp \tan\alpha\tan\beta}$
\end{itemize}

\vspace{3pt}
\noindent{\large\bfseries\color{uclblue} 4. Formules de Duplication, Carnot et Angle Demi}
\begin{itemize}[leftmargin=1.2em, itemsep=1.5pt, topsep=2pt]
    \item $\cos(2\alpha) = \cos^2\alpha - \sin^2\alpha = 2\cos^2\alpha - 1 = 1 - 2\sin^2\alpha$
    \item $\sin(2\alpha) = 2\sin\alpha\cos\alpha$ \quad et \quad $\tan(2\alpha) = \dfrac{2\tan\alpha}{1 - \tan^2\alpha}$
    \item \textbf{Carnot :} $\cos^2\alpha = \dfrac{1+\cos(2\alpha)}{2}$, \quad $\sin^2\alpha = \dfrac{1-\cos(2\alpha)}{2}$
    \item \textbf{Arc moitié ($t = \tan(\alpha/2)$) :}\\
    $\cos\alpha = \dfrac{1-t^2}{1+t^2}, \quad \sin\alpha = \dfrac{2t}{1+t^2}, \quad \tan\alpha = \dfrac{2t}{1-t^2}$
\end{itemize}

\vspace{3pt}
\noindent{\large\bfseries\color{uclblue} 5. Formules de Simpson (Factorisation)}
\begin{itemize}[leftmargin=1.2em, itemsep=2pt, topsep=2pt]
    \item $\cos p + \cos q = 2\cos\left(\dfrac{p+q}{2}\right)\cos\left(\dfrac{p-q}{2}\right)$
    \item $\cos p - \cos q = -2\sin\left(\dfrac{p+q}{2}\right)\sin\left(\dfrac{p-q}{2}\right)$
    \item $\sin p + \sin q = 2\sin\left(\dfrac{p+q}{2}\right)\cos\left(\dfrac{p-q}{2}\right)$
    \item $\sin p - \sin q = 2\sin\left(\dfrac{p-q}{2}\right)\cos\left(\dfrac{p+q}{2}\right)$
    \item $\tan p \pm \tan q = \dfrac{\sin(p \pm q)}{\cos p \cos q}$
\end{itemize}

\columnbreak

\noindent{\large\bfseries\color{uclblue} 6. Arcs, Secteurs et Triangles}

\vspace{2pt}
\noindent\textbf{Cercle de rayon $R$ ($\theta$ en radians) :}
\begin{itemize}[leftmargin=1.2em, itemsep=1.5pt, topsep=2pt]
    \item Longueur d'un arc : $L = R\theta$ \quad (Corde : $c = 2R\sin(\theta/2)$)
    \item Aire d'un secteur : $A_{\text{secteur}} = \dfrac{1}{2}R^2\theta$
    \item Aire d'un segment : $A_{\text{segment}} = \dfrac{1}{2}R^2(\theta - \sin\theta)$
\end{itemize}

\vspace{2pt}
\noindent\textbf{Triangle Rectangle (SOH-CAH-TOA) :}
\begin{itemize}[leftmargin=1.2em, itemsep=1.5pt, topsep=2pt]
    \item $\sin\theta = \dfrac{\text{Opposé}}{\text{Hypoténuse}}$
    \item $\cos\theta = \dfrac{\text{Adjacent}}{\text{Hypoténuse}}$
    \item $\tan\theta = \dfrac{\text{Opposé}}{\text{Adjacent}}$
\end{itemize}

\vspace{2pt}
\noindent\textbf{Triangle Quelconque ($A+B+C = \pi$) :}
\begin{itemize}[leftmargin=1.2em, itemsep=2pt, topsep=2pt]
    \item \textbf{Loi des sinus :} $\dfrac{a}{\sin A} = \dfrac{b}{\sin B} = \dfrac{c}{\sin C} = 2R$
    \item \textbf{Loi des cosinus (Al-Kashi) :}\\
    $a^2 = b^2 + c^2 - 2bc\cos A$
    \item \textbf{Aire :} $S = \dfrac{1}{2}bc\sin A = p\cdot r = \sqrt{p(p-a)(p-b)(p-c)}$\\
    avec $p = \dfrac{a+b+c}{2}$ (demi-périmètre) et $r$ rayon inscrit.
\end{itemize}

\vspace{4pt}
\noindent{\large\bfseries\color{uclblue} 7. Astuce d'Euler (Complexes)}
La formule d'Euler $e^{i\alpha} = \cos\alpha + i\sin\alpha$ permet de retrouver les formules d'addition, de duplication et de triplication en développant :
\[
    e^{i(\alpha+\beta)} = e^{i\alpha}e^{i\beta} \qquad \text{et} \qquad e^{in\alpha} = (e^{i\alpha})^n
\]

\vspace{2pt}
\noindent{\large\bfseries\color{uclblue} 8. Angles Remarquables}
\begin{center}
\vspace{-2pt}
\renewcommand{\arraystretch}{1.45}
\setlength{\tabcolsep}{4.5pt}
\begin{tabular}{|c|c|c|c|c|c|c|}
\hline
$\theta$ & $0$ & $\frac{\pi}{12}$ & $\frac{\pi}{6}$ & $\frac{\pi}{4}$ & $\frac{\pi}{3}$ & $\frac{\pi}{2}$ \\ \hline
$\sin$ & $0$ & $\frac{\sqrt{6}-\sqrt{2}}{4}$ & $\frac{1}{2}$ & $\frac{\sqrt{2}}{2}$ & $\frac{\sqrt{3}}{2}$ & $1$ \\ \hline
$\cos$ & $1$ & $\frac{\sqrt{6}+\sqrt{2}}{4}$ & $\frac{\sqrt{3}}{2}$ & $\frac{\sqrt{2}}{2}$ & $\frac{1}{2}$ & $0$ \\ \hline
$\tan$ & $0$ & $2-\sqrt{3}$ & $\frac{\sqrt{3}}{3}$ & $1$ & $\sqrt{3}$ & $-$ \\ \hline
\end{tabular}
\end{center}
\vspace{-2pt}
\noindent\emph{Note :} Pour $\theta = \frac{\pi}{8}$ ($22,5\degree$), $\tan\left(\frac{\pi}{8}\right) = \sqrt{2}-1$.

\end{multicols}
}

\newpage
% ============================================================================
% PAGE 2/4 : ÉQUATIONS TRIGONOMÉTRIQUES : TECHNIQUES ET EXEMPLES (1 PAGE A4)
% ============================================================================
\thispagestyle{plain}
\refstepcounter{section}
\addcontentsline{toc}{section}{\protect\numberline{\thesection}Équations Trigonométriques : Techniques et Exemples}

\begin{center}
    {\LARGE\bfseries\color{uclblue} Équations Trigonométriques : Techniques et Exemples}\\[2pt]
    {\small\color{darkgray} Examen d'entrée Ingénieur Civil -- UCLouvain (Fiche récapitulative imprimable 2/4)}
    \vspace{2pt}
    \hrule height 1pt
\end{center}
\vspace{-6pt}

{\fontsize{9.5pt}{11.5pt}\selectfont
\setlength{\columnsep}{18pt}
\begin{multicols}{2}

\noindent{\large\bfseries\color{uclblue} 1. Équations Élémentaires}

\vspace{2pt}
\noindent Règles de base ($k \in \mathbb{Z}$) :
\begin{itemize}[leftmargin=1.2em, itemsep=2pt, topsep=2pt]
    \item $\sin(\alpha) = \sin(\beta) \iff$\\
    $\alpha = \beta + 2k\pi \quad\text{ou}\quad \alpha = \pi - \beta + 2k\pi$
    \item $\cos(\alpha) = \cos(\beta) \iff$\\
    $\alpha = \pm\beta + 2k\pi$
    \item $\tan(\alpha) = \tan(\beta) \iff$\\
    $\alpha = \beta + k\pi \quad \left(\text{CE : } \alpha,\beta \neq \frac{\pi}{2}+m\pi\right)$
\end{itemize}
\textit{Astuce :} Pour $\sin(\alpha) = \cos(\beta)$, on transforme en $\sin(\alpha) = \sin\left(\frac{\pi}{2} - \beta\right)$ ou $\cos\left(\frac{\pi}{2} - \alpha\right) = \cos(\beta)$.

\vspace{6pt}
\noindent{\large\bfseries\color{uclblue} 2. Polynômes et Substitution}

\vspace{2pt}
\noindent Pour des équations avec des puissances paires, on peut remplacer et poser une variable.\\
\textbf{Exemple (UCLouvain, Juillet 2025) :}\\
$\sin^6 x - \cos^6 x = 3\sin^2 x \cos^2 x$
\begin{itemize}[leftmargin=1.2em, itemsep=2pt, topsep=2pt]
    \item Poser $y = \sin^2 x$, ce qui donne $\cos^2 x = 1 - y$.
    \item Remplacer : $y^3 - (1-y)^3 = 3y(1-y)$.
    \item Simplifier pour obtenir une équation algébrique en $y$ :
    \[
        2y^3 - 1 = 0 \implies \sin^2 x = \sqrt[3]{\frac{1}{2}}
    \]
\end{itemize}
\textit{Astuce :} Utiliser l'identité $a^3 \mp b^3 = (a\mp b)(a^2 \pm ab + b^2)$ pour factoriser quand c'est pertinent, ou les identités classiques $\sin^4 x + \cos^4 x = 1 - 2\sin^2 x\cos^2 x$ et $\sin^6 x + \cos^6 x = 1 - 3\sin^2 x\cos^2 x$.

\vspace{6pt}
\noindent{\large\bfseries\color{uclblue} 3. Linéaire en $\sin x, \cos x$}

\vspace{2pt}
\noindent Forme : $a\sin x + b\cos x = c \quad (\text{solutions ssi } c^2 \le a^2+b^2)$
\begin{itemize}[leftmargin=1.2em, itemsep=2pt, topsep=2pt]
    \item Diviser par $a$ : $\sin x + \dfrac{b}{a}\cos x = \dfrac{c}{a}$.
    \item Poser $\tan(y) = \dfrac{b}{a}$. L'équation devient :
    \[
        \sin(x)\cos(y) + \sin(y)\cos(x) = \frac{c}{a}\cos(y)
    \]
    \item Ce qui équivaut, avec la formule d'addition, à :
    \[
        \sin(x+y) = \frac{c}{a}\cos(y)
    \]
\end{itemize}

\columnbreak

\noindent{\large\bfseries\color{uclblue} 4. Homogène en $\sin x, \cos x$}

\vspace{2pt}
\noindent Polynôme en $\sin, \cos$ où la somme des exposants de chaque terme est constante.\\
\textbf{Exemple (UCLouvain, Juillet 2018) :}\\
$\sin^3 x + 2\cos^3 x = 3\sin^2 x \cos x$ (Degré 3)
\begin{itemize}[leftmargin=1.2em, itemsep=2pt, topsep=2pt]
    \item Vérifier qu'on ne peut pas mettre de terme en évidence.
    \item Vérifier si $\cos x = 0$ est solution, puis diviser les deux membres par $\cos^k x$ (ici $k=3$) pour se ramener à une équation en $\tan x$.
    \item On obtient l'équation algébrique :
    \[
        \tan^3 x + 2 = 3\tan^2 x
    \]
\end{itemize}

\vspace{6pt}
\noindent{\large\bfseries\color{uclblue} 5. Équation Symétrique}

\vspace{2pt}
\noindent L'équation ne change pas si on inverse $\sin x$ et $\cos x$.\\
\textbf{Exemple (UCLouvain, Juillet 2022) :}\\
$\cos x + \sin x = \sqrt{2}\cos x \sin x$
\begin{itemize}[leftmargin=1.2em, itemsep=2pt, topsep=2pt]
    \item Technique clé : poser $\mathbf{x = y + \dfrac{\pi}{4}}$.
    \item On sait via les formules d'addition que :
    \[
        \sin x = \frac{\sqrt{2}}{2}(\cos y + \sin y), \quad \cos x = \frac{\sqrt{2}}{2}(\cos y - \sin y)
    \]
    \item Remplacer dans l'équation ($\sin x + \cos x = \sqrt{2}\cos y$ et $\sin x\cos x = \cos^2 y - \frac{1}{2}$) pour obtenir une équation du second degré en $\cos y$ facile à résoudre.
\end{itemize}

\vspace{6pt}
\noindent{\large\bfseries\color{uclblue} 6. Formules de Simpson}

\vspace{2pt}
\noindent Pour les sommes/différences de sinus et cosinus avec des multiples ($2x, 3x\dots$).\\
\textbf{Exemple (UCLouvain, Septembre 2019) :}\\
$\sin(x) - 2\sin(3x) + \sin(5x) = -\cos(2x) + \cos(4x)$
\begin{itemize}[leftmargin=1.2em, itemsep=2pt, topsep=2pt]
    \item Regrouper stratégiquement les termes :
    \[
        \sin(5x) + \sin(x) = 2\sin(3x)\cos(2x)
    \]
    \item Regrouper :
    \[
        \cos(4x) - \cos(2x) = -2\sin(3x)\sin(x)
    \]
    \item L'équation devient factorisable par $\sin(3x)$ :
    \[
        2\sin(3x)\bigl(\cos(2x) - 1 + \sin(x)\bigr) = 0
    \]
    donc soit $\sin(3x) = 0$, soit l'autre facteur est nul.
\end{itemize}

\end{multicols}
}

\newpage
% ============================================================================
% PAGE 3/4 : GÉOMÉTRIE SYNTHÉTIQUE & ANALYTIQUE PLANE (CONIQUES) (1 PAGE A4)
% ============================================================================
\thispagestyle{plain}
\refstepcounter{section}
\addcontentsline{toc}{section}{\protect\numberline{\thesection}Formulaire de Géométrie Synthétique et Analytique Plane (Coniques)}

\begin{center}
    {\LARGE\bfseries\color{uclblue} Géométrie Synthétique \& Analytique Plane (Coniques)}\\[2pt]
    {\small\color{darkgray} Examen d'entrée Ingénieur Civil -- UCLouvain (Fiche récapitulative imprimable 3/4)}
    \vspace{2pt}
    \hrule height 1pt
\end{center}
\vspace{-6pt}

{\fontsize{9.2pt}{11.2pt}\selectfont
\setlength{\columnsep}{18pt}
\begin{multicols}{2}

\noindent{\large\bfseries\color{uclblue} 1. Géométrie Synthétique du Cercle}
\begin{itemize}[leftmargin=1.2em, itemsep=1.5pt, topsep=2pt]
    \item \textbf{Médiatrice d'une corde $[AB]$ :} La droite reliant le centre $O$ au milieu $M$ de $[AB]$ est perpendiculaire à $[AB]$ ($|OM|^2 + |AM|^2 = r^2$).
    \item \textbf{Angle inscrit et angle au centre} interceptant $\wideparen{AB}$ :
    \[
        |\widehat{AMB}| = \frac{1}{2}|\widehat{AOB}|
    \]
    \item \textbf{Demi-cercle :} Si $[AB]$ est un diamètre et $M \in \mathcal{C}$, alors $\widehat{AMB} = 90\degree$.
    \item \textbf{Quadrilatère inscriptible $ABCD$ :} Les angles opposés sont supplémentaires ($|\widehat{A}| + |\widehat{C}| = 180\degree = |\widehat{B}| + |\widehat{D}|$).
    \item \textbf{Puissance d'un point $P$ par rapport à $\mathcal{C}(O,r)$ :}
    \begin{itemize}[leftmargin=1em, itemsep=1pt, topsep=1pt]
        \item $P$ extérieur : $|PA|\cdot|PB| = |PT|^2 = |PO|^2 - r^2$
        \item $P$ intérieur : $|PA|\cdot|PB| = r^2 - |PO|^2$
    \end{itemize}
\end{itemize}

\vspace{3pt}
\noindent{\large\bfseries\color{uclblue} 2. Relations Métriques \& Aires}
\begin{itemize}[leftmargin=1.2em, itemsep=1.5pt, topsep=2pt]
    \item \textbf{Triangle rectangle en $C$} (hauteur $h = |CH|$ sur l'hypoténuse $c = p + q$) : $h^2 = p\cdot q$, \; $b^2 = c\cdot p$, \; $a^2 = c\cdot q$.
    \item \textbf{Figures semblables} de rapport $\rho = \frac{L_2}{L_1}$ : $\dfrac{\text{Aire}_2}{\text{Aire}_1} = \rho^2$.
    \item \textbf{Triangles de même hauteur :} $\dfrac{\text{Aire}_1}{\text{Aire}_2} = \dfrac{\text{base}_1}{\text{base}_2}$.
    \item \textbf{Identité du parallélogramme :}\\
    $|AC|^2 + |BD|^2 = 2\bigl(|AB|^2 + |BC|^2\bigr)$.
    \item \textbf{Solides :} $V_{\text{cyl}} = \pi R^2 h$, \; $V_{\text{cône}} = \frac{1}{3}\pi R^2 h$, \; $V_{\text{pyr}} = \frac{1}{3}\mathcal{B}h$, \; $V_{\text{sphère}} = \frac{4}{3}\pi r^3$, \; $A_{\text{sphère}} = 4\pi r^2$.
\end{itemize}

\vspace{3pt}
\noindent{\large\bfseries\color{uclblue} 3. Vecteurs, Droites et Cercles dans $\R^2$}
\begin{itemize}[leftmargin=1.2em, itemsep=1.5pt, topsep=2pt]
    \item \textbf{Produit scalaire :} $\vec{u}\cdot\vec{v} = u_x v_x + u_y v_y = \|\vec{u}\|\|\vec{v}\|\cos\theta$.
    \item \textbf{Orthogonalité :} $\vec{u} \perp \vec{v} \iff u_x v_x + u_y v_y = 0$.\\
    Vecteur orthogonal à $(a,b)$ : $(-b, a)$.
    \item \textbf{Droite $d \equiv ax + by + c = 0$ :} normale $\vec{n} = (a,b)$, vecteur directeur $\vec{d} = (-b,a)$, pente $m = -\frac{a}{b}$.\\
    Perpendicularité de $y = m_1 x + p_1$ et $y = m_2 x + p_2$ : $m_1 m_2 = -1$.
    \item \textbf{Distance d'un point $P(x_0,y_0)$ à $d \equiv ax+by+c=0$ :}
    \[
        d(P, d) = \frac{|ax_0 + by_0 + c|}{\sqrt{a^2 + b^2}}
    \]
    \item \textbf{Cercle de centre $\Omega(x_0,y_0)$ et rayon $R$ :}\\
    $(x - x_0)^2 + (y - y_0)^2 = R^2$.\\
    Le centre du cercle circonscrit à $ABC$ est l'intersection des médiatrices de $[AB]$ et $[AC]$.
\end{itemize}

\columnbreak

\noindent{\large\bfseries\color{uclblue} 4. Les Coniques dans le Plan}

\vspace{2pt}
\noindent\textbf{Définition par foyer $F$, directrice $d$ et excentricité $e$ :}
\[
    \frac{d(M, F)}{d(M, d)} = e \quad \begin{cases} 0 < e < 1 & \text{: Ellipse} \\ e = 1 & \text{: Parabole} \\ e > 1 & \text{: Hyperbole} \end{cases}
\]

\vspace{2pt}
\noindent\textbf{A. Ellipse de centre $C(x_0, y_0)$ :}
\[
    \frac{(x-x_0)^2}{a^2} + \frac{(y-y_0)^2}{b^2} = 1
\]
\begin{itemize}[leftmargin=1.2em, itemsep=1.5pt, topsep=2pt]
    \item \textbf{Si $a > b$ (grand axe horizontal $2a$) :}\\
    $c = \sqrt{a^2 - b^2}$, \; foyers $F_{1,2}(x_0 \pm c,\; y_0)$, \; $e = \dfrac{c}{a} < 1$,\\
    directrices $x = x_0 \pm \dfrac{a^2}{c}$. \quad ($|MF_1| + |MF_2| = 2a$)
    \item \textbf{Si $b > a$ (grand axe vertical $2b$) :}\\
    $c = \sqrt{b^2 - a^2}$, \; foyers $F_{1,2}(x_0,\; y_0 \pm c)$, \; $e = \dfrac{c}{b} < 1$,\\
    directrices $y = y_0 \pm \dfrac{b^2}{c}$.
\end{itemize}

\vspace{3pt}
\noindent\textbf{B. Hyperbole de centre $C(x_0, y_0)$ (axe horizontal) :}
\[
    \frac{(x-x_0)^2}{a^2} - \frac{(y-y_0)^2}{b^2} = 1
\]
\begin{itemize}[leftmargin=1.2em, itemsep=1.5pt, topsep=2pt]
    \item Sommets : $S_{1,2}(x_0 \pm a,\; y_0)$.
    \item Distance focale : $\mathbf{c = \sqrt{a^2 + b^2}}$, \; foyers $F_{1,2}(x_0 \pm c,\; y_0)$.
    \item Excentricité : $e = \dfrac{c}{a} > 1$, \; directrices : $x = x_0 \pm \dfrac{a^2}{c}$.
    \item \textbf{Asymptotes :} $y - y_0 = \pm \dfrac{b}{a}(x - x_0)$.
\end{itemize}

\vspace{3pt}
\noindent\textbf{C. Parabole de sommet $S(x_0, y_0)$ :}
\begin{itemize}[leftmargin=1.2em, itemsep=1.5pt, topsep=2pt]
    \item \textbf{Axe vertical :} $(x-x_0)^2 = 2p(y-y_0) \iff y - y_0 = a(x-x_0)^2$ (avec $p = \frac{1}{2a}$).\\
    Foyer $F\left(x_0,\; y_0 + \frac{p}{2}\right) = \left(x_0,\; y_0 + \frac{1}{4a}\right)$,\\
    Directrice $d \equiv y = y_0 - \frac{p}{2} = y_0 - \frac{1}{4a}$.
    \item \textbf{Axe horizontal :} $(y-y_0)^2 = 2p(x-x_0)$, foyer $F\left(x_0 + \frac{p}{2}, y_0\right)$, directrice $x = x_0 - \frac{p}{2}$.
\end{itemize}

\vspace{3pt}
\noindent{\large\bfseries\color{uclblue} 5. Tangentes à une Conique}
\begin{itemize}[leftmargin=1.2em, itemsep=1.5pt, topsep=2pt]
    \item \textbf{Condition de tangence d'une droite $y = mx+p$ :} Substituer $y$ dans l'équation de la conique et imposer que le discriminant de l'équation du $2^{\text{e}}$ degré en $x$ soit nul : $\mathbf{\Delta = B^2 - 4AC = 0}$.
    \item \textbf{Vecteur tangent à $y = f(x)$ en $x_1$ :} $\vec{t} = \bigl(1,\; f'(x_1)\bigr)$.
\end{itemize}

\end{multicols}
}

\newpage
% ============================================================================
% PAGE 4/4 : GÉOMÉTRIE ANALYTIQUE ET VECTORIELLE DANS L'ESPACE (1 PAGE A4)
% ============================================================================
\thispagestyle{plain}
\refstepcounter{section}
\addcontentsline{toc}{section}{\protect\numberline{\thesection}Formulaire de Géométrie Analytique et Vectorielle dans l'Espace}

\begin{center}
    {\LARGE\bfseries\color{uclblue} Géométrie Analytique et Vectorielle dans l'Espace ($\R^3$)}\\[2pt]
    {\small\color{darkgray} Examen d'entrée Ingénieur Civil -- UCLouvain (Fiche récapitulative imprimable 4/4)}
    \vspace{2pt}
    \hrule height 1pt
\end{center}
\vspace{-6pt}

{\fontsize{9.5pt}{11.5pt}\selectfont
\setlength{\columnsep}{18pt}
\begin{multicols}{2}

\noindent{\large\bfseries\color{uclblue} 1. Outils Vectoriels dans $\R^3$}
\begin{itemize}[leftmargin=1.2em, itemsep=2pt, topsep=2pt]
    \item \textbf{Distance et milieu de $A(x_A,y_A,z_A)$ et $B(x_B,y_B,z_B)$ :}
    \[
        |AB| = \sqrt{(x_B-x_A)^2 + (y_B-y_A)^2 + (z_B-z_A)^2}
    \]
    \[
        M = \left(\frac{x_A+x_B}{2},\; \frac{y_A+y_B}{2},\; \frac{z_A+z_B}{2}\right)
    \]
    \item \textbf{Produit scalaire et angle $\theta$ :}
    \[
        \vec{u}\cdot\vec{v} = u_x v_x + u_y v_y + u_z v_z = \|\vec{u}\|\,\|\vec{v}\|\cos\theta
    \]
    \textbf{Orthogonalité :} $\vec{u} \perp \vec{v} \iff \vec{u}\cdot\vec{v} = 0$.
    \item \textbf{Produit vectoriel $\vec{n} = \vec{u} \wedge \vec{v}$ ($\perp \vec{u}$ et $\perp \vec{v}$) :}
    \[
        \vec{u}\wedge\vec{v} = \bigl(u_y v_z - u_z v_y,\;\; u_z v_x - u_x v_z,\;\; u_x v_y - u_y v_x\bigr)
    \]
    \item \textbf{Coplanarité de 3 vecteurs $\vec{u}, \vec{v}, \vec{w}$ :}
    \[
        \det(\vec{u},\vec{v},\vec{w}) = \begin{vmatrix} u_x & v_x & w_x \\ u_y & v_y & w_y \\ u_z & v_z & w_z \end{vmatrix} = 0
    \]
\end{itemize}

\vspace{4pt}
\noindent{\large\bfseries\color{uclblue} 2. Droites et Plans dans l'Espace}
\begin{itemize}[leftmargin=1.2em, itemsep=2pt, topsep=2pt]
    \item \textbf{Droite $d$ passant par $A$ de vecteur directeur $\vec{u} = (a,b,c)$ :}
    \[
        d \equiv \begin{cases} x = x_A + t\,a \\ y = y_A + t\,b \\ z = z_A + t\,c \end{cases} \quad (t \in \R)
    \]
    Distance parcourue entre $A(0)$ et $M(t)$ : $|AM| = |t|\cdot\|\vec{u}\|$.
    \item \textbf{Plan $\Pi$ de vecteur normal $\vec{n} = (a,b,c)$ :}
    \[
        \Pi \equiv a(x-x_A) + b(y-y_A) + c(z-z_A) = 0
    \]
    soit $\mathbf{ax + by + cz + d = 0}$ avec $\vec{n} = (a,b,c) \perp \Pi$.
    \item \textbf{Plan défini par un point $A$ et deux vecteurs directeurs $\vec{u}, \vec{v}$ :} Prendre $\vec{n} = \vec{u} \wedge \vec{v}$ ou résoudre $\det(\overrightarrow{AM},\vec{u},\vec{v}) = 0$.
    \item \textbf{Plan médiateur d'un segment $[AB]$} (lieu des points équidistants de $A$ et $B$) : Plan passant par le milieu $M = \frac{A+B}{2}$ et de vecteur normal $\vec{n} = \overrightarrow{AB}$.
\end{itemize}

\vspace{4pt}
\noindent{\large\bfseries\color{uclblue} 3. Positions Relatives Droites / Plans}
Soit $d$ de vecteur directeur $\vec{u}$ et $\Pi_1, \Pi_2$ de normales $\vec{n}_1, \vec{n}_2$ :
\begin{itemize}[leftmargin=1.2em, itemsep=2pt, topsep=2pt]
    \item \textbf{Droite parallèle au plan ($d \parallel \Pi_1$) :} $\vec{u} \cdot \vec{n}_1 = 0$
    \item \textbf{Droite perpendiculaire au plan ($d \perp \Pi_1$) :} $\vec{u} \parallel \vec{n}_1$
    \item \textbf{Plans perpendiculaires ($\Pi_1 \perp \Pi_2$) :} $\vec{n}_1 \cdot \vec{n}_2 = 0$
    \item \textbf{Plans parallèles ($\Pi_1 \parallel \Pi_2$) :} $\vec{n}_1 \parallel \vec{n}_2$
\end{itemize}

\columnbreak

\noindent{\large\bfseries\color{uclblue} 4. Intersections, Projections et Distances}
\begin{itemize}[leftmargin=1.2em, itemsep=2pt, topsep=2pt]
    \item \textbf{Point de percée $B = d \cap \Pi$ :} Remplacer $x(t), y(t), z(t)$ de la droite $d$ dans l'équation cartésienne $ax+by+cz+d=0$ du plan $\Pi$ pour trouver $t$.
    \item \textbf{Projeté orthogonal $P$ d'un point $C$ sur une droite $d \equiv A + t\vec{u}$ :} Écrire $P(t) = A + t\vec{u}$ et résoudre l'équation d'orthogonalité :
    \[
        \overrightarrow{CP(t)} \cdot \vec{u} = 0
    \]
    \item \textbf{Projeté orthogonal $H$ d'un point $A$ sur un plan $\Pi$ de normale $\vec{n}$ :} $H$ est le point de percée dans $\Pi$ de la droite passant par $A$ et de vecteur directeur $\vec{n}$.
    \item \textbf{Droite $d_2$ parallèle à $d_1 \equiv A + t\vec{u}$ à distance $D$ dans la direction $\vec{v} \perp \vec{u}$ :}
    \[
        B = A + D\,\frac{\vec{v}}{\|\vec{v}\|} \quad \text{et} \quad d_2 \equiv B + t\,\vec{u}
    \]
    \item \textbf{Distance d'un point $A(x_A, y_A, z_A)$ au plan $\Pi \equiv ax + by + cz + d = 0$ :}
    \[
        d(A, \Pi) = \frac{|ax_A + by_A + cz_A + d|}{\sqrt{a^2 + b^2 + c^2}}
    \]
\end{itemize}

\vspace{4pt}
\noindent{\large\bfseries\color{uclblue} 5. Réflexion Optique sur un Miroir Plan}
Un rayon lumineux de vecteur directeur incident $\vec{v}_{\text{inc}}$ frappant un miroir plan $\Pi$ de vecteur normal $\vec{n}$ en un point $B$ fait un angle d'incidence $\theta$ avec la normale ($\cos\theta = \frac{|\vec{v}_{\text{inc}}\cdot\vec{n}|}{\|\vec{v}_{\text{inc}}\|\|\vec{n}\|}$) et repart selon le vecteur directeur réfléchi :
\[
    \vec{v}_{\text{réf}} = \vec{v}_{\text{inc}} - 2\,\frac{\vec{v}_{\text{inc}} \cdot \vec{n}}{\|\vec{n}\|^2}\,\vec{n}
\]

\vspace{4pt}
\noindent{\large\bfseries\color{uclblue} 6. Sphère et Section par un Plan}
\begin{itemize}[leftmargin=1.2em, itemsep=2pt, topsep=2pt]
    \item \textbf{Sphère $\mathcal{S}$ de centre $A(x_A, y_A, z_A)$ et de rayon $R$ :}
    \[
        (x - x_A)^2 + (y - y_A)^2 + (z - z_A)^2 = R^2
    \]
    \item \textbf{Intersection de $\mathcal{S}(A, R)$ avec un plan $\Pi$ :} Si $d = d(A, \Pi) < R$, l'intersection est un \textbf{cercle} (ou un disque pour la boule) centré au projeté orthogonal $H$ de $A$ sur $\Pi$, de rayon :
    \[
        r = \sqrt{R^2 - d^2} \quad \implies \quad \text{Aire} = \pi(R^2 - d^2)
    \]
    \item \textbf{Plan tangent à $\mathcal{S}(A,R)$ en un point $T \in \mathcal{S}$ :} C'est le plan passant par $T$ et de vecteur normal $\vec{n} = \overrightarrow{AT}$.
\end{itemize}

\end{multicols}
}

\restoregeometry

% ============================================================================
% ANNEXE A : ÉLÉMENTS DE LOGIQUE ET MÉTHODES DE DÉMONSTRATION
% ============================================================================
\chapter*{Annexe A : Éléments de Logique et Types de Démonstrations}
\addcontentsline{toc}{chapter}{Annexe A : Éléments de Logique et Types de Démonstrations}
\markboth{Annexe A : Logique et Démonstrations}{Annexe A : Logique et Démonstrations}

\emph{Cette annexe est directement adaptée des Annexes A et B de la « Synthèse de trigonométrie » d'Yvan Haine et Pierre Joris (ULiège). Bien que la logique formelle ne fasse pas explicitement partie du programme de l'examen d'admission, la maîtrise des connecteurs logiques (« et », « ou », $\implies$, $\iff$) est essentielle pour rédiger rigoureusement la résolution d'une équation (sans perdre de solutions ni introduire de solutions parasites) ou la démonstration d'une proposition géométrique.}

% ----------------------------------------------------------------------------
\section*{A.1 Assertions et opérateurs logiques (« non », « et », « ou »)}
% ----------------------------------------------------------------------------

\begin{itemize}
    \item Une \textbf{assertion} est un énoncé mathématique non ambigu qui possède une unique valeur de vérité : soit \textbf{Vrai (V)}, soit \textbf{Faux (F)}.
    \item Une \textbf{proposition} (ou théorème, lemme, corollaire) est une assertion dont on a démontré qu'elle est vraie.
    \item Deux assertions $P$ et $Q$ sont \textbf{logiquement équivalentes} (noté $P \iff Q$) si elles ont exactement la même table de vérité.
\end{itemize}

\begin{theoreme}[Négation, conjonction, disjonction et lois de De Morgan]
Soient $P$ et $Q$ deux assertions :
\begin{enumerate}
    \item \textbf{Négation ($\text{non } P$ ou $\neg P$) :} Vraie si $P$ est fausse, et fausse si $P$ est vraie. La double négation redonne $P$ : $\text{non}(\text{non } P) \iff P$.
    \item \textbf{Conjonction (« $P \text{ et } Q$ », notée $P \wedge Q$ ou par une accolade système $\begin{cases} P \\ Q \end{cases}$) :} Vraie uniquement lorsque $P$ et $Q$ sont \textbf{toutes les deux vraies} simultanément.
    \item \textbf{Disjonction (« $P \text{ ou } Q$ », notée $P \vee Q$) :} Vraie dès qu'\textbf{au moins une} des deux assertions $P$ ou $Q$ est vraie (le « ou » mathématique est toujours inclusif).
    \item \textbf{Lois de De Morgan (négation d'un « et » et d'un « ou ») :}
    \[
        \boxed{\text{non}(P \text{ ou } Q) \iff (\text{non } P) \text{ et } (\text{non } Q)} \qquad \text{et} \qquad \boxed{\text{non}(P \text{ et } Q) \iff (\text{non } P) \text{ ou } (\text{non } Q)}.
    \]
\end{enumerate}
\end{theoreme}

\begin{attention}[Erreur d'écriture très fréquente dans les équations et les conditions d'existence !]
\begin{itemize}
    \item Lorsqu'on résout une équation produit $A(x)\cdot B(x) = 0$ ou $\cos^2 x = 1$, les solutions sont reliées par \textbf{« ou »} (ce qui correspond à l'\textbf{union} $\cup$ des ensembles de solutions). On écrit :
    \[
        \cos^2 x = 1 \iff \cos x = 1 \text{ \textbf{ou} } \cos x = -1, \qquad \text{mais \textbf{JAMAIS} } \cancel{\cos^2 x = 1 \iff \begin{cases} \cos x = 1 \\ \cos x = -1 \end{cases}},
    \]
    car l'accolade signifie \textbf{« et »} (et aucun réel $x$ ne peut avoir son cosinus égal à $1$ et à $-1$ en même temps !).
    \item Inversement, pour les \textbf{conditions d'existence} d'une fraction $\frac{1}{A(x)\cdot B(x)}$ (où le dénominateur ne peut pas s'annuler), la loi de De Morgan transforme le « ou » en \textbf{« et »} :
    \[
        A(x)\cdot B(x) \neq 0 \iff \text{non}\bigl(A(x) = 0 \text{ ou } B(x) = 0\bigr) \iff A(x) \neq 0 \text{ \textbf{et} } B(x) \neq 0 \iff \begin{cases} A(x) \neq 0 \\ B(x) \neq 0 \end{cases}.
    \]
\end{itemize}
\end{attention}

% ----------------------------------------------------------------------------
\section*{A.2 Implication ($\implies$), réciproque, contraposée et équivalence ($\iff$)}
% ----------------------------------------------------------------------------

\begin{definition}[Implication, contraposée, réciproque et équivalence]
\begin{enumerate}
    \item \textbf{Implication ($P \implies Q$, lue « si $P$, alors $Q$ ») :} $P$ est l'\emph{hypothèse} (ou condition suffisante) et $Q$ est la \emph{thèse} (ou condition nécessaire). L'implication $P \implies Q$ n'est fausse que dans un seul cas : lorsque $P$ est vraie et $Q$ est fausse.
    \item \textbf{Contraposée ($\text{non } Q \implies \text{non } P$) :} L'implication contraposée a \textbf{toujours exactement la même valeur de vérité} que l'implication directe :
    \[
        \boxed{(P \implies Q) \iff (\text{non } Q \implies \text{non } P)}.
    \]
    \item \textbf{Réciproque ($Q \implies P$) :} Attention, la réciproque d'une proposition vraie \textbf{n'est pas nécessairement vraie} ! Par exemple, $x = \frac{\pi}{6} \implies \sin x = \frac{1}{2}$ est vraie, mais la réciproque $\sin x = \frac{1}{2} \implies x = \frac{\pi}{6}$ est fausse (car on pourrait avoir $x = \frac{5\pi}{6}$).
    \item \textbf{Équivalence ($P \iff Q$, lue « $P$ si et seulement si $Q$ », en abrégé « ssi ») :} $P \iff Q$ est vraie lorsque l'implication directe $P \implies Q$ \textbf{et} l'implication réciproque $Q \implies P$ sont toutes les deux vraies.
\end{enumerate}
\end{definition}

\begin{attention}[Piège lors de la résolution d'équations : élévation au carré et solutions parasites]
Lorsqu'on résout une équation $E(x)$, chaque étape devrait idéalement être une \textbf{équivalence} ($\iff$) pour garantir que l'ensemble des solutions finales est exactement celui de l'équation de départ.
Si vous effectuez une opération qui n'est qu'une simple \textbf{implication} ($\implies$), comme \textbf{élever les deux membres au carré} :
\[
    A(x) = B(x) \implies \bigl(A(x)\bigr)^2 = \bigl(B(x)\bigr)^2 \iff A(x) = B(x) \text{ ou } A(x) = -B(x),
\]
vous introduisez des \textbf{solutions parasites} (celles qui vérifient $A(x) = -B(x)$) ! Si vous élevez une équation au carré (comme dans la Méthode 3 de Juillet 2022 Q1 : $\cos x + \sin x = \sqrt{2}\cos x\sin x$), vous êtes \textbf{obligé} de tester vos solutions finales dans l'équation de départ (ou d'imposer la condition de signe $A(x)$ et $B(x)$ de même signe) pour éliminer les solutions parasites.
\end{attention}

% ----------------------------------------------------------------------------
\section*{A.3 Quantificateurs ($\forall$ et $\exists$) et méthodes de démonstration}
% ----------------------------------------------------------------------------

\begin{itemize}
    \item \textbf{Quantificateur universel ($\forall x \in E : P(x)$) :} « Pour tout $x$ appartenant à $E$, la propriété $P(x)$ est vraie ».
    \begin{itemize}
        \item Pour démontrer qu'une affirmation universelle $\forall x \in E : P(x)$ est \textbf{vraie}, il faut faire une démonstration générale valable pour une lettre $x$ quelconque de $E$ (un exemple numérique particulier ne prouve jamais une identité !).
        \item Pour démontrer qu'une affirmation universelle $\forall x \in E : P(x)$ est \textbf{fausse}, il suffit d'exhiber \textbf{un seul contre-exemple}, c'est-à-dire une valeur précise $x_0 \in E$ pour laquelle $P(x_0)$ est fausse, car :
        \[
            \text{non}\bigl(\forall x \in E : P(x)\bigr) \iff \exists x \in E : \text{non } P(x).
        \]
    \end{itemize}
    \item \textbf{Quantificateur existentiel ($\exists x \in E : P(x)$) :} « Il existe au moins un $x \in E$ tel que $P(x)$ est vraie ». Sa négation est $\forall x \in E : \text{non } P(x)$.
\end{itemize}

\begin{theoreme}[Récapitulatif des types de démonstrations utiles en trigonométrie]
\begin{enumerate}
    \item \textbf{Démonstration d'une identité $A(x) = B(x)$ :} Partir de $A(x)$ et appliquer des formules connues jusqu'à obtenir $B(x)$ (ou transformer $A(x)$ et $B(x)$ séparément vers une même expression $C(x)$).
    \item \textbf{Démonstration déductive directe de $P \implies Q$ :} On suppose l'hypothèse $P$ vraie, et on enchaîne des déductions logiques (« Or... », « Donc... », « Par conséquent... ») jusqu'à obtenir la thèse $Q$.
    \item \textbf{Démonstration d'une équivalence $P \iff Q$ :} Soit on procède par double implication (on démontre $P \implies Q$, puis réciproquement on démontre $Q \implies P$), soit on procède par une chaîne ininterrompue d'équivalences $P \iff P_1 \iff P_2 \iff \dots \iff Q$.
    \item \textbf{Démonstration par contraposition :} Pour prouver $P \implies Q$, on suppose $\text{non } Q$ et on démontre $\text{non } P$.
    \item \textbf{Démonstration par l'absurde :} Pour prouver une thèse $Q$ à partir d'une hypothèse $P$, on suppose que $P$ est vraie \textbf{et} que $Q$ est fausse ($\text{non } Q$), et on montre que cela conduit à une contradiction logique ou mathématique.
\end{enumerate}
\end{theoreme}

% ============================================================================
% CHAPITRE 6 : CORRIGÉS DÉTAILLÉS DES EXERCICES DE FIN DE CHAPITRE
% ============================================================================
\chapter{Corrigés Détaillés des Exercices}

Ce chapitre rassemble les solutions complètes et commentées de tous les exercices proposés à la fin des Chapitres 1 à 7.

% ----------------------------------------------------------------------------
\section{Corrigés du Chapitre 1 (Angles, Arcs et Cercle Trigonométrique)}
% ----------------------------------------------------------------------------

\paragraph{Corrigé de l'Exercice 1.1 (La rosace hexagonale -- Inspiré d'ULiège) :}
\begin{enumerate}[label=(\alph*)]
    \item Les six sommets $A, B, C, D, E, F$ d'un hexagone régulier inscrit dans un cercle de rayon $R = 3\text{ cm}$ divisent la circonférence en $6$ arcs égaux d'angle au centre $\theta = \frac{2\pi}{6} = \frac{\pi}{3}\text{ rad}$ ($60\degree$). La longueur du plus petit arc $\overset{\frown}{AB}$ vaut donc :
    \[
        L_{\overset{\frown}{AB}} = R\theta = 3 \cdot \frac{\pi}{3} = \boxed{\pi\text{ cm} \approx 3,14\text{ cm}}.
    \]
    \item Chaque pétale de la rosace est formé de deux arcs de cercle de rayon $R = 3\text{ cm}$ reliant le centre $O$ à un sommet (par exemple $A$ et $B$, centrés en $F$ et $C$). Dans les triangles équilatéraux $OAB$, $OBC$, etc., de côté $R$, chaque arc allant d'un sommet au sommet opposé à $120\degree$ (comme l'arc $\overset{\frown}{BF}$ centré en $A$, passant par $O$) correspond à un angle au centre de $\frac{2\pi}{3}\text{ rad}$ ($120\degree$). Comme la rosace complète à $6$ pétales est constituée de $6$ grands arcs d'angle $\frac{2\pi}{3}$ (ou $12$ petits demi-pétales d'angle $\frac{\pi}{3}$), sa longueur totale vaut :
    \[
        L_{\text{rosace}} = 12 \times \left(R\cdot\frac{\pi}{3}\right) = 4\pi R = \boxed{12\pi\text{ cm} \approx 37,70\text{ cm}}.
    \]
    \item Un demi-segment circulaire d'angle $\theta = \frac{\pi}{3}$ et de rayon $R$ a une aire $A_{\text{seg}} = \frac{1}{2}R^2\left(\frac{\pi}{3} - \sin\frac{\pi}{3}\right) = \frac{R^2}{2}\left(\frac{\pi}{3} - \frac{\sqrt{3}}{2}\right)$. Un pétale complet est la réunion de deux tels segments accolés le long de leur corde de longueur $R$, et la rosace comporte $6$ pétales (soit $12$ segments circulaires d'angle $\frac{\pi}{3}$) :
    \[
        A_{\text{rosace}} = 12 \times \frac{R^2}{2}\left(\frac{\pi}{3} - \frac{\sqrt{3}}{2}\right) = R^2(2\pi - 3\sqrt{3}) = \boxed{9(2\pi - 3\sqrt{3})\text{ cm}^2 \approx 9,78\text{ cm}^2}.
    \]
\end{enumerate}

\paragraph{Corrigé de l'Exercice 1.2 (Rayon de la Terre par Ératosthène -- Inspiré d'ULiège) :}
Les rayons du Soleil étant parallèles, l'angle $\theta$ que font les rayons solaires avec la verticale de l'obélisque à Alexandrie est égal (angles correspondants) à l'angle au centre de la Terre entre Syène et Alexandrie. Dans le triangle rectangle formé par l'obélisque ($h = 1\text{ m} = 100\text{ cm}$) et son ombre ($d = 12,6\text{ cm}$), on a $\tan\theta = \frac{12,6}{100} = 0,126$. Pour un petit angle exprimé en radians, $\theta \approx \tan\theta \approx 0,126\text{ rad}$ (plus précisément $\theta = \arctan(0,126) \approx 0,1253\text{ rad}$). La relation de l'arc de méridien $L = R\theta$ avec $L = 800\text{ km}$ donne alors :
\[
    R = \frac{L}{\theta} \approx \frac{800}{0,1253} \approx \boxed{6383\text{ km}} \quad \text{(ou } \frac{800}{0,126} \approx 6349\text{ km)}.
\]

\paragraph{Corrigé de l'Exercice 1.3 (UCLouvain -- Mini-Blocus 2026, Q1b) :}
La perpendiculaire abaissée du centre $O$ du cercle sur la corde $[AB]$ coupe celle-ci en son milieu $M$, avec $\overline{OM} = 3\text{ cm}$ et $\overline{OA} = R = 5\text{ cm}$. Dans le triangle $OMA$ rectangle en $M$, le théorème de Pythagore donne $\overline{AM} = \sqrt{R^2 - \overline{OM}^2} = \sqrt{5^2 - 3^2} = \sqrt{16} = 4\text{ cm}$. La longueur totale de la corde vaut donc $L = 2\overline{AM} = \boxed{8\text{ cm}}$.

\paragraph{Corrigé de l'Exercice 1.4 (Simplifications et identités fondamentales -- Inspiré d'ULiège) :}
\begin{enumerate}[label=(\alph*)]
    \item \textbf{CE :} $\cos x \neq 0$ et $\sin x \neq 0 \iff x \neq k\frac{\pi}{2}$ ($k\in\Z$). On met au même dénominateur :
    \[
        \tan x + \cot x = \frac{\sin x}{\cos x} + \frac{\cos x}{\sin x} = \frac{\sin^2 x + \cos^2 x}{\sin x \cos x} = \frac{1}{\sin x \cos x}.
    \]
    \item \textbf{CE :} $x \neq \frac{\pi}{2} + k\pi$ ($k\in\Z$). En utilisant $1 + \tan^2 x = \frac{1}{\cos^2 x}$ :
    \[
        \frac{\tan x}{1 + \tan^2 x} = \frac{\sin x / \cos x}{1 / \cos^2 x} = \frac{\sin x}{\cos x} \cdot \cos^2 x = \sin x \cos x.
    \]
    \item \textbf{CE :} $\forall x \in \R$. Par différence de deux carrés $a^4 - b^4 = (a^2-b^2)(a^2+b^2)$ et $\sin^2 x + \cos^2 x = 1$ :
    \[
        \sin^4 x - \cos^4 x = \sin^2 x - \cos^2 x = (1 - \cos^2 x) - \cos^2 x = 1 - 2\cos^2 x = 2\sin^2 x - 1.
    \]
    \item \textbf{CE :} $\forall x \in \R$. En développant $1 = (\sin^2 x + \cos^2 x)^3 = \sin^6 x + \cos^6 x + 3\sin^2 x\cos^2 x(\sin^2 x + \cos^2 x)$, il vient immédiatement $\sin^6 x + \cos^6 x = 1 - 3\sin^2 x\cos^2 x$.
\end{enumerate}

\paragraph{Corrigé de l'Exercice 1.5 (UCLouvain -- Juillet 2022, Q3a) :}
Le triangle $ABC$ étant rectangle, l'un de ses trois angles est droit (supposons $\alpha = \frac{\pi}{2}$, donc $\sin\alpha = 1$), et les deux autres angles aigus $\beta$ et $\gamma$ sont complémentaires ($\gamma = \frac{\pi}{2} - \beta$), ce qui entraîne $\sin\gamma = \sin\left(\frac{\pi}{2}-\beta\right) = \cos\beta$. Par conséquent :
\[
    \sin^2(\alpha) + \sin^2(\beta) + \sin^2(\gamma) = 1^2 + \sin^2(\beta) + \cos^2(\beta) = 1 + 1 = \boxed{2}.
\]

\paragraph{Corrigé de l'Exercice 1.6 (UCLouvain -- Juillet 2022, Q3d) :}
Posons $y = \sin x - \cos x$. Remarquons que :
\[
    (\sin x + \cos x)^2 + (\sin x - \cos x)^2 = 2(\sin^2 x + \cos^2 x) = 2.
\]
En remplaçant $\sin x + \cos x = \frac{\sqrt{2}}{3}$, on obtient $\left(\frac{\sqrt{2}}{3}\right)^2 + y^2 = 2 \iff \frac{2}{9} + y^2 = 2 \iff y^2 = \frac{16}{9} \iff y = \pm\frac{4}{3}$.
Or, dans le deuxième quadrant ($\frac{\pi}{2} < x < \pi$), on a $\sin x > 0$ et $\cos x < 0$, donc $\sin x - \cos x > 0$. La valeur cherchée est donc positive : $\sin(x) - \cos(x) = \boxed{\dfrac{4}{3}}$.

\paragraph{Corrigé de l'Exercice 1.7 (Angles associés) :}
On simplifie chaque terme grâce aux angles associés : $\cos(\pi-x) = -\cos x$, $\sin\left(\frac{\pi}{2}+x\right) = \cos x$, $\cos(3\pi+x) = \cos(\pi+x) = -\cos x$, et $\sin\left(\frac{3\pi}{2}-x\right) = \sin\left(\pi + \left(\frac{\pi}{2}-x\right)\right) = -\sin\left(\frac{\pi}{2}-x\right) = -\cos x$. En additionnant :
\[
    E(x) = -\cos x + \cos x - \cos x - (-\cos x) = \boxed{0}.
\]

\paragraph{Corrigé de l'Exercice 1.8 (UCLouvain -- Juillet 2026, Q3d) :}
Puisque $\sin\theta = x$ et $-\frac{\pi}{2} \le \theta \le \frac{\pi}{2}$ (quadrants IV et I), le cosinus est positif ou nul : $\cos\theta = +\sqrt{1 - \sin^2\theta} = \sqrt{1 - x^2}$. Par la formule de duplication du sinus :
\[
    \sin(2\theta) = 2\sin\theta\cos\theta = \boxed{2x\sqrt{1 - x^2}}.
\]

% ----------------------------------------------------------------------------
\section{Corrigés du Chapitre 2 (Formules Trigonométriques et Identités)}
% ----------------------------------------------------------------------------

\paragraph{Corrigé de l'Exercice 2.1 (UCLouvain -- Juillet 2022 Q3f et Juillet 2023 Q3d) :}
En lisant la formule de Simpson $\sin(a+b) + \sin(a-b) = 2\sin a \cos b$ de droite à gauche ($\sin a\cos b = \frac{1}{2}[\sin(a+b) + \sin(a-b)]$) :
\begin{align*}
    \sin(70\degree)\cos(50\degree) &= \frac{1}{2}\bigl(\sin(120\degree) + \sin(20\degree)\bigr) = \frac{1}{2}\left(\frac{\sqrt{3}}{2} + \sin(20\degree)\right), \\
    \sin(260\degree)\cos(280\degree) &= \frac{1}{2}\bigl(\sin(540\degree) + \sin(-20\degree)\bigr) = \frac{1}{2}\bigl(\sin(180\degree) - \sin(20\degree)\bigr) = -\frac{1}{2}\sin(20\degree).
\end{align*}
En additionnant les deux termes, $\pm\frac{1}{2}\sin(20\degree)$ s'éliminent et il reste $E = \boxed{\dfrac{\sqrt{3}}{4}}$.
\emph{Alternative rapide :} Par angles associés, $\sin(260\degree) = \sin(180\degree+80\degree) = -\sin(80\degree) = -\cos(10\degree)$ et $\cos(280\degree) = \cos(360\degree-80\degree) = \cos(80\degree) = \sin(10\degree)$, ce qui ramène au même calcul.

\paragraph{Corrigé de l'Exercice 2.2 (UCLouvain -- Septembre 2022, Q3a et Q3b) :}
\begin{itemize}
    \item \textbf{Pour $S_1$ :} Appliquons Simpson aux deux premiers termes $\sin(50\degree) + \sin(10\degree)$ (moyenne $30\degree$, demi-écart $20\degree$) :
    \[
        S_1 = 2\sin(30\degree)\cos(20\degree) - \sin(70\degree) = \cos(20\degree) - \sin(90\degree - 20\degree) = \cos(20\degree) - \cos(20\degree) = \boxed{0}.
    \]
    \item \textbf{Pour $S_2$ :} Appliquons Simpson aux premier et troisième termes $\cos\left(\frac{\pi}{9}\right) + \cos\left(\frac{7\pi}{9}\right)$ (moyenne $\frac{4\pi}{9}$, demi-écart $\frac{3\pi}{9} = \frac{\pi}{3}$) :
    \[
        S_2 = 2\cos\left(\frac{4\pi}{9}\right)\cos\left(\frac{\pi}{3}\right) + \cos\left(\pi - \frac{4\pi}{9}\right) = \cos\left(\frac{4\pi}{9}\right) - \cos\left(\frac{4\pi}{9}\right) = \boxed{0}.
    \]
\end{itemize}

\paragraph{Corrigé de l'Exercice 2.3 (UCLouvain -- Septembre 2022 Q3d et Juillet 2024 Q3b) :}
\begin{enumerate}[label=(\alph*)]
    \item Comme $\cos x = \frac{1}{3}$ et $x \in \left]0, \pi\right[$, on a $\sin x = +\sqrt{1 - \cos^2 x} = \sqrt{1 - \frac{1}{9}} = \frac{\sqrt{8}}{3} = \frac{2\sqrt{2}}{3}$. Par les formules de duplication :
    \[
        \sin(2x) = 2\sin x\cos x = \frac{4\sqrt{2}}{9} \quad \text{et} \quad \cos(2x) = 2\cos^2 x - 1 = \frac{2}{9} - 1 = -\frac{7}{9},
    \]
    d'où $\tan(2x) = \frac{\sin(2x)}{\cos(2x)} = -\frac{4\sqrt{2}}{7}$ et $7\times\tan(2x) = \boxed{-4\sqrt{2}}$.
    \item Par la formule de $\sin(2\alpha)$ en fonction de $t = \tan(\alpha)$ (avec $\alpha = 10\degree$ et $t \simeq \frac{9}{50}$) :
    \[
        \sin(20\degree) = \frac{2\tan(10\degree)}{1 + \tan^2(10\degree)} \simeq \frac{2 \cdot \frac{9}{50}}{1 + \left(\frac{9}{50}\right)^2} = \frac{\frac{18}{50}}{\frac{2500 + 81}{2500}} = \frac{18 \times 50}{2581} = \boxed{\frac{900}{2581}}.
    \]
\end{enumerate}

\paragraph{Corrigé de l'Exercice 2.4 (UCLouvain -- Septembre 2024, Q3a, Q3b et Q3d) :}
\begin{enumerate}[label=(\alph*)]
    \item Par Carnot, $\cos(2x) = 1 - 2\sin^2 x = 1 - 2\left(\frac{\sqrt{5}-1}{4}\right)^2 = 1 - 2\cdot\frac{6 - 2\sqrt{5}}{16} = 1 - \frac{3-\sqrt{5}}{4} = \frac{1+\sqrt{5}}{4}$ (qui est la valeur de $\cos 36\degree$, liée au nombre d'or !). Puis en appliquant une seconde fois la duplication :
    \[
        \cos(4x) + \frac{1}{4} = 2\cos^2(2x) - 1 + \frac{1}{4} = 2\left(\frac{1+\sqrt{5}}{4}\right)^2 - \frac{3}{4} = \frac{6+2\sqrt{5}}{8} - \frac{3}{4} = \boxed{\frac{\sqrt{5}}{4}}.
    \]
    \item Appliquons Simpson à $\sin(36\degree+x) - \sin(36\degree-x) = 2\sin(x)\cos(36\degree)$ et à $\sin(72\degree-x) - \sin(72\degree+x) = -2\sin(x)\cos(72\degree)$. En factorisant $2\sin(x)$ et en utilisant $\cos(72\degree) = 2\cos^2(36\degree) - 1$ :
    \begin{align*}
        E(x) &= 2\sin(x)\bigl(\cos(36\degree) - 2\cos^2(36\degree) + 1\bigr) \\
        &= 2\sin(x)\left(\frac{1+\sqrt{5}}{4} - \frac{3+\sqrt{5}}{4} + 1\right) = 2\sin(x)\cdot\frac{1}{2} = \boxed{\sin(x)}.
    \end{align*}
    \item Utilisons la formule de l'angle moitié $\tan\left(\frac{\alpha}{2}\right) = \frac{1 - \cos\alpha}{\sin\alpha}$ pour $\frac{\alpha}{2} = \frac{\pi}{24}$ (donc $\alpha = \frac{\pi}{12}$). Sachant que $\cos\frac{\pi}{12} = \frac{\sqrt{2}(\sqrt{3}+1)}{4}$ et $\sin\frac{\pi}{12} = \frac{\sqrt{2}(\sqrt{3}-1)}{4}$ :
    \[
        \sqrt{2}(\sqrt{3}-1)\tan\left(\frac{\pi}{24}\right) = 4\sin\left(\frac{\pi}{12}\right)\cdot\frac{1 - \cos(\pi/12)}{\sin(\pi/12)} = 4 - 4\cos\left(\frac{\pi}{12}\right) = 4 - \sqrt{2}(\sqrt{3}+1).
    \]
    Par conséquent, $4 - \sqrt{2}(\sqrt{3}-1)\tan\left(\frac{\pi}{24}\right) = \boxed{\sqrt{2}(\sqrt{3}+1) = \sqrt{6} + \sqrt{2}}$.
\end{enumerate}

\paragraph{Corrigé de l'Exercice 2.5 (UCLouvain -- Juillet 2021 Q3d et Septembre 2023 Q3d) :}
\begin{itemize}
    \item \textbf{Pour $P_1$ :} Remarquons que $\frac{\pi}{24} + \frac{11\pi}{24} = \frac{\pi}{2}$ et $\frac{5\pi}{24} + \frac{7\pi}{24} = \frac{\pi}{2}$ (angles complémentaires !). Donc $\sin\left(\frac{11\pi}{24}\right) = \cos\left(\frac{\pi}{24}\right)$ et $\sin\left(\frac{7\pi}{24}\right) = \cos\left(\frac{5\pi}{24}\right)$. En utilisant deux fois $\sin\alpha\cos\alpha = \frac{1}{2}\sin(2\alpha)$ :
    \[
        P_1 = \left(\sin\frac{\pi}{24}\cos\frac{\pi}{24}\right)\left(\sin\frac{5\pi}{24}\cos\frac{5\pi}{24}\right) = \frac{1}{4}\sin\left(\frac{\pi}{12}\right)\cos\left(\frac{\pi}{12}\right) = \frac{1}{8}\sin\left(\frac{\pi}{6}\right) = \boxed{\frac{1}{16}}.
    \]
    \item \textbf{Pour $P_2$ :} On connaît $\sin\left(\frac{3\pi}{9}\right) = \sin\left(\frac{\pi}{3}\right) = \frac{\sqrt{3}}{2}$. Calculons $Q = \sin\left(\frac{\pi}{9}\right)\sin\left(\frac{2\pi}{9}\right)\sin\left(\frac{4\pi}{9}\right)$ (soit $\sin 20\degree\sin 40\degree\sin 80\degree$) en appliquant Werner $2\sin 20\degree\sin 40\degree = \cos 20\degree - \cos 60\degree = \cos 20\degree - \frac{1}{2}$ :
    \[
        Q = \frac{1}{2}\left(\cos 20\degree - \frac{1}{2}\right)\sin 80\degree = \frac{1}{4}\bigl(\sin 100\degree + \sin 60\degree - \sin 80\degree\bigr) = \frac{\sin 60\degree}{4} = \frac{\sqrt{3}}{8},
    \]
    car $\sin 100\degree = \sin(180\degree-80\degree) = \sin 80\degree$. D'où $P_2 = \frac{\sqrt{3}}{2} \cdot \frac{\sqrt{3}}{8} = \boxed{\dfrac{3}{16}}$.
\end{itemize}

\paragraph{Corrigé de l'Exercice 2.6 (UCLouvain -- Juillet 2021, Q3e) :}
L'hypothèse $\tan x + \tan y + \tan x\tan y = 1$ se réécrit $\tan x + \tan y = 1 - \tan x\tan y$. Comme $1 - \tan x\tan y \neq 0$ (sinon $\tan x + \tan y = 0$ et $\tan x\tan y = 1$, impossible dans $\R$), on divise par $1 - \tan x\tan y$ :
\[
    \tan(x + y) = \frac{\tan x + \tan y}{1 - \tan x\tan y} = 1.
\]
Pour $(x+y) \in [0, \pi]$, l'unique solution est $\boxed{x + y = \dfrac{\pi}{4}}$.

\paragraph{Corrigé de l'Exercice 2.7 (ULiège -- Juillet 2002 Q1b, Juillet 2003 Q1 et Juillet 2004 Q1) :}
\begin{enumerate}[label=(\alph*)]
    \item \textbf{CE :} $\cos(2x)+\cos(2y) \neq 0$ et $\tan(x+y)$ existe $\iff \cos(x+y) \neq 0$ et $\cos(x-y) \neq 0 \iff x \neq \pm y + \frac{\pi}{2} + k\pi$. Par Simpson :
    \[
        \frac{\sin(2x) + \sin(2y)}{\cos(2x) + \cos(2y)} = \frac{2\sin(x+y)\cos(x-y)}{2\cos(x+y)\cos(x-y)} = \frac{\sin(x+y)}{\cos(x+y)} = \tan(x+y).
    \]
    \item \textbf{CE :} $a \neq \frac{\pi}{2} + k\pi$ et $a \neq -\frac{\pi}{4} + k\pi$. Par Carnot ($1-\cos 2a = 2\sin^2 a$, $1+\cos 2a = 2\cos^2 a$) et duplication ($\sin 2a = 2\sin a\cos a$) :
    \[
        \frac{1 - \cos(2a) + \sin(2a)}{1 + \cos(2a) + \sin(2a)} = \frac{2\sin^2 a + 2\sin a\cos a}{2\cos^2 a + 2\sin a\cos a} = \frac{2\sin a(\sin a + \cos a)}{2\cos a(\cos a + \sin a)} = \tan a.
    \]
    \item Par Carnot $\sin^2 u = \frac{1-\cos(2u)}{2}$ puis Simpson sur la somme des deux derniers cosinus, qui vaut $2\cos(2x)\cos\left(\frac{4\pi}{3}\right) = -\cos(2x)$ :
    \[
        \text{MG} = \frac{1}{2}\left[3 - \cos(2x) - \cos\left(2x+\frac{4\pi}{3}\right) - \cos\left(2x-\frac{4\pi}{3}\right)\right] = \frac{1}{2}\bigl(3 - \cos(2x) + \cos(2x)\bigr) = \frac{3}{2}.
    \]
\end{enumerate}

\paragraph{Corrigé de l'Exercice 2.8 (ULiège -- Septembre 2002 Q1 et Juillet 2010 Q4) :}
\begin{enumerate}[label=(\alph*)]
    \item Pour $a \neq k\pi$, par Carnot et duplication sur l'angle demi $\frac{a}{2}$ :
    \[
        \csc a + \cot a = \frac{1 + \cos a}{\sin a} = \frac{2\cos^2(a/2)}{2\sin(a/2)\cos(a/2)} = \cot\left(\frac{a}{2}\right) \iff \csc a = \cot\left(\frac{a}{2}\right) - \cot a.
    \]
    Par télescopage, pour $a \neq \frac{k\pi}{8}$ :
    \begin{align*}
        \csc a + \csc(2a) + \csc(4a) + \csc(8a) &= \left(\cot\frac{a}{2} - \cot a\right) + (\cot a - \cot 2a) \\
        &\quad + (\cot 2a - \cot 4a) + (\cot 4a - \cot 8a) \\
        &= \boxed{\cot\left(\frac{a}{2}\right) - \cot(8a)}.
    \end{align*}
    \item Regroupons les angles complémentaires deux à deux ($9\degree + 81\degree = 90\degree$ et $27\degree + 63\degree = 90\degree$) grâce à $\tan a + \tan b = \frac{\sin(a+b)}{\cos a\cos b}$ et Simpson $2\cos a\cos b = \cos(a+b) + \cos(a-b)$ :
    \begin{align*}
        E &= (\tan 9\degree + \tan 81\degree) - (\tan 27\degree + \tan 63\degree) = \frac{\sin 90\degree}{\cos 9\degree\cos 81\degree} - \frac{\sin 90\degree}{\cos 27\degree\cos 63\degree} \\
        &= \frac{2}{\cos 90\degree + \cos 72\degree} - \frac{2}{\cos 90\degree + \cos 36\degree} = 2\left(\frac{1}{\cos 72\degree} - \frac{1}{\cos 36\degree}\right) = 2\,\frac{\cos 36\degree - \cos 72\degree}{\cos 36\degree\cos 72\degree} \\
        &= 2\,\frac{-2\sin 54\degree\sin(-18\degree)}{\cos 36\degree\cos 72\degree} = 4\,\frac{\cos 36\degree\cos 72\degree}{\cos 36\degree\cos 72\degree} = \boxed{4}.
    \end{align*}
\end{enumerate}

\paragraph{Corrigé de l'Exercice 2.9 (Identités dans un triangle -- ULiège) :}
\begin{enumerate}[label=(\alph*)]
    \item Comme $A+B+C = \pi$, on a $\sin C = \sin(A+B) = 2\sin\left(\frac{A+B}{2}\right)\cos\left(\frac{A+B}{2}\right)$ et $\sin\left(\frac{A+B}{2}\right) = \cos\left(\frac{C}{2}\right)$, $\cos\left(\frac{A+B}{2}\right) = \sin\left(\frac{C}{2}\right)$. Par Simpson :
    \begin{align*}
        \sin A + \sin B + \sin C &= 2\sin\frac{A+B}{2}\cos\frac{A-B}{2} + 2\sin\frac{A+B}{2}\cos\frac{A+B}{2} \\
        &= 2\cos\frac{C}{2}\left(\cos\frac{A-B}{2} + \cos\frac{A+B}{2}\right) = 4\cos\frac{A}{2}\cos\frac{B}{2}\cos\frac{C}{2}.
    \end{align*}
    \item Par Simpson et Carnot ($\cos(2C) = 2\cos^2 C - 1 = 2\cos^2(A+B) - 1$ et $\cos(A+B) = -\cos C$) :
    \begin{align*}
        \cos(2A) + \cos(2B) + \cos(2C) &= 2\cos(A+B)\cos(A-B) + 2\cos^2(A+B) - 1 \\
        &= -1 + 2\cos(A+B)\bigl(\cos(A-B) + \cos(A+B)\bigr) \\
        &= -1 - 4\cos A\cos B\cos C.
    \end{align*}
    \item Par Simpson sur le membre de droite et sachant que $\frac{B+C}{2} = \frac{\pi}{2} - \frac{A}{2}$ :
    \[
        \sin A = \frac{2\sin\frac{B+C}{2}\cos\frac{B-C}{2}}{2\cos\frac{B+C}{2}\cos\frac{B-C}{2}} = \tan\left(\frac{B+C}{2}\right) = \cot\left(\frac{A}{2}\right) \iff 2\sin\frac{A}{2}\cos\frac{A}{2} = \frac{\cos(A/2)}{\sin(A/2)}.
    \]
    Comme $A \in \left]0, \pi\right[$, $\cos(A/2) \neq 0$, d'où $\sin^2(A/2) = \frac{1}{2} \implies \sin(A/2) = \frac{\sqrt{2}}{2} \implies \frac{A}{2} = \frac{\pi}{4} \implies \boxed{A = \frac{\pi}{2}}$ (triangle rectangle en $A$).
\end{enumerate}

\paragraph{Corrigé de l'Exercice 2.10 (ULiège -- Septembre 1983 Q2 / Septembre 2010 Q3) :}
\begin{enumerate}[label=(\alph*)]
    \item Si $a = b-r$ et $c = b+r$, appliquons Simpson à $\sin(b-r)+\sin(b+r) = 2\sin b\cos r$ et $\cos(b-r)+\cos(b+r) = 2\cos b\cos r$ :
    \[
        \frac{\sin a + \sin b + \sin c}{\cos a + \cos b + \cos c} = \frac{\sin b(2\cos r + 1)}{\cos b(2\cos r + 1)} = \tan b,
    \]
    qui est valide sauf si $\cos b = 0$ ($b = \frac{\pi}{2} + k\pi$) ou $2\cos r + 1 = 0$ ($r = \pm\frac{2\pi}{3} + 2k\pi$).
    \item Réciproquement, $\frac{\sin a + \sin b + \sin c}{\cos a + \cos b + \cos c} = \frac{\sin b}{\cos b} \iff (\sin a\cos b - \cos a\sin b) + (\sin c\cos b - \cos c\sin b) = 0 \iff \sin(a-b) + \sin(c-b) = 0 \iff \sin(a-b) = \sin(b-c)$. Cela équivaut à $a - b = b - c + 2k\pi$ (soit $b = \frac{a+c}{2} + k\pi$) ou $a - b = \pi - (b - c) + 2k\pi$ (soit $a = c + (2k+1)\pi$, cas dégénéré où $\sin a + \sin c = 0$ et $\cos a + \cos c = 0$). La réciproque n'est donc vraie qu'à un multiple de $\pi$ près (ou exactement si l'on impose $|b - \frac{a+c}{2}| < \pi$ et $a \neq c + (2k+1)\pi$).
\end{enumerate}

% ----------------------------------------------------------------------------
\section{Corrigés du Chapitre 3 (Équations et Inéquations Trigonométriques)}
% ----------------------------------------------------------------------------

\paragraph{Corrigé de l'Exercice 3.1 (UCLouvain -- Mini-Blocus 2026 Q1a et Juillet 2024 Q3a) :}
\begin{enumerate}[label=(\alph*)]
    \item Par la méthode des équations linéaires (ou symétriques), $\sin x + \cos x = \sqrt{2}\sin\left(x + \frac{\pi}{4}\right) = \sqrt{2} \iff \sin\left(x + \frac{\pi}{4}\right) = 1 \iff x + \frac{\pi}{4} = \frac{\pi}{2} + 2k\pi \iff x = \frac{\pi}{4} + 2k\pi$. Dans $[0, 2\pi[$, l'unique solution est $\boxed{x = \dfrac{\pi}{4}}$.
    \item Par Simpson sur $\sin(x) + \sin(3x) = 2\sin(2x)\cos(x)$ :
    \[
        \sin(x) + \sin(2x) + \sin(3x) = \sin(2x)\bigl(2\cos x + 1\bigr) = 0 \iff \sin(2x) = 0 \quad \text{ou} \quad \cos x = -\frac{1}{2}.
    \]
    Dans l'intervalle ouvert $]0, \pi[$, $\sin(2x) = 0 \iff x = \frac{k\pi}{2}$ donne $x = \frac{\pi}{2}$, et $\cos x = -\frac{1}{2}$ donne $x = \frac{2\pi}{3}$. D'où $\boxed{x \in \left\{\frac{\pi}{2}, \frac{2\pi}{3}\right\}}$.
\end{enumerate}

\paragraph{Corrigé de l'Exercice 3.2 (UCLouvain -- Septembre 2019, Q1) :}
\begin{enumerate}[label=(\alph*)]
    \item Par Simpson, $\sin(5x)+\sin(x) = 2\sin(3x)\cos(2x)$ et $\cos(4x)-\cos(2x) = -2\sin(3x)\sin(x)$. L'équation devient :
    \[
        2\sin(3x)\cos(2x) - 2\sin(3x) = -2\sin(3x)\sin(x) \iff 2\sin(3x)\bigl(\cos(2x) + \sin(x) - 1\bigr) = 0.
    \]
    En remplaçant $\cos(2x) - 1 = -2\sin^2 x$, le second facteur s'écrit $\sin x(1 - 2\sin x) = 0$. On obtient donc $\sin(3x) = 0$ ($x = \frac{k\pi}{3}$, qui inclut $\sin x = 0$) ou $\sin x = \frac{1}{2}$ ($x = \frac{\pi}{6} + 2k\pi$ ou $x = \frac{5\pi}{6} + 2k\pi$), soit :
    \[
        \boxed{S = \left\{\frac{k\pi}{3} \;\middle|\; k\in\Z\right\} \cup \left\{\frac{\pi}{6} + 2k\pi,\; \frac{5\pi}{6} + 2k\pi \;\middle|\; k\in\Z\right\}}.
    \]
    \item \textbf{CE :} $\sin x \neq 0 \iff x \neq k\pi$ ($k\in\Z$). Par Simpson au numérateur, $\cos(x) - \cos(3x) = -2\sin(2x)\sin(-x) = 2\sin(2x)\sin(x)$. En simplifiant par $\sin x \neq 0$ :
    \[
        2\sin(2x) = \cos(x) \iff 4\sin x\cos x - \cos x = 0 \iff \cos x(4\sin x - 1) = 0.
    \]
    D'où $\cos x = 0 \iff x = \frac{\pi}{2} + k\pi$ ou $\sin x = \frac{1}{4} \iff x = \arcsin\left(\frac{1}{4}\right) + 2k\pi$ ou $x = \pi - \arcsin\left(\frac{1}{4}\right) + 2k\pi$ ($k\in\Z$), toutes compatibles avec la CE $x \neq k\pi$.
\end{enumerate}

\paragraph{Corrigé de l'Exercice 3.3 (UCLouvain -- Juillet 2019 Q1 et Juillet 2020 Q1) :}
\begin{enumerate}[label=(\alph*)]
    \item Équation symétrique : posons $u = \sin x + \cos x = \sqrt{2}\cos\left(x - \frac{\pi}{4}\right)$, d'où $2\sin x\cos x = u^2 - 1$. L'équation devient :
    \[
        4u - 4(u^2 - 1) = 5 \iff 4u^2 - 4u + 1 = 0 \iff (2u - 1)^2 = 0 \iff u = \frac{1}{2}.
    \]
    On résout donc $\cos\left(x - \frac{\pi}{4}\right) = \frac{\sqrt{2}}{4}$, soit $\boxed{x = \frac{\pi}{4} \pm \arccos\left(\frac{\sqrt{2}}{4}\right) + 2k\pi}$ ($k\in\Z$).
    \item Par l'identité remarquable $a^3 + b^3 = (a+b)(a^2 - ab + b^2)$, on a $\sin^3 x + \cos^3 x = (\sin x + \cos x)(1 - \sin x\cos x)$. En mettant $(\sin x + \cos x)$ en évidence :
    \[
        \bigl(\sin x + \cos x\bigr)\bigl(3 - 4(1 - \sin x\cos x)\bigr) = 0 \iff \bigl(\sin x + \cos x\bigr)\bigl(4\sin x\cos x - 1\bigr) = 0.
    \]
    \begin{itemize}
        \item $\sin x + \cos x = 0 \iff \tan x = -1 \iff x = -\dfrac{\pi}{4} + k\pi$ ($k\in\Z$).
        \item $4\sin x\cos x - 1 = 0 \iff 2\sin(2x) = 1 \iff \sin(2x) = \frac{1}{2} \iff x = \frac{\pi}{12} + k\pi$ ou $x = \frac{5\pi}{12} + k\pi$ ($k\in\Z$).
    \end{itemize}
    D'où $\boxed{S = \left\{-\frac{\pi}{4} + k\pi,\; \frac{\pi}{12} + k\pi,\; \frac{5\pi}{12} + k\pi \;\middle|\; k\in\Z\right\}}$.
\end{enumerate}

\paragraph{Corrigé de l'Exercice 3.4 (UCLouvain -- Septembre 2020, Q1) :}
\textbf{CE :} $\cos x \neq 0 \iff x \neq \frac{\pi}{2} + k\pi$ ($k\in\Z$). En multipliant toute l'équation par $\cos^2 x \neq 0$ et en remplaçant $\cos^2 x = 1 - \sin^2 x$ :
\[
    \sin^2 x - 3\sin x - 1 = 1 - \sin^2 x \iff 2\sin^2 x - 3\sin x - 2 = 0 \iff (2\sin x + 1)(\sin x - 2) = 0.
\]
Comme $\sin x = 2$ est impossible, il reste $\sin x = -\frac{1}{2} \iff \boxed{x = -\frac{\pi}{6} + 2k\pi \text{ ou } x = \frac{7\pi}{6} + 2k\pi}$ ($k\in\Z$), compatibles avec la CE.

\paragraph{Corrigé de l'Exercice 3.5 (ULiège -- Septembre 2002 Q2, Juillet 2003 Q3 et Juillet 2008 Q2) :}
\begin{enumerate}[label=(\alph*)]
    \item \textbf{CE :} $\sin x \neq 0$ et $\cot x + \cos x = \cos x\left(\frac{1+\sin x}{\sin x}\right) \neq 0 \iff x \neq k\frac{\pi}{2}$ ($k\in\Z$). Sous ces CE, on simplifie la fraction par $\frac{\cos x}{\sin x}$ :
    \[
        \frac{1 - \sin x}{1 + \sin x} = 2(1 - \sin x) \iff (1 - \sin x)\frac{-1-2\sin x}{1+\sin x} = 0 \iff \sin x = 1 \text{ ou } \sin x = -\frac{1}{2}.
    \]
    La solution $\sin x = 1$ ($x = \frac{\pi}{2}+2k\pi$) est \textbf{rejetée} par les CE ! Il reste $\sin x = -\frac{1}{2} \iff \boxed{x = -\frac{\pi}{6} + 2k\pi \text{ ou } x = \frac{7\pi}{6} + 2k\pi}$ ($k\in\Z$).
    \item Par Simpson à gauche ($\cos 2x + \cos 6x = 2\cos 4x\cos 2x$) et Carnot à droite ($1 + \cos 8x = 2\cos^2 4x$) :
    \begin{align*}
        2\cos(4x)\bigl(\cos(2x) - \cos(4x)\bigr) = 0 &\iff \cos(4x) = 0 \quad \text{ou} \quad \cos(4x) = \cos(2x) \\
        &\iff \boxed{x = \frac{\pi}{8} + \frac{k\pi}{4} \quad \text{ou} \quad x = \frac{k\pi}{3}} \qquad (k\in\Z).
    \end{align*}
    \item À gauche, par Simpson sur le 1\up{er} et le 3\up{e} terme : $\sin x + \sin 3x + \sin 2x = 2\sin 2x\cos x + \sin 2x = \sin(2x)(2\cos x + 1)$. À droite, par Carnot sur $1 + \cos 2x = 2\cos^2 x$ : $\sqrt{2}(\cos x + 2\cos^2 x) = \sqrt{2}\cos x(2\cos x + 1)$. En factorisant :
    \[
        (2\cos x + 1)\bigl(\sin(2x) - \sqrt{2}\cos x\bigr) = \cos x(2\cos x + 1)(2\sin x - \sqrt{2}) = 0.
    \]
    D'où $\cos x = 0$ ($x = \frac{\pi}{2}+k\pi$), $\cos x = -\frac{1}{2}$ ($x = \pm\frac{2\pi}{3}+2k\pi$) ou $\sin x = \frac{\sqrt{2}}{2}$ ($x = \frac{\pi}{4}+2k\pi$ ou $x = \frac{3\pi}{4}+2k\pi$).
\end{enumerate}

\paragraph{Corrigé de l'Exercice 3.6 (ULiège -- Septembre 2005 Q2, Juillet 2006 Q4 et Septembre 2009 Q2) :}
\begin{enumerate}[label=(\alph*)]
    \item \textbf{CE :} $x \neq \frac{\pi}{2} + k\pi$. En posant $t = \tan x$, $t^3 - 3t^2 - 3t + 1 = (t^3 + 1) - 3t(t + 1) = (t + 1)(t^2 - 4t + 1) = 0$.
    \begin{itemize}
        \item $t = -1 \iff \tan x = -1 \iff x = -\frac{\pi}{4} + k\pi$.
        \item $t^2 - 4t + 1 = 0 \iff t = 2 \pm \sqrt{3} \iff x = \frac{5\pi}{12} + k\pi$ ou $x = \frac{\pi}{12} + k\pi$.
    \end{itemize}
    Ces trois familles d'écart $\frac{\pi}{3}$ se regroupent en une seule : $\boxed{x = \frac{\pi}{12} + \frac{k\pi}{3}}$ ($k\in\Z$).
    \item \textbf{CE :} $x \neq (2k+1)\pi$ et $x \neq \frac{\pi}{2} + 2k\pi$. En exprimant $\tan(x/2) = \frac{\sin(x/2)}{\cos(x/2)}$ :
    \[
        \left(\frac{1 + \tan(x/2)}{1 - \tan(x/2)}\right)^2 = \left(\frac{\cos(x/2) + \sin(x/2)}{\cos(x/2) - \sin(x/2)}\right)^2 = \frac{1 + \sin x}{1 - \sin x}.
    \]
    Sous la CE $\sin x \neq 1$, l'équation devient $1 + \sin x + \sin x(1 - \sin x) = 0 \iff \sin^2 x - 2\sin x - 1 = 0 \iff \sin x = 1 \pm \sqrt{2}$. Comme $1+\sqrt{2} > 1$, il reste $\sin x = 1 - \sqrt{2} \in [-1, 1]$, soit $\boxed{x = \arcsin(1-\sqrt{2}) + 2k\pi \text{ ou } x = \pi - \arcsin(1-\sqrt{2}) + 2k\pi}$ ($k\in\Z$).
    \item \textbf{CE :} $x \neq \frac{\pi}{2}+k\pi$, $x \neq \frac{\pi}{4}+\frac{k\pi}{2}$, $x \neq \frac{\pi}{6}+\frac{k\pi}{3}$. En développant $\tan(3x) = \tan(x + 2x) = \frac{\tan x + \tan 2x}{1 - \tan x\tan 2x}$ et en mettant $(\tan x + \tan 2x)$ en évidence :
    \[
        (\tan x + \tan 2x)\left(1 + \frac{1}{1 - \tan x\tan 2x}\right) = 0 \iff \tan(2x) = \tan(-x) \text{ ou } \tan x\tan 2x = 2.
    \]
    La première donne $x = \frac{k\pi}{3}$, et la seconde donne $\frac{2\tan^2 x}{1-\tan^2 x} = 2 \iff 4\tan^2 x = 2 \iff \tan x = \pm\frac{\sqrt{2}}{2} \iff x = \pm\arctan\left(\frac{\sqrt{2}}{2}\right) + k\pi$ ($k\in\Z$).
\end{enumerate}

\paragraph{Corrigé de l'Exercice 3.7 (ULiège -- Juillet 2005 Q2, Juillet 2007 Q2 et Juillet 2010 Q2) :}
En utilisant les identités classiques $\sin^4 x + \cos^4 x = 1 - \frac{1}{2}\sin^2(2x)$ et $\sin^6 x + \cos^6 x = 1 - 3\sin^2 x\cos^2 x = 1 - \frac{3}{4}\sin^2(2x)$ :
\begin{enumerate}[label=(\alph*)]
    \item $1 - \frac{3}{4}\sin^2(2x) = \frac{5}{8} \iff \sin^2(2x) = \frac{1}{2} \iff \sin(2x) = \pm\frac{\sqrt{2}}{2} \iff 2x = \frac{\pi}{4} + \frac{k\pi}{2} \iff \boxed{x = \frac{\pi}{8} + \frac{k\pi}{4}}$ ($k\in\Z$).
    \item $1 - \frac{3}{4}\sin^2(2x) = \frac{7}{16} \iff \sin^2(2x) = \frac{3}{4} \iff \sin(2x) = \pm\frac{\sqrt{3}}{2} \iff 2x = \pm\frac{\pi}{3} + k\pi \iff \boxed{x = \pm\frac{\pi}{6} + \frac{k\pi}{2}}$ ($k\in\Z$).
    \item $1 - \frac{1}{2}\sin^2(2x) = \frac{1}{2}\sin(2x) \iff \sin^2(2x) + \sin(2x) - 2 = 0 \iff \bigl(\sin(2x) - 1\bigr)\bigl(\sin(2x) + 2\bigr) = 0 \iff \sin(2x) = 1 \iff 2x = \frac{\pi}{2} + 2k\pi \iff \boxed{x = \frac{\pi}{4} + k\pi}$ ($k\in\Z$).
\end{enumerate}

\paragraph{Corrigé de l'Exercice 3.8 (Inéquations trigonométriques -- Inspiré d'ULiège) :}
\begin{enumerate}[label=(\alph*)]
    \item Sur le cercle trigonométrique, $\sin u \le -\frac{\sqrt{3}}{2} \iff -\frac{2\pi}{3} + 2k\pi \le u \le -\frac{\pi}{3} + 2k\pi$. Avec $u = 2x$, on divise par $2$ :
    \[
        \boxed{-\frac{\pi}{3} + k\pi \le x \le -\frac{\pi}{6} + k\pi \qquad (k \in \Z)}.
    \]
    \item Par Simpson : $\sin\left(x - \frac{\pi}{3}\right) - \sin x = 2\sin\left(-\frac{\pi}{6}\right)\cos\left(x - \frac{\pi}{6}\right) = -\cos\left(x - \frac{\pi}{6}\right) > 0 \iff \cos\left(x - \frac{\pi}{6}\right) < 0$. Or $\cos u < 0 \iff \frac{\pi}{2} + 2k\pi < u < \frac{3\pi}{2} + 2k\pi$. En remplaçant $u = x - \frac{\pi}{6}$ :
    \[
        \boxed{\frac{2\pi}{3} + 2k\pi < x < \frac{5\pi}{3} + 2k\pi \qquad (k \in \Z)}.
    \]
    \item En factorisant le trinôme $2y^2 - y - 1 = (y - 1)(2y + 1)$ avec $y = \cos x$, l'inéquation s'écrit $(\cos x - 1)(2\cos x + 1) \ge 0$. Comme $\cos x \le 1$ pour tout $x$, le facteur $(\cos x - 1)$ est toujours négatif ou nul (nul en $x = 2k\pi$). Pour que le produit soit positif ou nul, il faut donc soit $\cos x = 1$ ($x = 2k\pi$), soit $2\cos x + 1 \le 0 \iff \cos x \le -\frac{1}{2}$ :
    \[
        \boxed{x \in \{2k\pi \mid k\in\Z\} \cup \bigcup_{k\in\Z} \left[\frac{2\pi}{3} + 2k\pi,\; \frac{4\pi}{3} + 2k\pi\right]}.
    \]
\end{enumerate}

% ----------------------------------------------------------------------------
\section{Corrigés du Chapitre 4 (Résolution de Triangles et Problèmes)}
% ----------------------------------------------------------------------------

\paragraph{Corrigé de l'Exercice 4.1 (UCLouvain -- Septembre 2024, Q2 : Théorème de la bissectrice) :}
Notons $\theta = \widehat{AXB}$, de sorte que $\widehat{AXC} = \pi - \theta$ et $\sin(\widehat{AXC}) = \sin(\pi-\theta) = \sin\theta$. Par la loi des sinus dans les triangles $ABX$ (de côtés $x = \overline{BX}$ et $c = \overline{AB}$) et $ACX$ (de côtés $a-x = \overline{XC}$ et $b = \overline{AC}$) :
\[
    \frac{x}{\sin\alpha} = \frac{c}{\sin\theta} \quad \text{et} \quad \frac{a-x}{\sin\alpha} = \frac{b}{\sin(\pi-\theta)} = \frac{b}{\sin\theta}.
\]
En divisant membre à membre ces deux égalités ($\sin\alpha$ et $\sin\theta$ se simplifient), on obtient le théorème de la bissectrice :
\[
    \frac{x}{a-x} = \frac{c}{b} \iff bx = c(a-x) \iff (b+c)x = ac \iff \boxed{x = \frac{ac}{b+c}}.
\]

\paragraph{Corrigé de l'Exercice 4.2 (UCLouvain -- Mini-Blocus 2026 Q1e, Septembre 2026 Q3e et Septembre 2023 Q3f) :}
\begin{enumerate}[label=(\alph*)]
    \item Considérons un triangle rectangle isocèle d'hypoténuse $c$ et de cathètes $a = b = \frac{c}{\sqrt{2}}$.
    \begin{itemize}
        \item Le rayon du cercle circonscrit à tout triangle rectangle est la moitié de son hypoténuse : $R = \dfrac{c}{2}$.
        \item L'aire du triangle vaut $S = \frac{1}{2}ab = \frac{c^2}{4}$, et son demi-périmètre vaut $p = \frac{a+b+c}{2} = \frac{c(\sqrt{2}+1)}{2}$. Par la formule $r = \frac{S}{p}$ (ou la formule du triangle rectangle $r = \frac{a+b-c}{2}$), on obtient $r = \frac{c(\sqrt{2}-1)}{2}$.
    \end{itemize}
    Le rapport vaut donc $\dfrac{r}{R} = \boxed{\sqrt{2} - 1}$.
    \item Plaçons l'angle droit à l'origine $(0,0)$ avec les deux cathètes de longueur $a = \frac{1}{\sqrt{2}}$ sur les axes $Ox$ et $Oy$. Le centre du cercle inscrit est en $I = (r, r)$, et le centre du cercle circonscrit (milieu de l'hypoténuse) est en $O = \left(\frac{a}{2}, \frac{a}{2}\right) = \left(\frac{\sqrt{2}}{4}, \frac{\sqrt{2}}{4}\right)$. Comme ces deux points sont sur la bissectrice $y=x$, leur distance vaut :
    \[
        d = \sqrt{2}\left(\frac{\sqrt{2}}{4} - r\right) = \frac{1}{2} - \sqrt{2}\cdot\frac{\sqrt{2}-1}{2} = \frac{1 - 2 + \sqrt{2}}{2} = \boxed{\frac{\sqrt{2}-1}{2}}.
    \]
\end{enumerate}

\paragraph{Corrigé de l'Exercice 4.3 (UCLouvain -- Juillet 2026 Q3e et Septembre 2026 Q3a : Trapèzes) :}
\begin{enumerate}[label=(\alph*)]
    \item Dans le trapèze rectangle $ABCD$ (rectangle en $A$ et $D$, de hauteur $h = \overline{AD}$), notons $\alpha = \widehat{ABD}$ l'angle que fait la diagonale $[BD]$ avec la base $[AB]$. Dans le triangle $ABD$ rectangle en $A$, $\tan\alpha = \frac{h}{x}$. Comme $[AC] \perp [BD]$, l'angle $\widehat{DAC}$ est complémentaire à $\widehat{ADB}$, donc égal à $\alpha$ ! Dans le triangle $ADC$ rectangle en $D$, on a alors $\tan\alpha = \frac{y}{h}$. En égalant les deux expressions de $\tan\alpha$ :
    \[
        \frac{h}{x} = \frac{y}{h} \implies h^2 = xy \implies \boxed{h = \overline{AD} = \sqrt{xy}} \quad \text{(moyenne géométrique des bases)}.
    \]
    \item Dans le trapèze \textbf{isocèle} $ABCD$ dont les diagonales $[AC]$ et $[BD]$ se coupent perpendiculairement en $O$, par symétrie les triangles $OAB$ et $OCD$ sont tous deux \textbf{rectangles isocèles en $O$}. La hauteur de $OAB$ issue de $O$ (médiane relative à l'hypoténuse $x = \overline{AB}$) vaut donc $h_1 = \frac{x}{2}$, et celle de $OCD$ vaut $h_2 = \frac{y}{2}$. La hauteur totale du trapèze vaut $h = h_1 + h_2 = \frac{x+y}{2}$, et son aire vaut :
    \[
        \text{Aire} = \frac{x+y}{2} \cdot h = \boxed{\frac{(x+y)^2}{4}}.
    \]
\end{enumerate}

\paragraph{Corrigé de l'Exercice 4.4 (UCLouvain -- Juillet 2022, Q3b) :}
Soient $A, B, C$ les trois points consécutifs sur le cercle de centre $O$ et de rayon $R$ tels que $\overline{AB} = 2$ (angle au centre $\alpha$), $\overline{BC} = 3$ (angle au centre $\beta$) et $\overline{AC} = 4$ (angle au centre $\alpha+\beta$). D'après le théorème de l'angle inscrit, l'angle inscrit $\widehat{ACB}$ interceptant le même arc $\overset{\frown}{AB}$ que l'angle au centre $\alpha$ vaut exactement $\widehat{ACB} = \dfrac{\alpha}{2}$ !
En appliquant la \textbf{loi des cosinus} dans le triangle $ABC$ (de côtés $2, 3, 4$) pour l'angle $\widehat{ACB} = \frac{\alpha}{2}$ :
\[
    2^2 = 3^2 + 4^2 - 2\cdot 3\cdot 4\cos\left(\frac{\alpha}{2}\right) \implies 4 = 25 - 24\cos\left(\frac{\alpha}{2}\right) \implies \cos\left(\frac{\alpha}{2}\right) = \frac{21}{24} = \frac{7}{8}.
\]
Enfin, par la formule de Carnot / duplication :
\[
    \cos(\alpha) = 2\cos^2\left(\frac{\alpha}{2}\right) - 1 = 2\left(\frac{7}{8}\right)^2 - 1 = \frac{98}{64} - 1 = \frac{34}{64} = \boxed{\frac{17}{32}}.
\]

\paragraph{Corrigé de l'Exercice 4.5 (UCLouvain -- Juillet 2021, Q3a, Q3c et Q3f) :}
\begin{enumerate}[label=(\alph*)]
    \item Chaque côté d'un polygone régulier à $n$ côtés est vu depuis le centre $O$ du cercle de rayon $R$ sous un angle $\theta = \frac{2\pi}{n}$. Le côté du polygone inscrit est une corde de longueur $c_I = 2R\sin\left(\frac{\pi}{n}\right)$, tandis que le côté du polygone circonscrit (tangent au cercle) vaut $c_C = 2R\tan\left(\frac{\pi}{n}\right)$. Le rapport des périmètres vaut donc :
    \[
        \frac{P_C}{P_I} = \frac{n \cdot 2R\tan(\pi/n)}{n \cdot 2R\sin(\pi/n)} = \boxed{\frac{1}{\cos(\pi/n)}}.
    \]
    \item Par Carnot ($\cos^2(A-B) = \frac{1+\cos(2A-2B)}{2}$) et Simpson ($\cos(2A-2B) - 2\cos(2A)\cos(2B) = -\cos(2A+2B)$) :
    \[
        \cos^2(A-B) - \cos(2A)\cos(2B) = \frac{1 - \cos(2A+2B)}{2} = \sin^2(A+B) = \sin^2(\pi - C) = \sin^2 C.
    \]
    On a donc $\sin^2 C = \frac{1}{2} \implies \sin C = \frac{\sqrt{2}}{2}$ (car $C \in \left]0, \pi\right[$). Comme $\widehat{C}$ est obtus, $\boxed{\widehat{C} = 135\degree}$.
    \item Le triangle $ABC$ étant rectangle en $A$, son hypoténuse est le diamètre du cercle circonscrit : $a = \overline{BC} = 2R = 4$. Par Pythagore, $\overline{AC} = \sqrt{4^2 - (\sqrt{7})^2} = \sqrt{9} = 3$. De plus, comme $\widehat{B} + \widehat{C} = \frac{\pi}{2}$, on a $\cos(\widehat{B}-\widehat{C}) = \cos\left(2\widehat{B} - \frac{\pi}{2}\right) = \sin(2\widehat{B}) = 2\sin\widehat{B}\cos\widehat{B}$. Or $\cos\widehat{B} = \frac{\overline{AB}}{\overline{BC}} = \frac{\sqrt{7}}{4}$ et $\sin\widehat{B} = \frac{\overline{AC}}{\overline{BC}} = \frac{3}{4}$, d'où :
    \[
        \cos(\widehat{B} - \widehat{C}) = 2 \cdot \frac{3}{4} \cdot \frac{\sqrt{7}}{4} = \boxed{\frac{3\sqrt{7}}{8}}.
    \]
\end{enumerate}

\paragraph{Corrigé de l'Exercice 4.6 (UCLouvain -- Juillet 2024 Q3c et Septembre 2024 Q3c) :}
\begin{enumerate}[label=(\alph*)]
    \item Par la formule de l'aire $S = \frac{1}{2}bc\sin\widehat{A}$ avec $S = 3$, $b = \overline{AC} = 2$ et $\sin\widehat{A} = \sin\frac{\pi}{6} = \frac{1}{2}$, on trouve $3 = \frac{1}{2}\cdot 2 \cdot c \cdot \frac{1}{2} \implies c = \overline{AB} = 6\text{ cm}$. Puis par la loi des cosinus :
    \[
        \overline{BC} = a = \sqrt{b^2 + c^2 - 2bc\cos\left(\frac{\pi}{6}\right)} = \sqrt{4 + 36 - 24\cdot\frac{\sqrt{3}}{2}} = \boxed{\sqrt{40 - 12\sqrt{3}}\text{ cm}}.
    \]
    \item Comme $\widehat{A} = 60\degree$ et $\widehat{B} = 45\degree$, on a $\widehat{C} = 180\degree - 105\degree = 75\degree$, dont le sinus vaut $\sin 75\degree = \frac{\sqrt{2}(\sqrt{3}+1)}{4}$. Par la loi des sinus $\frac{a}{\sin A} = \frac{b}{\sin B} = \frac{c}{\sin C}$, il vient $\frac{a+b}{c} = \frac{\sin 60\degree + \sin 45\degree}{\sin 75\degree} = \frac{2(\sqrt{3}+\sqrt{2})}{\sqrt{2}(\sqrt{3}+1)}$, d'où :
    \[
        \sqrt{2}(\sqrt{3}+1)\frac{a+b}{c} = \boxed{2(\sqrt{3}+\sqrt{2})}.
    \]
\end{enumerate}

\paragraph{Corrigé de l'Exercice 4.7 (UCLouvain -- Juillet 2025 / Recueil 2025, Q3.5) :}
Le triangle $OAB$ a deux côtés égaux au rayon $\overline{OA} = \overline{OB} = 1$ et un angle au sommet $\widehat{AOB} = \frac{\pi}{3}$ ($60\degree$) : il est donc \textbf{équilatéral} ! Ainsi, $\overline{AB} = 1$ et $\widehat{OBA} = \frac{\pi}{3} = 60\degree$.
Puisque $C$ est sur la droite $(AB)$ au-delà de $B$ avec $\overline{BC} = \overline{AB} = 1$, on a $\overline{OB} = \overline{BC} = 1$, donc le triangle $OBC$ est \textbf{isocèle en $B$}. Son angle au sommet $B$ est le supplément de $\widehat{OBA}$ : $\widehat{OBC} = \pi - \frac{\pi}{3} = \frac{2\pi}{3}$ ($120\degree$). Les deux angles à la base du triangle isocèle $OBC$ sont donc égaux à :
\[
    \widehat{BOC} = \widehat{BCO} = \frac{\pi - \frac{2\pi}{3}}{2} = \boxed{\frac{\pi}{6} \quad (30\degree)}.
\]

\paragraph{Corrigé de l'Exercice 4.8 (UCLouvain -- Juillet 2022, Q2) :}
\begin{enumerate}[label=(\alph*)]
    \item La diagonale du carré $ABCD$ inscrit dans le cercle de rayon $1$ vaut $2R = 2$. Par Pythagore, $2\overline{AB}^2 = 2^2 = 4 \implies \overline{AB} = \boxed{\sqrt{2}}$, et la distance du centre $O$ à chaque côté du carré vaut $\frac{\sqrt{2}}{2}$.
    \item Posons $x = \overline{HE}$ le côté du petit carré $EFGH$. Le point $H$ a pour coordonnées (par rapport à $O$) une abscisse de $\frac{x}{2}$ et une ordonnée de $\frac{\sqrt{2}}{2} + x$. Comme $H$ appartient au cercle de rayon $1$ :
    \[
        \left(\frac{x}{2}\right)^2 + \left(\frac{\sqrt{2}}{2} + x\right)^2 = 1 \iff \frac{5}{4}x^2 + \sqrt{2}x - \frac{1}{2} = 0 \implies x = \frac{\sqrt{2}}{5} \implies \boxed{\frac{\overline{AB}}{\overline{HE}} = \frac{\sqrt{2}}{\sqrt{2}/5} = 5}.
    \]
    \item Posons $y = \overline{JK}$. Le point $K$ sur le cercle a pour ordonnée $y + \frac{\sqrt{2}}{2}$ et pour abscisse $\frac{\overline{IJ}}{2} = \sqrt{1 - \left(y + \frac{\sqrt{2}}{2}\right)^2}$. L'aire du rectangle vaut donc :
    \[
        A(y) = \overline{IJ}\cdot\overline{JK} = 2y\sqrt{1 - \left(y + \frac{\sqrt{2}}{2}\right)^2}.
    \]
    Pour maximiser $A(y)$, on maximise $A^2(y) = 4\left(-y^4 - \sqrt{2}y^3 + \frac{1}{2}y^2\right)$ en annulant sa dérivée : $-4y^3 - 3\sqrt{2}y^2 + y = 0 \implies y^2 + \frac{3\sqrt{2}}{4}y - \frac{1}{4} = 0 \implies y = \frac{\sqrt{2}(\sqrt{17}-3)}{8}$. Comme $\overline{DC} = \sqrt{2}$, le rapport optimal vaut $\boxed{\dfrac{\overline{DC}}{\overline{JK}} = \dfrac{8}{\sqrt{17}-3} = \sqrt{17}+3}$.
\end{enumerate}

\paragraph{Corrigé de l'Exercice 4.9 (ULiège -- Juillet 2008, Q3 : Trois cordes dans un demi-cercle) :}
Notons $d$ la distance entre le centre $O$ du demi-cercle et la corde la plus basse $C_3 = 20\text{ m}$ (demi-corde $10\text{ m}$). Les distances de $O$ aux cordes $C_2 = 16\text{ m}$ (demi-corde $8\text{ m}$) et $C_1 = 8\text{ m}$ (demi-corde $4\text{ m}$) sont alors $d+h$ et $d+2h$. Par le théorème de Pythagore dans les trois triangles rectangles d'hypoténuse $R$ :
\[
    \begin{cases}
        R^2 = d^2 + 10^2 = d^2 + 100 & (1) \\
        R^2 = (d+h)^2 + 8^2 = d^2 + 2dh + h^2 + 64 & (2) \\
        R^2 = (d+2h)^2 + 4^2 = d^2 + 4dh + 4h^2 + 16 & (3)
    \end{cases}
\]
Les différences $(2)-(1)$ et $(3)-(1)$ donnent $2dh + h^2 = 36$ et $4dh + 4h^2 = 84$. En combinant $(4dh+4h^2) - 2(2dh+h^2) = 84 - 72$, on élimine $dh$ :
\[
    2h^2 = 12 \implies \boxed{h = \sqrt{6}\text{ m}}, \qquad \text{puis } d = \frac{36 - 6}{2\sqrt{6}} = \frac{15}{\sqrt{6}}\text{ m},
\]
et enfin l'équation $(1)$ fournit le rayon :
\[
    R = \sqrt{d^2 + 100} = \sqrt{\frac{225}{6} + 100} = \sqrt{37,5 + 100} = \boxed{\sqrt{137,5}\text{ m} = \frac{\sqrt{550}}{2}\text{ m} \approx 11,73\text{ m}}.
\]

\paragraph{Corrigé de l'Exercice 4.10 (ULiège -- Septembre 2006, Q2) :}
Appliquons la loi des sinus dans les trois triangles formés par la sécante $\mathcal{D}$ :
\begin{itemize}
    \item Dans le triangle $DCE$ : $\dfrac{\overline{DC}}{\sin\widehat{DEC}} = \dfrac{\overline{EC}}{\sin\widehat{CDE}} \implies \overline{DC} = \overline{EC}\,\dfrac{\sin\widehat{DEC}}{\sin\widehat{CDE}}$.
    \item Dans le triangle $DBF$ : $\dfrac{\overline{DB}}{\sin\widehat{DFB}} = \dfrac{\overline{FB}}{\sin\widehat{BDF}} \implies \overline{DB} = \overline{FB}\,\dfrac{\sin\widehat{DFB}}{\sin\widehat{BDF}}$.
\end{itemize}
Comme $D, C, B$ sont alignés, $\widehat{CDE} = \widehat{BDF}$. De plus, $\widehat{DEC} = \widehat{AEF}$ (angles opposés par le sommet) et $\widehat{DFB} = \pi - \widehat{AFE}$ (angles supplémentaires, donc $\sin\widehat{DFB} = \sin\widehat{AFE}$). En faisant le quotient :
\[
    \frac{\overline{DC}}{\overline{DB}} = \frac{\overline{EC}}{\overline{FB}} \cdot \frac{\sin\widehat{AEF}}{\sin\widehat{AFE}}.
\]
Or, par la loi des sinus dans le triangle $AEF$, on a $\dfrac{\sin\widehat{AEF}}{\sin\widehat{AFE}} = \dfrac{\overline{FA}}{\overline{EA}}$. On retrouve ainsi le célèbre théorème de Ménélaüs :
\[
    \frac{\overline{DC}}{\overline{DB}} = \frac{\overline{EC}}{\overline{EA}} \cdot \frac{\overline{FA}}{\overline{FB}} = \frac{3}{2} \cdot \frac{1}{3} = \boxed{\frac{1}{2}}.
\]


% ----------------------------------------------------------------------------
\section{Corrigés du Chapitre 5 (Géométrie Synthétique Plane et Spatiale)}
% ----------------------------------------------------------------------------

\paragraph{Corrigé de l'Exercice 5.1 (UCLouvain -- Juillet 2026, Q3 : Carré et arc de cercle) :}
Tracons par le point $X$ (milieu de $[DY]$) la droite parallèle à $(AD)$, qui coupe $[AB]$ en $M$ et $[DC]$ en $N$.
\begin{enumerate}
    \item Les triangles rectangles $DYC$ (rectangle en $C$) et $DXN$ (rectangle en $N$) partagent l'angle aigu $\widehat{YDC}$ : ils sont semblables. Comme $X$ est le milieu de l'hypoténuse $[DY]$, $N$ est le milieu de $[DC]$, et par conséquent $M$ est le milieu de $[AB]$, soit $|AM| = \frac{c}{2}$.
    \item Le point $X$ appartient au quart de cercle de centre $A$ et de rayon $c$, donc $|AX| = c$. Dans le triangle $AMX$ rectangle en $M$, on a $\sin\widehat{MAX} = \frac{|MX|}{|AX|}$... Attention : le côté de longueur $\frac{c}{2}$ est $|AM|$ (adjacent à $\widehat{MAX}$, donc opposé à $\widehat{AXM} = \widehat{DAX}$ car $(MN) \parallel (AD)$ !). Ainsi, $\sin\widehat{DAX} = \frac{|AM|}{|AX|} = \frac{c/2}{c} = \frac{1}{2}$, d'où $|\widehat{DAX}| = 30\degree$.
    \item Enfin, comme $|AD| = |AX| = c$, le triangle $ADX$ est \textbf{isocèle en $A$} ! Ses deux angles à la base $[DX]$ (qui porte $[DY]$) sont donc égaux à :
    \[
        |\widehat{ADY}| = |\widehat{ADX}| = \frac{180\degree - |\widehat{DAX}|}{2} = \frac{180\degree - 30\degree}{2} = \boxed{75\degree}.
    \]
\end{enumerate}

\paragraph{Corrigé de l'Exercice 5.2 (UCLouvain -- Juillet 2024, Q3 : Tangence inscrite) :}
\begin{enumerate}[label=(\alph*)]
    \item La projection horizontale du cercle $\mathcal{C}$ de rayon $r$ sur le diamètre $[AB]$ est le segment $[PQ]$ de longueur $|PQ| = 2r$ (dont le milieu est le point de tangence $S$). Le diamètre total du demi-cercle $\mathcal{D}$ vaut donc :
    \[
        2R = |AB| = |AP| + |PQ| + |QB| = 2 + 2r + 8 = 2r + 10 \implies \boxed{R(r) = r + 5}.
    \]
    \item Le demi-cercle $\mathcal{D}(O, R)$ et le cercle $\mathcal{C}(C, r)$ sont tangents intérieurement au point $T$, donc les points $O, C, T$ sont alignés et $|OC| = R - r = 5$. De plus, $S$ est le point de tangence de $\mathcal{C}$ avec $[AB]$, donc $|CS| = r$, $(CS) \perp (AB)$ et $|AS| = |AP| + |PS| = 2 + r$, d'où $|OS| = |OA| - |AS| = R - (2+r) = (r+5) - (r+2) = 3$.\\
    Dans le triangle $OSC$ rectangle en $S$, le théorème de Pythagore donne :
    \[
        |OC|^2 = |OS|^2 + |CS|^2 \iff 5^2 = 3^2 + r^2 \implies r = 4, \quad \text{et donc } \boxed{R = r + 5 = 9}.
    \]
\end{enumerate}

\paragraph{Corrigé de l'Exercice 5.3 (UCLouvain -- Juillet 2023, Q3 : Surface cachée de l'hexagone) :}
\begin{enumerate}[label=(\alph*)]
    \item Dans un hexagone régulier $ABCDEF$ de côté $c$, le grand diamètre $[AD]$ mesure $2c$ et chaque angle intérieur vaut $120\degree$ ($|\widehat{BAD}| = 60\degree$, $|\widehat{BCD}| = 120\degree$). Dans le triangle isocèle $BCD$ ($|BC|=|CD|=c$), les angles à la base valent $30\degree$, donc $|\widehat{ABD}| = |\widehat{ABC}| - |\widehat{CBD}| = 120\degree - 30\degree = 90\degree$.\\
    Comme $\widehat{ABD} = 90\degree$ et que $M, B, N$ sont alignés, les angles $\widehat{MBA}$ et $\widehat{NBD}$ sont complémentaires ($|\widehat{MBA}| + |\widehat{NBD}| = 90\degree$), donc $|\widehat{MAB}| = 90\degree - |\widehat{MBA}| = |\widehat{NBD}|$. Les triangles rectangles $MAB$ et $NBD$ ont les mêmes angles aigus : ils sont \textbf{semblables}.
    \item Dans le triangle $ABD$ rectangle en $B$, on a $|AB| = c$ et $|AD| = 2c$, donc par Pythagore $|BD| = \sqrt{(2c)^2 - c^2} = c\sqrt{3}$. Le rapport de similitude de $MAB$ (d'hypoténuse $|AB| = c$) vers $NBD$ (d'hypoténuse $|BD| = c\sqrt{3}$) vaut $\rho = \frac{|BD|}{|AB|} = \sqrt{3}$. Le rapport de leurs aires vaut donc :
    \[
        r = \frac{\text{Aire}(\triangle NBD)}{\text{Aire}(\triangle MAB)} = \rho^2 = (\sqrt{3})^2 = \boxed{3}.
    \]
    \item On a $\text{Aire}(\triangle NBD) = 3 \times \text{Aire}(\triangle MAB) = 3 \times 20 = 60$. En retranchant l'aire du quadrilatère concave $NBCD$ ($24$), on obtient l'aire du triangle $BCD$ :
    \[
        \text{Aire}(\triangle BCD) = \text{Aire}(\triangle NBD) - \text{Aire}(NBCD) = 60 - 24 = 36.
    \]
    Or, un hexagone régulier de côté $c$ est composé de $6$ triangles équilatéraux d'aire $\frac{\sqrt{3}}{4}c^2$, tandis que $\text{Aire}(\triangle BCD) = \frac{1}{2}c^2\sin 120\degree = \frac{\sqrt{3}}{4}c^2$ (exactement $\frac{1}{6}$ de l'hexagone!). L'aire totale de l'hexagone vaut donc $6 \times 36 = \boxed{216}$.
\end{enumerate}

\paragraph{Corrigé de l'Exercice 5.4 (UCLouvain -- Septembre 2024, Q3 : Rapport d'aires paramétré) :}
\begin{enumerate}[label=(\alph*)]
    \item Le quadrilatère $NBOD$ a ses côtés opposés parallèles deux à deux ($(NB) \parallel (DO)$ par construction du rectangle $BPDM$, et $(ND) \parallel (BO)$ côtés opposés du carré $ABCD$) : c'est un \textbf{parallélogramme} de base $|ND| = |AD| - |AN| = (1-\lambda)c$ et de hauteur $|DC| = c$. Son aire vaut donc :
    \[
        \mathcal{Q} = |ND| \cdot |DC| = \boxed{(1-\lambda)c^2}.
    \]
    \item Les triangles $ABN$ (rectangle en $A$) et $MDN$ (rectangle en $M$) ont leurs angles en $N$ opposés par le sommet ($\widehat{ANB} = \widehat{MND}$) : ils sont donc \textbf{semblables}. Dans $ABN$, l'hypoténuse vaut $|NB| = \sqrt{|AB|^2 + |AN|^2} = c\sqrt{1+\lambda^2}$, tandis que dans $MDN$ l'hypoténuse est $|ND| = (1-\lambda)c$. Le rapport de similitude vaut donc :
    \[
        \rho(\lambda) = \frac{|ND|}{|NB|} = \boxed{\frac{1-\lambda}{\sqrt{1+\lambda^2}}}.
    \]
    \item Le rectangle $BPDM$ est la réunion du parallélogramme $NBOD$ et des deux triangles rectangles isométriques $MDN$ et $PBO$. Or $\text{Aire}(\triangle MDN) = \rho^2 \text{Aire}(\triangle ABN) = \frac{(1-\lambda)^2}{1+\lambda^2}\cdot\frac{1}{2}\lambda c^2$. Dès lors :
    \begin{align*}
        \text{Aire}(BPDM) &= \mathcal{Q} + 2\,\text{Aire}(\triangle MDN) = (1-\lambda)c^2 + \frac{(1-\lambda)^2\lambda}{1+\lambda^2}c^2 \\
        &= (1-\lambda)\left(\frac{1+\lambda^2+\lambda-\lambda^2}{1+\lambda^2}\right)c^2 = \frac{1-\lambda^2}{1+\lambda^2}c^2.
    \end{align*}
    En divisant par $\text{Aire}(ABCD) = c^2$, on obtient $R(\lambda) = \boxed{\dfrac{1-\lambda^2}{1+\lambda^2}}$.
\end{enumerate}

\paragraph{Corrigé de l'Exercice 5.5 (UCLouvain -- Juillet 2024, Q4b : Carré et quarts de cercle) :}
\begin{enumerate}[label=(\alph*)]
    \item Puisque $(DP) \parallel (QB)$ et $|DP| = |QB| = d$, le quadrilatère $DQBP$ est un \textbf{parallélogramme} dont les côtés mesurent $|DP| = |QB| = d$ et $|DQ| = |BP| = c$ (car $Q$ est sur le cercle de centre $D$ et de rayon $c$, et $P$ sur le cercle de centre $B$ et de rayon $c$). Ses deux diagonales sont $[DB]$ (diagonale du carré, de longueur $c\sqrt{2}$) et $[PQ]$ (de longueur $\ell$). Par l'identité du parallélogramme :
    \[
        |DB|^2 + |PQ|^2 = 2(|DQ|^2 + |DP|^2) \iff 2c^2 + \ell^2 = 2(c^2 + d^2) \implies \boxed{\ell(d) = \sqrt{2}\,d}.
    \]
    \item La longueur $\ell = \sqrt{2}\,d$ est minimale lorsque la distance $d = |DP|$ entre $D$ et un point $P$ de l'arc de centre $B$ et de rayon $c$ est minimale, ce qui se produit lorsque $P$ (et $Q$) se situe sur la diagonale $[BD]$ ! En ce point, $d_{\min} = |BD| - c = (\sqrt{2}-1)c$, et la distance entre $P$ et $Q$ sur la diagonale $[BD]$ vaut $m(c) = |BD| - 2d_{\min} = \sqrt{2}c - 2(\sqrt{2}-1)c = \boxed{(2-\sqrt{2})c}$ (note : par continuité de la formule $\ell = \sqrt{2}d_{\min} = (2-\sqrt{2})c$, les deux approches donnent exactement le même résultat !).
\end{enumerate}

\paragraph{Corrigé de l'Exercice 5.6 (UCLouvain -- Juillet 2022, Q4 : Sélection « Top Cheffe ») :}
\begin{enumerate}[label=(\alph*)]
    \item En coupe transversale, les centres des $4$ petits cercles de rayon $r$ forment un carré de côté $2r$, dont le centre $O$ est le centre du grand cercle de rayon $R$. La distance de $O$ au centre $C$ d'un petit cercle est la demi-diagonale du carré de côté $2r$, soit $|OC| = r\sqrt{2}$. Le rayon du grand cercle tangent extérieurement aux petits cercles vaut donc $R = |OC| + r = (\sqrt{2}+1)r$, d'où :
    \[
        \frac{r}{R} = \frac{1}{\sqrt{2}+1} = \boxed{\sqrt{2} - 1}.
    \]
    \item La largeur $L$ de la pâte fait deux fois le tour du cylindre de rayon $R$ : $L = 2 \times (2\pi R) = 4\pi R \implies R = \frac{L}{4\pi}$. Le volume vide à l'intérieur du cannelloni de longueur $\ell$ vaut la différence entre le volume du grand cylindre et celui des $4$ petits cylindres :
    \[
        V = \pi R^2 \ell - 4(\pi r^2 \ell) = \pi R^2 \ell\left(1 - 4\left(\frac{r}{R}\right)^2\right) = \pi\ell\frac{L^2}{16\pi^2}\bigl(1 - 4(\sqrt{2}-1)^2\bigr) = \boxed{\frac{\ell L^2}{16\pi}(8\sqrt{2} - 11)}.
    \]
\end{enumerate}

\paragraph{Corrigé de l'Exercice 5.7 (UCLouvain -- Septembre 2022, Q4 : Tangence équilatérale) :}
\begin{enumerate}[label=(\alph*)]
    \item Notons $A$ et $B$ les milieux de deux côtés du petit triangle équilatéral de côté $2a$ (centres de deux demi-cercles de rayon $a$). Le segment $[AB]$ qui joint les milieux de deux côtés est parallèle au troisième côté et mesure la moitié de celui-ci : $|AB| = a$. Les rayons $[BC]$ et $[AD]$ perpendiculaires au côté supérieur $[FE]$ du grand triangle ont pour longueur $a$, donc $ABCD$ est un \textbf{carré} de côté $|CD| = a$.\\
    Au sommet $E$ du grand triangle équilatéral ($\widehat{E} = 60\degree$), le triangle $ADE$ est rectangle en $D$ ($[AD] \perp [FE]$) avec $|AD| = a$ et $|\widehat{DAE}| = 30\degree$, donc $|DE| = a\tan 60\degree = a\sqrt{3}$. Par symétrie ($|FC| = |DE| = a\sqrt{3}$), le côté du grand triangle vaut :
    \[
        L = |FC| + |CD| + |DE| = a\sqrt{3} + a + a\sqrt{3} = \boxed{(2\sqrt{3} + 1)a}.
    \]
    \item À chacun des $3$ sommets du grand triangle, la région grisée est symétrique et formée de $2$ triangles rectangles isométriques à $ADE$ (d'aire $\frac{1}{2}\cdot a \cdot a\sqrt{3} = \frac{\sqrt{3}}{2}a^2$) dont on retire $2$ secteurs circulaires de rayon $a$ et d'angle $|\widehat{DAE}| = 60\degree = \frac{\pi}{3}$ (soit au total $6$ triangles rectangles et $6$ secteurs de $60\degree$ qui forment un disque complet d'aire $\pi a^2$ !). L'aire totale des $3$ régions vaut donc :
    \[
        S = 6 \times \frac{\sqrt{3}}{2}a^2 - \pi a^2 = \boxed{(3\sqrt{3} - \pi)a^2}.
    \]
\end{enumerate}

\paragraph{Corrigé de l'Exercice 5.8 (UCLouvain -- Juillet 2021, Q1 et Q2) :}
\begin{enumerate}[label=(\alph*)]
    \item Le rectangle $AGDH$ a ses diagonales $[AD]$ et $[GH]$ de même longueur qui se coupent en leur milieu $O$, qui est le centre du cercle $\mathcal{C}$ circonscrit à l'hexagone régulier $ABCDEF$. Par conséquent, le sommet $G$ (où $\widehat{AGD} = 90\degree$) appartient au cercle $\mathcal{C}$! Sur ce cercle $\mathcal{C}$, l'angle inscrit $\widehat{AGC}$ intercepte le même arc $\wideparen{AC}$ que l'angle inscrit $\widehat{ABC} = 120\degree$ (ou l'angle au centre $\widehat{AOC} = 120\degree$ situé du même côté), donc $|\widehat{AGC}| = 120\degree$. Dès lors :
    \[
        \alpha = |\widehat{DGC}| = |\widehat{AGC}| - |\widehat{AGD}| = 120\degree - 90\degree = \boxed{30\degree \quad \left(\frac{\pi}{6}\text{ rad}\right)}.
    \]
    \item Soit $T$ le point de tangence de la corde $[AB]$ avec le petit cercle de rayon $r$. Comme $(OT) \perp (AB)$, $T$ est le milieu de $[AB]$, donc $|AT| = L$. Par Pythagore dans $OTA$, $R^2 - r^2 = L^2$. L'aire de l'anneau vaut donc $\mathcal{A} = \pi R^2 - \pi r^2 = \pi(R^2 - r^2) = \boxed{\pi L^2}$.
    \item Dans le carré de côté $4$, la droite $(NM)$ est parallèle à $(AB)$ à distance $h = 2$ de $(AB)$, donc $\text{Aire}(\triangle ABE) = \frac{1}{2} \times 4 \times 2 = 4$. Comme $F$ est le milieu de $[EB]$, la médiane $[AF]$ partage le triangle $ABE$ en deux triangles $AEF$ et $ABF$ de même aire : $\mathcal{T} = \frac{1}{2}\text{Aire}(\triangle ABE) = \boxed{2}$.
\end{enumerate}

\paragraph{Corrigé de l'Exercice 5.9 (UCLouvain -- Septembre 2021, Q3 et Q4) :}
\begin{enumerate}[label=(\alph*)]
    \item $[AM]$ est la hauteur du triangle équilatéral de côté $2$, donc $|AM| = \boxed{\sqrt{3}}$. Les deux petits carrés ont pour côté $\frac{|AM|}{2} = \frac{\sqrt{3}}{2}$, donc $|AD| = \frac{\sqrt{3}}{2}$ et $|MB| = 1$. Les triangles rectangles $AMB$ (de cathètes $\sqrt{3}$ et $1$) et $DAM$ (de cathètes $\frac{\sqrt{3}}{2}$ et $\sqrt{3}$) ne sont pas directement ceux de $H$, mais comme $(AD) \perp (AM)$ et $(MB) \perp (AM)$, on a $(AD) \parallel (MB)$ et le rectangle $ADGM$ a pour diagonale $[DG]$... Attention : dans l'énoncé original, les points $D, H, B$ sont alignés (la droite $(DB)$ passe par $H \in [AB]$ ? Non, $H \in [AB]$ est sur la droite $(DG)$ ou $(DM)$ : dans le triangle $AMB$, $\tan\widehat{BAM} = \frac{1}{\sqrt{3}} \implies \widehat{BAM} = 30\degree$, et dans le rectangle $ADGM$ de côtés $\frac{\sqrt{3}}{2}$ et $\sqrt{3}$, $\tan\widehat{ADM} = \frac{\sqrt{3}}{\sqrt{3}/2} = 2$). En fait, sur le dessin de Septembre 2021 Q3, les triangles $AMB$ et $ADH$ sont semblables avec $|\widehat{ADH}| = |\widehat{AMB}| = 90\degree$ et $|\widehat{AHD}| = |\widehat{MAH}| = 30\degree$, ce qui donne $|AH| = |AB|\cdot\frac{|AD|}{|MB|} = 2\cdot\frac{\sqrt{3}/2}{1} = \boxed{\sqrt{3}}$. Le triangle $AMH$ est alors isocèle en $A$ ($|AM| = |AH| = \sqrt{3}$) d'angle au sommet $\widehat{MAH} = 30\degree$, d'où $\alpha = \widehat{MHA} = \frac{180\degree - 30\degree}{2} = \boxed{75\degree}$.
    \item Les triangles $AEG$ et $GEM$ ont le sommet $G$ en commun et leurs bases $[AE]$ et $[EM]$ alignées et égales ($E$ milieu de $[AM]$), donc le rapport de leurs aires vaut $r = \boxed{1}$. Par symétrie centrale du carré autour de $M$, $(BG) \parallel (DH)$. Dans le triangle $GBC$, la droite $(FH)$ parallèle à $(BG)$ passe par $F$ qui est au tiers de $[EC]$ ($\frac{|CF|}{|CE|} = \frac{1}{3}$), donc par Thalès $\frac{|CH|}{|CB|} = \frac{1}{3}$, soit $\frac{|BH|}{|HC|} = \boxed{2}$. Enfin, l'aire du parallélogramme $GBHD$ de base $|BH| = \frac{2}{3}c$ et de hauteur $c$ occupe la fraction $f = \boxed{\frac{2}{3}}$ de l'aire du carré.
\end{enumerate}

\paragraph{Corrigé de l'Exercice 5.10 (UCLouvain -- Juillet 2023, Q4a et Q4b) :}
\begin{enumerate}[label=(\alph*)]
    \item En projetant orthogonalement $A$ et $C$ sur les directions parallèle et perpendiculaire à $(XY)$, la diagonale $[AC]$ du carré $ABCD$ est l'hypoténuse d'un triangle rectangle de côtés $|AX| + |YC| = 13 + 8 = 21$ et $|XY| = 3$. Par Pythagore, $|AC|^2 = 21^2 + 3^2 = 441 + 9 = 450$. Or, dans un carré de côté $c$, $|AC|^2 = 2c^2$, d'où $c^2 = 225 \implies \boxed{c = 15}$.
    \item Notons $C$ et $D$ les deux régions latérales du rectangle (chacune dans un seul quart de disque). L'hypothèse $\text{Aire}(A) = \text{Aire}(B)$ entraîne que l'aire totale du rectangle $\ell \cdot L = \text{Aire}(A) + \text{Aire}(B) + \text{Aire}(C) + \text{Aire}(D) = \bigl(\text{Aire}(A) + \text{Aire}(C)\bigr) + \bigl(\text{Aire}(A) + \text{Aire}(D)\bigr)$ est exactement égale à la somme des aires des deux quarts de disque de rayon $\ell$ ! Ainsi :
    \[
        \ell \cdot L = 2 \times \frac{\pi\ell^2}{4} = \frac{\pi\ell^2}{2} \implies \boxed{L = \frac{\pi\ell}{2}}.
    \]
\end{enumerate}

% ----------------------------------------------------------------------------
\section{Corrigés du Chapitre 6 (Géométrie Analytique Plane et Coniques)}
% ----------------------------------------------------------------------------

\paragraph{Corrigé de l'Exercice 6.1 (UCLouvain -- Septembre 2023, Q1b) :}
\begin{enumerate}[label=(\alph*)]
    \item On a $\overrightarrow{AB} = (1-4,\; -2-(-3)) = (-3, 1)$. Tout point $C$ aligné avec $A$ et $B$ s'écrit $C = A + r\overrightarrow{AB} = (4-3r,\; -3+r)$. En imposant $C \in d \equiv y = 1+x$ :
    \[
        -3 + r = 1 + (4 - 3r) \iff 4r = 8 \implies r = 2 \implies \boxed{C(-2, -1)}.
    \]
    \item Un vecteur directeur orthogonal à $\overrightarrow{AB} = (-3, 1)$ est $\vec{v} = (1, 3)$. Tout point $D$ tel que $(BD) \perp (AB)$ s'écrit $D = B + s\vec{v} = (1+s,\; -2+3s)$. En imposant $D \in d \equiv y = 1+x$ :
    \[
        -2 + 3s = 1 + (1 + s) \iff 2s = 4 \implies s = 2 \implies \boxed{D(3, 4)}.
    \]
\end{enumerate}

\paragraph{Corrigé de l'Exercice 6.2 (UCLouvain -- Septembre 2021, Q1) :}
\begin{enumerate}[label=(\alph*)]
    \item La droite $\mathcal{A} \equiv 3x + y + 12 = 0$ a pour vecteur normal $\vec{n}_{\mathcal{A}} = (3, 1)$, qui est donc un vecteur directeur de toute perpendiculaire $\mathcal{B}$. Un vecteur normal à $\mathcal{B}$ est $\vec{n}_{\mathcal{B}} = (1, -3)$. Comme $P(2, -2) \in \mathcal{B}$, on a $1(x-2) - 3(y+2) = 0 \iff \boxed{\mathcal{B} \equiv x - 3y - 8 = 0}$.
    \item Soit $\Omega(x_c, y_c)$ le centre du cercle de rayon $R = \sqrt{10}$. Sa distance à $\mathcal{A}$ et à $\mathcal{B}$ vaut $\sqrt{10}$ :
    \[
        \frac{|3x_c + y_c + 12|}{\sqrt{10}} = \sqrt{10} \quad \text{et} \quad \frac{|x_c - 3y_c - 8|}{\sqrt{10}} = \sqrt{10}.
    \]
    Pour le quadrant contenant l'origine $(0,0)$ (où $3(0)+0+12 > 0$ et $0-3(0)-8 < 0$), on a $3x_c + y_c + 12 = 10$ et $x_c - 3y_c - 8 = -10$, soit le système :
    \[
        \begin{cases} 3x_c + y_c = -2 \\ x_c - 3y_c = -2 \end{cases} \implies 10x_c = -8 \implies \boxed{\Omega(-0,8\;;\; 0,4) = \left(-\frac{4}{5},\; \frac{2}{5}\right)}.
    \]
\end{enumerate}

\paragraph{Corrigé de l'Exercice 6.3 (UCLouvain -- Septembre 2022, Q1) :}
\begin{enumerate}[label=(\alph*)]
    \item Le milieu de $E(-2,-4)$ et $G(-5,5)$ est $M\left(-\frac{7}{2}, \frac{1}{2}\right)$, et $\overrightarrow{EG} = (-3, 9) = -3(1, -3)$. La médiatrice $m$ de $[EG]$ a donc pour vecteur normal $(1, -3)$ et passe par $M$ :
    \[
        1\left(x + \frac{7}{2}\right) - 3\left(y - \frac{1}{2}\right) = 0 \iff \boxed{m \equiv x - 3y + 5 = 0}.
    \]
    \item Le milieu de $E(-2,-4)$ et $F(3,1)$ est $N\left(\frac{1}{2}, -\frac{3}{2}\right)$ et $\overrightarrow{EF} = (5, 5) = 5(1, 1)$, donc la médiatrice de $[EF]$ est $(x - \frac{1}{2}) + (y + \frac{3}{2}) = 0 \iff x + y + 1 = 0$. L'intersection avec $x - 3y + 5 = 0$ donne $4y = 4 \implies y = 1$ et $x = -2$, soit le centre $\Omega(-2, 1)$. Le rayon vaut $R = |\Omega E| = \sqrt{0^2 + (-5)^2} = 5$, d'où l'équation du cercle circonscrit :
    \[
        \boxed{\mathcal{C} \equiv (x + 2)^2 + (y - 1)^2 = 25}.
    \]
    \item La distance du centre $\Omega(-2,1)$ au milieu $M(-3,5\;;\; 0,5)$ de la corde $[EG]$ vaut $|\Omega M| = \sqrt{(-1,5)^2 + (-0,5)^2} = \sqrt{2,5} = \sqrt{\frac{5}{2}}$. La flèche (largeur maximale du segment circulaire entre $[EG]$ et l'arc $\wideparen{EG}$) vaut $R - |\Omega M| = 5 - \sqrt{\frac{5}{2}}$, donc le rayon du plus grand cercle inscrit dans cette lunule vaut :
    \[
        r = \frac{R - |\Omega M|}{2} = \boxed{\frac{1}{2}\left(5 - \sqrt{\frac{5}{2}}\right) \approx 1,71}.
    \]
\end{enumerate}

\paragraph{Corrigé de l'Exercice 6.4 (UCLouvain -- Juillet 2021, Q3) :}
\begin{enumerate}[label=(\alph*)]
    \item La droite $d \equiv 4x + 3y - 12 = 0$ est tangente au cercle $x^2 + y^2 = a$ (centré à l'origine $O(0,0)$ et de rayon $R = \sqrt{a}$) ssi la distance de $O$ à $d$ est égale à $\sqrt{a}$ :
    \[
        a = \bigl(d(O, d)\bigr)^2 = \left(\frac{|0 + 0 - 12|}{\sqrt{4^2 + 3^2}}\right)^2 = \left(\frac{12}{5}\right)^2 = \boxed{\frac{144}{25}}.
    \]
    \item La droite $d$ a pour pente $-\frac{4}{3}$, donc toute droite $m \perp d$ a pour pente $\frac{3}{4}$ et s'écrit $y = \frac{3}{4}x + c \iff x = \frac{4}{3}(y - c)$. Substituons $x$ dans la parabole $\mathcal{P} \equiv \frac{y^2}{27} + \frac{x}{6} + \frac{1}{2} = 0$ :
    \[
        \frac{y^2}{27} + \frac{2}{9}(y - c) + \frac{1}{2} = 0 \iff y^2 + 6y + \left(\frac{27}{2} - 6c\right) = 0.
    \]
    La droite $m$ est tangente à $\mathcal{P}$ ssi le discriminant $\Delta$ est nul :
    \[
        \Delta = 6^2 - 4\left(\frac{27}{2} - 6c\right) = 36 - 54 + 24c = 0 \implies c = \frac{18}{24} = \frac{3}{4}.
    \]
    D'où $\boxed{m \equiv y = \frac{3}{4}x + \frac{3}{4} \iff 3x - 4y + 3 = 0}$, et le point de tangence (racine double $y = -3$) est $\boxed{T(-5, -3)}$.
\end{enumerate}

\paragraph{Corrigé de l'Exercice 6.5 (UCLouvain -- Juillet 2022, Q2) :}
\begin{enumerate}[label=(\alph*)]
    \item Le milieu de $C\left(-\frac{12}{5}, \frac{6}{5}\right)$ et $D(5, -3)$ est $M\left(\frac{13}{10}, -\frac{9}{10}\right)$, et $\overrightarrow{CD} = \left(\frac{37}{5}, -\frac{21}{5}\right) = \frac{1}{5}(37, -21)$. Un vecteur directeur de la médiatrice $m \perp [CD]$ est donc $\vec{u} = (21, 37)$, d'où :
    \[
        \boxed{m \equiv \begin{cases} x = \frac{13}{10} + 21k \\ y = -\frac{9}{10} + 37k \end{cases} \quad (k \in \R) \qquad \Longleftrightarrow \qquad 37x - 21y - 67 = 0}.
    \]
    \item L'ellipse $\mathcal{E} \equiv \frac{x^2}{a^2} + \frac{y^2}{b^2} = 1$ passe par $C\left(-\frac{12}{5}, \frac{6}{5}\right)$, donc $\frac{144}{a^2} + \frac{36}{b^2} = 25$. D'autre part, par la règle de dédoublement (ou en annulant le discriminant de l'intersection avec $d \equiv 8x + 9y = 30 \iff \frac{4}{15}x + \frac{3}{10}y = 1$), une droite $ux + vy = 1$ est tangente à $\frac{x^2}{a^2} + \frac{y^2}{b^2} = 1$ ssi $a^2 u^2 + b^2 v^2 = 1$, soit ici :
    \[
        \frac{16}{225}a^2 + \frac{9}{100}b^2 = 1.
    \]
    En posant $X = \frac{1}{a^2}$ et $Y = \frac{1}{b^2}$, la première équation donne $144X + 36Y = 25$, et la condition de tangence $\Delta = 0$ se réduit à $(4Y - 1)^2 = 0 \implies Y = \frac{1}{b^2} = \frac{1}{4}$ ($b^2 = 4$), puis $144X = 25 - 9 = 16 \implies X = \frac{1}{a^2} = \frac{1}{9}$ ($a^2 = 9$). D'où l'ellipse :
    \[
        \boxed{\mathcal{E} \equiv \frac{x^2}{9} + \frac{y^2}{4} = 1}.
    \]
\end{enumerate}

\paragraph{Corrigé de l'Exercice 6.6 (UCLouvain -- Juillet 2024, Q1a) :}
En complétant le carré en $x$ : $4(x^2 - 2x) + y^2 = 0 \iff 4\bigl((x-1)^2 - 1\bigr) + y^2 = 0 \iff 4(x-1)^2 + y^2 = 4$, soit :
\[
    \boxed{\frac{(x-1)^2}{1} + \frac{y^2}{4} = 1}.
\]
C'est une \textbf{ellipse} de centre $C(1, 0)$, de petit axe horizontal $2a = 2$ ($a=1$) et de \textbf{grand axe vertical} $2b = 4$ ($b=2$). La distance focale vaut $c = \sqrt{b^2 - a^2} = \sqrt{4-1} = \sqrt{3}$, les foyers sont $\boxed{F_1(1, -\sqrt{3}) \text{ et } F_2(1, \sqrt{3})}$, l'excentricité vaut $e = \frac{c}{b} = \frac{\sqrt{3}}{2}$ et les directrices sont $y = \pm \frac{b^2}{c} = \pm \frac{4}{\sqrt{3}}$.

\paragraph{Corrigé de l'Exercice 6.7 (UCLouvain -- Septembre 2025, Q1b) :}
Le segment reliant $A(2,0)$ et $B(2,2)$ est vertical et a pour milieu $M(2,1)$, donc sa médiatrice est la droite horizontale d'équation $\boxed{y = 1}$. En remplaçant $y = 1$ dans l'hyperbole $\frac{x^2}{4} - y^2 = 1$, on obtient $\frac{x^2}{4} - 1 = 1 \iff x^2 = 8 \implies x = \pm 2\sqrt{2}$, soit les deux points d'intersection $\boxed{C(-2\sqrt{2}, 1) \text{ et } D(2\sqrt{2}, 1)}$.

\paragraph{Corrigé de l'Exercice 6.8 (UCLouvain -- Juillet 2023, Q1a) :}
L'équation s'écrit $(x+2)^2 = 2(y+4) \iff y = \frac{1}{2}(x+2)^2 - 4$. C'est une parabole d'axe vertical $x = -2$, concave vers le haut ($2p = 2 \implies p = 1$ et $\frac{p}{2} = \frac{1}{2}$), de \textbf{sommet} $\boxed{S(-2, -4)}$, de \textbf{foyer} $\boxed{F(-2, -3,5)}$ et de \textbf{directrice} $\boxed{d \equiv y = -4,5}$. Ses intersections avec $Ox$ ($y=0$) vérifient $(x+2)^2 = 8 \implies \boxed{x = -2 \pm 2\sqrt{2}}$.

\paragraph{Corrigé de l'Exercice 6.9 (Classification et réduction de coniques) :}
\begin{enumerate}[label=(\alph*)]
    \item $4(x^2-4x) - 9(y^2+2y) = 29 \iff 4(x-2)^2 - 16 - 9(y+1)^2 + 9 = 29 \iff 4(x-2)^2 - 9(y+1)^2 = 36$, soit :
    \[
        \boxed{\frac{(x-2)^2}{9} - \frac{(y+1)^2}{4} = 1}.
    \]
    C'est une \textbf{hyperbole} de centre $C(2, -1)$, d'axe focal horizontal ($a=3, b=2$), de distance focale $c = \sqrt{9+4} = \sqrt{13}$, de foyers $F_{1,2}(2 \pm \sqrt{13}, -1)$, d'excentricité $e = \frac{\sqrt{13}}{3}$ et d'asymptotes $y + 1 = \pm \frac{2}{3}(x-2)$.
    \item $(y-3)^2 - 9 - 8x + 17 = 0 \iff \boxed{(y-3)^2 = 8(x-1)}$. C'est une \textbf{parabole} de sommet $S(1, 3)$, d'axe horizontal $y = 3$, avec $2p = 8 \implies p = 4$ ($\frac{p}{2} = 2$), de foyer $F(3, 3)$ et de directrice $x = -1$.
\end{enumerate}

\paragraph{Corrigé de l'Exercice 6.10 (Lieu bifocal et tangentes parallèles) :}
Le centre est $O(0,0)$, la distance focale est $c = 3$ (foyers sur $Ox$) et $2a = 10 \implies a = 5$ ($a^2 = 25$), d'où $b^2 = a^2 - c^2 = 25 - 9 = 16$ et l'ellipse $\boxed{\frac{x^2}{25} + \frac{y^2}{16} = 1}$.\\
Toute droite parallèle à $4x + 3y = 0$ s'écrit $y = -\frac{4}{3}x + k$. En substituant dans $16x^2 + 25y^2 = 400$ :
\[
    16x^2 + 25\left(\frac{16}{9}x^2 - \frac{8}{3}kx + k^2\right) = 400 \iff 544x^2 - 600kx + 9(25k^2 - 400) = 0.
\]
L'annulation du discriminant réduit $\Delta' = (-300k)^2 - 544 \cdot 9(25k^2 - 400) = 0$ (ou la formule $k^2 = a^2 m^2 + b^2 = 25\cdot\frac{16}{9} + 16 = \frac{544}{9}$) donne $k = \pm \frac{4\sqrt{34}}{3}$, soit les deux tangentes $\boxed{4x + 3y = \pm 4\sqrt{34}}$.

% ----------------------------------------------------------------------------
\section{Corrigés du Chapitre 7 (Géométrie Analytique dans l'Espace)}
% ----------------------------------------------------------------------------

\paragraph{Corrigé de l'Exercice 7.1 (UCLouvain -- Juillet 2025, Q2 : Éboulement rocheux) :}
\begin{enumerate}[label=(\alph*)]
    \item La trajectoire a pour équations paramétriques $x(t) = 10 + 4t$, $y(t) = 20 - 2t$, $z(t) = 100 - 4t$. L'altitude $z(t) = 80$ est atteinte lorsque $100 - 4t = 80 \implies t = 5$, ce qui donne le point $\boxed{M(30, 10, 80)}$. Comme $\|\vec{v}\| = \sqrt{4^2 + (-2)^2 + (-4)^2} = \sqrt{36} = 6$, la distance parcourue vaut $|AM| = 5 \times 6 = \boxed{30}$.
    \item Le plan du versant s'écrit $\Pi \equiv 0,6x - 0,8y + z - 85 = 0$, de vecteur normal $\vec{n} = (0,6\;;\; -0,8\;;\; 1)$. Pour que la trajectoire de vecteur directeur $\vec{w} = (a, b, c)$ suive le plan $\Pi$ (soit parallèle à $\Pi$), il faut que $\vec{w} \cdot \vec{n} = 0$, soit :
    \[
        0,6a - 0,8b + c = 0 \iff \boxed{c = 0,8b - 0,6a}.
    \]
    \item Pour le vecteur observé $\vec{v} = (4, -2, -4)$, on calcule $0,8b - 0,6a = 0,8(-2) - 0,6(4) = -1,6 - 2,4 = -4 = c$ ! La condition de parallélisme au plan $\Pi$ est donc vérifiée (et de plus $A(10,20,100) \in \Pi$ car $-0,6(10)+0,8(20)+85 = 95 \neq 100$ : le point de départ $A$ est $5\text{ m}$ au-dessus du plan, ou selon le barème officiel de Juillet 2025 Q2c, il suffisait de tester si $(4,-2,-4)$ vérifie la relation 2.b).
\end{enumerate}

\paragraph{Corrigé de l'Exercice 7.2 (UCLouvain -- Septembre 2024, Q2) :}
\begin{enumerate}[label=(\alph*)]
    \item On a $\overrightarrow{AB} = (2 - (-7),\; 4 - 10,\; 2 - 5) = (9, -6, -3)$. Tout point $P \in (AB)$ s'écrit $P(s) = A + s\overrightarrow{AB} = (-7+9s,\; 10-6s,\; 5-3s)$, de sorte que $\overrightarrow{CP} = P(s) - C = (-10+9s,\; -6s,\; -2-3s)$. La condition $\overrightarrow{CP} \perp \overrightarrow{AB}$ donne :
    \[
        \overrightarrow{AB} \cdot \overrightarrow{CP} = 9(-10+9s) - 6(-6s) - 3(-2-3s) = -84 + 126s = 0 \implies s = \frac{84}{126} = \frac{2}{3}.
    \]
    En remplaçant $s = \frac{2}{3}$, on obtient $\boxed{P(-1, 6, 3)}$.
    \item Le plan médiateur du segment $[AC]$ passe par le milieu $M\left(-2, 10, 6\right)$ de $A(-7,10,5)$ et $C(3,10,7)$ et a pour vecteur normal $\overrightarrow{AC} = (10, 0, 2) = 2(5, 0, 1)$. Son équation cartésienne est donc $5(x+2) + 0 + 1(z-6) = 0 \iff 5x + z + 4 = 0$.\\
    Or tout point $Q \in (BC)$ (où $\overrightarrow{BC} = (1, 6, 5)$) s'écrit $Q(t) = B + t\overrightarrow{BC} = (2+t,\; 4+6t,\; 2+5t)$. En substituant dans $5x + z + 4 = 0$ :
    \[
        5(2+t) + (2+5t) + 4 = 0 \iff 10t + 16 = 0 \implies t = -\frac{8}{5} \implies \boxed{Q\left(\frac{2}{5},\; -\frac{28}{5},\; -6\right)}.
    \]
\end{enumerate}

\paragraph{Corrigé de l'Exercice 7.3 (UCLouvain -- Septembre 2023, Q2) :}
\begin{enumerate}[label=(\alph*)]
    \item La droite $d_1 \perp \Pi$ a pour vecteur directeur le vecteur normal $\vec{n} = (-2, 1, -3)$ du plan $\Pi$ et passe par $A(5, -1, 7)$ :
    \[
        \boxed{d_1 \equiv \begin{cases} x = 5 - 2t \\ y = -1 + t \\ z = 7 - 3t \end{cases} \quad (t \in \R)}.
    \]
    Son point de percée $B \in d_1 \cap \Pi$ vérifie $-2(5-2t) + (-1+t) - 3(7-3t) + 4 = 0 \iff 14t - 28 = 0 \implies t = 2$, d'où $\boxed{B(1, 1, 1)}$.
    \item La droite $d_2 \parallel d_1$ passant par $C(-6, 5, -7)$ est elle aussi perpendiculaire au plan $\Pi$! Ses équations paramétriques sont $x = -6-2u, y = 5+u, z = -7-3u$, et son point de percée $D \in d_2 \cap \Pi$ vérifie $-2(-6-2u) + (5+u) - 3(-7-3u) + 4 = 0 \iff 14u + 42 = 0 \implies u = -3$, soit $D(0, 2, 2)$.\\
    Puisque le plan $\Pi$ est perpendiculaire aux deux droites parallèles $d_1$ et $d_2$, la distance entre $d_1$ et $d_2$ est exactement la distance entre leurs points de percée $B(1,1,1)$ et $D(0,2,2)$ :
    \[
        d(d_1, d_2) = |BD| = \sqrt{(0-1)^2 + (2-1)^2 + (2-1)^2} = \boxed{\sqrt{3}}.
    \]
\end{enumerate}

\paragraph{Corrigé de l'Exercice 7.4 (UCLouvain -- Juillet 2023, Q2) :}
\begin{enumerate}[label=(\alph*)]
    \item Le point d'intersection $B \in d_1$ s'écrit $B(t-3,\; 3t-4,\; -t+5)$, et le vecteur $\overrightarrow{BA} = A - B = (4-t,\; 5-3t,\; t-3)$ doit être orthogonal au vecteur directeur $\vec{d}_1 = (1, 3, -1)$ de $d_1$ :
    \[
        \vec{d}_1 \cdot \overrightarrow{BA} = 1(4-t) + 3(5-3t) - 1(t-3) = 22 - 11t = 0 \implies t = 2.
    \]
    On a donc $B(-1, 2, 3)$ et $\vec{d}_2 = \overrightarrow{BA} = (2, -1, -1)$, d'où $\boxed{d_2 \equiv \begin{cases} x = -1 + 2s \\ y = 2 - s \\ z = 3 - s \end{cases}}$.
    \item Le plan $\mathcal{P}$ contenant $d_1$ et $d_2$ a pour vecteur normal :
    \[
        \vec{n} = \vec{d}_1 \wedge \vec{d}_2 = (1, 3, -1) \wedge (2, -1, -1) = (-4, -1, -7) = -(4, 1, 7),
    \]
    et passe par $A(1, 1, 2)$ ($4(1)+1+7(2)=19$), d'où $\boxed{\mathcal{P} \equiv 4x + y + 7z = 19}$.
    \item Dans le plan $\mathcal{P}$, le vecteur $\vec{d}_2 = (2, -1, -1)$ est perpendiculaire à $d_1$ et sa norme vaut justement $\|\vec{d}_2\| = \sqrt{4+1+1} = \sqrt{6}$ ! Par conséquent, l'ensemble des points de $\mathcal{P}$ situés à distance $\sqrt{6}$ de $d_1$ est formé des deux droites $B + r\vec{d}_1 \pm \vec{d}_2$ :
    \begin{itemize}
        \item \emph{1\ier{} cas ($+\vec{d}_2$) :} $C = (-1+r+2,\; 2+3r-1,\; 3-r-1) = (1+r,\; 1+3r,\; 2-r)$. L'intersection avec $x = 2z$ donne $1+r = 2(2-r) \implies r = 1 \implies \boxed{C_1(2, 4, 1)}$.
        \item \emph{2\ieme{} cas ($-\vec{d}_2$) :} $C = (-1+r-2,\; 2+3r+1,\; 3-r+1) = (-3+r,\; 3+3r,\; 4-r)$. L'intersection avec $x = 2z$ donne $-3+r = 2(4-r) \implies r = \frac{11}{3} \implies \boxed{C_2\left(\frac{2}{3},\; 14,\; \frac{1}{3}\right)}$.
    \end{itemize}
\end{enumerate}

\paragraph{Corrigé de l'Exercice 7.5 (UCLouvain -- Juillet 2022, Q1) :}
\begin{enumerate}[label=(\alph*)]
    \item La boule ouverte de centre $A(7, -2, 8)$ et de rayon $R = 10$ a pour inéquation :
    \[
        \boxed{\mathcal{L} \equiv (x - 7)^2 + (y + 2)^2 + (z - 8)^2 < 100}.
    \]
    \item La distance de $A(7, -2, 8)$ au plan $p \equiv 3x - 6y + 2z + 7 = 0$ vaut :
    \[
        d = d(A, p) = \frac{|3(7) - 6(-2) + 2(8) + 7|}{\sqrt{3^2 + (-6)^2 + 2^2}} = \frac{|21 + 12 + 16 + 7|}{\sqrt{49}} = \frac{56}{7} = 8\text{ m}.
    \]
    Le rayon du disque de section vaut $r = \sqrt{R^2 - d^2} = \sqrt{100 - 64} = 6\text{ m}$, et son aire vaut $\pi r^2 = \boxed{36\pi\text{ m}^2}$.
    \item Le plan $q \equiv x + ay + bz = c$ a pour normale $\vec{n}_q = (1, a, b)$.
    \begin{itemize}
        \item $q \perp p \implies \vec{n}_q \cdot \vec{n}_p = 3 - 6a + 2b = 0 \implies b = 3a - \frac{3}{2}$.
        \item $q$ fait un angle de $45\degree$ avec le plan horizontal (de normale $\vec{k} = (0,0,1)$) $\implies \frac{|b|}{\sqrt{1+a^2+b^2}} = \cos 45\degree = \frac{\sqrt{2}}{2} \iff b^2 = 1 + a^2$.
    \end{itemize}
    En substituant $b = 3a - \frac{3}{2}$ dans $b^2 = 1 + a^2$, on obtient $8a^2 - 9a + \frac{5}{4} = 0 \implies a = \frac{9 \pm \sqrt{41}}{16}$. Pour que $B(5,5,5)$ soit en dessous de $q$ ($z_B = 5 < z_q(5,5)$), on retient $a = \frac{9+\sqrt{41}}{16} \approx 0,96$, $b = \frac{3\sqrt{41}+3}{16} \approx 1,39$ et $c = 7 - 2a + 8b \approx 16,18$.
\end{enumerate}

\paragraph{Corrigé de l'Exercice 7.6 (UCLouvain -- Juillet 2021, Q4 : Alpiniste) :}
\begin{enumerate}[label=(\alph*)]
    \item L'alpiniste touche le sol $z(t) = 0$ lorsque $-18t + 144 = 0 \implies t = 8\text{ min}$. Le vecteur déplacement entre $t=0$ et $t=8$ est $\Delta\vec{r} = 8\vec{u} = 8(-1, 6, -18)$, et la distance parcourue vaut :
    \[
        d = 8\sqrt{(-1)^2 + 6^2 + (-18)^2} = 8\sqrt{1 + 36 + 324} = 8\sqrt{361} = 8 \times 19 = \boxed{152\text{ m}}.
    \]
    \item Une droite horizontale $h$ a un vecteur directeur de composante verticale nulle $\vec{h} = (h_x, h_y, 0)$. Pour qu'elle soit perpendiculaire à la trajectoire de vecteur directeur $\vec{u} = (-1, 6, -18)$, il faut $\vec{h} \cdot \vec{u} = -h_x + 6h_y + 0 = 0$, ce qui donne par exemple $\boxed{\vec{h} = (6, 1, 0)}$.
    \item Le plan de la paroi passe par le point initial $A(-50, -100, 144)$ et est dirigé par $\vec{u} = (-1, 6, -18)$ et $\vec{h} = (6, 1, 0)$. Leur produit vectoriel donne la normale $\vec{n} = \vec{u} \wedge \vec{h} = (18, -108, -37)$, d'où l'équation cartésienne du plan :
    \[
        \boxed{18(x + 50) - 108(y + 100) - 37(z - 144) = 0 \iff 18x - 108y - 37z - 4572 = 0}.
    \]
\end{enumerate}

\paragraph{Corrigé de l'Exercice 7.7 (UCLouvain -- Septembre 2021, Q2 : Pyramide) :}
\begin{enumerate}[label=(\alph*)]
    \item Le plan médiateur $\mathcal{L}$ de $B(-3,-1,0)$ et $C(4,0,0)$ passe par leur milieu $M\left(\frac{1}{2}, -\frac{1}{2}, 0\right)$ et a pour vecteur normal $\overrightarrow{BC} = (7, 1, 0)$ :
    \[
        7\left(x - \frac{1}{2}\right) + 1\left(y + \frac{1}{2}\right) = 0 \iff \boxed{\mathcal{L} \equiv 7x + y - 3 = 0}.
    \]
    \item Comme $C(4,0,0)$ ne vérifie pas l'équation du plan $\Pi \equiv -3x + y + z - 8 = 0$ ($-12 - 8 = -20 \neq 0$), la face contenue dans $\Pi$ est nécessairement la face $ABD$ ! Le sommet $D$ appartient donc à la fois à $z = 5$, au plan médiateur $\mathcal{L} \equiv 7x + y - 3 = 0$ (car $|DB| = |DC|$) et au plan $\Pi \equiv -3x + y + z - 8 = 0$. Avec $z = 5$, le système :
    \[
        \begin{cases} 7x + y = 3 \\ -3x + y = 3 \end{cases} \implies x = 0 \text{ et } y = 3 \implies \boxed{D(0, 3, 5)}.
    \]
    \item Le point $A(x_A, y_A, 0)$ appartient à $z = 0$ et à $\Pi \equiv -3x + y + z - 8 = 0$, donc $y_A = 3x_A + 8$. De plus, $|DA|^2 = |DC|^2 = (4-0)^2 + (0-3)^2 + (0-5)^2 = 16 + 9 + 25 = 50$. Ainsi :
    \[
        x_A^2 + (3x_A + 8 - 3)^2 + 5^2 = 50 \iff 10x_A^2 + 30x_A = 0 \iff 10x_A(x_A + 3) = 0.
    \]
    La solution $x_A = -3$ correspond au point $B(-3,-1,0)$, et la solution $x_A = 0$ donne $y_A = 8$, soit $\boxed{A(0, 8, 0)}$.
\end{enumerate}

\paragraph{Corrigé de l'Exercice 7.8 (Symétrique d'un point par rapport à un plan) :}
\begin{enumerate}[label=(\alph*)]
    \item La distance de $A(5, -3, 7)$ au plan $\Pi \equiv 2x - y + 2z - 6 = 0$ vaut :
    \[
        d(A, \Pi) = \frac{|2(5) - (-3) + 2(7) - 6|}{\sqrt{2^2 + (-1)^2 + 2^2}} = \frac{|10 + 3 + 14 - 6|}{3} = \frac{21}{3} = \boxed{7}.
    \]
    \item La perpendiculaire à $\Pi$ passant par $A$ a pour équations $x = 5 + 2t, y = -3 - t, z = 7 + 2t$. Son point de percée $H \in \Pi$ vérifie $2(5+2t) - (-3-t) + 2(7+2t) - 6 = 0 \iff 9t + 21 = 0 \implies t = -\frac{7}{3}$, d'où $\boxed{H\left(\frac{1}{3},\; -\frac{2}{3},\; \frac{7}{3}\right)}$. Le symétrique $A'$ correspond à $2t = -\frac{14}{3}$, soit $\boxed{A'\left(-\frac{13}{3},\; \frac{5}{3},\; -\frac{7}{3}\right)}$.
\end{enumerate}

\paragraph{Corrigé de l'Exercice 7.9 (Intersection de deux plans et distance point--droite) :}
\begin{enumerate}[label=(\alph*)]
    \item En additionnant $\Pi_1 \equiv x + y - z = 1$ et $\Pi_2 \equiv 2x - y + z = 2$, on obtient $3x = 3 \implies x = 1$. En posant $z = t$, la première équation donne $y = z + 1 - x = t$. D'où la droite :
    \[
        \boxed{d \equiv \begin{cases} x = 1 \\ y = t \\ z = t \end{cases} \quad (t \in \R), \quad \text{passant par } A(1,0,0) \text{ et de vecteur directeur } \vec{u} = (0, 1, 1)}.
    \]
    \item Un point $H(1, t, t) \in d$ est le projeté orthogonal de $P(3, 1, 0)$ ssi $\overrightarrow{PH} \cdot \vec{u} = (-2, t-1, t)\cdot(0, 1, 1) = 2t - 1 = 0 \implies t = \frac{1}{2}$, d'où $\boxed{H\left(1, \frac{1}{2}, \frac{1}{2}\right)}$ et $d(P, d) = |PH| = \sqrt{(-2)^2 + \left(-\frac{1}{2}\right)^2 + \left(\frac{1}{2}\right)^2} = \boxed{\frac{3\sqrt{2}}{2}}$.
\end{enumerate}

\paragraph{Corrigé de l'Exercice 7.10 (Sphère tangente et plan tangent) :}
\begin{enumerate}[label=(\alph*)]
    \item Le rayon $R$ de la sphère centrée en $\Omega(1, -2, 3)$ et tangente à $\Pi \equiv 2x + 3y - 6z + 5 = 0$ est la distance de $\Omega$ à $\Pi$ :
    \[
        R = d(\Omega, \Pi) = \frac{|2(1) + 3(-2) - 6(3) + 5|}{\sqrt{2^2 + 3^2 + (-6)^2}} = \frac{|-17|}{7} = \frac{17}{7},
    \]
    d'où l'équation de la sphère : $\boxed{(x-1)^2 + (y+2)^2 + (z-3)^2 = \frac{289}{49}}$.
    \item Pour $T(3, 3, 2)$ et la sphère $\mathcal{S}'$ de centre $C(2, 1, -1)$, on a $(3-2)^2 + (3-1)^2 + (2+1)^2 = 1 + 4 + 9 = 14$, donc $T \in \mathcal{S}'$. Le plan tangent en $T$ passe par $T(3,3,2)$ et a pour vecteur normal le rayon $\overrightarrow{CT} = (1, 2, 3)$ :
    \[
        1(x - 3) + 2(y - 3) + 3(z - 2) = 0 \iff \boxed{x + 2y + 3z - 15 = 0}.
    \]
\end{enumerate}


\end{document}
